% Generated mechanically from frozen input inputs/p05/ja/descent.tex.
% Do not edit this generated file; edit the bound locale source instead.

% OK, start here.
%

\setcounter{chapter}{34}
\chapter{降下}


\phantomsection
\label{descent-section-phantom}


\section{はじめに}
\label{descent-section-introduction}

\noindent
スキーム上の位相を扱う章
（Topologies, Section \ref{topologies-section-introduction} 参照）では、
スキームの Zariski, \'etale, fppf, 平滑, syntomic および fpqc 被覆を導入した。
本章では、そのような被覆に沿ってスキーム上のどのような構造を
降下できるかを論じる。
たとえば \cite{Gr-I}, \cite{Gr-II}, \cite{Gr-III},
\cite{Gr-IV}, \cite{Gr-V}, \cite{Gr-VI} を見よ。
本章はまた、準連接層やスキームなどの降下問題という比較的形式化されて
いない設定において、降下、降下データ、有効降下データの概念を導入する
ことも目的とする。サイト上のスタックという形式的な概念については、
スタックを扱う章で論じる
（Stacks, Section \ref{stacks-section-introduction} 参照）。

\section{準連接層の降下データ}
\label{descent-section-equivalence}

\noindent
本章では、射影
$\text{pr}_i : X \times \ldots \times X \to X$
の添字を $i = 0$ から付ける規約を用いる。したがって、
$\text{pr}_0, \text{pr}_1 : X \times X  \to X$,
$\text{pr}_0, \text{pr}_1, \text{pr}_2 : X \times X \times X  \to X$,
などとなる。


\begin{definition}
\label{descent-definition-descent-datum-quasi-coherent}
$S$ をスキームとし、$\{f_i : S_i \to S\}_{i \in I}$ を
$S$ を終域とする射の族とする。
\begin{enumerate}
\item 与えられた族に関する{\it 準連接層の降下データ
$(\mathcal{F}_i, \varphi_{ij})$}とは、各 $i \in I$ に対する
$S_i$ 上の準連接層 $\mathcal{F}_i$ と、各組
$(i, j) \in I^2$ に対する準連接
$\mathcal{O}_{S_i \times_S S_j}$-加群の同型
$\varphi_{ij} : \text{pr}_0^*\mathcal{F}_i \to \text{pr}_1^*\mathcal{F}_j$
からなり、任意の添字の三つ組 $(i, j, k) \in I^3$ に対して
次の図式
$$
\xymatrix{
\text{pr}_0^*\mathcal{F}_i \ar[rd]_{\text{pr}_{01}^*\varphi_{ij}}
\ar[rr]_{\text{pr}_{02}^*\varphi_{ik}} & &
\text{pr}_2^*\mathcal{F}_k \\
& \text{pr}_1^*\mathcal{F}_j \ar[ru]_{\text{pr}_{12}^*\varphi_{jk}} &
}
$$
が $\mathcal{O}_{S_i \times_S S_j \times_S S_k}$-加群の図式として
可換となるものをいう。これを{\it コサイクル条件}という。
\item {\it 降下データの射
$\psi : (\mathcal{F}_i, \varphi_{ij}) \to
(\mathcal{F}'_i, \varphi'_{ij})$}とは、
$\mathcal{O}_{S_i}$-加群の射
$\psi_i : \mathcal{F}_i \to \mathcal{F}'_i$ の族
$\psi = (\psi_i)_{i\in I}$ であって、すべての図式
$$
\xymatrix{
\text{pr}_0^*\mathcal{F}_i \ar[r]_{\varphi_{ij}} \ar[d]_{\text{pr}_0^*\psi_i}
& \text{pr}_1^*\mathcal{F}_j \ar[d]^{\text{pr}_1^*\psi_j} \\
\text{pr}_0^*\mathcal{F}'_i \ar[r]^{\varphi'_{ij}} &
\text{pr}_1^*\mathcal{F}'_j \\
}
$$
が可換となるものをいう。
\end{enumerate}
\end{definition}

\noindent
念頭に置くべき好例は次のものである。
$S = \bigcup S_i$ が開被覆であると仮定する。
この場合、Sheaves, Section \ref{sheaves-section-glueing-sheaves} で
集合の層の降下データをすでに見ており、そこではこれを
「与えられた被覆に関する集合の層の貼り合わせデータ」と呼んだ。
さらに、貼り合わせデータの圏が $S$ 上の層の圏と同値であることを証明した。
以下では、$\{S_i \to S\}_{i\in I}$ が fpqc 被覆であるときに、
上の設定における同様の結果を示す。

\medskip\noindent
極端な場合として、被覆 $\{S \to S\}$ が $\text{id}_S$ で
与えられるとき、降下データは必ず
$(\mathcal{F}, \text{id}_\mathcal{F})$ の形をもつ。
コサイクル条件により、この場合に許される射は
$\mathcal{F}$ 上の恒等射だけである。
次の補題は、とくに任意の準連接 $\mathcal{O}_S$-加群の層に対し、
任意に与えられた族に関する降下データが一意に付随することを示す。

\begin{lemma}
\label{descent-lemma-refine-descent-datum}
$\mathcal{U} = \{U_i \to U\}_{i \in I}$ および
$\mathcal{V} = \{V_j \to V\}_{j \in J}$
を、それぞれ固定した終域をもつスキーム射の族とする。
$(g, \alpha : I \to J, (g_i)) : \mathcal{U} \to \mathcal{V}$
を、固定した終域をもつ射の族の射とする
（Sites, Definition \ref{sites-definition-morphism-coverings} 参照）。
$(\mathcal{F}_j, \varphi_{jj'})$ を、族
$\{V_j \to V\}_{j \in J}$ に関する準連接層の降下データとする。このとき、
\begin{enumerate}
\item 系
$$
\left(g_i^*\mathcal{F}_{\alpha(i)},
(g_i \times g_{i'})^*\varphi_{\alpha(i)\alpha(i')}\right)
$$
は族 $\{U_i \to U\}_{i \in I}$ に関する降下データである。
\item この構成は降下データ
$(\mathcal{F}_j, \varphi_{jj'})$ について函手的である。
\item $g = g'$ を満たす、固定した終域をもつ射の族の第二の射
$(g', \alpha' : I \to J, (g'_i))$ が与えられたとき、
降下データの函手的同型
$$
(g_i^*\mathcal{F}_{\alpha(i)},
(g_i \times g_{i'})^*\varphi_{\alpha(i)\alpha(i')})
\cong
((g'_i)^*\mathcal{F}_{\alpha'(i)},
(g'_i \times g'_{i'})^*\varphi_{\alpha'(i)\alpha'(i')}).
$$
\end{enumerate}
\end{lemma}

\begin{proof}
省略する。ヒント：部分 (3) の降下データの同型を与える射
$g_i^*\mathcal{F}_{\alpha(i)} \to (g'_i)^*\mathcal{F}_{\alpha'(i)}$
は、射
$(g_i, g'_i) : U_i \to V_{\alpha(i)} \times_V V_{\alpha'(i)}$
による射 $\varphi_{\alpha(i)\alpha'(i)}$ の引き戻しである。
\end{proof}

\noindent
任意の族 $\mathcal{U} = \{S_i \to S\}_{i \in I}$ は、一意な仕方で
自明被覆 $\{S \to S\}$ の細分となる。
$S$ 上の準連接層 $\mathcal{F}$ に対し、上の手順で得られる
$\mathcal{U}$ に関する降下データを単に
$(\mathcal{F}|_{S_i}, can)$ と書く。

\begin{definition}
\label{descent-definition-descent-datum-effective-quasi-coherent}
$S$ をスキームとする。
$\{S_i \to S\}_{i \in I}$ を $S$ を終域とする射の族とする。
\begin{enumerate}
\item $\mathcal{F}$ を準連接 $\mathcal{O}_S$-加群とする。
被覆 $\{S \to S\}$ に関する $\mathcal{F}$ 上の一意な降下データを
{\it 自明降下データ}と呼ぶ。
\item 自明降下データの $\{S_i \to S\}$ への引き戻しを
{\it 標準降下データ}と呼び、$(\mathcal{F}|_{S_i}, can)$ と書く。
\item 与えられた被覆に関する準連接層の降下データ
$(\mathcal{F}_i, \varphi_{ij})$ が{\it 有効}であるとは、
$S$ 上の準連接層 $\mathcal{F}$ が存在して
$(\mathcal{F}_i, \varphi_{ij})$ が
$(\mathcal{F}|_{S_i}, can)$ と同型になることをいう。
\end{enumerate}
\end{definition}

\begin{lemma}
\label{descent-lemma-zariski-descent-effective}
$S$ をスキームとし、$S = \bigcup U_i$ を開被覆とする。
族 $\mathcal{U} = \{U_i \to S\}$ に関する準連接層の
任意の降下データは有効である。さらに、準連接
$\mathcal{O}_S$-加群の圏から $\mathcal{U}$ に関する
降下データの圏への函手は充満忠実である。
\end{lemma}

\begin{proof}
これは Sheaves, Section \ref{sheaves-section-glueing-sheaves} と、
準連接性が局所的性質であること
（Modules, Definition \ref{modules-definition-quasi-coherent} 参照）
から直ちに従う。
\end{proof}

\noindent
さらに進むには、まず環上の加群の場合を調べる必要がある。










\section{加群の降下}
\label{descent-section-descent-modules}

\noindent
$R \to A$ を環準同型とする。
Simplicial, Example \ref{simplicial-example-push-outs-simplicial-object}
により、これから余単体的 $R$-代数
$$
\xymatrix{
A
\ar@<1ex>[r]
\ar@<-1ex>[r]
&
A \otimes_R A
\ar@<0ex>[l]
\ar@<2ex>[r]
\ar@<0ex>[r]
\ar@<-2ex>[r]
&
A \otimes_R A \otimes_R A
\ar@<1ex>[l]
\ar@<-1ex>[l]
}
$$
が得られる。これを $(A/R)_\bullet$ と書く。したがって
$(A/R)_n$ は $R$ 上の $A$ の $(n + 1)$ 重テンソル積である。
写像 $\varphi : [n] \to [m]$ に対し、$R$-代数準同型
$(A/R)_\bullet(\varphi)$ は
$$
a_0 \otimes \ldots \otimes a_n
\longmapsto
\prod\nolimits_{\varphi(i) = 0} a_i
\otimes
\prod\nolimits_{\varphi(i) = 1} a_i
\otimes \ldots \otimes
\prod\nolimits_{\varphi(i) = m} a_i
$$
で与えられる。ここでは空積を $1$ とする規約を用いる。
したがって、Simplicial, Section
\ref{simplicial-section-cosimplicial-object} の記法による最初のいくつかの写像は
$$
\begin{matrix}
\delta^1_0 & : & a_0 & \mapsto & 1 \otimes a_0 \\
\delta^1_1 & : & a_0 & \mapsto & a_0 \otimes 1 \\
\sigma^0_0 & : & a_0 \otimes a_1 & \mapsto & a_0a_1 \\
\delta^2_0 & : & a_0 \otimes a_1 & \mapsto & 1 \otimes a_0 \otimes a_1 \\
\delta^2_1 & : & a_0 \otimes a_1 & \mapsto & a_0 \otimes 1 \otimes a_1 \\
\delta^2_2 & : & a_0 \otimes a_1 & \mapsto & a_0 \otimes a_1 \otimes 1 \\
\sigma^1_0 & : & a_0 \otimes a_1 \otimes a_2 & \mapsto & a_0a_1 \otimes a_2 \\
\sigma^1_1 & : & a_0 \otimes a_1 \otimes a_2 & \mapsto & a_0 \otimes a_1a_2
\end{matrix}
$$
などである。

\medskip\noindent
$R$-加群 $M$ から余単体的 $(A/R)_\bullet$-加群
$(A/R)_\bullet \otimes_R M$ が得られる。言い換えると、
$M_n = (A/R)_n \otimes_R M$ と置き、$R$-代数準同型
$(A/R)_n \to (A/R)_m$ を用いて
$M \otimes_R (A/R)_\bullet$ 上の対応する写像を定める。

\medskip\noindent
加群の設定における準連接層の降下データの類似物は次のものである。

\begin{definition}
\label{descent-definition-descent-datum-modules}
$R \to A$ を環準同型とする。
\begin{enumerate}
\item {\it $R \to A$ に関する加群の降下データ $(N, \varphi)$}
とは、$A$-加群 $N$ と $A \otimes_R A$-加群の同型
$$
\varphi : N \otimes_R A \to A \otimes_R N
$$
からなり、{\it コサイクル条件}、すなわち
$A \otimes_R A \otimes_R A$-加群準同型の図式
$$
\xymatrix{
N \otimes_R A \otimes_R A \ar[rr]_{\varphi_{02}}
\ar[rd]_{\varphi_{01}}
& &
A \otimes_R A \otimes_R N \\
& A \otimes_R N \otimes_R A \ar[ru]_{\varphi_{12}} &
}
$$
が可換となるものをいう（記法については下を見よ）。
\item {\it 降下データの射
$(N, \varphi) \to (N', \varphi')$}とは、図式
$$
\xymatrix{
N \otimes_R A \ar[r]_\varphi \ar[d]_{\psi \otimes \text{id}_A} &
A \otimes_R N \ar[d]^{\text{id}_A \otimes \psi} \\
N' \otimes_R A \ar[r]^{\varphi'} &
A \otimes_R N'
}
$$
を可換にする $A$-加群準同型 $\psi : N \to N'$ のことである。
\end{enumerate}
\end{definition}

\noindent
定義では、
$\varphi_{01} = \varphi \otimes \text{id}_A$,
$\varphi_{12} = \text{id}_A \otimes \varphi$ と書き、
$\varphi(n \otimes 1) = \sum a_i \otimes n_i$ のとき
$\varphi_{02}(n \otimes 1 \otimes 1) = \sum a_i \otimes 1 \otimes n_i$
と定める。これら三つはいずれも
$A \otimes_R A \otimes_R A$-加群準同型である。同値な記述として
$$
\varphi_{ij}
=
\varphi \otimes_{(A/R)_1, \ (A/R)_\bullet(\tau^2_{ij})} (A/R)_2
$$
である。ここで $\tau^2_{ij} : [1] \to [2]$ は
$0 \mapsto i$, $1 \mapsto j$ で与えられる写像である。実際、
$(A/R)_{\bullet}(\tau^2_{02})(a_0 \otimes a_1) =
a_0 \otimes 1 \otimes a_1$,
であり、他も同様である\footnote{
$\tau^2_{ij} = \delta^2_k$ であることに注意せよ。ただし
$\{i, j, k\} = [2] = \{0, 1, 2\}$ である。
Simplicial, Definition \ref{simplicial-definition-face-degeneracy} を見よ。}。

\medskip\noindent
次の補題を述べるため、さらに記法を導入する。
$(N, \varphi)$ を環準同型 $R \to A$ に関する降下データとする。
$n \geq 0$ および $i \in [n]$ に対し、
$$
N_{n, i} =
A \otimes_R
\ldots
\otimes_R A \otimes_R N \otimes_R A \otimes_R
\ldots
\otimes_R A
$$
と置く。ここで因子 $N$ は $i$ 番目の位置にあり、これは
$(A/R)_n$-加群である。写像
$\tau^n_i : [0] \to [n]$, $0 \mapsto i$ を導入すると、
$$
N_{n, i} = N \otimes_{(A/R)_0, \ (A/R)_\bullet(\tau^n_i)} (A/R)_n
$$
となる。$0 \leq i \leq j \leq n$ に対し、
$\tau^n_{ij} : [1] \to [n]$ を $0$ を $i$ に、$1$ を $j$ に
写す写像とする。上と同様に、準同型 $\varphi$ は $i < j$ のとき
$(A/R)_n$-加群の同型
$$
\varphi^n_{ij}
=
\varphi \otimes_{(A/R)_1, \ (A/R)_\bullet(\tau^n_{ij})} (A/R)_n :
N_{n, i} \longrightarrow N_{n, j}
$$
を誘導する。$i = j$ のときは
$\varphi^n_{ij} = \text{id}$ と置く。これらはすべて同型なので、
因子 $N$ を任意の位置へ移せる。コサイクル条件はまさに、
その移し方（二つの同型を合成するか、一度に移すかなど）に
結果が依存しないことを意味する。最後に、任意の
$\beta : [n] \to [m]$ に対し、射
$$
N_{\beta, i} : N_{n, i} \to N_{m, \beta(i)}
$$
を、すべての $n \in N$ に対して
$$
N_{\beta, i}(1 \otimes \ldots \otimes n \otimes \ldots \otimes 1)
=
1 \otimes \ldots \otimes n \otimes \ldots \otimes 1
$$
を満たす一意な $(A/R)_\bullet(\beta)$-半線形写像として定める。
以上から次の補題が示唆される。

\begin{lemma}
\label{descent-lemma-descent-datum-cosimplicial}
$R \to A$ を環準同型とする。
降下データ $(N, \varphi)$ に対し、規則 $N_n = N_{n, n}$ と、
$\beta : [n] \to [m]$ が与えられたときの定義
$$
N_\bullet(\beta) = (\varphi^m_{\beta(n)m}) \circ N_{\beta, n} :
N_{n, n} \longrightarrow N_{m, m}.
$$
により、余単体的 $(A/R)_\bullet$-加群
$N_\bullet$\footnote{正確には $(N, \varphi)_\bullet$ と書くべきである。}
を対応させることができる。この手順は降下データについて函手的である。
\end{lemma}

\begin{proof}
$\varphi(n \otimes 1) = \sum \alpha_i \otimes x_i$ と書くと、
最初のいくつかの写像は
$$
\begin{matrix}
\delta^1_0 & : & N & \to & A \otimes N & n & \mapsto & 1 \otimes n \\
\delta^1_1 & : & N & \to & A \otimes N & n & \mapsto &
\sum \alpha_i \otimes x_i\\
\sigma^0_0 & : & A \otimes N & \to & N & a_0 \otimes n & \mapsto & a_0n \\
\delta^2_0 & : & A \otimes N & \to & A \otimes A \otimes N &
a_0 \otimes n & \mapsto & 1 \otimes a_0 \otimes n \\
\delta^2_1 & : & A \otimes N & \to & A \otimes A \otimes N &
a_0 \otimes n & \mapsto & a_0 \otimes 1 \otimes n \\
\delta^2_2 & : & A \otimes N & \to & A \otimes A \otimes N &
a_0 \otimes n & \mapsto & \sum a_0 \otimes \alpha_i \otimes x_i \\
\sigma^1_0 & : & A \otimes A \otimes N & \to & A \otimes N &
a_0 \otimes a_1 \otimes n & \mapsto & a_0a_1 \otimes n \\
\sigma^1_1 & : & A \otimes A \otimes N & \to & A \otimes N &
a_0 \otimes a_1 \otimes n & \mapsto & a_0 \otimes a_1n
\end{matrix}
$$
となる。記法は Simplicial, Section
\ref{simplicial-section-cosimplicial-object} に従う。
まず二つの性質
$\sigma^0_0 \circ \delta^1_0 = \text{id}$ および
$\sigma^0_0 \circ \delta^1_1 = \text{id}$ を確かめる。
第一の等式 $\sigma^0_0 \circ \delta^1_0 = \text{id}$ は、
上の射の明示的記述から明らかである。
第二の関係を証明するにはコサイクル条件を使う必要がある
（任意の同型 $\varphi : N \otimes_R A \to A \otimes_R N$
に対しては成り立たないからである）。
$p = \sigma^0_0 \circ \delta^1_1 : N \to N$ と置く。
上の写像の記述から、$p$ は
$$
p = \varphi \otimes \text{id} :
N = (N \otimes_R A) \otimes_{(A \otimes_R A)} A
\longrightarrow
(A \otimes_R N) \otimes_{(A \otimes_R A)} A = N
$$
にも等しいことが分かる。$\varphi$ は同型なので、$p$ も同型である。
ある $\alpha_i \in A$ と $x_i \in N$ により
$\varphi(n \otimes 1) = \sum \alpha_i \otimes x_i$ と書く。
すると $p(n) = \sum \alpha_ix_i$ である。次に、ある
$\alpha_{ij} \in A$ と $y_j \in N$ により
$\varphi(x_i \otimes 1) = \sum \alpha_{ij} \otimes y_j$ と書く。
このときコサイクル条件は
$$
\sum \alpha_i \otimes \alpha_{ij} \otimes y_j
=
\sum \alpha_i \otimes 1 \otimes x_i.
$$
を意味する。したがって
$p(n) = \sum \alpha_ix_i = \sum \alpha_i\alpha_{ij}y_j =
\sum \alpha_i p(x_i) = p(p(n))$ である。ゆえに $p$ は射影子であり、
しかも同型なので恒等射である。

\medskip\noindent
$N_\bullet$ が余単体的加群であることを完全に証明するには、
Simplicial, Remark \ref{simplicial-remark-relations-cosimplicial} の
五種類の関係をすべて確かめる必要がある。
$\sigma$ 同士の合成に関する関係は明らかである。
$\delta$ 同士の合成に関する関係は $\varphi$ のコサイクル条件に帰着する。
上とまったく同じ方法で
$\sigma_j \circ \delta_j = \sigma_j \circ \delta_{j + 1} = \text{id}$
も確かめられる。最後に、$\delta$ と $\sigma$ の合成に関する
その他の関係は、どのような $\varphi$ に対しても成り立つ。
\end{proof}

\noindent
$R$-加群 $M$ には標準降下データ $(M \otimes_R A, can)$ を
対応させられることに注意せよ。ここで
$can : (M \otimes_R A) \otimes_R A \to A \otimes_R (M \otimes_R A)$
は明らかな写像
$(m \otimes a) \otimes a' \mapsto a \otimes (m \otimes a')$.
である。

\begin{lemma}
\label{descent-lemma-canonical-descent-datum-cosimplicial}
$R \to A$ を環準同型とし、$M$ を $R$-加群とする。
標準降下データに付随する余単体的 $(A/R)_\bullet$-加群は、
余単体的加群 $(A/R)_\bullet \otimes_R M$ と同型である。
\end{lemma}

\begin{proof}
省略する。
\end{proof}

\begin{definition}
\label{descent-definition-descent-datum-effective-module}
$R \to A$ を環準同型とする。$R$-加群 $M$ と、降下データ
$(M \otimes_R A, can)$ から $(N, \varphi)$ への同型が存在するとき、
降下データ $(N, \varphi)$ は{\it 有効}であるという。
\end{definition}

\noindent
$R \to A$ を環準同型とし、$(N, \varphi)$ を降下データとする。
$N_\bullet$ に付随するコチェイン複体 $s(N_\bullet)$ を取ることができる
（Simplicial, Section \ref{simplicial-section-dold-kan-cosimplicial} 参照）。
これは次の形をもつ：
$$
N \to A \otimes_R N \to A \otimes_R A \otimes_R N \to \ldots
$$
各写像を記述できる。第一の写像は
$$
n \longmapsto 1 \otimes n - \varphi(n \otimes 1).
$$
であり、第二の写像は純テンソル上で
$$
a \otimes n \longmapsto 1 \otimes a \otimes n
- a \otimes 1 \otimes n + a \otimes \varphi(n \otimes 1).
$$
と定まる。このパターンがどのように続くかは明らかである。

\medskip\noindent
特別な場合 $N = A \otimes_R M$ には、任意の $m \in M$ に対し
元 $1 \otimes m$ が、余単体的加群
$(A/R)_\bullet \otimes_R M$ に付随するコチェイン複体の
第一の写像の核に属する。したがって拡大コチェイン複体
\begin{equation}
\label{descent-equation-extended-complex}
0 \to M \to A \otimes_R M \to A \otimes_R A \otimes_R M \to \ldots
\end{equation}
を得る。ここで $0$ は次数 $-2$、加群 $M$ は次数 $-1$、
加群 $A \otimes_R M$ は次数 $0$ にあるものと考え、以下同様とする。
$M = R$ のとき、この複体は
$$
0 \to R \to A \to A \otimes_R A \to A \otimes_R A \otimes_R A \to \ldots
$$
という形をもつことに注意せよ。

\begin{lemma}
\label{descent-lemma-with-section-exact}
$R \to A$ が切断をもつと仮定する。このとき任意の $R$-加群 $M$ に対し、
拡大コチェイン複体 (\ref{descent-equation-extended-complex}) は完全である。
\end{lemma}

\begin{proof}
Simplicial, Lemma
\ref{simplicial-lemma-push-outs-simplicial-object-w-section} により、
写像 $R \to (A/R)_\bullet$ は余単体的 $R$-代数のホモトピー同値である
（ここで $R$ は定値余単体的 $R$-代数を表す）。
また $\otimes_R M$ は $R$-代数の圏から $R$-加群の圏への函手なので、
$M \to (A/R)_\bullet \otimes_R M$ は余単体的 $R$-加群の圏における
ホモトピー同値である。Simplicial, Lemma
\ref{simplicial-lemma-functorial-homotopy} を見よ。
したがって付随する複体の間に誘導される写像もホモトピー同値である。
Simplicial, Lemma \ref{simplicial-lemma-homotopy-s-Q} を見よ。
定値余単体的 $R$-加群 $M$ に付随する複体は
$$
\xymatrix{
M \ar[r]^0 & M \ar[r]^1 & M \ar[r]^0 & M \ar[r]^1 & M \ldots
}
$$
であるから結論を得る
（拡大版は先頭に $M$ を一つ余分に置くだけである）。
\end{proof}

\begin{lemma}
\label{descent-lemma-ff-exact}
$R \to A$ が忠実平坦であると仮定する
（Algebra, Definition \ref{algebra-definition-flat} 参照）。
このとき任意の $R$-加群 $M$ に対し、拡大コチェイン複体
(\ref{descent-equation-extended-complex}) は完全である。
\end{lemma}

\begin{proof}
忠実平坦な環準同型 $R \to R'$ が存在し、環準同型
$R' \to A' = R' \otimes_R A$ に対して結論が成り立つことを示せたとする。
すると $R \to A$ に対しても結論が従う。実際、任意の $R$-加群 $M$ に対し、
余単体的加群 $(M \otimes_R R') \otimes_{R'} (A'/R')_\bullet$ は
余単体的加群 $R' \otimes_R (M \otimes_R (A/R)_\bullet)$ にほかならない。
したがって $(M \otimes_R R') \otimes_{R'} (A'/R')_\bullet$ に付随する
複体のコホモロジーの消滅から、$R \to R'$ の忠実平坦性により
$M \otimes_R (A/R)_\bullet$ に付随する複体のコホモロジーの消滅が従う。
拡大複体の次数 $-1$ および $0$ のコホモロジー群の消滅についても
同様である（証明は省略する）。

\medskip\noindent
そのような忠実平坦拡大は実際に存在する。すなわち $R' = A$ と取ればよい。
環準同型 $R' = A \to A' = A \otimes_R A$ は切断
$a \otimes a' \mapsto aa'$ をもち、
Lemma \ref{descent-lemma-with-section-exact} を適用できるからである。
\end{proof}

\noindent
この複体と有効性の問題との関係は次のとおりである。

\begin{lemma}
\label{descent-lemma-recognize-effective}
$R \to A$ を忠実平坦な環準同型とし、
$(N, \varphi)$ を降下データとする。このとき $(N, \varphi)$ が
有効であることと、標準写像
$$
A \otimes_R H^0(s(N_\bullet)) \longrightarrow N
$$
が同型であることは同値である。
\end{lemma}

\begin{proof}
$(N, \varphi)$ が有効ならば、$\varphi = can$ として
$N = A \otimes_R M$ と書ける。Lemmas
\ref{descent-lemma-canonical-descent-datum-cosimplicial} および
\ref{descent-lemma-ff-exact} により $H^0(s(N_\bullet)) = M$ が従う。
逆に、補題の写像が同型であると仮定し、
$M = H^0(s(N_\bullet))$ と置く。これは $N$ の $R$-部分加群であり、
$M = \{n \in N \mid 1 \otimes n = \varphi(n \otimes 1)\}$ である。
確かめるべきことは、同型 $A \otimes_R M \to N$ を介して
標準降下データが $\varphi$ と一致することだけである。
この確認は省略する。
\end{proof}

\begin{lemma}
\label{descent-lemma-descent-descends}
$R \to A$ を忠実平坦な環準同型とし、$R \to R'$ も忠実平坦とする。
$A' = R' \otimes_R A$ と置く。$R' \to A'$ に関するすべての
降下データが有効ならば、$R \to A$ に関するすべての降下データも有効である。
\end{lemma}

\begin{proof}
$(N, \varphi)$ を $R \to A$ に関する降下データとする。
$N' = R' \otimes_R N = A' \otimes_A N$ と置き、降下データ
$\varphi$ の基底変換を
$\varphi' = \text{id}_{R'} \otimes \varphi$ と書く。
すると $(N', \varphi')$ は $R' \to A'$ に関する降下データであり、
$H^0(s(N'_\bullet)) = R' \otimes_R H^0(s(N_\bullet))$.
さらに、写像
$A' \otimes_{R'} H^0(s(N'_\bullet)) \to N'$ は、忠実平坦な写像
$A \to A'$ による $A$-加群準同型
$A \otimes_R H^0(s(N)) \to N$ の基底変換と同一視される。
したがって Lemma \ref{descent-lemma-recognize-effective} から結論を得る。
\end{proof}

\noindent
次が本節の主結果である。
証明はいささか不器用に見えるかもしれない。より高い観点からの方法については、
下の Remark \ref{descent-remark-homotopy-equivalent-cosimplicial-algebras} を見よ。

\begin{proposition}
\label{descent-proposition-descent-module}
\begin{slogan}
忠実平坦な環準同型に沿う加群の有効降下。
\end{slogan}
$R \to A$ を忠実平坦な環準同型とする。このとき、
\begin{enumerate}
\item $R \to A$ に関する加群の任意の降下データは有効である。
\item $R$-加群から降下データの圏への函手
$M \mapsto (A \otimes_R M, can)$ は圏同値である。
\item 逆函手は
$(N, \varphi) \mapsto H^0(s(N_\bullet))$ で与えられる。
\end{enumerate}
\end{proposition}

\begin{proof}
(1) だけを証明し、(2) と (3) の証明は省略する。
$R \to A$ は忠実平坦なので、忠実平坦な基底変換 $R \to R'$ であって
$R' \to A' = R' \otimes_R A$ が切断をもつものが存在する
（Lemma \ref{descent-lemma-ff-exact} の証明と同様に $R' = A$ と取ればよい）。
したがって Lemma \ref{descent-lemma-descent-descends} により、
$R \to A$ が切断 $\sigma : A \to R$ をもつと仮定してよい。
$(N, \varphi)$ を $R \to A$ に関する降下データとし、
$$
M = H^0(s(N_\bullet)) = \{n \in N \mid 1 \otimes n = \varphi(n \otimes 1)\}
\subset
N
$$
と置く。Lemma \ref{descent-lemma-recognize-effective} により、
$A \otimes_R M \to N$ が同型であることを示せば十分である。

\medskip\noindent
元 $n \in N$ を取る。ある $a_i \in A$ と $x_i \in N$ により
$\varphi(n \otimes 1) = \sum a_i \otimes x_i$ と書く。
Lemma \ref{descent-lemma-descent-datum-cosimplicial} により、
$N$ において $n = \sum a_i x_i$ である
（任意の余単体的対象で
$\sigma^0_0 \circ \delta^1_1 = \text{id}$ だからである）。
次に、ある $a_{ij} \in A$ と $y_j \in N$ により
$\varphi(x_i \otimes 1) = \sum a_{ij} \otimes y_j$ と書く。
コサイクル条件は
$$
\sum a_i \otimes a_{ij} \otimes y_j = \sum a_i \otimes 1 \otimes x_i
$$
が $A \otimes_R A \otimes_R N$ において成り立つことを意味する。
ここから二つのことが従う：
\begin{enumerate}
\item 第一の $A$ に $\sigma$ を適用すると、
$\sum \sigma(a_i) \varphi(x_i \otimes 1) = \sum \sigma(a_i) \otimes x_i$,
\item 中央の $A$ に $\sigma$ を適用すると、
$\sum_i a_i \otimes \sum_j \sigma(a_{ij}) y_j = \sum a_i \otimes x_i$.
\end{enumerate}
(1) より $\sum \sigma(a_i) x_i \in M$ である。これを $x_i$ に
適用すると、すべての $i$ に対し
$\sum \sigma(a_{ij})y_i \in M$ を得る。
(2) の等式を掛け合わせると、
$\sum_i a_i (\sum_j \sigma(a_{ij}) y_j) = \sum a_i x_i = n$ を得る。
したがって $A \otimes_R M \to N$ は全射である。
最後に $m_i \in M$ かつ $\sum a_i m_i = 0$ と仮定する。
$\sum a_im_i \otimes 1$ に $\varphi$ を適用すると
$\sum a_i \otimes m_i = 0$ が分かる。
すなわち $A \otimes_R M \to N$ は単射でもあり、結論を得る。
\end{proof}

\begin{remark}
\label{descent-remark-standard-covering}
$R$ を環とし、$f_1, \ldots, f_n\in R$ が単位イデアルを生成するとする。
環 $A = \prod_i R_{f_i}$ は忠実平坦な $R$-代数である。
余単体的環 $(A/R)_\bullet$ の次数 $n$ の環は
$$
\prod\nolimits_{i_0, \ldots, i_n} R_{f_{i_0}\ldots f_{i_n}}
$$
であることに注意する。したがって上の結果から
Algebra, Lemmas \ref{algebra-lemma-standard-covering},
\ref{algebra-lemma-cover-module}, \ref{algebra-lemma-glue-modules}
が再び得られる。しかし高次でも完全であるため、上の結果は実際には
さらに多くを述べている。すなわち、アフィンスキーム上の準連接層の
{\v C}ech コホモロジーが自明であることを意味する。
これにより Cohomology of Schemes, Lemma
\ref{coherent-lemma-cech-cohomology-quasi-coherent-trivial}
の第二の証明を得る。
\end{remark}

\begin{remark}
\label{descent-remark-homotopy-equivalent-cosimplicial-algebras}
$R$ を環とし、$A_\bullet$ を余単体的 $R$-代数とする。
この設定では、降下データは次の性質をもつ余単体的
$A_\bullet$-加群 $M_\bullet$ に対応する。任意の
$n, m \geq 0$ と任意の $\varphi : [n] \to [m]$ に対し、
写像 $M(\varphi) : M_n \to M_m$ は同型
$$
M_n \otimes_{A_n, A(\varphi)} A_m \longrightarrow M_m.
$$
を誘導する。このような余単体的加群を{\it Cartesian 加群}と呼ぶ。
この設定では、Proposition \ref{descent-proposition-descent-module} の証明を
次の段階に分けられる：
\begin{enumerate}
\item $R \to R'$ と $R \to A$ が忠実平坦であるとする。
$(R' \otimes_R A)/R'$ に関する降下データが有効ならば、
$A/R$ に関する降下データも有効である。
\item $A$ を $R$-代数とする。$A/R$ に関する降下データは、
Cartesian $(A/R)_\bullet$-加群に対応する。
\item $R \to A$ が切断をもつならば、$(A/R)_\bullet$ は
$R$ を値とする定値余単体的 $R$-代数 $R$ とホモトピー同値である。
\item $A_\bullet \to B_\bullet$ が余単体的 $R$-代数の
ホモトピー同値ならば、函手
$M_\bullet \mapsto M_\bullet \otimes_{A_\bullet} B_\bullet$
は Cartesian $A_\bullet$-加群の圏と Cartesian
$B_\bullet$-加群の圏の間の圏同値を誘導する。
\end{enumerate}
(1) については Lemma \ref{descent-lemma-descent-descends} を見よ。
(2) では Lemma \ref{descent-lemma-descent-datum-cosimplicial} を用いる。
(3) は Lemma \ref{descent-lemma-with-section-exact} の証明ですでに見た
（Simplicial, Lemma
\ref{simplicial-lemma-push-outs-simplicial-object-w-section} に依拠する）。
さらに、(4) は正しく考えれば直ちに分かる。
\end{remark}








\section{普遍単射な射に関する降下}
\label{descent-section-descent-universally-injective}

\noindent
代数幾何における多くの構成は{\it 降下}の技法を用いて行われる。
たとえば、与えられた空間へ射影する少し大きな空間上でまず作業して
その空間上の対象を構成したり、被覆写像に沿って引き戻すことで
空間や射の性質を確認したりする。このような技法の有用性はもちろん、
広いクラスの{\it 有効降下射}を特定できるかどうかに依存する。
現代代数幾何が Grothendieck 流に発展した初期には、
{\it 準コンパクト}かつ{\it 忠実平坦}な射のクラスが、
対象、射、およびそれらの多くの性質を降下させるうえで有効であることが示された。

\medskip\noindent
通常どおり、この主張は環と加群の性質に帰着する。
Grothendieck は、準同型 $f: R \to S$ が加群の有効降下射となるためには、
$f$ が忠実平坦であれば十分であることを示した。
しかしこれは多くの自然な例を除外してしまう。たとえば任意の
分裂環準同型は有効降下射である。この自然な例は忠実平坦降下の証明にも現れる。
任意の環準同型 $f: R \to S$ に対し、
$1_S \otimes f: S \to S \otimes_R S$ は、平坦性の有無にかかわらず
乗法写像によって分裂する。

\medskip\noindent
そこで、忠実平坦な射の場合と分裂環準同型の場合の双方を含み、
加群の有効降下を導く自然な環論的条件があるかを問うことができる。
この問いに対する完全な答えが 1970 年頃から知られていたと知れば、
読者は驚くかもしれない（少なくとも本節の著者は驚いた）。
すなわち、$f: R \to S$ が加群の有効降下射となるための必要条件は、
$f$ が $R$-加群の圏で{\it 普遍単射}となること、つまり任意の
$R$-加群 $M$ に対して写像
$1_M  \otimes f: M \to M \otimes_R S$ が単射となることであると
容易に確認できる。そしてこれは十分条件でもあることが分かる。
たとえば $f$ が $R$-加群の圏で分裂するならば
（環の圏で分裂する必要はない）、$f$ は加群の有効降下射である。
  
\medskip\noindent
この結果の歴史はいささか込み入っている。もともとは Olivier
\cite{olivier} が主張し、普遍単射な射を{\it 純粋}と呼んだが、
明確な証明は示さなかった。Joyal と Tierney の仕事
\cite{joyal-tierney} から結果を取り出すこともできるが、
我々の知るかぎり、文献に現れた最初の独立した証明は
Mesablishvili \cite{mesablishvili1} によるものである。
本節の第一の目的は Mesablishvili の証明を詳しく示すことである。
誤植の修正を除けば彼の元の提示からほとんど変更を要しないが、
一つだけ例外があり、代数幾何で通常用いられる降下データの定義と
\cite{mesablishvili1} で用いられる定義との関係を明示する。
この証明は完全に圏論的であり、したがってモナドの言葉で述べられる
（ゆえに他の状況にも適用できる）。\cite{janelidze-tholen} を見よ。

\medskip\noindent
本節の第二の目的は、加群、代数、および射のどの性質が
普遍単射な環準同型に沿って降下できるかについて情報を集めることである。
有限加群と平坦加群の場合は Mesablishvili
\cite{mesablishvili2} により扱われた。


\subsection{圏論的準備}
\label{descent-subsection-category-prelims}

\noindent
記述を簡潔にするため、まず圏論のいくつかの基本概念を復習する。
本小節では周囲の圏を一つ固定する。

\medskip\noindent
二つの射 $g_1, g_2: B \to C$ に対し、$g_1$ と $g_2$ の
{\it 等化子}とは、$g_1 \circ f = g_2 \circ f$ を満たし、
この性質について普遍的な射 $f: A \to B$ のことであった。
後半の主張は、破線の矢印を除いた任意の可換図式
$$
\xymatrix{A' \ar[rd]^e \ar@/^1.5pc/[rrd] \ar@{-->}[d] & & \\
A \ar[r]^f & B \ar@<1ex>[r]^{g_1} \ar@<-1ex>[r]_{g_2} & 
C
}
$$
が一意に補完できることを意味する。この状況では図式
\begin{equation}
\label{descent-equation-equalizer}
\xymatrix{
A \ar[r]^f  & B \ar@<1ex>[r]^{g_1} \ar@<-1ex>[r]_{g_2} & C
}
\end{equation}
も等化子であるという。矢印を逆にすると{\it 余等化子}の定義が得られる。
Categories, Sections \ref{categories-section-equalizers} および
\ref{categories-section-coequalizers} を見よ。

\medskip\noindent
等化子であるという性質は普遍性を含むので、通常は共変函手を適用しても
保たれるとは限らない。単射や全射の場合と同様、分裂を明示することで
この問題を回避できる場合がある。

\begin{definition}
\label{descent-definition-split-equalizer}
{\it 分裂等化子}とは、$g_1 \circ f = g_2 \circ f$ を満たす図式
(\ref{descent-equation-equalizer}) であって、補助的な射
$h : B \to A$ と $i : C \to B$ が存在し、
\begin{equation}
\label{descent-equation-split-equalizer-conditions}
h \circ f = 1_A, \quad f \circ h = i \circ g_1, \quad i \circ g_2 = 1_B.
\end{equation}
を満たすものをいう。
\end{definition}

\noindent
要点は、矢印の間の等式によって
(\ref{descent-equation-equalizer}) が等化子となることである。実際、
$e = f \circ (h \circ e)$ と書けば、写像 $e$ は $f$ を介して
一意に分解する。したがって分裂等化子に共変函手を適用すると
分裂等化子が得られ、反変函手を適用すると{\it 分裂余等化子}が得られる。
後者の定義は明らかである。

\subsection{普遍単射な射}
\label{descent-subsection-universally-injective}

\noindent
$\textit{Rings}$ は単位元 $1$ をもつ可換環の圏を表すものとする。
$\textit{Rings}$ の対象 $R$ に対し、$R$-加群の圏を
$\text{Mod}_R$ と書く。

\begin{remark}
\label{descent-remark-reflects}
アーベル圏の間の函手 $F : \mathcal{A} \to \mathcal{B}$ が完全で、
非零対象を非零対象へ送るならば、この函手は単射と全射を反映する。
実際、完全性により $F$ は核と余核を保つ
（Homology, Section \ref{homology-section-functors} と比較せよ）。
たとえば $f : R \to S$ が忠実平坦な環準同型ならば、
$\bullet \otimes_R S: \text{Mod}_R \to \text{Mod}_S$ はこれらの性質をもつ。
\end{remark}

\noindent
$R$ を環とする。$\text{Mod}_R$ における射 $f : M \to N$ が
{\it 普遍単射}であるとは、すべての $P \in \text{Mod}_R$ に対して射
$f \otimes 1_P: M \otimes_R P \to N \otimes_R P$ が単射となることである。
Algebra, Definition \ref{algebra-definition-universally-injective} を見よ。

\begin{definition}
\label{descent-definition-universally-injective}
環準同型 $f: R \to S$ が{\it 普遍単射}であるとは、
$\text{Mod}_R$ の射として普遍単射であることをいう。
\end{definition}

\begin{example}
\label{descent-example-split-injection-universally-injective}
$\text{Mod}_R$ における任意の分裂単射は普遍単射である。
とくに $\textit{Rings}$ における任意の分裂単射は普遍単射である。
\end{example}

\begin{example}
\label{descent-example-cover-universally-injective}
$R$ を環とし、$f_1, \ldots, f_n \in R$ が単位イデアルを生成するとする。
射 $R \to R_{f_1} \oplus \ldots \oplus R_{f_n}$ は普遍単射である。
これは Lemma \ref{descent-lemma-faithfully-flat-universally-injective} から
直ちに従うが、直接確認することも有益である。すぐに $R$ が局所環の場合へ
帰着できる。この場合、ある $f_i$ は単元でなければならず、
したがって写像 $R \to R_{f_i}$ は同型である。
\end{example}

\begin{lemma}
\label{descent-lemma-faithfully-flat-universally-injective}
任意の忠実平坦な環準同型は普遍単射である。
\end{lemma}

\begin{proof}
これは Algebra, Lemma
\ref{algebra-lemma-faithfully-flat-universally-injective} の言い換えである。
\end{proof}

\noindent
\cite{mesablishvili1} の鍵となる観察は、$\text{Mod}_R$ における
単射余生成子の通常の構成を用いると、普遍単射性を分裂によって
有用な形に言い換えられることである。

\begin{definition}
\label{descent-definition-C}
$R$ を環とする。反変函手
{\it $C$} $ : \text{Mod}_R \to \text{Mod}_R$ を
$$
C(M) = \Hom_{\textit{Ab}}(M, \mathbf{Q}/\mathbf{Z}),
$$
と定める。$C(M)$ 上の $R$-作用は $rf(s) = f(rs)$ で与える。
\end{definition}

\noindent
この函手は More on Algebra, Section
\ref{more-algebra-section-injectives-modules} では
$M \mapsto M^\vee$ と書かれていた。

\begin{lemma}
\label{descent-lemma-C-is-faithful}
環 $R$ に対し、函手 $C : \text{Mod}_R \to \text{Mod}_R$ は完全であり、
単射と全射を反映する。
\end{lemma}

\begin{proof}
完全性は More on Algebra, Lemma \ref{more-algebra-lemma-vee-exact}
である。他の性質はこれから従う。Remark \ref{descent-remark-reflects} を見よ。
\end{proof}

\begin{remark}
\label{descent-remark-adjunction}
$\Hom$ とテンソル積の間の標準随伴を、反変函手の自然同型
\begin{equation}
\label{descent-equation-adjunction}
C(\bullet_1 \otimes_R \bullet_2) \cong \Hom_R(\bullet_1, C(\bullet_2)): 
\text{Mod}_R \times \text{Mod}_R \to \text{Mod}_R
\end{equation}
の形で頻繁に用いる。この同型は
$f: M_1 \otimes_R M_2 \to \mathbf{Q}/\mathbf{Z}$ を写像
$m_1 \mapsto (m_2 \mapsto f(m_1 \otimes m_2))$ へ送る。
Algebra, Lemma \ref{algebra-lemma-hom-from-tensor-product-variant} を見よ。
この観察の系として、
$$
\xymatrix@C=9pc{
C(M) \ar@<1ex>[r] \ar@<-1ex>[r] & C(N) \ar[r] & C(P)
}
$$
が $\text{Mod}_R$ における分裂余等化子図式ならば、任意の
$Q \in \text{Mod}_R$ に対して
$$
\xymatrix@C=9pc{
C(M \otimes_R Q) \ar@<1ex>[r] \ar@<-1ex>[r] & C(N \otimes_R Q) \ar[r] & C(P 
\otimes_R Q)
}
$$
も分裂余等化子図式となる。
\end{remark}

\begin{lemma}
\label{descent-lemma-split-surjection}
$R$ を環とする。$\text{Mod}_R$ における射 $f: M \to N$ が
普遍単射であることと、
$C(f): C(N) \to C(M)$ が分裂全射であることは同値である。
\end{lemma}

\begin{proof}
(\ref{descent-equation-adjunction}) により、任意の
$P \in \text{Mod}_R$ に対して可換図式
$$
\xymatrix@C=9pc{
\Hom_R( P, C(N)) \ar[r]_{\Hom_R(P,C(f))} \ar[d]^{\cong} &
\Hom_R(P,C(M)) \ar[d]^{\cong} \\
C(P \otimes_R N ) \ar[r]^{C(1_{P} \otimes f)} & C(P \otimes_R M ).
}
$$
を得る。$f$ が普遍単射ならば、
$1_{C(M)} \otimes f: C(M) \otimes_R M \to C(M) \otimes_R N$ は単射なので、
$P = C(M)$ とした上の図式の両行は全射である。したがって
$1_{C(M)} \in \Hom_R(C(M), C(M))$ を、
$C(f)$ を分裂させるある
$g \in \Hom_R(C(N), C(M))$ へ持ち上げられる。
逆に $C(f)$ が分裂全射ならば、上の図式の両行は全射である。
したがって Lemma \ref{descent-lemma-C-is-faithful} により
$1_{P} \otimes f$ は単射である。
\end{proof}

\begin{remark}
\label{descent-remark-functorial-splitting}
$f: M \to N$ を $\text{Mod}_R$ における普遍単射な射とする。
$C(f)$ の分裂 $g$ を選ぶと、各 $P \in \text{Mod}_R$ に対して
$C(1_P \otimes f)$ の分裂を函手的に構成できる。
実際、(\ref{descent-equation-adjunction}) により、これは
$\Hom_R(P, C(f))$ を $P$ について函手的に分裂させることと同値であり、
写像 $g \circ \bullet$ によって実現される。
\end{remark}


\subsection{加群とその射の降下}
\label{descent-subsection-descent-modules-morphisms}

\noindent
本小節を通じて環準同型 $f: R \to S$ を固定する。
Section \ref{descent-section-descent-modules} で見たように、
余単体的代数の言葉を使って加群の降下データを論じることもできるが、
本小節ではより具体的な用語を用いる。

\medskip\noindent
$i = 1, 2, 3$ に対し、$S_i$ を $R$ 上の $S$ の $i$ 重テンソル積とする。
環準同型 $\delta_0^1, \delta_1^1: S_1 \to S_2$,
$\delta_{01}^1, \delta_{02}^1, \delta_{12}^1: S_1 \to S_3$ および
$\delta_0^2, \delta_1^2, \delta_2^2: S_2 \to S_3$ を次の式で定める：
\begin{align*}
\delta^1_0  (a_0) & =  1 \otimes a_0 \\
\delta^1_1  (a_0) & = a_0 \otimes 1 \\
\delta^2_0  (a_0 \otimes a_1) & =  1 \otimes a_0 \otimes a_1 \\
\delta^2_1  (a_0 \otimes a_1) & =  a_0 \otimes 1 \otimes a_1 \\
\delta^2_2  (a_0 \otimes a_1) & =  a_0 \otimes a_1 \otimes 1 \\
\delta_{01}^1(a_0) & = 1 \otimes 1 \otimes a_0 \\
\delta_{02}^1(a_0) & = 1 \otimes a_0 \otimes 1 \\
\delta_{12}^1(a_0) & = a_0 \otimes 1 \otimes 1.
\end{align*}
言い換えると、上付き添字は始域の環を表し、下付き添字は因子 1 を
挿入する位置を表す
（この記法は Section \ref{descent-section-descent-modules} で導入した記法と整合する）。

\medskip\noindent
Definition \ref{descent-definition-descent-datum-modules} を思い出すと、
$M \in \text{Mod}_S$ に対する $f$ に関する $M$ 上の
{\it 降下データ}とは、$S_2$-加群の同型
$$
\theta :
M \otimes_{S,\delta^1_0} S_2
\longrightarrow
M \otimes_{S,\delta^1_1} S_2
$$
であって、{\it コサイクル条件}
\begin{equation}
\label{descent-equation-cocycle-condition}
(\theta \otimes \delta_2^2) \circ (\theta \otimes \delta_2^0) = (\theta \otimes 
\delta_2^1):
M \otimes_{S, \delta^1_{01}} S_3 \to M \otimes_{S,\delta^1_{12}} S_3.
\end{equation}
を満たすものをいう\footnote{正確には、ここでの $\theta$ は
Definition \ref{descent-definition-descent-datum-modules} の $\varphi$ の逆である。}。
$f$ に関する降下データを備えた $S$-加群の圏を $DD_{S/R}$ と書く。

\medskip\noindent
たとえば $M_0 \in \text{Mod}_R$ と同型
$M \cong M_0 \otimes_R S$ を選ぶと、
$M \otimes_{S,\delta^1_0} S_2$ と
$M \otimes_{S,\delta^1_1} S_2$ をともに
$M_0 \otimes_R S_2$ と自然に同一視することにより降下データが得られる。
とくにこの構成は函手 $f^*: \text{Mod}_R \to DD_{S/R}$ を定める。

\begin{definition}
\label{descent-definition-effective-descent}
函手 $f^*: \text{Mod}_R \to DD_{S/R}$ を
{\it $f$ に沿う基底拡大}と呼ぶ。$f^*$ が充満忠実であるとき、
$f$ を{\it 加群の降下射}という。$f^*$ が圏同値であるとき、
$f$ を{\it 加群の有効降下射}という。
\end{definition}

\noindent
目標は、$f$ が普遍単射であるとき、$\theta$ を用いて
$M$ の中に $M_0$ を見いだせることを示すことである。
この過程ではいくつかの等化子図式が本質的な役割を果たす。

\begin{lemma}
\label{descent-lemma-equalizer-M}
$(M,\theta) \in DD_{S/R}$ に対し、図式
\begin{equation}
\label{descent-equation-equalizer-M}
\xymatrix@C=8pc{
M \ar[r]^{\theta \circ (1_M \otimes \delta_0^1)} &
M \otimes_{S, \delta_1^1} S_2
\ar@<1ex>[r]^{(\theta \otimes \delta_2^2) \circ (1_M \otimes \delta^2_0)}
\ar@<-1ex>[r]_{1_{M \otimes S_2} \otimes \delta^2_1} & 
M \otimes_{S, \delta_{12}^1} S_3
}
\end{equation}
は分裂等化子である。
\end{lemma}

\begin{proof}
環準同型 $\sigma^0_0: S_2 \to S_1$ および
$\sigma_0^1, \sigma_1^1: S_3 \to S_2$ を次の式で定める：
\begin{align*}
\sigma^0_0 (a_0 \otimes a_1) & = a_0a_1 \\
\sigma^1_0 (a_0 \otimes a_1 \otimes a_2) & = a_0a_1 \otimes a_2 \\
\sigma^1_1 (a_0 \otimes a_1 \otimes a_2) & = a_0 \otimes a_1a_2.
\end{align*}
補助的な射として
$1_M \otimes \sigma_0^0: M \otimes_{S, \delta_1^1} S_2 \to M$
および
$1_M \otimes \sigma_0^1: M \otimes_{S,\delta_{12}^1} S_3 \to
M \otimes_{S, \delta_1^1} S_2$ を取る。
(\ref{descent-equation-split-equalizer-conditions}) で必要とされる整合性のうち、
第一のものはコサイクル条件
(\ref{descent-equation-cocycle-condition}) を $\sigma_1^1$ とテンソルすることで
従い、他は直ちに従う。
\end{proof}

\begin{lemma}
\label{descent-lemma-equalizer-CM}
$(M, \theta) \in DD_{S/R}$ に対し、
(\ref{descent-equation-equalizer-M}) に $C$ を適用して得られる図式
\begin{equation}
\label{descent-equation-coequalizer-CM}
\xymatrix@C=8pc{
C(M \otimes_{S, \delta_{12}^1} S_3)
\ar@<1ex>[r]^{C((\theta \otimes \delta_2^2) \circ (1_M \otimes \delta^2_0))}
\ar@<-1ex>[r]_{C(1_{M \otimes S_2} \otimes \delta^2_1)} &
C(M \otimes_{S, \delta_1^1} S_2 )
\ar[r]^{C(\theta \circ (1_M \otimes \delta_0^1))} & C(M).
}
\end{equation}
は分裂余等化子である。
\end{lemma}

\begin{proof}
省略する。
\end{proof}

\begin{lemma}
\label{descent-lemma-equalizer-S}
図式
\begin{equation}
\label{descent-equation-equalizer-S}
\xymatrix@C=8pc{
S_1 \ar[r]^{\delta^1_1} &
S_2 \ar@<1ex>[r]^{\delta^2_2} \ar@<-1ex>[r]_{\delta^2_1} & 
S_3
}
\end{equation}
は分裂等化子である。
\end{lemma}

\begin{proof}
Lemma \ref{descent-lemma-equalizer-M} で
$(M, \theta) = f^*(S)$ と取ればよい。
\end{proof}

\noindent
これは $f^*$ の逆函手となる候補の定義を示唆する。

\begin{definition}
\label{descent-definition-pushforward}
函手 {\it $f_*$} $: DD_{S/R} \to \text{Mod}_R$ を次のように定める。
$f_*(M, \theta)$ は、図式
\begin{equation}
\label{descent-equation-equalizer-f}
\xymatrix@C=8pc{f_*(M,\theta) \ar[r] & M \ar@<1ex>^{\theta \circ (1_M \otimes 
\delta_0^1)}[r] \ar@<-1ex>_{1_M \otimes \delta_1^1}[r] & 
M \otimes_{S, \delta_1^1} S_2 
}
\end{equation}
を等化子にする $M$ の $R$-部分加群とする。
\end{definition}

\noindent
Lemma \ref{descent-lemma-equalizer-M} と、制限函手
$\text{Mod}_S \to \text{Mod}_R$ が基底拡大函手
$\bullet \otimes_R S: \text{Mod}_R \to \text{Mod}_S$ の右随伴であることから、
$f_*$ が $f^*$ の右随伴であることが従う。

\medskip\noindent
これで鍵となる補題を述べる準備ができた。忠実平坦の場合には自明である
（Remark \ref{descent-remark-descent-lemma} 参照）が、
一般の場合には若干の議論が必要である。

\begin{lemma}
\label{descent-lemma-descent-lemma}
$f$ が普遍単射ならば、(\ref{descent-equation-equalizer-f}) を
$R$ 上で $S$ とテンソルして得られる図式
\begin{equation}
\label{descent-equation-equalizer-f2}
\xymatrix@C=8pc{
f_*(M, \theta) \otimes_R S
\ar[r]^{\theta \circ (1_M \otimes \delta_0^1)} &
M \otimes_{S, \delta_1^1} S_2 
\ar@<1ex>[r]^{(\theta \otimes \delta_2^2) \circ (1_M \otimes \delta^2_0)}
\ar@<-1ex>[r]_{1_{M \otimes S_2} \otimes \delta^2_1} &
M \otimes_{S, \delta_{12}^1} S_3
}
\end{equation}
は等化子である。
\end{lemma}

\begin{proof}
Lemma \ref{descent-lemma-split-surjection} と
Remark \ref{descent-remark-functorial-splitting} により、写像
$C(1_N \otimes f): C(N \otimes_R S) \to C(N)$ は $N$ について
函手的に分裂できる。これにより、次の可換図式の上側の縦矢印が得られる：
$$
\xymatrix@C=8pc{
C(M \otimes_{S, \delta_1^1} S_2)
\ar@<1ex>^{C(\theta \circ (1_M \otimes \delta_0^1))}[r]
\ar@<-1ex>_{C(1_M \otimes \delta_1^1)}[r] \ar[d] &
C(M) \ar[r]\ar[d] & C(f_*(M,\theta)) \ar@{-->}[d] \\
C(M \otimes_{S,\delta_{12}^1} S_3)
\ar@<1ex>^{C((\theta \otimes \delta_2^2) \circ (1_M \otimes \delta^2_0))}[r]
\ar@<-1ex>_{C(1_{M \otimes S_2} \otimes \delta^2_1)}[r] \ar[d] &
C(M \otimes_{S, \delta_1^1} S_2 )
\ar[r]^{C(\theta \circ (1_M \otimes \delta_0^1))}
\ar[d]^{C(1_M \otimes \delta_1^1)} &
C(M) \ar[d] \ar@{=}[dl] \\
C(M \otimes_{S, \delta_1^1} S_2)
\ar@<1ex>[r]^{C(\theta \circ (1_M \otimes \delta_0^1))}
\ar@<-1ex>[r]_{C(1_M \otimes \delta_1^1)} &
C(M) \ar[r] &
C(f_*(M,\theta))
}
$$
ここで各列に沿う合成は恒等射である。
第二行は余等化子図式 (\ref{descent-equation-coequalizer-CM}) であり、
これによって破線の矢印が得られる。右上の正方形から補助的な射
$C(f_*(M,\theta)) \to C(M)$ および
$C(M) \to C(M\otimes_{S,\delta_1^1} S_2)$ を得る。
これらにより第一行は分裂余等化子図式となる。
Remark \ref{descent-remark-adjunction} により、$C$ の内部で $S$ と
テンソルして、分裂余等化子図式
$$
\xymatrix@C=8pc{
C(M \otimes_{S,\delta_2^2 \circ \delta_1^1} S_3)
\ar@<1ex>^{C((\theta \otimes \delta_2^2) \circ (1_M \otimes \delta^2_0))}[r] 
\ar@<-1ex>_{C(1_{M \otimes S_2} \otimes \delta^2_1)}[r] &
C(M \otimes_{S, \delta_1^1} S_2 )
\ar[r]^{C(\theta \circ (1_M \otimes \delta_0^1))} &
C(f_*(M,\theta) \otimes_R S).
}
$$
を得る。Lemma \ref{descent-lemma-C-is-faithful} により、
(\ref{descent-equation-equalizer-f2}) も等化子でなければならない。
\end{proof}

\begin{remark}
\label{descent-remark-descent-lemma}
$f$ が $\text{Mod}_R$ における分裂単射ならば、$C$ を使わずに
$f$ を直接分裂させることで議論を簡単にできる。
$f$ が忠実平坦ならばさらに簡単である。この場合、
$R$ 上で $S$ とテンソルする操作はすべての等化子を保つので、
Lemma \ref{descent-lemma-descent-lemma} の結論は直ちに従う。
\end{remark}

\begin{theorem}
\label{descent-theorem-descent}
次の条件は同値である。
\begin{enumerate}
\item[(a)] 射 $f$ は加群の降下射である。
\item[(b)] 射 $f$ は加群の有効降下射である。
\item[(c)] 射 $f$ は普遍単射である。
\end{enumerate}
\end{theorem}

\begin{proof}
(b) が (a) を含意することは明らかである。
次に (a) が (c) を含意することを確かめる。
$f$ が普遍単射でなければ、写像
$1_M \otimes f: M \to M \otimes_R S$ が非自明な核 $N$ をもつような
$M \in \text{Mod}_R$ を取れる。自然な射影 $M \to M/N$ は
同型ではないが、$DD_{S/R}$ におけるその像は同型である。
したがって $f^*$ は充満忠実ではない。

\medskip\noindent
最後に (c) が (b) を含意することを確かめる。
Lemmas \ref{descent-lemma-equalizer-M} および
\ref{descent-lemma-descent-lemma} により、
$(M, \theta) \in DD_{S/R}$,
に対して自然な写像 $f^* f_*(M,\theta) \to M$ は $S$-加群の同型である。
$M_0 \in \text{Mod}_R$ に対し、
$M_0 \to f_*f^*M_0$ が同型であることを示したい。
随伴により、合成
$f^*M_0 \to f^*f_*f^*M_0 \to f^*M_0$ は恒等射であり、
一方、上の議論により第二の矢印は同型である。
したがって $f^*M_0 \to f^*f_*f^*M_0$ は同型である。
Algebra, Lemma
\ref{algebra-lemma-base-change-along-universally-injective-ring-map}
により $M_0 \to f_* f^* M_0$ は同型であると結論できる。
ゆえに $f_*$ と $f^*$ は互いに準逆な函手であり、主張が従う。
\end{proof}

\subsection{加群の性質の降下}
\label{descent-subsection-descent-properties-modules}

\noindent
この小節を通じて、普遍単射な環準同型 $f : R \to S$、対象
$M \in \text{Mod}_R$、および環準同型 $R \to A$ を固定する。ここでは、
$M$ または $A$ のどの性質が $f$ に沿う基底拡大の後で判定できるかを
調べる。まず \cite{mesablishvili2} の結果をいくつか示す。

\begin{lemma}
\label{descent-lemma-flat-to-injective}
$M \in \text{Mod}_R$ が平坦ならば、$C(M)$ は入射 $R$-加群である。
\end{lemma}

\begin{proof}
$0 \to N \to P \to Q \to 0$ を $\text{Mod}_R$ における完全列とする。
$M$ は平坦なので、
$$
0 \to N \otimes_R M \to P \otimes_R M \to Q \otimes_R M \to 0
$$
は完全である。補題 \ref{descent-lemma-C-is-faithful} により、
$$
0 \to C(Q \otimes_R M) \to C(P \otimes_R M) \to C(N \otimes_R M) \to 0
$$
は完全である。(\ref{descent-equation-adjunction}) により、この最後の列は
$$
0 \to \Hom_R(Q, C(M)) \to \Hom_R(P, C(M)) \to \Hom_R(N, C(M)) \to 0.
$$
と書き直せる。したがって $C(M)$ は $\text{Mod}_R$ の入射対象である。
\end{proof}

\begin{theorem}
\label{descent-theorem-descend-module-properties}
$M \otimes_R S$ が $S$-加群として次のいずれかの性質をもつとする：
\begin{enumerate}
\item[(a)]
有限生成である；
\item[(b)]
有限表示である；
\item[(c)]
平坦である；
\item[(d)]
忠実平坦である；
\item[(e)]
有限射影である。
\end{enumerate}
このとき $M$ も $R$-加群として同じ性質をもち、逆も成り立つ。
\end{theorem}

\begin{proof}
(a) を示すため、$M \otimes_R S$ の $\text{Mod}_S$ における生成元の有限集合
$\{n_i\}$ を選ぶ。各 $n_i$ を $\sum_j m_{ij} \otimes s_{ij}$ と書く。ただし
$m_{ij} \in M$ かつ $s_{ij} \in S$ である。$e_{ij}$ を基底にもつ有限自由
$R$-加群を $F$ とし、$e_{ij}$ を $m_{ij}$ に写す $R$-加群準同型
$F \to M$ をとる。このとき $F \otimes_R S\to M \otimes_R S$ は全射なので、
$\Coker(F \to M) \otimes_R S$ は零であり、したがって
$\Coker(F \to M)$ も零である。これで (a) が示された。

\medskip\noindent
(b) を示すため、$M \otimes_R S$ が有限表示であると仮定する。(a) により
$M$ は有限生成である。核が $K$ である全射 $R^{\oplus n} \to M$ を選ぶ。
すると $K \otimes_R S \to S^{\oplus r} \to M \otimes_R S \to 0$ は完全である。
Algebra, Lemma \ref{algebra-lemma-extension} により、
$S^{\oplus r} \to M \otimes_R S$ の核は有限 $S$-加群である。したがって、
$S^{\oplus r}$ における $k_i \otimes 1$ の像が
$S^{\oplus r} \to M \otimes_R S$ の核を生成するような有限個の元
$k_1, \ldots, k_t \in K$ をとれる。$k_1, \ldots, k_t$ が生成する部分加群を
$K' \subset K$ とする。このとき $M' = R^{\oplus r}/K'$ は有限表示
$R$-加群であり、射 $M' \to M$ をもち、
$M' \otimes_R S \to M \otimes_R S$ は同型である。ゆえに所望のとおり
$M' \cong M$ となる。

\medskip\noindent
(c) を示す。$0 \to M' \to M'' \to M \to 0$ を $\text{Mod}_R$ における
短完全列とする。$\bullet \otimes_R S$ は右完全関手なので、
$M'' \otimes_R S \to M \otimes_R S$ は全射である。したがって補題
\ref{descent-lemma-C-is-faithful} により、写像
$C(M \otimes_R S) \to C(M'' \otimes_R S)$ は単射である。
$M \otimes_R S$ が平坦ならば、補題 \ref{descent-lemma-flat-to-injective} より
$C(M \otimes_R S)$ は $\text{Mod}_S$ の入射対象であるから、単射
$C(M \otimes_R S) \to C(M'' \otimes_R S)$ は $\text{Mod}_S$ で分裂し、
したがって $\text{Mod}_R$ でも分裂する。また補題
\ref{descent-lemma-split-surjection} により
$C(M \otimes_R S) \to C(M)$ は分裂全射なので、
$C(M) \to C(M'')$ は $\text{Mod}_R$ における分裂単射である。すなわち列
$$
0 \to C(M) \to C(M'') \to C(M') \to 0
$$
は分裂完全である。$N \in \text{Mod}_R$ に対し、
(\ref{descent-equation-adjunction}) により列
$$
0 \to C(M \otimes_R N) \to C(M'' \otimes_R N) \to C(M' \otimes_R N) \to 0
$$
も分裂完全である。補題 \ref{descent-lemma-C-is-faithful} により、
$$
0 \to M' \otimes_R N \to M'' \otimes_R N \to M \otimes_R N \to 0
$$
は完全である。これは $M$ が $R$ 上平坦であることを意味する。実際、$M'$ を
$M$ へ全射する自由加群としてとれば、任意の加群 $N$ に対し
$\text{Tor}_1^R(M, N) = 0$ であると結論でき、Algebra, Lemma
\ref{algebra-lemma-characterize-flat} を適用できる。これで (c) が示された。

\medskip\noindent
(c) から (d) を導く。$N \in \text{Mod}_R$ かつ $M \otimes_R N$ が零ならば、
$M \otimes_R S \otimes_S (N \otimes_R S) \cong (M \otimes_R N) \otimes_R 
S$ は零である。ゆえに $N \otimes_R S$ は零であり、したがって $N$ も零である。

\medskip\noindent
最後に (e) を導くには、$M$ が有限生成かつ射影であることと、有限表示かつ
平坦であることが同値だと想起すればよい。Algebra, Lemma
\ref{algebra-lemma-finite-projective} を参照せよ。
\end{proof}

\noindent
$R$-代数については次の変種がある。

\begin{theorem}
\label{descent-theorem-descend-algebra-properties}
$A \otimes_R S$ が $S$-代数として次のいずれかの性質をもつとする：
\begin{enumerate}
\item[(a)]
有限型である；
\item[(b)]
有限表示である；
\item[(c)]
形式的不分岐である；
\item[(d)]
不分岐である；
\item[(e)]
エタールである。
\end{enumerate}
このとき $A$ も $R$-代数として同じ性質をもち、もちろん逆も成り立つ。
\end{theorem}

\begin{proof}
(a) を示すため、$A \otimes_R S$ の $S$ 上の生成元の有限集合 $\{x_i\}$ を
選ぶ。各 $x_i$ を $\sum_j y_{ij} \otimes s_{ij}$ と書く。ただし
$y_{ij} \in A$ かつ $s_{ij} \in S$ である。変数 $e_{ij}$ 上の多項式
$R$-代数を $F$ とし、$e_{ij}$ を $y_{ij}$ に写す $R$-代数準同型
$F \to M$ をとる。このとき $F \otimes_R S\to A \otimes_R S$ は全射なので、
$\Coker(F \to A) \otimes_R S$ は零であり、したがって
$\Coker(F \to A)$ も零である。これで (a) が示された。

\medskip\noindent
(b) を示すため、$A \otimes_R S$ が有限表示 $S$-代数であると仮定する。
(a) により $A$ は $R$ 上有限型である。核が $I$ である全射
$R[x_1, \ldots, x_n] \to A$ を選ぶ。すると
$I \otimes_R S \to S[x_1, \ldots, x_n] \to A \otimes_R S \to 0$ は完全である。
Algebra, Lemma \ref{algebra-lemma-finite-presentation-independent} により、
$S[x_1, \ldots, x_n] \to A \otimes_R S$ の核は有限生成イデアルである。
したがって、$S[x_1, \ldots, x_n]$ における $y_i \otimes 1$ の像が
$S[x_1, \ldots, x_n] \to A \otimes_R S$ の核を生成するような有限個の元
$y_1, \ldots, y_t \in I$ をとれる。$y_1, \ldots, y_t$ が生成するイデアルを
$I' \subset I$ とする。このとき
$A' = R[x_1, \ldots, x_n]/I'$ は有限表示 $R$-代数であり、射
$A' \to A$ をもち、$A' \otimes_R S \to A \otimes_R S$ は同型である。
ゆえに所望のとおり $A' \cong A$ となる。

\medskip\noindent
(c) を示す。$A$ が $R$ 上形式的不分岐であることと、相対微分加群
$\Omega_{A/R}$ が零であることは同値である。Algebra, Lemma
\ref{algebra-lemma-characterize-formally-unramified} または
\cite[Proposition~17.2.1]{EGA4} を参照せよ。
$\Omega_{(A \otimes_R S)/S} = \Omega_{A/R} \otimes_R S$ なので、
定理 \ref{descent-theorem-descent} によりこの消滅性は降下する。

\medskip\noindent
先の各場合から (d) を導く。$A$ が $R$ 上不分岐であることと、$A$ が $R$ 上
形式的不分岐かつ有限型であることは同値である。Algebra, Lemma
\ref{algebra-lemma-formally-unramified-unramified} を参照せよ。

\medskip\noindent
(e) を示す。Algebra, Lemma
\ref{algebra-lemma-etale-flat-unramified-finite-presentation} または
\cite[Th\'eor\`eme~17.6.1]{EGA4} により、代数 $A$ が $R$ 上エタールである
ことと、$A$ が $R$ 上平坦、不分岐、かつ有限表示であることは同値である。
\end{proof}

\begin{remark}
\label{descent-remark-when-locally-split}
$T \to S \otimes_R T$ が何らかの追加構造をもつような忠実平坦環準同型
$g: R \to T$ があれば、議論はより容易になる。たとえば、
$T \to S \otimes_R T$ が $\textit{Rings}$ で分裂することを保証できれば、
基底拡大で安定であり、かつ忠実平坦射に沿って降下する加群または代数の
任意の性質は、普遍単射な射に沿っても降下することになる。自然な予想は、
$T$ が忠実平坦であるだけでなく $\text{Mod}_R$ で入射的でもあるような
$g$ を見つけることであろう。しかし $R = \mathbf{Z}$ の場合でさえ、
そのような準同型は存在しえない。
\end{remark}
















\section{準連接層の fpqc 降下}
\label{descent-section-fpqc-descent-quasi-coherent}

\noindent
加群の平坦降下の主要な応用は、fpqc 被覆に関する準連接層の
対応する降下命題である。

\begin{lemma}
\label{descent-lemma-standard-fpqc-covering}
$S$ をアフィンスキームとする。
$\mathcal{U} = \{f_i : U_i \to S\}_{i = 1, \ldots, n}$ を $S$ の標準 fpqc 被覆とする。
Topologies, Definition \ref{topologies-definition-standard-fpqc} を参照せよ。
$\mathcal{U} = \{U_i \to S\}$ に関する準連接層上の任意の降下データは有効である。
さらに、準連接 $\mathcal{O}_S$-加群の圏から $\mathcal{U}$ に関する降下データの
圏への関手は充満忠実である。
\end{lemma}

\begin{proof}
これは命題 \ref{descent-proposition-descent-module} をスキームの言葉で言い換えたものである。
まず、$\mathcal{U}$ に関する準連接層上の降下データ $\xi$ は、被覆
$\mathcal{U}' = \{\coprod_{i = 1, \ldots, n} U_i \to S\}$ に関する準連接層上の
降下データ $\xi'$ とまったく同じものである。また、$\xi$ の有効性と $\xi'$ の
有効性も同じである。したがって $n = 1$、すなわち
$\mathcal{U} = \{U \to S\}$ で $U$ と $S$ がアフィンである場合を仮定してよい。
この場合、降下データは環準同型
$$
\Gamma(S, \mathcal{O})
\longrightarrow
\Gamma(U, \mathcal{O}).
$$
に関する加群上の降下データに対応する。$U \to S$ は全射かつ平坦なので、
この環準同型は忠実平坦である。すなわち、命題
\ref{descent-proposition-descent-module} を適用すればよい。
\end{proof}

\begin{proposition}
\label{descent-proposition-fpqc-descent-quasi-coherent}
$S$ をスキームとする。
$\mathcal{U} = \{\varphi_i : U_i \to S\}$ を fpqc 被覆とする。Topologies,
Definition \ref{topologies-definition-fpqc-covering} を参照せよ。
$\mathcal{U} = \{U_i \to S\}$ に関する準連接層上の任意の降下データは有効である。
さらに、準連接 $\mathcal{O}_S$-加群の圏から $\mathcal{U}$ に関する降下データの
圏への関手は充満忠実である。
\end{proposition}

\begin{proof}
$S = \bigcup_{j \in J} V_j$ をアフィン開被覆とする。$j, j' \in J$ に対し、
交わり（アフィンとは限らない）を $V_{jj'} = V_j \cap V_{j'}$ と書く。
開集合 $V \subset S$ に対し、
$\mathcal{U}_V = \{V \times_S U_i \to V\}_{i \in I}$ と書く。これは fpqc 被覆である
（Topologies, Lemma \ref{topologies-lemma-fpqc}）。fpqc 被覆の定義により、
各 $j \in J$ に対して有限集合 $K_j$、写像 $\underline{i} : K_j \to I$、および
$k \in K_j$ ごとのアフィン開集合
$U_{\underline{i}(k), k} \subset U_{\underline{i}(k)}$ で、
$\mathcal{V}_j = \{U_{\underline{i}(k), k} \to V_j\}_{k \in K_j}$ が $V_j$ の標準 fpqc
被覆となるものをとれる。もちろん $\mathcal{V}_j$ は $\mathcal{U}_{V_j}$ の細分である。
図式で表すと
$$
\xymatrix{
\mathcal{V}_j \ar[r] \ar@{~>}[d] &
\mathcal{U}_{V_j} \ar[r] \ar@{~>}[d] &
\mathcal{U} \ar@{~>}[d] \\
V_j \ar@{=}[r] & V_j \ar[r] & S
}
$$
となる。上段の水平矢印は、終域を固定した射の族の間の射である
（Sites, Definition \ref{sites-definition-morphism-coverings} を参照せよ）。

\medskip\noindent
命題を示すには、準連接層から降下データへの関手の忠実性、充満性、
本質的全射性を順に示す。

\medskip\noindent
忠実性。$\mathcal{F}$、$\mathcal{G}$ を $S$ 上の準連接層とし、
$a, b : \mathcal{F} \to \mathcal{G}$ を $\mathcal{O}_S$-加群準同型とする。
すべての $i$ に対し $\varphi_i^*(a) = \varphi_i^*(b)$ と仮定する。
$s \in S$ をとる。ある $i \in I$ と $u \in U_i$ に対し
$s = \varphi_i(u)$ である。$\mathcal{O}_{S, s} \to \mathcal{O}_{U_i, u}$ は平坦であり、
したがって忠実平坦なので（Algebra, Lemma
\ref{algebra-lemma-local-flat-ff}）、
$a_s = b_s : \mathcal{F}_s \to \mathcal{G}_s$ である。ゆえに $a = b$ である。

\medskip\noindent
充満忠実性。$\mathcal{F}$、$\mathcal{G}$ を $S$ 上の準連接層とし、
$a_i : \varphi_i^*\mathcal{F} \to \varphi_i^*\mathcal{G}$ を
$\mathcal{O}_{U_i}$-加群準同型で、$U_i \times_U U_j$ 上
$\text{pr}_0^*a_i = \text{pr}_1^*a_j$ を満たすものとする。
これらの射を引き戻して、射
$$
a_k :
\varphi_{i(k)}^*\mathcal{F}|_{U_{\underline{i}(k), k}}
\longrightarrow
\varphi_{i(k)}^*\mathcal{G}|_{U_{\underline{i}(k), k}}
$$
$k \in K_j$ を得る（記号は上のとおり）。さらに補題
\ref{descent-lemma-refine-descent-datum} により、これらは $\mathcal{V}_j$ 上の
（標準）降下データの間の射を定める。したがって補題
\ref{descent-lemma-standard-fpqc-covering} により、対応する一意な射
$a_j : \mathcal{F}|_{V_j} \to \mathcal{G}|_{V_j}$ を得る。
$a_j|_{V_{jj'}} = a_{j'}|_{V_{jj'}}$ を見るには、$a_j$ と $a_{j'}$ のいずれも、
$\mathcal{V}_{j, V_{jj'}}$ と $\mathcal{V}_{j', V_{jj'}}$ の双方を細分する任意の被覆への
（標準）降下データの射 $(a_i)_{i \in I}$ の引き戻しと一致すること、および既に
示した忠実性を用いる。たとえば被覆
$\mathcal{V}_{jj'} =
\{V_k \times_S V_{k'} \to V_{jj'}\}_{k \in K_j, k' \in K_{j'}}$
を用いればよい。

\medskip\noindent
本質的全射性。$\xi = (\mathcal{F}_i, \varphi_{ii'})$ を被覆 $\mathcal{U}$ に関する
準連接層上の降下データとする。この降下データを引き戻すと、$V_j$ の被覆
$\mathcal{V}_j$ に関する準連接層上の降下データ $\xi_j$ を得る。補題
\ref{descent-lemma-standard-fpqc-covering} により、付随する標準降下データが $\xi_j$ と
同型になる $V_j$ 上の準連接層 $\mathcal{F}_j$ が存在する。上で示した充満忠実性に
より、同型
$$
\phi_{jj'} :
\mathcal{F}_j|_{V_{jj'}}
\longrightarrow
\mathcal{F}_{j'}|_{V_{jj'}}
$$
が存在する。これは、$\xi_j$ と $\xi_{j'}$ の $\mathcal{V}_{jj'}$ への引き戻しの間の
降下データの同型に対応する。これらの写像 $\phi_{jj'}$ がコサイクル条件を
満たすことを見るには、三重交叉 $V_{jj'j''}$ 上で上に示した忠実性を用いる。
したがって補題 \ref{descent-lemma-zariski-descent-effective} により、層
$\mathcal{F}_j$ は所望の準連接層 $\mathcal{F}$ に貼り合わされる。
なお、$\mathcal{F}$ に付随する $\mathcal{U}$ に関する標準降下データが、初めに
与えた降下データと同型であることを確かめなければならない。この確認は省略する。
\end{proof}









\section{準連接層のガロア降下}
\label{descent-section-galois-descent}

\noindent
準連接層のガロア降下は、準連接層の fpqc 降下の特別な場合にすぎない。
この節では、ガロア降下を fpqc 降下に移す方法を説明し、先の結果を
適用して結論を得る。

\medskip\noindent
$k'/k$ を体拡大とする。このとき $\{\Spec(k') \to \Spec(k)\}$ は fpqc 被覆である。
$X$ を $k$ 上のスキームとする。$k$-代数 $A$ に対し
$X_A = X \times_{\Spec(k)} \Spec(A)$ とおく。Topologies, Lemma
\ref{topologies-lemma-fpqc} により、$\{X_{k'} \to X\}$ は fpqc 被覆である。また
$$
X_{k'} \times_X X_{k'} = X_{k' \otimes_k k'}
\quad\text{および}\quad
X_{k'} \times_X X_{k'} \times_X X_{k'} = X_{k' \otimes_k k' \otimes_k k'}
$$
である。したがって $\{X_{k'} \to X\}$ に関する準連接層上の降下データは、
$X_{k'}$ 上の準連接層 $\mathcal{F}$ と、$X_{k' \otimes_k k'}$ 上の同型
$\varphi : \text{pr}_0^*\mathcal{F} \to \text{pr}_1^*\mathcal{F}$ で、
$X_{k' \otimes_k k' \otimes_k k'}$ 上の明らかなコサイクル条件を満たすものにより
与えられる。以下では、ガロア拡大の場合にこれが何を意味するかを具体的に求める。

\medskip\noindent
$k'/k$ をガロア群 $G = \text{Gal}(k'/k)$ をもつ有限ガロア拡大とする。このとき
$k$-代数の同型
$$
k' \otimes_k k' \longrightarrow \prod\nolimits_{\sigma \in G} k',\quad
a \otimes b \longrightarrow \prod a\sigma(b)
$$
および
$$
k' \otimes_k k' \otimes_k k' \longrightarrow
\prod\nolimits_{(\sigma, \tau) \in G \times G} k',\quad
a \otimes b \otimes c \longrightarrow \prod a\sigma(b)\sigma(\tau(c))
$$
がある。ここで $a\sigma(b)\tau(c)$ ではなく
$a\sigma(b)\sigma(\tau(c))$ を選んだのは、以下の式が簡潔になるためであり、
本質的な必要性はない。$\sigma \in G$ に対し
$$
f_\sigma = \text{id}_X \times \Spec(\sigma) :
X_{k'} \longrightarrow X_{k'}
$$
と書く。$\Spec(-)$ は反変関手なので
$f_{\sigma \tau} = f_\tau \circ f_\sigma$ であり、逆の順序ではないことに注意せよ。
上の第一の同型から同一視
$$
X_{k' \otimes_k k'} = \coprod\nolimits_{\sigma \in G} X_{k'}
$$
を得る。この同一視のもとで $\text{pr}_0$ は写像
$$
\coprod\nolimits_{\sigma \in G} X_{k'}
\xrightarrow{\coprod \text{id}}
X_{k'}
$$
に対応し、$\text{pr}_1$ は写像
$$
\coprod\nolimits_{\sigma \in G} X_{k'}
\xrightarrow{\coprod f_\sigma}
X_{k'}
$$
に対応する。したがって、$X_{k'}$ 上の $\mathcal{F}$ の降下データ $\varphi$ は
同型の族 $\varphi_\sigma : \mathcal{F} \to f_\sigma^*\mathcal{F}$ に対応する。
コサイクル条件を求めるため、上の代数の同型から得られる同一視
$$
X_{k' \otimes_k k' \otimes_k k'} =
\coprod\nolimits_{(\sigma, \tau) \in G \times G} X_{k'}.
$$
を用いる。この同一視のもとで写像 $\text{pr}_{01}$ は写像
$$
\coprod\nolimits_{(\sigma, \tau) \in G \times G} X_{k'}
\longrightarrow
\coprod\nolimits_{\sigma \in G} X_{k'}
$$
に対応し、これは添字 $(\sigma, \tau)$ の直和因子を恒等射により添字 $\sigma$ の
直和因子へ写す。写像 $\text{pr}_{12}$ は写像
$$
\coprod\nolimits_{(\sigma, \tau) \in G \times G} X_{k'}
\longrightarrow
\coprod\nolimits_{\sigma \in G} X_{k'}
$$
に対応し、これは添字 $(\sigma, \tau)$ の直和因子を射 $f_\sigma$ により添字
$\tau$ の直和因子へ写す。最後に、写像 $\text{pr}_{02}$ は写像
$$
\coprod\nolimits_{(\sigma, \tau) \in G \times G} X_{k'}
\longrightarrow
\coprod\nolimits_{\sigma \in G} X_{k'}
$$
に対応し、これは添字 $(\sigma, \tau)$ の直和因子を恒等射により添字
$\sigma\tau$ の直和因子へ写す。したがってコサイクル条件
$$
\text{pr}_{02}^*\varphi = \text{pr}_{12}^*\varphi \circ \text{pr}_{01}^*\varphi
$$
は、各組 $(\sigma, \tau)$ に対する条件
$$
\varphi_{\sigma\tau} = f_\sigma^*\varphi_\tau \circ \varphi_\sigma
$$
に移る。ここで両辺は写像 $\mathcal{F} \to f_{\sigma\tau}^*\mathcal{F}$ である。
（すべて整合している。たとえば $\varphi_\sigma$ の終域は
$f_\sigma^*\mathcal{F}$ であり、$f_\sigma^*\varphi_\tau$ の始域も同じく
$f_\sigma^*\mathcal{F}$ である。）

\begin{lemma}
\label{descent-lemma-galois-descent}
$k'/k$ をガロア群 $G$ をもつ（有限）ガロア拡大とし、$X$ を $k$ 上の
スキームとする。準連接 $\mathcal{O}_X$-加群の圏は、次を満たす系
$(\mathcal{F}, (\varphi_\sigma)_{\sigma \in G})$ の圏と同値である：
\begin{enumerate}
\item $\mathcal{F}$ は $X_{k'}$ 上の準連接加群である；
\item $\varphi_\sigma : \mathcal{F} \to f_\sigma^*\mathcal{F}$ は加群の同型である；
\item すべての $\sigma, \tau \in G$ に対し
$\varphi_{\sigma\tau} = f_\sigma^*\varphi_\tau \circ \varphi_\sigma$ である。
\end{enumerate}
ここで $f_\sigma = \text{id}_X \times \Spec(\sigma) : X_{k'} \to X_{k'}$ である。
\end{lemma}

\begin{proof}
上で見たように、補題にあるデータ $(\mathcal{F}, (\varphi_\sigma)_{\sigma \in G})$ は、
fpqc 被覆 $\{X_{k'} \to X\}$ に関する降下データと同じものである。したがって
命題 \ref{descent-proposition-fpqc-descent-quasi-coherent} から従う。
\end{proof}

\noindent
上の議論を少し一般化した場合が次である。全射有限エタール射 $X \to Y$ と
有限群 $G$、および群準同型
$G^{opp} \to \text{Aut}_Y(X), \sigma \mapsto f_\sigma$ が与えられ、写像
$$
G \times X \longrightarrow X \times_Y X,\quad
(\sigma, x) \longmapsto (x, f_\sigma(x))
$$
が同型であると仮定する。このとき上と同じ結果が成り立つ。

\begin{lemma}
\label{descent-lemma-galois-descent-more-general}
$X \to Y$、$G$、および $f_\sigma : X \to X$ を上のとおりとする。
準連接 $\mathcal{O}_Y$-加群の圏は、次を満たす系
$(\mathcal{F}, (\varphi_\sigma)_{\sigma \in G})$ の圏と同値である：
\begin{enumerate}
\item $\mathcal{F}$ は準連接 $\mathcal{O}_X$-加群である；
\item $\varphi_\sigma : \mathcal{F} \to f_\sigma^*\mathcal{F}$ は加群の同型である；
\item すべての $\sigma, \tau \in G$ に対し
$\varphi_{\sigma\tau} = f_\sigma^*\varphi_\tau \circ \varphi_\sigma$ である。
\end{enumerate}
\end{lemma}

\begin{proof}
$X \to Y$ は全射有限エタールなので、$\{X \to Y\}$ は fpqc 被覆である。
$G \times X \to X \times_Y X$, $(\sigma, x) \mapsto (x, f_\sigma(x))$ は同型なので、
$G \times G \times X \to X \times_Y X \times_Y X$,
$(\sigma, \tau, x) \mapsto (x, f_\sigma(x), f_{\sigma\tau}(x))$ も同型である。
これらの同一視のもとで、補題にあるデータの圏は、被覆 $\{x \to Y\}$ に関する
準連接層上の降下データの圏と同じである。したがって命題
\ref{descent-proposition-fpqc-descent-quasi-coherent} から従う。
\end{proof}









\section{加群の有限性の降下}
\label{descent-section-descent-finiteness}

\noindent
この節では、準連接加群がある種の有限性条件を満たすかどうかは、
被覆の各要素上で確かめればよいことを示す。

\begin{lemma}
\label{descent-lemma-finite-type-descends}
$X$ をスキーム、$\mathcal{F}$ を準連接 $\mathcal{O}_X$-加群とする。
$\{f_i : X_i \to X\}_{i \in I}$ を fpqc 被覆とし、各
$f_i^*\mathcal{F}$ が有限型 $\mathcal{O}_{X_i}$-加群であると仮定する。
このとき $\mathcal{F}$ は有限型 $\mathcal{O}_X$-加群である。
\end{lemma}

\begin{proof}
省略する。アフィンの場合については Algebra, Lemma
\ref{algebra-lemma-descend-properties-modules} を参照せよ。
\end{proof}

\begin{lemma}
\label{descent-lemma-finite-type-descends-fppf}
$f : (X, \mathcal{O}_X) \to (Y, \mathcal{O}_Y)$ を局所環付き空間の射とし、
$\mathcal{F}$ を $\mathcal{O}_Y$-加群の層とする。次を仮定する：
\begin{enumerate}
\item $f$ は位相空間の写像として開である；
\item $f$ は全射かつ平坦である；
\item $f^*\mathcal{F}$ は有限型である。
\end{enumerate}
このとき $\mathcal{F}$ は有限型である。
\end{lemma}

\begin{proof}
$y \in Y$ を点とし、$y$ に写る点 $x \in X$ を選ぶ。開集合
$x \in U \subset X$ と、$U$ 上で $f^*\mathcal{F}$ を生成する
$f^*\mathcal{F}(U)$ の元 $s_1, \ldots, s_n$ を選ぶ。
$f^*\mathcal{F} =
f^{-1}\mathcal{F} \otimes_{f^{-1}\mathcal{O}_Y} \mathcal{O}_X$ なので、$U$ を縮小して
$s_i = \sum t_{ij} \otimes a_{ij}$ と仮定できる。ただし
$t_{ij} \in f^{-1}\mathcal{F}(U)$ かつ $a_{ij} \in \mathcal{O}_X(U)$ である。
$U$ をさらに縮小して、$f(U) \subset V$ である開集合 $V \subset Y$ 上の切断
$s_{ij} \in \mathcal{F}(V)$ から $t_{ij}$ が来ると仮定できる。切断 $s_{ij}$ の
個数を $N$ とし、写像
$$
\sigma = (s_{ij}) : \mathcal{O}_V^{\oplus N} \to \mathcal{F}|_V
$$
を考える。切断の選び方により $f^*\sigma|_U$ は全射である。したがって
すべての $u \in U$ に対して写像
$$
\sigma_{f(u)} \otimes_{\mathcal{O}_{Y, f(u)}} \mathcal{O}_{X, u} :
\mathcal{O}_{X, u}^{\oplus N}
\longrightarrow
\mathcal{F}_{f(u)} \otimes_{\mathcal{O}_{Y, f(u)}} \mathcal{O}_{X, u}
$$
は全射である。$f$ は平坦なので、局所環準同型
$\mathcal{O}_{Y, f(u)} \to \mathcal{O}_{X, u}$ は平坦であり、したがって忠実平坦である
（Algebra, Lemma \ref{algebra-lemma-local-flat-ff}）。ゆえに
$\sigma_{f(u)}$ は全射である。$f$ は開なので、$f(U)$ は $y$ の開近傍であり、
証明が完了する。
\end{proof}

\begin{lemma}
\label{descent-lemma-finite-presentation-descends}
$X$ をスキーム、$\mathcal{F}$ を準連接 $\mathcal{O}_X$-加群とする。
$\{f_i : X_i \to X\}_{i \in I}$ を fpqc 被覆とし、各
$f_i^*\mathcal{F}$ が有限表示 $\mathcal{O}_{X_i}$-加群であると仮定する。
このとき $\mathcal{F}$ は有限表示 $\mathcal{O}_X$-加群である。
\end{lemma}

\begin{proof}
省略する。アフィンの場合については Algebra, Lemma
\ref{algebra-lemma-descend-properties-modules} を参照せよ。
\end{proof}

\begin{lemma}
\label{descent-lemma-locally-generated-by-r-sections-descends}
$X$ をスキーム、$\mathcal{F}$ を準連接 $\mathcal{O}_X$-加群とする。
$\{f_i : X_i \to X\}_{i \in I}$ を fpqc 被覆とし、各
$f_i^*\mathcal{F}$ が $\mathcal{O}_{X_i}$-加群として局所的に $r$ 個の切断で
生成されると仮定する。このとき $\mathcal{F}$ は $\mathcal{O}_X$-加群として
局所的に $r$ 個の切断で生成される。
\end{lemma}

\begin{proof}
補題 \ref{descent-lemma-finite-type-descends} により $\mathcal{F}$ は有限型である。
したがって中山の補題（Algebra, Lemma \ref{algebra-lemma-NAK}）により、
$\mathcal{F}$ が点 $x \in X$ の近傍で $r$ 個の切断により生成されることと、
$\dim_{\kappa(x)} \mathcal{F}_x \otimes \kappa(x) \leq r$ であることは同値である。
$i$ と、$x$ に写る点 $x_i \in X_i$ を選ぶ。このとき
$\dim_{\kappa(x)} \mathcal{F}_x \otimes \kappa(x) = 
\dim_{\kappa(x_i)} (f_i^*\mathcal{F})_{x_i} \otimes \kappa(x_i)$
であり、$f_i^*\mathcal{F}$ は局所的に $r$ 個の切断で生成されるので、これは
$\leq r$ である。
\end{proof}

\begin{lemma}
\label{descent-lemma-flat-descends}
$X$ をスキーム、$\mathcal{F}$ を準連接 $\mathcal{O}_X$-加群とする。
$\{f_i : X_i \to X\}_{i \in I}$ を fpqc 被覆とし、各
$f_i^*\mathcal{F}$ が平坦 $\mathcal{O}_{X_i}$-加群であると仮定する。
このとき $\mathcal{F}$ は平坦 $\mathcal{O}_X$-加群である。
\end{lemma}

\begin{proof}
省略する。アフィンの場合については Algebra, Lemma
\ref{algebra-lemma-descend-properties-modules} を参照せよ。
\end{proof}

\begin{lemma}
\label{descent-lemma-finite-locally-free-descends}
$X$ をスキーム、$\mathcal{F}$ を準連接 $\mathcal{O}_X$-加群とする。
$\{f_i : X_i \to X\}_{i \in I}$ を fpqc 被覆とし、各
$f_i^*\mathcal{F}$ が有限局所自由 $\mathcal{O}_{X_i}$-加群であると仮定する。
このとき $\mathcal{F}$ は有限局所自由 $\mathcal{O}_X$-加群である。
\end{lemma}

\begin{proof}
準連接層が有限局所自由であることと、有限表示かつ平坦であることは同値なので
従う。Algebra, Lemma \ref{algebra-lemma-finite-projective} を参照せよ。
実際、各 $f_i^*\mathcal{F}$ が平坦かつ有限表示ならば、補題
\ref{descent-lemma-flat-descends} および
\ref{descent-lemma-finite-presentation-descends} により $\mathcal{F}$ も同様である。
\end{proof}

\noindent
局所射影準連接層の定義は Properties, Section
\ref{properties-section-locally-projective} にある。

\begin{lemma}
\label{descent-lemma-locally-projective-descends}
$X$ をスキーム、$\mathcal{F}$ を準連接 $\mathcal{O}_X$-加群とする。
$\{f_i : X_i \to X\}_{i \in I}$ を fpqc 被覆とし、各
$f_i^*\mathcal{F}$ が局所射影 $\mathcal{O}_{X_i}$-加群であると仮定する。
このとき $\mathcal{F}$ は局所射影 $\mathcal{O}_X$-加群である。
\end{lemma}

\begin{proof}
省略する。Zariski 被覆については Properties, Lemma
\ref{properties-lemma-locally-projective} であり、アフィンの場合については
Algebra, Theorem \ref{algebra-theorem-ffdescent-projectivity} である。
\end{proof}

\begin{remark}
\label{descent-remark-locally-free-descends}
局所自由であるという準連接加群の性質は、fpqc 位相では降下しない。
実際、$R$ を環とし、$M$ を階数 1 の局所自由加群の可算直和
$M = \bigoplus L_n$ である射影 $R$-加群で、局所自由ではないものとする。
Examples, Lemma \ref{examples-lemma-projective-not-locally-free} を参照せよ。
忠実平坦基底拡大
$$
R \longrightarrow
\bigoplus\nolimits_{m \geq 1}
\bigoplus\nolimits_{(i_1, \ldots, i_m) \in \mathbf{Z}^{\oplus m}}
L_1^{\otimes i_1} \otimes_R \ldots \otimes_R L_m^{\otimes i_m}
$$
を行うと $M$ は自由になる。しかし fppf 被覆の場合に何が起こるかは分からない。
言い換えると、次の問いへの答えは分かっていない。$A \to B$ を有限表示な
忠実平坦環準同型とする。$M$ を $A$-加群とし、$M \otimes_A B$ が自由であるとする。
$M$ は局所自由 $A$-加群か。$A$ が Noether 環ならば答えは肯定的であり、これは
\cite{Bass} の結果から従う。しかし一般の場合の答えは分かっていない。答えまたは
参考文献を知っている場合は、次へ電子メールを送られたい：
\href{mailto:stacks.project@gmail.com}{stacks.project@gmail.com}.
\end{remark}

\noindent
ここで、上の結果に関連するが少し性質の異なる二つの結果も加える。

\begin{lemma}
\label{descent-lemma-finite-over-finite-module}
$f : X \to Y$ をスキームの射、$\mathcal{F}$ を準連接
$\mathcal{O}_X$-加群とし、$f$ は有限射であると仮定する。このとき
$\mathcal{F}$ が有限型 $\mathcal{O}_X$-加群であることと、
$f_*\mathcal{F}$ が有限型 $\mathcal{O}_Y$-加群であることは同値である。
\end{lemma}

\begin{proof}
$f$ は有限なのでアフィンである。したがって、有限環準同型 $A \to B$ が定める射
$f$ が $\Spec(B) \to \Spec(A)$ である場合に帰着される。このときさらに、
$\mathcal{F} = \widetilde{M}$ は $B$-加群 $M$ に付随する加群の層である。
$M$ が有限 $B$-加群であることと $M$ が有限 $A$-加群であることは同値である。
Algebra, Lemma \ref{algebra-lemma-finite-module-over-finite-extension} を参照せよ。
これと Properties, Lemma \ref{properties-lemma-finite-type-module} を合わせれば
補題が示される。
\end{proof}

\begin{lemma}
\label{descent-lemma-finite-finitely-presented-module}
$f : X \to Y$ をスキームの射、$\mathcal{F}$ を準連接
$\mathcal{O}_X$-加群とし、$f$ は有限かつ有限表示であると仮定する。このとき
$\mathcal{F}$ が有限表示 $\mathcal{O}_X$-加群であることと、
$f_*\mathcal{F}$ が有限表示 $\mathcal{O}_Y$-加群であることは同値である。
\end{lemma}

\begin{proof}
$f$ は有限なのでアフィンである。したがって、有限かつ有限表示な環準同型
$A \to B$ が定める射 $f$ が $\Spec(B) \to \Spec(A)$ である場合に帰着される。
このときさらに、$\mathcal{F} = \widetilde{M}$ は $B$-加群 $M$ に付随する
加群の層である。$M$ が有限表示 $B$-加群であることと $M$ が有限表示 $A$-加群である
ことは同値である。Algebra, Lemma
\ref{algebra-lemma-finite-finitely-presented-extension} を参照せよ。
これと Properties, Lemma \ref{properties-lemma-finite-presentation-module} を
合わせれば補題が示される。
\end{proof}
















\section{準連接層と位相、I}
\label{descent-section-quasi-coherent-sheaves}

\noindent
この節の結果は、スキーム $S$ 上の準連接加群の圏と、位相の章で $S$ に
付随させた多くのサイト上の準連接加群の圏との間に自然な同値があることを述べる。

\medskip\noindent
$S$ をスキーム、$\mathcal{F}$ を準連接 $\mathcal{O}_S$-加群とする。関手
\begin{equation}
\label{descent-equation-quasi-coherent-presheaf}
(\Sch/S)^{opp} \longrightarrow \textit{Ab},
\quad
(f : T \to S) \longmapsto \Gamma(T, f^*\mathcal{F}).
\end{equation}
を考える。

\begin{lemma}
\label{descent-lemma-sheaf-condition-holds}
$S$ をスキーム、$\mathcal{F}$ を準連接 $\mathcal{O}_S$-加群とし、
$\tau \in \{Zariski, \linebreak[0] \etale, \linebreak[0] smooth,
\linebreak[0] syntomic, \linebreak[0] fppf, \linebreak[0] fpqc\}$ とする。
(\ref{descent-equation-quasi-coherent-presheaf}) で定義した関手は、$S$ 上の任意のスキーム
$T$ の任意の $\tau$-被覆 $\{T_i \to T\}_{i \in I}$ に関して層条件を満たす。
\end{lemma}

\begin{proof}
$\tau \in \{Zariski, \linebreak[0] \etale, \linebreak[0] smooth,
\linebreak[0] syntomic, \linebreak[0] fppf\}$ の場合、$\tau$-被覆は fpqc 被覆でもある。
Topologies, Lemmas
\ref{topologies-lemma-zariski-etale},
\ref{topologies-lemma-zariski-etale-smooth},
\ref{topologies-lemma-zariski-etale-smooth-syntomic},
\ref{topologies-lemma-zariski-etale-smooth-syntomic-fppf}, および
\ref{topologies-lemma-zariski-etale-smooth-syntomic-fppf-fpqc} の結果を参照せよ。
したがって fpqc 被覆について示せば十分である。$f : T \to S$ が与えられ、
$\{f_i : T_i \to T\}_{i \in I}$ が fpqc 被覆であると仮定する。切断の族
$s_i \in \Gamma(T_i , f_i^*f^*\mathcal{F})$ が
$s_i|_{T_i \times_T T_j} = s_j|_{T_i \times_T T_j}$ を満たすとする。
対応する切断 $s \in \Gamma(T, f^*\mathcal{F})$ を見つけなければならない。
$s_i$ を、$T$ 上の準連接層 $\mathcal{O}_T$ と $f^*\mathcal{F}$ に付随する
標準降下データと両立する写像の族
$\varphi_i : f_i^*\mathcal{O}_T = \mathcal{O}_{T_i} \to f_i^*f^*\mathcal{F}$
と解釈し直せる。したがって命題 \ref{descent-proposition-fpqc-descent-quasi-coherent} により、
これらは写像 $\mathcal{O}_T \to f^*\mathcal{F}$ へ（一意に）降下し、この写像が
切断 $s$ を与える。
\end{proof}

\noindent
特に、次の定義を置ける。

\begin{definition}
\label{descent-definition-structure-sheaf}
$\tau \in \{Zariski, \linebreak[0] \etale, \linebreak[0]
smooth, \linebreak[0] syntomic, \linebreak[0] fppf\}$ とする。
$S$ をスキーム、$\Sch_\tau$ を $S$ を含む大サイト、$\mathcal{F}$ を準連接
$\mathcal{O}_S$-加群とする。
\begin{enumerate}
\item 大サイト $(\Sch/S)_\tau$ の {\it 構造層}とは、環の層
$T/S \mapsto \Gamma(T, \mathcal{O}_T)$ であり、$\mathcal{O}$ または
$\mathcal{O}_S$ と書く。
\item $\tau = Zariski$ または $\tau = \etale$ の場合、小サイト $S_{Zar}$ または
$S_\etale$ の {\it 構造層}とは、環の層
$T/S \mapsto \Gamma(T, \mathcal{O}_T)$ であり、$\mathcal{O}$ または
$\mathcal{O}_S$ と書く。
\item 大サイト $(\Sch/S)_\tau$ 上で $\mathcal{F}$ に付随する
{\it $\mathcal{O}$-加群の層}とは、$\mathcal{O}$-加群の層
$(f : T \to S) \mapsto \Gamma(T, f^*\mathcal{F})$ であり、
$\mathcal{F}^a$（しばしば単に $\mathcal{F}$）と書く。
\item $\tau = Zariski$ または $\tau = \etale$ の場合、小サイト $S_{Zar}$ または
$S_\etale$ 上で $\mathcal{F}$ に付随する {\it $\mathcal{O}$-加群の層}とは、
$\mathcal{O}$-加群の層 $(f : T \to S) \mapsto \Gamma(T, f^*\mathcal{F})$ であり、
$\mathcal{F}^a$（しばしば単に $\mathcal{F}$）と書く。
\end{enumerate}
\end{definition}

\noindent
各場合に同じ記号 $\mathcal{F}^a$ を用いていることに注意せよ。定義上、層
$\mathcal{F}^a$ を定める規則は考察するサイトによらないので、実際に混乱は生じない。

\begin{remark}
\label{descent-remark-Zariski-site-space}
Topologies, Lemma \ref{topologies-lemma-Zariski-usual} で、スキーム $S$ の
小 Zariski サイトは、層の圏が自然に同値であるという意味で、位相空間としての
$S$ と同値であることを見た。いま $S_{Zar}$ にも構造層 $\mathcal{O}$ が備わるので、
環付きサイト $(S_{Zar}, \mathcal{O})$ 上の加群の層は、環付き空間
$(S, \mathcal{O}_S)$ 上の加群の層と一致する。
\end{remark}

\begin{remark}
\label{descent-remark-change-topologies-ringed}
$f : T \to S$ をスキームの射とする。Topologies, Section
\ref{topologies-section-change-topologies} に列挙したサイトの各射 $f_{sites}$ は、
環付きサイトの射となる。実際、サイトのこれらの射
$f_{sites} : (\Sch/T)_\tau \to (\Sch/S)_{\tau'}$ または
$f_{sites} : (\Sch/S)_\tau \to S_{\tau'}$ は、連続関手
$S'/S \mapsto T \times_S S'/S$ により与えられる。そこで $S'/S$ が与えられたとき、
$$
f_{sites}^\sharp :
\mathcal{O}(S'/S)
\longrightarrow
f_{sites, *}\mathcal{O}(S'/S) =
\mathcal{O}(T \times_S S'/T)
$$
を通常の写像
$\text{pr}_{S'}^\sharp : \mathcal{O}(S') \to \mathcal{O}(T \times_S S')$ とする。
同様に、$\tau \in \{Zar, \etale\}$ に対する射
$i_f : \Sh(T_\tau) \to \Sh((\Sch/S)_\tau)$
も環付きトポスの射となる。Topologies, Lemmas
\ref{topologies-lemma-put-in-T} および
\ref{topologies-lemma-put-in-T-etale} を参照せよ。これは
$i_f^{-1}\mathcal{O} = \mathcal{O}$ だからである。いくつかの特別な場合を挙げる：
\begin{enumerate}
\item 大サイトの射
$f_{big} : (\Sch/X)_{fppf} \to (\Sch/Y)_{fppf}$ は、環付きサイトの射
$$
(f_{big}, f_{big}^\sharp) :
((\Sch/X)_{fppf}, \mathcal{O}_X)
\longrightarrow
((\Sch/Y)_{fppf}, \mathcal{O}_Y)
$$
となる。Modules on Sites, Definition \ref{sites-modules-definition-ringed-site} を
参照せよ。大 syntomic、smooth、\'etale、および Zariski サイトについても同様である。
\item 小サイトの射 $f_{small} : X_\etale \to Y_\etale$ は、環付きサイトの射
$$
(f_{small}, f_{small}^\sharp) :
(X_\etale, \mathcal{O}_X)
\longrightarrow
(Y_\etale, \mathcal{O}_Y)
$$
となる。Modules on Sites, Definition \ref{sites-modules-definition-ringed-site} を
参照せよ。小 Zariski サイトについても同様である。
\end{enumerate}
\end{remark}

\noindent
$S$ をスキームとする。（たとえば）$(\Sch/S)_{Zar}$ 上の $\mathcal{O}$-加群が
与えられたとき、その（たとえば）$(\Sch/S)_{fppf}$ への引き戻しが単に fppf
層化であることは明らかである。大サイトと小サイトを比較するときに何が起こるかを
見るため、次を示す。

\begin{lemma}
\label{descent-lemma-compare-sites}
$S$ をスキームとする。
$$
\begin{matrix}
\text{id}_{\tau, Zar} & : & (\Sch/S)_\tau \to S_{Zar}, &
\tau \in \{Zar, \etale, smooth, syntomic, fppf\} \\
\text{id}_{\tau, \etale} & : &
(\Sch/S)_\tau \to S_\etale, &
\tau \in \{\etale, smooth, syntomic, fppf\} \\
\text{id}_{small, \etale, Zar} & : & S_\etale \to S_{Zar},
\end{matrix}
$$
を Remark \ref{descent-remark-change-topologies-ringed} の環付きサイトの射とする。
$\mathcal{F}$ を $\mathcal{O}_S$-加群の層とし、$S_{Zar}$ 上の
$\mathcal{O}$-加群の層とみなす。このとき
\begin{enumerate}
\item $(\text{id}_{\tau, Zar})^*\mathcal{F}$ は、$(\Sch/S)_\tau$ 上の Zariski 層
$$
(f : T \to S) \longmapsto \Gamma(T, f^*\mathcal{F})
$$
の $\tau$-層化である；
\item $(\text{id}_{small, \etale, Zar})^*\mathcal{F}$ は、$S_\etale$ 上の Zariski 層
$$
(f : T \to S) \longmapsto \Gamma(T, f^*\mathcal{F})
$$
の \'etale 層化である。
\end{enumerate}
$\mathcal{G}$ を $S_\etale$ 上の $\mathcal{O}$-加群の層とする。このとき
\begin{enumerate}
\item[(3)] $(\text{id}_{\tau, \etale})^*\mathcal{G}$ は、\'etale 層
$$
(f : T \to S) \longmapsto \Gamma(T, f_{small}^*\mathcal{G})
$$
の $\tau$-層化である。ここで $f_{small} : T_\etale \to S_\etale$ は Remark
\ref{descent-remark-change-topologies-ringed} の環付き小 \'etale サイトの射である。
\end{enumerate}
\end{lemma}

\begin{proof}
(1) の証明。まず $\tau = Zar$ のとき結果が成り立つことに注意する。この場合、
トポスの射
$i_f : \Sh(T_{Zar}) \to \Sh((\Sch/S)_{Zar})$
があり、$T_{Zar} \to S_{Zar}$ の射として
$\text{id}_{\tau, Zar} \circ i_f = f_{small}$ を満たす。Topologies, Lemmas
\ref{topologies-lemma-put-in-T} および
\ref{topologies-lemma-morphism-big-small} を参照せよ。引き戻しは推移的なので
（Modules on Sites, Lemma
\ref{sites-modules-lemma-push-pull-composition-modules} を参照せよ）、
$i_f^*(\text{id}_{\tau, Zar})^*\mathcal{F} = f_{small}^*\mathcal{F}$
となり、所望の結果を得る。したがって、この補題の直前の注意により、
$(\text{id}_{\tau, Zar})^*\mathcal{F}$ は前層
$T \mapsto \Gamma(T, f^*\mathcal{F})$ の $\tau$-層化である。

\medskip\noindent
(3) の証明は、Topologies, Lemmas \ref{topologies-lemma-put-in-T-etale} および
\ref{topologies-lemma-morphism-big-small-etale} を用いる点を除き、(1) の証明と
まったく同じである。(2) の証明は省略する。
\end{proof}

\begin{remark}
\label{descent-remark-change-topologies-ringed-sites}
Remark \ref{descent-remark-change-topologies-ringed} と
Lemma \ref{descent-lemma-compare-sites} には次の応用がある：
\begin{enumerate}
\item $S$ をスキームとする。構成 $\mathcal{F} \mapsto \mathcal{F}^a$ は、
環付きサイトの射
$\text{id}_{\tau, Zar} : ((\Sch/S)_\tau, \mathcal{O})
\to (S_{Zar}, \mathcal{O})$
または射
$\text{id}_{small, \etale, Zar} :
(S_\etale, \mathcal{O}) \to (S_{Zar}, \mathcal{O})$ による引き戻しである。
\item $f : X \to Y$ をスキームの射とする。Remark
\ref{descent-remark-change-topologies-ringed} の環付きサイトの任意の射 $f_{sites}$ に対し
$$
(f^*\mathcal{F})^a = f_{sites}^*\mathcal{F}^a.
$$
である。これは (1) と、引き戻しが環付きサイトの射の合成と両立することから従う。
Modules on Sites, Lemma
\ref{sites-modules-lemma-push-pull-composition-modules} を参照せよ。
\end{enumerate}
\end{remark}

\begin{lemma}
\label{descent-lemma-quasi-coherent-gives-quasi-coherent}
$S$ をスキーム、$\mathcal{F}$ を準連接 $\mathcal{O}_S$-加群とし、
$\tau \in \{Zariski, \linebreak[0] \etale, \linebreak[0]
smooth, \linebreak[0] syntomic, \linebreak[0] fppf\}$ とする。
\begin{enumerate}
\item 層 $\mathcal{F}^a$ は、Modules on Sites, Definition
\ref{sites-modules-definition-site-local} で定義される $(\Sch/S)_\tau$ 上の
準連接 $\mathcal{O}$-加群である。
\item $\tau = Zariski$ または $\tau = \etale$ の場合、層 $\mathcal{F}^a$ は、
Modules on Sites, Definition \ref{sites-modules-definition-site-local} で定義される
$S_{Zar}$ または $S_\etale$ 上の準連接 $\mathcal{O}$-加群である。
\end{enumerate}
\end{lemma}

\begin{proof}
$\{S_i \to S\}$ を Zariski 被覆で、ある添字集合 $K_i$ と $J_i$ に対して完全列
$$
\bigoplus\nolimits_{k \in K_i} \mathcal{O}_{S_i} \longrightarrow
\bigoplus\nolimits_{j \in J_i} \mathcal{O}_{S_i} \longrightarrow
\mathcal{F}|_{S_i} \longrightarrow 0
$$
をもつものとする。環付き空間上の準連接層の定義により、このような被覆をとれる
（Modules, Definition \ref{modules-definition-quasi-coherent} を参照せよ）。

\medskip\noindent
(1) の証明。$\tau \in \{Zariski, \linebreak[0] fppf, \linebreak[0]
\etale, \linebreak[0] smooth, \linebreak[0] syntomic\}$ とする。
$\mathcal{F}^a|_{(\Sch/S_i)_\tau}$ も完全列
$$
\bigoplus\nolimits_{k \in K_i} \mathcal{O}|_{(\Sch/S_i)_\tau}
\longrightarrow
\bigoplus\nolimits_{j \in J_i} \mathcal{O}|_{(\Sch/S_i)_\tau}
\longrightarrow
\mathcal{F}^a|_{(\Sch/S_i)_\tau} \longrightarrow 0
$$
に入ることは明らかである。したがって Modules on Sites, Lemma
\ref{sites-modules-lemma-local-final-object} により $\mathcal{F}^a$ は準連接である。

\medskip\noindent
(2) の証明。$\tau = \etale$ とする。
$\mathcal{F}^a|_{(S_i)_\etale}$ も完全列
$$
\bigoplus\nolimits_{k \in K_i} \mathcal{O}|_{(S_i)_\etale}
\longrightarrow
\bigoplus\nolimits_{j \in J_i} \mathcal{O}|_{(S_i)_\etale}
\longrightarrow
\mathcal{F}^a|_{(S_i)_\etale} \longrightarrow 0
$$
に入ることは明らかである。したがって Modules on Sites, Lemma
\ref{sites-modules-lemma-local-final-object} により $\mathcal{F}^a$ は準連接である。
$\tau = Zariski$ の場合も同様である（実際、対応する環付きトポスが一致するので、
これはまったくの恒真命題である）。
\end{proof}

\begin{lemma}
\label{descent-lemma-fully-faithful-associated}
$S$ をスキームとし、
$\tau \in \{Zariski, \linebreak[0] \etale, \linebreak[0]
smooth, \linebreak[0] syntomic, \linebreak[0] fppf\}$ とする。Definition
\ref{descent-definition-structure-sheaf} の各関手 $\mathcal{F} \mapsto \mathcal{F}^a$
$$
\QCoh(\mathcal{O}_S) \to \QCoh((\Sch/S)_\tau, \mathcal{O})
\quad\text{または}\quad
\QCoh(\mathcal{O}_S) \to \QCoh(S_\tau, \mathcal{O})
$$
は充満忠実である。
\end{lemma}

\begin{proof}
（補題 \ref{descent-lemma-quasi-coherent-gives-quasi-coherent} により、実際に表示した
関手が得られる。）$S$ 上の $\mathcal{O}_S$-加群を
$(S_{Zar}, \mathcal{O}_S)$ 上の加群と同一視する。準連接加群 $\mathcal{F}$ 上の
関手 $\mathcal{F} \mapsto \mathcal{F}^a$ は、環付きサイトの射 $f$ による引き戻しで
与えられる。Remark \ref{descent-remark-change-topologies-ringed-sites} を参照せよ。
各場合に関手 $f_*$ は、包含関手 $S_{Zar} \to S_\tau$ または
$S_{Zar} \to (\Sch/S)_\tau$ に沿う制限で与えられる（これらのサイトの射の定義に
ついては Topologies, Section \ref{topologies-section-change-topologies} の議論を
参照せよ）。これを $f^*\mathcal{F} = \mathcal{F}^a$ という記述と合わせると、
$\mathcal{F}$ が準連接ならば $f_*f^*\mathcal{F} = \mathcal{F}$ である。したがって
$$
\Hom_\mathcal{O}(\mathcal{F}^a, \mathcal{G}^a) =
\Hom_\mathcal{O}(f^*\mathcal{F}, f^*\mathcal{G}) =
\Hom_{\mathcal{O}_S}(\mathcal{F}, f_*f^*\mathcal{G}) =
\Hom_{\mathcal{O}_S}(\mathcal{F}, \mathcal{G})
$$
となり、所望の結果を得る。
\end{proof}

\begin{proposition}
\label{descent-proposition-equivalence-quasi-coherent}
$S$ をスキームとし、
$\tau \in \{Zariski, \linebreak[0] \etale, \linebreak[0]
smooth, \linebreak[0] syntomic, \linebreak[0] fppf\}$ とする。
\begin{enumerate}
\item 関手 $\mathcal{F} \mapsto \mathcal{F}^a$ は圏の同値
$$
\QCoh(\mathcal{O}_S)
\longrightarrow
\QCoh((\Sch/S)_\tau, \mathcal{O})
$$
を定め、$S$ 上の準連接層の圏と $S$ の大 $\tau$ サイト上の準連接
$\mathcal{O}$-加群の圏とを同値にする。
\item $\tau = Zariski$ または $\tau = \etale$ とする。関手
$\mathcal{F} \mapsto \mathcal{F}^a$ は圏の同値
$$
\QCoh(\mathcal{O}_S)
\longrightarrow
\QCoh(S_\tau, \mathcal{O})
$$
を定め、$S$ 上の準連接層の圏と $S$ の小 $\tau$ サイト上の準連接
$\mathcal{O}$-加群の圏とを同値にする。
\end{enumerate}
\end{proposition}

\begin{proof}
補題 \ref{descent-lemma-quasi-coherent-gives-quasi-coherent} で関手が良定義であることを見た。
補題 \ref{descent-lemma-fully-faithful-associated} により関手は充満忠実である。証明を
終えるため、$(\Sch/S)_\tau$ 上の準連接 $\mathcal{O}$-加群が、$S$ の
$\tau$-被覆に関する準連接層上の降下データを定めることを示す。この降下データを
構成した後、命題 \ref{descent-proposition-fpqc-descent-quasi-coherent} を用いて対応する
$S$ 上の準連接層を得る。

\medskip\noindent
$\mathcal{G}$ を $S$ の大 $\tau$ サイト上の準連接 $\mathcal{O}$-加群とする。
Modules on Sites, Definition \ref{sites-modules-definition-site-local} により、
$S$ の $\tau$-被覆 $\{S_i \to S\}_{i \in I}$ で、各制限
$\mathcal{G}|_{(\Sch/S_i)_\tau}$ が大域表示
$$
\bigoplus\nolimits_{k \in K_i} \mathcal{O}|_{(\Sch/S_i)_\tau}
\longrightarrow
\bigoplus\nolimits_{j \in J_i} \mathcal{O}|_{(\Sch/S_i)_\tau}
\longrightarrow
\mathcal{G}|_{(\Sch/S_i)_\tau} \longrightarrow 0
$$
をもつものが存在する。ただし $J_i$ と $K_i$ はある添字集合である。これは、
$S_i$ 上のある準連接層 $\mathcal{F}_i$ に対し
$\mathcal{G}|_{(\Sch/S_i)_\tau}$ が $\mathcal{F}_i^a$ であることを意味すると主張する。
実際、直和
$\bigoplus\nolimits_{k \in K_i} \mathcal{O}|_{(\Sch/S_i)_\tau}$
および
$\bigoplus\nolimits_{j \in J_i} \mathcal{O}|_{(\Sch/S_i)_\tau}$.
については明らかである。したがって、$S_i$ 上のある準連接層
$\mathcal{K}_i$、$\mathcal{L}_i$ に対し、$\mathcal{G}|_{(\Sch/S_i)_\tau}$ は写像
$\varphi : \mathcal{K}_i^a \to \mathcal{L}_i^a$ の余核である。$(\ )^a$ の
充満忠実性により、$S_i$ 上の準連接層のある写像
$\phi : \mathcal{K}_i \to \mathcal{L}_i$ に対し $\varphi = \phi^a$ である。
すると
$\mathcal{G}|_{(\Sch/S_i)_\tau} \cong \Coker(\phi)^a$
であることは明らかであり、主張が従う。

\medskip\noindent
$\mathcal{G}$ は圏 $(\Sch/S)_\tau$ 全体上にあるので、
$$
(\text{pr}_0^*\mathcal{F}_i)^a
\cong
\mathcal{G}|_{(\Sch/(S_i \times_S S_j))_\tau}
\cong
(\text{pr}_1^*\mathcal{F})^a
$$
が $(\Sch/(S_i \times_S S_j))_\tau$ 上の $\mathcal{O}$-加群として成り立つ。
したがって再び充満忠実性を用いると、$S_i \times_S S_j$ 上の準連接加群の
標準同型
$$
\phi_{ij} :
\text{pr}_0^*\mathcal{F}_i
\longrightarrow
\text{pr}_1^*\mathcal{F}_j
$$
を得る。これらがコサイクル条件を満たすことの確認は省略する。実際に満たすので、
準連接層の降下の有効性と被覆 $\{S_i \to S\}$（命題
\ref{descent-proposition-fpqc-descent-quasi-coherent}）により、与えられた降下データと両立し、
$\mathcal{F}|_{S_i} \cong \mathcal{F}_i$ を満たす $S$ 上の準連接層
$\mathcal{F}$ が存在する。言い換えると、$\mathcal{O}$-加群の同型
$$
\phi_i :
\mathcal{F}^a|_{(\Sch/S_i)_\tau}
\longrightarrow
\mathcal{G}|_{(\Sch/S_i)_\tau}
$$
が与えられ、$S_i \times_S S_j$ 上で一致する。したがって
$\SheafHom_\mathcal{O}(\mathcal{F}^a, \mathcal{G})$ は層なので
（Modules on Sites, Lemma \ref{sites-modules-lemma-internal-hom}）、上の同型
$\phi_i$ を復元する $\mathcal{O}$-加群の射 $\mathcal{F}^a \to \mathcal{G}$ が
存在する。ゆえにこれは同型であり、証明が完了する。

\medskip\noindent
サイト $S_\etale$ と $S_{Zar}$ の場合もまったく同じ方法で証明される。
\end{proof}

\begin{lemma}
\label{descent-lemma-equivalence-quasi-coherent-properties}
$S$ をスキームとし、
$\tau \in \{Zariski, \linebreak[0] \etale, \linebreak[0]
smooth, \linebreak[0] syntomic, \linebreak[0] fppf\}$ とする。
$\mathcal{P}$ を Modules on Sites, Definitions
\ref{sites-modules-definition-global}、
\ref{sites-modules-definition-site-local}、および
\ref{sites-modules-definition-flat} で定義された加群の性質の一つとする\footnote{一覧は
次のとおりである：自由、有限自由、大域切断により生成、$r$ 個の大域切断により生成、
有限個の大域切断により生成、大域表示をもつ、大域有限表示をもつ、局所自由、
有限局所自由、局所的に切断により生成、局所的に $r$ 個の切断により生成、有限型、
有限表示、連接、または平坦。}。規則 $\mathcal{F} \mapsto \mathcal{F}^a$ により
定義される圏の同値
$$
\QCoh(\mathcal{O}_S)
\longrightarrow
\QCoh((\Sch/S)_\tau, \mathcal{O})
\quad\text{および}\quad
\QCoh(\mathcal{O}_S)
\longrightarrow
\QCoh(S_\tau, \mathcal{O})
$$
は、Proposition \ref{descent-proposition-equivalence-quasi-coherent} で見たものであり、性質
$$
\mathcal{F}\text{ は }\mathcal{P}\text{ をもつ}
\Leftrightarrow
\mathcal{F}^a\text{ は }\mathcal{O}\text{-加群として }\mathcal{P}\text{ をもつ}
$$
をもつ。ただし $\mathcal{P}$ が「局所自由」または「連接」の場合は（おそらく）
例外である。$\mathcal{P}=$「連接」の場合、$S$ が局所 Noether スキームで、
$\tau$ が Zariski または \'etale ならば、同値
$\QCoh(\mathcal{O}_S) \to \QCoh(S_\tau, \mathcal{O})$ についてこの主張は成り立つ。
\end{lemma}

\begin{proof}
大域的性質、すなわち Modules on Sites, Definition
\ref{sites-modules-definition-global} で定義された性質については直ちに従う。
局所的性質については Modules on Sites, Lemma
\ref{sites-modules-lemma-local-final-object} を用い、「$\mathcal{F}^a$ が
$\mathcal{P}$ をもつ」ことを $X$ の被覆の各要素上の性質へ移せる。したがって結果は
Lemmas \ref{descent-lemma-finite-type-descends},
\ref{descent-lemma-finite-presentation-descends},
\ref{descent-lemma-locally-generated-by-r-sections-descends},
\ref{descent-lemma-flat-descends}、および
\ref{descent-lemma-finite-locally-free-descends} から従う。準連接加群が連接であることは、
局所 Noether スキーム上で有限型であることと同じである（Cohomology of Schemes,
Lemma \ref{coherent-lemma-coherent-Noetherian} を参照せよ）。したがって有限型加群の
場合に帰着される（詳細は省略する）。
\end{proof}










\section{準連接加群のコホモロジーと位相}
\label{descent-section-quasi-coherent-cohomology}

\noindent
この節では、準連接加群のコホモロジーが位相の選択によらないことを示す。

\begin{lemma}
\label{descent-lemma-standard-covering-Cech}
$S$ をスキームとする。次のいずれかの場合を考える：
\begin{enumerate}
\item[(a)] $\tau \in \{Zariski, \linebreak[0] fppf, \linebreak[0]
\etale, \linebreak[0] smooth, \linebreak[0] syntomic\}$
かつ $\mathcal{C} = (\Sch/S)_\tau$；
\item[(b)] $\tau = \etale$ かつ $\mathcal{C} = S_\etale$；
\item[(c)] $\tau = Zariski$ かつ $\mathcal{C} = S_{Zar}$。
\end{enumerate}
$\mathcal{F}$ を $\mathcal{C}$ 上のアーベル層、$U \in \Ob(\mathcal{C})$ をアフィンとし、
$\mathcal{U} = \{U_i \to U\}_{i = 1, \ldots, n}$ を $\mathcal{C}$ における標準アフィン
$\tau$-被覆とする。このとき
\begin{enumerate}
\item $\mathcal{V} = \{\coprod_{i = 1, \ldots, n} U_i \to U\}$ は $U$ の
$\tau$-被覆である；
\item $\mathcal{U}$ は $\mathcal{V}$ の細分である；
\item 誘導される {\v C}ech 複体の写像（Cohomology on Sites,
Equation (\ref{sites-cohomology-equation-map-cech-complexes})）
$$
\check{\mathcal{C}}^\bullet(\mathcal{V}, \mathcal{F})
\longrightarrow
\check{\mathcal{C}}^\bullet(\mathcal{U}, \mathcal{F})
$$
は複体の同型である。
\end{enumerate}
\end{lemma}

\begin{proof}
これは
$$
(\coprod\nolimits_{i_0 = 1, \ldots, n} U_{i_0}) \times_U
\ldots \times_U
(\coprod\nolimits_{i_p = 1, \ldots, n} U_{i_p})
=
\coprod\nolimits_{i_0, \ldots, i_p \in \{1, \ldots, n\}}
U_{i_0} \times_U \ldots \times_U U_{i_p}
$$
であること、および非交和は $\tau$-被覆なので
$\mathcal{F}(\coprod_a V_a) = \prod_a \mathcal{F}(V_a)$ であることから従う。
\end{proof}

\begin{lemma}
\label{descent-lemma-standard-covering-Cech-quasi-coherent}
$S$ をスキーム、$\mathcal{F}$ を $S$ 上の準連接層とする。$\tau$、$\mathcal{C}$、
$U$、$\mathcal{U}$ を Lemma \ref{descent-lemma-standard-covering-Cech} のとおりとする。
このとき複体の同型
$$
\check{\mathcal{C}}^\bullet(\mathcal{U}, \mathcal{F}^a)
\cong
s((A/R)_\bullet \otimes_R M)
$$
がある（Section \ref{descent-section-descent-modules} を参照せよ）。ここで
$R = \Gamma(U, \mathcal{O}_U)$、$M = \Gamma(U, \mathcal{F}^a)$ であり、
$R \to A$ は忠実平坦環準同型である。特に、すべての $p \geq 1$ に対し
$$
\check{H}^p(\mathcal{U}, \mathcal{F}^a) = 0
$$
である。
\end{lemma}

\begin{proof}
補題 \ref{descent-lemma-standard-covering-Cech} により、
$\check{\mathcal{C}}^\bullet(\mathcal{U}, \mathcal{F}^a)$ は
$\check{\mathcal{C}}^\bullet(\mathcal{V}, \mathcal{F}^a)$ と同型である。ここで
$\mathcal{V} = \{V \to U\}$、$V = \coprod_{i = 1, \ldots n} U_i$ であり、これも
アフィンである。$A = \Gamma(V, \mathcal{O}_V)$ とおく。$\{V \to U\}$ は
$\tau$-被覆なので、$R \to A$ は忠実平坦である。一方、$\mathcal{F}^a$ の定義により、
次数 $p$ の項 $\check{\mathcal{C}}^p(\mathcal{V}, \mathcal{F}^a)$ は
$$
\Gamma(V \times_U \ldots \times_U V, \mathcal{F}^a)
=
\Gamma(\Spec(A \otimes_R \ldots \otimes_R A), \mathcal{F}^a)
=
A \otimes_R \ldots \otimes_R A \otimes_R M
$$
である。{\v C}ech 複体の写像が複体 $s((A/R)_\bullet \otimes_R M)$ の写像と
一致することの確認は省略する。コホモロジーの消滅は補題
\ref{descent-lemma-ff-exact} である。
\end{proof}

\begin{proposition}
\label{descent-proposition-same-cohomology-quasi-coherent}
\begin{slogan}
準連接層のコホモロジーは、どの位相を用いても同じである。
\end{slogan}
$S$ をスキーム、$\mathcal{F}$ を $S$ 上の準連接層とし、
$\tau \in \{Zariski, \linebreak[0] \etale, \linebreak[0]
smooth, \linebreak[0] syntomic, \linebreak[0] fppf\}$ とする。
\begin{enumerate}
\item 標準同型
$$
H^q(S, \mathcal{F}) = H^q((\Sch/S)_\tau, \mathcal{F}^a).
$$
がある。
\item 標準同型
$$
H^q(S, \mathcal{F}) =
H^q(S_{Zar}, \mathcal{F}^a) =
H^q(S_\etale, \mathcal{F}^a).
$$
がある。
\end{enumerate}
\end{proposition}

\begin{proof}
$q = 0$ の場合は $\mathcal{F}^a$ の定義から明らかである。
$\mathcal{C} = (\Sch/S)_\tau$、$\mathcal{C} = S_\etale$、または
$\mathcal{C} = S_{Zar}$ とする。

\medskip\noindent
Cohomology on Sites, Lemma
\ref{sites-cohomology-lemma-cech-vanish-collection} を、
$\mathcal{F} = \mathcal{F}^a$、$\mathcal{B} \subset \Ob(\mathcal{C})$ を
$\mathcal{C}$ におけるアフィンスキームの集合、
$\text{Cov} \subset \text{Cov}_\mathcal{C}$ を標準アフィン $\tau$-被覆の集合として
適用する。補題の仮定 (3) は補題
\ref{descent-lemma-standard-covering-Cech-quasi-coherent} により満たされる。したがって、
$\mathcal{C}$ のすべてのアフィン対象 $U$ に対し
$H^p(U, \mathcal{F}^a) = 0$ である。

\medskip\noindent
次に、$U \in \Ob(\mathcal{C})$ を任意の分離対象とし、構造射を
$f : U \to S$ と書く。$U = \bigcup U_i$ をアフィン開被覆とする。これを
$\mathcal{C}$ における $U$ の $\tau$-被覆
$\mathcal{U} = \{U_i \to U\}$ とみなすこともできる。
$U$ は分離的と仮定したので
$U_{i_0} \times_U \ldots \times_U U_{i_p} =
U_{i_0} \cap \ldots \cap U_{i_p}$ はアフィンである。Cohomology on Sites, Lemma
\ref{sites-cohomology-lemma-cech-spectral-sequence-application}
と上の結果により
$$
H^p(U, \mathcal{F}^a) = \check{H}^p(\mathcal{U}, \mathcal{F}^a)
= H^p(U, f^*\mathcal{F})
$$
である。最後の等号は Cohomology of Schemes, Lemma
\ref{coherent-lemma-cech-cohomology-quasi-coherent} による。特に $S$ が分離的ならば
$U = S$、$f = \text{id}_S$ ととれて、命題が証明される。読者には証明の残りを
飛ばす（または、より明瞭な叙述となるよう書き直す）ことを勧める。

\medskip\noindent
$S$ 上の入射分解 $\mathcal{F} \to \mathcal{I}^\bullet$ と、$\mathcal{C}$ 上の入射分解
$\mathcal{F}^a \to \mathcal{J}^\bullet$ を選ぶ。$\mathcal{J}^n$ の $S$ の開集合への
制限を $\mathcal{J}^n|_S$ と書く。開被覆は $\tau$-被覆なので、これは位相空間
$S$ 上の層である。複体
$$
0 \to \mathcal{F} \to \mathcal{J}^0|_S \to \mathcal{J}^1|_S \to \ldots
$$
を得る。これは、任意のアフィン開集合 $U \subset S$ 上の切断の複体が完全だから
完全である（上で見た $p > 0$ に対する $H^p(U, \mathcal{F}^a)$ の消滅による）。
したがって Derived Categories, Lemma \ref{derived-lemma-morphisms-lift} により、
複体の写像 $\mathcal{J}^\bullet|_S \to \mathcal{I}^\bullet$ が存在し、特に写像
$$
R\Gamma(\mathcal{C}, \mathcal{F}^a)
=
\Gamma(S, \mathcal{J}^\bullet)
\longrightarrow
\Gamma(S, \mathcal{I}^\bullet)
=
R\Gamma(S, \mathcal{F}).
$$
を誘導する。コホモロジーをとると写像
$H^n(\mathcal{C}, \mathcal{F}^a) \to H^n(S, \mathcal{F})$ を得る。これが同型で
あることを示さなければならない。$\mathcal{U} : S = \bigcup U_i$ をアフィン開被覆と
する。これを $\tau$-被覆とみなすこともできる。上の構成から二重複体の写像
$$
\check{\mathcal{C}}^\bullet(\mathcal{U}, \mathcal{J})
=
\check{\mathcal{C}}^\bullet(\mathcal{U}, \mathcal{J}|_S)
\longrightarrow
\check{\mathcal{C}}^\bullet(\mathcal{U}, \mathcal{I}).
$$
を得る。この写像はスペクトル系列の写像
$$
{}^\tau\! E_2^{p, q} = \check{H}^p(\mathcal{U}, \underline{H}^q(\mathcal{F}^a))
\longrightarrow
E_2^{p, q} = \check{H}^p(\mathcal{U}, \underline{H}^q(\mathcal{F}))
$$
を誘導する。第一のスペクトル系列は $H^{p + q}(\mathcal{C}, \mathcal{F})$ に、
第二のスペクトル系列は $H^{p + q}(S, \mathcal{F})$ に収束する。一方、誘導写像
${}^\tau\! E_2^{p, q} \to E_2^{p, q}$ は全単射であることを既に見た
（すべての交叉はアフィン内の開集合なので分離的である）。したがって写像
$H^n(\mathcal{C}, \mathcal{F}^a) \to H^n(S, \mathcal{F})$ も同型であり、
証明が完了する。
\end{proof}

\begin{proposition}
\label{descent-proposition-equivalence-quasi-coherent-functorial}
$f : T \to S$ をスキームの射とする。
\begin{enumerate}
\item Proposition \ref{descent-proposition-equivalence-quasi-coherent} の圏の同値は
引き戻しと両立する。より正確には、$S$ 上の任意の準連接層 $\mathcal{G}$ に対し
$f^*(\mathcal{G}^a) = (f^*\mathcal{G})^a$ である。
\item Proposition \ref{descent-proposition-equivalence-quasi-coherent} の (1) の圏の同値は、
一般には直像と {\bf 両立しない}。
\item $f$ が準コンパクトかつ準分離で、$\tau \in \{Zariski, \etale\}$ ならば、
$f_*$ と $f_{small, *}$ は準連接層を保ち、図式
$$
\xymatrix{
\QCoh(\mathcal{O}_T)
\ar[rr]_{f_*} \ar[d]_{\mathcal{F} \mapsto \mathcal{F}^a} & &
\QCoh(\mathcal{O}_S)
\ar[d]^{\mathcal{G} \mapsto \mathcal{G}^a} \\
\QCoh(T_\tau, \mathcal{O}) \ar[rr]^{f_{small, *}} & &
\QCoh(S_\tau, \mathcal{O})
}
$$
は可換である。すなわち
$f_{small, *}(\mathcal{F}^a) = (f_*\mathcal{F})^a$ である。
\end{enumerate}
\end{proposition}

\begin{proof}
(1) は Remark \ref{descent-remark-change-topologies-ringed-sites} の議論から従う。
(2) は単なる警告であり、次のように説明できる。まず $f_*$ は一般には準連接層を
準連接層へ写さないので、主張を正確に定式化できない。これが $f$（および $f$ の
任意の基底変換）について成り立つ場合でさえ、大サイト上での両立性は
$f_*\mathcal{F}$ の形成が任意の基底変換と可換であることを意味するが、これは一般に
成り立たない。明示的な例は準コンパクト開埋め込み
$j : X = \mathbf{A}^2_k \setminus \{0\} \to \mathbf{A}^2_k = Y$
である。ここで $k$ は体である。$j_*\mathcal{O}_X = \mathcal{O}_Y$ だが、$0$ 写像に
よって $\Spec(k)$ へ基底変換すると、直像が零になることが分かる。

\medskip\noindent
$\tau = \etale$ の場合に (3) を示す。$f$ および $f$ の任意の基底変換は準連接層を
準連接層へ写す。Schemes, Lemma
\ref{schemes-lemma-push-forward-quasi-coherent} を参照せよ。等式
$f_{small, *}(\mathcal{F}^a) = (f_*\mathcal{F})^a$ は、任意の \'etale 射
$g : U \to S$ に対し
$\Gamma(U, g^*f_*\mathcal{F}) = \Gamma(U \times_S T, (g')^*\mathcal{F})$
であることを意味する。ここで $g' : U \times_S T \to T$ は射影である。これは
Cohomology of Schemes, Lemma
\ref{coherent-lemma-flat-base-change-cohomology} により成り立つ。
\end{proof}

\begin{lemma}
\label{descent-lemma-higher-direct-images-small-etale}
$f : T \to S$ を準コンパクトかつ準分離なスキームの射とし、$\mathcal{F}$ を
$T$ 上の準連接層とする。\'etale 位相と Zariski 位相のいずれについても標準同型
$R^if_{small, *}(\mathcal{F}^a) = (R^if_*\mathcal{F})^a$.
がある。
\end{lemma}

\begin{proof}
\'etale 位相について示し、Zariski 位相の場合の証明は省略する。Cohomology of Schemes,
Lemma \ref{coherent-lemma-quasi-coherence-higher-direct-images} により層
$R^if_*\mathcal{F}$ は準連接なので、主張には意味がある。層
$R^if_{small, *}\mathcal{F}^a$ は前層
$$
U \longmapsto H^i(U \times_S T, \mathcal{F}^a)
$$
に付随する層である。ここで $g : U \to S$ は $S_\etale$ の対象である。Cohomology on
Sites, Lemma \ref{sites-cohomology-lemma-higher-direct-images} を参照せよ。
規約により右辺は、$\mathcal{F}^a$ を局所化 $T_\etale/U \times_S T$ へ制限したものの
\'etale コホモロジーであり、この局所化は $(U \times_S T)_\etale$ に等しい。
命題 \ref{descent-proposition-same-cohomology-quasi-coherent} により、この前層は前層
$$
U \longmapsto
H^i(U \times_S T, (g')^*\mathcal{F}),
$$
と同じである。ここで $g' : U \times_S T \to T$ は射影である。$U$ がアフィンならば、
これは $H^0(U, R^if'_*(g')^*\mathcal{F})$ と同じである。Cohomology of Schemes,
Lemma \ref{coherent-lemma-quasi-coherence-higher-direct-images-application} を参照せよ。
Cohomology of Schemes, Lemma \ref{coherent-lemma-flat-base-change-cohomology} により、
これは $H^0(U, g^*R^if_*\mathcal{F})$ に等しく、$(R^if_*\mathcal{F})^a$ の $U$ 上の
値である。したがって、加群の層 $R^if_{small, *}(\mathcal{F}^a)$ と
$(R^if_*\mathcal{F})^a$ の $S_\etale$ の各アフィン対象上の値は標準的に同型であり、
よって両者は標準的に同型である。
\end{proof}




\section{準連接層と位相、II}
\label{descent-section-quasi-coherent-sheaves-bis}

\noindent
位相の章で $S$ に付随させた各サイト上の準連接加群と、スキーム $S$ 上の
準連接加群との比較を引き続き論じる。


\begin{lemma}
\label{descent-lemma-compare-etale-zariski-flat}
Lemma \ref{descent-lemma-compare-sites} の環付きサイトの射
$\text{id}_{small, \etale, Zar} : S_\etale \to S_{Zar}$ は平坦である。
\end{lemma}

\begin{proof}
ここで $\epsilon = \text{id}_{small, \etale, Zar}$ と書き、
$S_\etale$ および $S_{Zar}$ 上の構造層をそれぞれ
$\mathcal{O}_\etale$ および $\mathcal{O}_{Zar}$ と書く。
$\mathcal{O}_\etale$ が平坦な
$\epsilon^{-1}\mathcal{O}_{Zar}$-加群であることを示せばよい。
\'etale 射は開写像であることを思い出そう。Morphisms, Lemma
\ref{morphisms-lemma-etale-open} を参照せよ。
したがって（層の引き戻しの構成から）、
$\epsilon^{-1}\mathcal{O}_{Zar}$ は $S_\etale$ 上の前層
$\mathcal{O}'$ の層化である。この前層は \'etale 射
$f : V \to S$ を $\mathcal{O}_S(f(V))$ へ写す。
$V$ と $U = f(V) \subset S$ がともにアフィンならば、
$V \to U$ はアフィン間の \'etale 射であり、したがって
\'etale 環準同型に対応する。\'etale 環準同型は平坦なので、
$\mathcal{O}_S(U) = \mathcal{O}'(V) \to
\mathcal{O}_\etale(V) = \mathcal{O}_V(V)$ は平坦である。
最後に、$S_\etale$ の任意の対象、すなわち任意の \'etale 射
$f : V \to S$ に対し、アフィン開被覆 $V = \bigcup V_i$ であって、
すべての $i$ について $f(V_i)$ が $S$ のアフィン開部分となるものが存在する
\footnote{実際、$y \in V$ に対し、$f(V')$ があるアフィン開部分
$U \subset S$ に含まれるようなアフィン開部分 $y \in V' \subset V$ を取る。
次にアフィン開部分 $f(y) \in U' \subset f(V')$ を取る。
すると $V'' = f^{-1}(U') \subset V'$ は $U' \times_U V'$ に等しいので
アフィンであり、$f(V'') = U'$ もアフィンである。}。
ゆえに主張は Modules on Sites, Lemma
\ref{sites-modules-lemma-flatness-sheafification-refined} から従う。
\end{proof}

\begin{lemma}
\label{descent-lemma-equivalence-quasi-coherent-limits}
スキーム $S$ を取り、
$\tau \in \{Zariski, \linebreak[0] \etale,
\linebreak[0] smooth, \linebreak[0] syntomic, \linebreak[0] fppf\}$.
とする。規則 $\mathcal{F} \mapsto \mathcal{F}^a$ によって定まり、
Proposition \ref{descent-proposition-equivalence-quasi-coherent} で見た関手
$$
\QCoh(\mathcal{O}_S)
\longrightarrow
\textit{Mod}((\Sch/S)_\tau, \mathcal{O})
\quad\text{および}\quad
\QCoh(\mathcal{O}_S)
\longrightarrow
\textit{Mod}(S_\tau, \mathcal{O})
$$
は次の性質をもつ。
\begin{enumerate}
\item 充満忠実である。
\item 直和と可換である。
\item 余極限と可換である。
\item 右完全である。
\item 関手
$\QCoh(\mathcal{O}_S) \to \textit{Mod}(S_\etale, \mathcal{O})$
として完全である。
\item 関手
$\QCoh(\mathcal{O}_S) \to
\textit{Mod}((\Sch/S)_\tau, \mathcal{O})$
としては一般に {\bf 完全でない}。
\item 二つの準連接 $\mathcal{O}_S$-加群
$\mathcal{F}$, $\mathcal{G}$ に対し
$(\mathcal{F} \otimes_{\mathcal{O}_S} \mathcal{G})^a =
\mathcal{F}^a \otimes_\mathcal{O} \mathcal{G}^a$ が成り立つ。
\item $\tau = \etale$ または $\tau = Zariski$ とする。
二つの準連接 $\mathcal{O}_S$-加群 $\mathcal{F}$, $\mathcal{G}$ について
$\mathcal{F}$ が有限表示ならば
$(\SheafHom_{\mathcal{O}_S}(\mathcal{F}, \mathcal{G}))^a =
\SheafHom_\mathcal{O}(\mathcal{F}^a, \mathcal{G}^a)$ が
$\textit{Mod}(S_\tau, \mathcal{O})$ で成り立つ。
\item 二つの準連接 $\mathcal{O}_S$-加群
$\mathcal{F}$, $\mathcal{G}$ に対しては、$\mathcal{F}$ が有限表示であっても
$(\SheafHom_{\mathcal{O}_S}(\mathcal{F}, \mathcal{G}))^a =
\SheafHom_\mathcal{O}(\mathcal{F}^a, \mathcal{G}^a)$
は $\textit{Mod}((\Sch/S)_\tau, \mathcal{O})$ で一般には
{\bf 成り立たない}。
\item $\mathcal{O}$-加群の短完全列
$0 \to \mathcal{F}_1^a \to \mathcal{E} \to \mathcal{F}_2^a \to 0$
が与えられれば、$\mathcal{E}$ は準連接である
\footnote{注意：これは誤解を招きうる。(6) を参照せよ。}。すなわち
$\mathcal{E}$ はこの関手の本質的像に属する。
\end{enumerate}
\end{lemma}

\begin{proof}
(1) は Proposition \ref{descent-proposition-equivalence-quasi-coherent} で示した。

\medskip\noindent
Schemes, Section \ref{schemes-section-quasi-coherent} で、
スキーム上の準連接層の余極限は準連接層であることを見た。
さらに Remark \ref{descent-remark-change-topologies-ringed-sites} で、
$\mathcal{F} \mapsto \mathcal{F}^a$ は環付きサイトの射に対する
引き戻し関手であり、したがってすべての余極限と可換であることを見た。
Modules on Sites, Lemma
\ref{sites-modules-lemma-exactness-pushforward-pullback} を参照せよ。
よって (3) と、その特別な場合である (2) が成り立つ。

\medskip\noindent
これはまた、この関手が右完全（すなわち有限余極限と可換）であることも示すので、
(4) が従う。

\medskip\noindent
関手 $\QCoh(\mathcal{O}_S) \to
\textit{Mod}(S_\etale, \mathcal{O})$,
$\mathcal{F} \mapsto \mathcal{F}^a$
は左完全である。実際、\'etale 射は平坦である。Morphisms, Lemma
\ref{morphisms-lemma-etale-flat} を参照せよ。これで (5) が示された。

\medskip\noindent
(6) を示すため $S = \Spec(\mathbf{Z})$ とする。このとき
$2 : \mathcal{O}_S \to \mathcal{O}_S$ は単射だが、付随する
$(\Sch/S)_\tau$ 上の $\mathcal{O}$-加群の写像は単射でない。
実際、$2 : \mathbf{F}_2 \to \mathbf{F}_2$ は単射でなく、
$\Spec(\mathbf{F}_2)$ は $(\Sch/S)_\tau$ の対象である。

\medskip\noindent
(7) が成り立つのは、上で述べたように、関手
$\mathcal{F} \mapsto \mathcal{F}^a$ が環付きサイトの射に対する
引き戻し関手であり、そのような関手はテンソル積と可換だからである。
Modules on Sites, Lemma
\ref{sites-modules-lemma-tensor-product-pullback} を参照せよ。

\medskip\noindent
(8) は $\tau = Zariski$ なら明らかである。というのも、
$S_{Zar}$ 上の $\mathcal{O}$-加群の圏は、位相空間 $S$ 上の
$\mathcal{O}_S$-加群の圏と同じだからである。
$\tau = \etale$ ならば、上で述べたように関手
$\mathcal{F} \mapsto \mathcal{F}^a$ は環付きサイトの平坦射
$(S_\etale, \mathcal{O}) \to (S_{Zar}, \mathcal{O}_S)$
に対する引き戻し関手なので、(8) が成り立つ。Lemma
\ref{descent-lemma-compare-etale-zariski-flat} を参照せよ。
環付きサイトの平坦射による引き戻しは、有限表示加群を始域とする内部 Hom と可換する。
Modules on Sites, Lemma \ref{sites-modules-lemma-pullback-internal-hom}
を参照せよ。

\medskip\noindent
(9) を示すため $S = \Spec(\mathbf{Z})$ とし、
$\mathcal{F} = \Coker(2 : \mathcal{O}_S \to \mathcal{O}_S)$
$\mathcal{G} = \mathcal{O}_S$ と置く。
このとき $\mathcal{F}^a = \Coker(2 : \mathcal{O} \to \mathcal{O})$,
$\mathcal{G}^a = \mathcal{O}$ である。したがって
$\SheafHom_\mathcal{O}(\mathcal{F}^a, \mathcal{G}^a) = \mathcal{O}[2]$
は $\mathcal{O}$ の $2$-捩れに等しく、零でない。(6) の証明を参照せよ。
一方、加群
$\SheafHom_{\mathcal{O}_S}(\mathcal{F}, \mathcal{G})$
は零である。

\medskip\noindent
(10) を示す。
$0 \to \mathcal{F}_1^a \to \mathcal{E} \to \mathcal{F}_2^a \to 0$
を $\mathcal{O}$-加群の短完全列とし、$\mathcal{F}_1$ と
$\mathcal{F}_2$ は $S$ 上準連接であるとする。$S_{Zar}$ への制限
$$
0 \to \mathcal{F}_1 \to \mathcal{E}|_{S_{Zar}} \to \mathcal{F}_2
$$
を考える。Proposition \ref{descent-proposition-same-cohomology-quasi-coherent}
により、任意のアフィン $U \subset S$ について
$H^1(U, \mathcal{F}_1^a) = H^1(U, \mathcal{F}_1) = 0$.
である。したがって上の列は右でも完全である。Schemes, Section
\ref{schemes-section-quasi-coherent} により、
$\mathcal{F} = \mathcal{E}|_{S_{Zar}}$ は準連接である。
ゆえに可換図式
$$
\xymatrix{
& \mathcal{F}_1^a \ar[r] \ar[d] &
\mathcal{F}^a \ar[r] \ar[d] &
\mathcal{F}_2^a \ar[r] \ar[d] & 0 \\
0 \ar[r] &
\mathcal{F}_1^a \ar[r] &
\mathcal{E} \ar[r] &
\mathcal{F}_2^a \ar[r] & 0
}
$$
を得る。証明を終えるには、上段も右完全であることを示せば十分である。
そのため、再び $S$ のアフィン開部分を
$U = \Spec(A) \subset S$ と書く。上で
$0 \to \mathcal{F}_1(U) \to \mathcal{E}(U) \to \mathcal{F}_2(U) \to 0$
が完全であることを見た。任意のアフィン・スキーム $V/U$,
$V = \Spec(B)$ に対し、写像
$\mathcal{F}_1^a(V) \to \mathcal{E}(V)$ は単射である。
定義により $\mathcal{F}_1^a(V) = \mathcal{F}_1(U) \otimes_A B$ である。
単射 $\mathcal{F}_1^a(V) \to \mathcal{E}(V)$ は
$$
\mathcal{F}_1(U) \otimes_A B \to
\mathcal{E}(U) \otimes_A B \to \mathcal{E}(V)
$$
と分解する。$B = A \oplus M$ という形の $A$-代数 $B$ を考えると、
$\mathcal{F}_1(U) \to \mathcal{E}(U)$ は普遍単射であることが分かる
（Algebra, Definition \ref{algebra-definition-universally-injective} を参照せよ）。
$\mathcal{E}(U) = \mathcal{F}(U)$ なので、
$\mathcal{F}_1 \to \mathcal{F}$ は任意の基底変換後も単射のままである。
同値なこととして、$\mathcal{F}_1^a \to \mathcal{F}^a$ は単射である。
\end{proof}

\begin{lemma}
\label{descent-lemma-properties-quasi-coherent}
スキーム $S$ を取る。$S_\etale$ 上の準連接加群の圏
$\QCoh(S_\etale, \mathcal{O})$ は次の性質をもつ。
\begin{enumerate}
\item 準連接層の任意の直和は準連接である。
\item 準連接層の任意の余極限は準連接である。
\item 準連接層の射の核と余核は準連接である。
\item $\mathcal{O}$-加群の短完全列
$0 \to \mathcal{F}_1 \to \mathcal{F}_2 \to \mathcal{F}_3 \to 0$
において、三項のうち二項が準連接ならば残りの一項も準連接である。
\item 二つの準連接 $\mathcal{O}$-加群のテンソル積は準連接である。
\item 二つの準連接 $\mathcal{O}$-加群
$\mathcal{F}$, $\mathcal{G}$ について $\mathcal{F}$ が有限表示ならば、
内部 Hom
$\SheafHom_\mathcal{O}(\mathcal{F}, \mathcal{G})$
は準連接である。
\end{enumerate}
\end{lemma}

\begin{proof}
対応する事実はスキーム $S$ 上の準連接加群について成り立つ。
Schemes, Section \ref{schemes-section-quasi-coherent} を参照せよ。
Lemma \ref{descent-lemma-equivalence-quasi-coherent-limits} を用いて
これらの事実を $S_\etale$ へ移すことにより証明する。

\medskip\noindent
(1) を示す。$\mathcal{F}_i$, $i \in I$ を
$\QCoh(S_\etale, \mathcal{O})$ の対象の族とする。
ある $S$ 上の準連接加群 $\mathcal{G}_i$ を用いて
$\mathcal{F}_i = \mathcal{G}_i^a$ と書く。上で引用した補題により
$\bigoplus \mathcal{F}_i = (\bigoplus \mathcal{G}_i)^a$ なので結論を得る。

\medskip\noindent
(2) を示す。$\mathcal{I} \to \QCoh(S_\etale, \mathcal{O})$,
$i \mapsto \mathcal{F}_i$ を図式とする。
$\mathcal{F}_i = \mathcal{G}_i^a$ と書けば、図式
$\mathcal{I} \to \QCoh(\mathcal{O}_S)$.
を得る。上で引用した補題により
$\colim \mathcal{F}_i = (\colim \mathcal{G}_i)^a$ なので結論を得る。

\medskip\noindent
(3) を示す。$a : \mathcal{F} \to \mathcal{F}'$ を
$\QCoh(S_\etale, \mathcal{O})$ の射とする。
ある $S$ 上の準連接加群の写像
$b : \mathcal{G} \to \mathcal{G}'$ によって $a = b^a$ と書く。
上で引用した補題により
$\Ker(a) = \Ker(b)^a$ および $\Coker(a) = \Coker(b)^a$ なので結論を得る。

\medskip\noindent
(4) を示す。$\mathcal{F}_1$ と $\mathcal{F}_3$ が準連接であることだけが
分かっている場合を除けば、これは (3) から従う。この残る場合には、
$S$ 上準連接な $\mathcal{G}_i$ を用いて
$\mathcal{F}_1 = \mathcal{G}_1^a$ および
$\mathcal{F}_3 = \mathcal{G}_3^a$ と書く。
Lemma \ref{descent-lemma-equivalence-quasi-coherent-limits} の (10) から結論を得る。

\medskip\noindent
(5) を示す。$\mathcal{F}$ と $\mathcal{F}'$ を
$\QCoh(S_\etale, \mathcal{O})$ の対象とする。
$S$ 上準連接な $\mathcal{G}$ と $\mathcal{G}'$ を用いて
$\mathcal{F} = \mathcal{G}^a$ および
$\mathcal{F}' = (\mathcal{G}')^a$ と書く。
上で引用した補題により
$\mathcal{F} \otimes_\mathcal{O} \mathcal{F}' =
(\mathcal{G} \otimes_{\mathcal{O}_S} \mathcal{G}')^a$
なので結論を得る。

\medskip\noindent
(6) を示す。$\mathcal{F}$ と $\mathcal{G}$ を
$\QCoh(S_\etale, \mathcal{O})$ の対象とし、
$\mathcal{F}$ は有限表示であるとする。
$S$ 上準連接な $\mathcal{H}$ と $\mathcal{I}$ を用いて
$\mathcal{F} = \mathcal{H}^a$ および
$\mathcal{G} = (\mathcal{I})^a$ と書く。
Lemma \ref{descent-lemma-equivalence-quasi-coherent-properties} により
$\mathcal{H}$ は有限表示である。Lemma
\ref{descent-lemma-equivalence-quasi-coherent-limits} の (8) により
$\SheafHom_\mathcal{O}(\mathcal{F}, \mathcal{G}) =
(\SheafHom_{\mathcal{O}_S}(\mathcal{H}, \mathcal{I}))^a$
なので結論を得る。
\end{proof}

\begin{lemma}
\label{descent-lemma-properties-quasi-coherent-on-big}
スキーム $S$ を取り、
$\tau \in \{Zariski, \linebreak[0] \etale, \linebreak[0]
smooth, \linebreak[0] syntomic, \linebreak[0] fppf\}$.
とする。$(\Sch/S)_\tau$ 上の準連接加群の圏
$\QCoh((\Sch/S)_\tau, \mathcal{O})$ は次の性質をもつ。
\begin{enumerate}
\item 準連接層の任意の直和は準連接である。
\item 準連接層の任意の余極限は準連接である。
\item 準連接層の射の余核は準連接である。
\item $\mathcal{O}$-加群の短完全列
$0 \to \mathcal{F}_1 \to \mathcal{F}_2 \to \mathcal{F}_3 \to 0$
において $\mathcal{F}_1$ と $\mathcal{F}_3$ が準連接ならば
$\mathcal{F}_2$ も準連接である。
\item 二つの準連接 $\mathcal{O}$-加群のテンソル積は準連接である。
\item 二つの準連接 $\mathcal{O}$-加群
$\mathcal{F}$, $\mathcal{G}$ について $\mathcal{F}$ が有限局所自由ならば、
内部 Hom
$\SheafHom_\mathcal{O}(\mathcal{F}, \mathcal{G})$
は準連接である。
\end{enumerate}
\end{lemma}

\begin{proof}
対応する事実はスキーム $S$ 上の準連接加群について成り立つ。
Schemes, Section \ref{schemes-section-quasi-coherent} を参照せよ。
Lemma \ref{descent-lemma-equivalence-quasi-coherent-limits} を用いて
これらの事実を $(\Sch/S)_\tau$ へ移すことにより証明する。

\medskip\noindent
(1) を示す。$\mathcal{F}_i$, $i \in I$ を
$\QCoh((\Sch/S)_\tau, \mathcal{O})$ の対象の族とする。
ある $S$ 上の準連接加群 $\mathcal{G}_i$ を用いて
$\mathcal{F}_i = \mathcal{G}_i^a$ と書く。上で引用した補題により
$\bigoplus \mathcal{F}_i = (\bigoplus \mathcal{G}_i)^a$ なので結論を得る。

\medskip\noindent
(2) を示す。$\mathcal{I} \to \QCoh((\Sch/S)_\tau, \mathcal{O})$,
$i \mapsto \mathcal{F}_i$ を図式とする。
$\mathcal{F}_i = \mathcal{G}_i^a$ と書けば、図式
$\mathcal{I} \to \QCoh(\mathcal{O}_S)$.
を得る。上で引用した補題により
$\colim \mathcal{F}_i = (\colim \mathcal{G}_i)^a$ なので結論を得る。

\medskip\noindent
(3) を示す。$a : \mathcal{F} \to \mathcal{F}'$ を
$\QCoh((\Sch/S)_\tau, \mathcal{O})$ の射とする。
ある $S$ 上の準連接加群の写像
$b : \mathcal{G} \to \mathcal{G}'$ によって $a = b^a$ と書く。
上で引用した補題により
$\Coker(a) = \Coker(b)^a$（余核は余極限だから）なので結論を得る。

\medskip\noindent
(4) を示す。$S$ 上準連接な $\mathcal{G}_i$ を用いて
$\mathcal{F}_1 = \mathcal{G}_1^a$ および
$\mathcal{F}_3 = \mathcal{G}_3^a$ と書く。
Lemma \ref{descent-lemma-equivalence-quasi-coherent-limits} の (10) から結論を得る。

\medskip\noindent
(5) を示す。$\mathcal{F}$ と $\mathcal{F}'$ を
$\QCoh((\Sch/S)_\tau, \mathcal{O})$ の対象とする。
$S$ 上準連接な $\mathcal{G}$ と $\mathcal{G}'$ を用いて
$\mathcal{F} = \mathcal{G}^a$ および
$\mathcal{F}' = (\mathcal{G}')^a$ と書く。
上で引用した補題により
$\mathcal{F} \otimes_\mathcal{O} \mathcal{F}' =
(\mathcal{G} \otimes_{\mathcal{O}_S} \mathcal{G}')^a$
なので結論を得る。

\medskip\noindent
(6) を示す。ある準連接 $\mathcal{O}_S$-加群を用いて
$\mathcal{F} = \mathcal{H}^a$ と書く。Lemma
\ref{descent-lemma-equivalence-quasi-coherent-properties} により
$\mathcal{H}$ は有限局所自由である。この問題は $S$ 上 Zariski 局所的
（詳細は省略する）なので、$\mathcal{H} = \mathcal{O}_S^{\oplus n}$ が
有限自由であると仮定してよい。このとき $\mathcal{F} = \mathcal{O}^{\oplus n}$ であり、
$\SheafHom_\mathcal{O}(\mathcal{F}, \mathcal{G}) = \mathcal{G}^{\oplus n}$
は準連接である。
\end{proof}

\begin{example}
\label{descent-example-internal-hom-not-qcoh}
スキーム $S$ を取る。Lemma \ref{descent-lemma-properties-quasi-coherent-on-big}
で考えた位相 $\tau$ の一つについて、$\mathcal{F}$ と $\mathcal{G}$ を
$(\Sch/S)_\tau$ 上の準連接加群とする。一般には
$\SheafHom_\mathcal{O}(\mathcal{F}, \mathcal{G})$
は、$\mathcal{F}$ が有限表示であっても準連接とは限らない。
実際、$S = \Spec(\mathbf{Z})$,
$\mathcal{F} = \Coker(2 : \mathcal{O} \to \mathcal{O})$,
$\mathcal{G} = \mathcal{O}$ とする。このとき
$\SheafHom_\mathcal{O}(\mathcal{F}, \mathcal{G}) = \mathcal{O}[2]$
は $\mathcal{O}$ の $2$-捩れに等しく、これは準連接でない。
\end{example}

\begin{lemma}
\label{descent-lemma-qc-colimits}
スキーム $S$ を取る。
\begin{enumerate}
\item 圏 $\QCoh((\Sch/S)_{fppf}, \mathcal{O})$ は余極限をもち、
それらは各圏
$\textit{Mod}((\Sch/S)_{Zar}, \mathcal{O})$,
$\textit{Mod}((\Sch/S)_\etale, \mathcal{O})$, および
$\textit{Mod}((\Sch/S)_{fppf}, \mathcal{O})$ における余極限と一致する。
\item $\mathcal{F}, \mathcal{G}$ を
$\QCoh((\Sch/S)_{fppf}, \mathcal{O})$ の対象とする。\newline
各圏 $\textit{Mod}((\Sch/S)_{Zar}, \mathcal{O})$,
$\textit{Mod}((\Sch/S)_\etale, \mathcal{O})$, または
$\textit{Mod}((\Sch/S)_{fppf}, \mathcal{O})$ で計算したテンソル積
$\mathcal{F} \otimes_\mathcal{O} \mathcal{G}$ は一致し、その共通の値は
$\QCoh((\Sch/S)_{fppf}, \mathcal{O})$ の対象である。
\item $\mathcal{F}, \mathcal{G}$ を
$\QCoh((\Sch/S)_{fppf}, \mathcal{O})$ の対象とし、$\mathcal{F}$ は
有限局所自由（fppf 位相に関して、または同値なこととして \'etale 位相に関して、
または同値なこととして Zariski 位相に関して）とする。各圏
$\textit{Mod}((\Sch/S)_{Zar}, \mathcal{O})$,
$\textit{Mod}((\Sch/S)_\etale, \mathcal{O})$, または
$\textit{Mod}((\Sch/S)_{fppf}, \mathcal{O})$ で計算した内部 Hom
$\SheafHom_\mathcal{O}(\mathcal{F}, \mathcal{G})$
は一致し、その共通の値は
$\QCoh((\Sch/S)_{fppf}, \mathcal{O})$ の対象である。
\end{enumerate}
\end{lemma}

\begin{proof}
この補題は上で示した結果を少し異なる形にまとめたものである。
まず Lemma \ref{descent-lemma-properties-quasi-coherent-on-big} により、
(1), (2), (3) の各構成の出力が
$\QCoh((\Sch/S)_\tau, \mathcal{O})$ に属することは既に分かっている。
残るのは、各場合に結果が用いる位相に依存しないことを示すことである。
鍵となるのは、Proposition \ref{descent-proposition-equivalence-quasi-coherent} の同値
$\QCoh(\mathcal{O}_S) \to \QCoh((\Sch/S)_\tau, \mathcal{O})$,
$\mathcal{F} \mapsto \mathcal{F}^a$
が、位相 $\tau \in \{Zariski, \etale, fppf\}$ の選択によらず
同じ公式で与えられることである。

\medskip\noindent
(1) を示す。$\mathcal{I} \to \QCoh((\Sch/S)_{fppf}, \mathcal{O})$,
$i \mapsto \mathcal{F}_i$ を図式とする。
$\mathcal{F}_i = \mathcal{G}_i^a$ と書けば、図式
$\mathcal{I} \to \QCoh(\mathcal{O}_S)$.
を得る。このとき
$\colim \mathcal{F}_i = (\colim \mathcal{G}_i)^a$ は
$\textit{Mod}((\Sch/S)_\tau, \mathcal{O})$ において
$\tau \in \{Zariski, \etale, fppf\}$
の各場合に成り立つ。これは Lemma
\ref{descent-lemma-equivalence-quasi-coherent-limits} による。
これで (1) が示された。

\medskip\noindent
(2) を示す。$S$ 上準連接な $\mathcal{H}$ と $\mathcal{I}$ を用いて
$\mathcal{F} = \mathcal{H}^a$ および
$\mathcal{G} = (\mathcal{I})^a$ と書く。このとき
$\mathcal{F} \otimes_\mathcal{O} \mathcal{G} =
(\mathcal{H} \otimes_\mathcal{O} \mathcal{I})^a$ は
$\textit{Mod}((\Sch/S)_\tau, \mathcal{O})$ において
$\tau \in \{Zariski, \etale, fppf\}$
の各場合に成り立つ。これは Lemma
\ref{descent-lemma-equivalence-quasi-coherent-limits} による。
これで (2) が示された。

\medskip\noindent
(3) を示す。$\mathcal{F}$ と $\mathcal{G}$ を
$\QCoh((\Sch/S)_{fppf}, \mathcal{O})$ の対象とする。
$S$ 上準連接な $\mathcal{H}$ を用いて
$\mathcal{F} = \mathcal{H}^a$ と書く。Lemma
\ref{descent-lemma-equivalence-quasi-coherent-properties} により
\begin{align*}
\mathcal{F}\text{ は fppf 位相に関して有限局所自由}
& \Leftrightarrow
\mathcal{H}\text{ は }S\text{ 上有限局所自由} \\
& \Leftrightarrow
\mathcal{F}\text{ は \'etale 位相に関して有限局所自由} \\
& \Leftrightarrow
\mathcal{H}\text{ は }S\text{ 上有限局所自由} \\
& \Leftrightarrow
\mathcal{F}\text{ は Zariski 位相に関して有限局所自由}
\end{align*}
を得る。これが (3) の括弧内の主張を説明する。
さて、これらの同値な条件が成り立つならば、$\mathcal{H}$ は有限局所自由である。
$\SheafHom_\mathcal{O}(\mathcal{F}, \mathcal{G})$
の構成は Modules on Sites, Section
\ref{sites-modules-section-internal-hom} において、
加群の前層としての $\mathcal{F}$ と $\mathcal{G}$ のみに依存する
（出力 $\SheafHom$ が層であるかどうかだけが、
$\mathcal{F}$ と $\mathcal{G}$ が層であるかどうかに依存する）。
\end{proof}










\section{準連接加群とアフィン}
\label{descent-section-alternative-quasi-coherent}

\noindent
スキーム $S$ を取る\footnote{この節では Topologies, Section
\ref{topologies-section-change-topologies} と同様に、
$\tau \in \{Zariski, \etale, smooth, syntomic, fppf\}$ の各場合について
サイト $(\Sch/S)_\tau$ の台となる圏を同一に選ぶ。このときサイト
$(\textit{Aff}/S)_\tau$ の台となる圏も同一である。}。
$\tau \in \{Zariski, \etale, smooth, syntomic, fppf\}$ とする。
$(\textit{Aff}/S)_\tau$ は $(\Sch/S)_\tau$ の、対象をアフィン・スキームとする
充満部分圏に、被覆を標準 $\tau$-被覆と宣言してサイトの構造を入れたものであった。
Topologies, Lemmas
\ref{topologies-lemma-affine-big-site-Zariski},
\ref{topologies-lemma-affine-big-site-etale},
\ref{topologies-lemma-affine-big-site-smooth},
\ref{topologies-lemma-affine-big-site-syntomic}, および
\ref{topologies-lemma-affine-big-site-fppf}
により、トポスの同値
$g : \Sh((\textit{Aff}/S)_\tau) \to \Sh((\Sch/S)_\tau)$
を得る。その引き戻し関手は制限によって与えられる。
$\mathcal{O}$ が $(\Sch/S)_\tau$ 上の構造層を表すことを思い出し、
一時的に厳密さを期して、$\mathcal{O}$ の $(\textit{Aff}/S)_\tau$ への制限を
$\mathcal{O}_{\textit{Aff}}$ と書く。このとき環付きトポスの同値
\begin{equation}
\label{descent-equation-alternative-ringed}
(\Sh((\textit{Aff}/S)_\tau), \mathcal{O}_{\textit{Aff}})
\longrightarrow
(\Sh((\Sch/S)_\tau), \mathcal{O})
\end{equation}
を得る。これにより準連接加群も比較できる。

\begin{lemma}
\label{descent-lemma-quasi-coherent-alternative}
スキーム $S$ を取り、$\mathcal{F}$ を
$(\textit{Aff}/S)_{fppf}$ 上の $\mathcal{O}_{\textit{Aff}}$-加群の前層とする。
次の条件は同値である。
\begin{enumerate}
\item $(\textit{Aff}/S)_{fppf}$ の任意の射 $U \to U'$ に対し、写像
$\mathcal{F}(U') \otimes_{\mathcal{O}(U')} \mathcal{O}(U) \to \mathcal{F}(U)$
は同型である。
\item $\mathcal{F}$ は $(\textit{Aff}/S)_{Zar}$ 上の層であり、かつ
Modules on Sites, Definition \ref{sites-modules-definition-site-local}
の意味で環付きサイト
$((\textit{Aff}/S)_{Zar}, \mathcal{O}_{\textit{Aff}})$ 上の準連接加群である。
\item \'etale 位相について (2) と同じ条件が成り立つ。
\item smooth 位相について (2) と同じ条件が成り立つ。
\item syntomic 位相について (2) と同じ条件が成り立つ。
\item fppf 位相について (2) と同じ条件が成り立つ。
\item 同値 (\ref{descent-equation-alternative-ringed}) を介して、$\mathcal{F}$ は
$(\Sch/S)_{Zar}$,
$(\Sch/S)_\etale$,
$(\Sch/S)_{smooth}$,
$(\Sch/S)_{syntomic}$, または
$(\Sch/S)_{fppf}$
上の準連接加群に対応する。
\item $\mathcal{F}$ は、Section \ref{descent-section-quasi-coherent-sheaves}
で述べた手続きによって一意な準連接 $\mathcal{O}_S$-加群
$\mathcal{G}$ から得られる。
\end{enumerate}
\end{lemma}

\begin{proof}
準連接加群の概念は内在的なので
（Modules on Sites, Lemma \ref{sites-modules-lemma-special-locally-free}）、
同値 (\ref{descent-equation-alternative-ringed}) は準連接加群の圏の間の同値を誘導する。
Proposition \ref{descent-proposition-equivalence-quasi-coherent} は、
$\Sch/S$ 上の準連接加群を調べるために用いる位相は関係しないことを述べ、
さらに (8) が (7) と同じであることも示している。
したがって (2) -- (8) はすべて同値である。

\medskip\noindent
同値な条件 (2) -- (8) が成り立つと仮定し、
$\mathcal{G}$ を (8) のものとする。$h : U \to U' \to S$ を
$\textit{Aff}/S$ の射とする。構造射を
$f : U \to S$ および $f' : U' \to S$ と書けば、
$f = f' \circ h$ である。また
$\mathcal{F}(U') = \Gamma(U', (f')^*\mathcal{G})$ および
$\mathcal{F}(U) = \Gamma(U, f^*\mathcal{G}) = \Gamma(U, h^*(f')^*\mathcal{G})$.
である。したがって Schemes, Lemma
\ref{schemes-lemma-widetilde-pullback} により (1) が成り立つ。

\medskip\noindent
(1) が成り立つと仮定する。証明を終えるには (2) を示せば十分である。
$U$ を $(\textit{Aff}/S)_{Zar}$ の対象とし、
$U = \Spec(R)$ とする。標準開被覆 $U = U_1 \cup \ldots \cup U_n$ は、
$R$ の単位イデアルを生成する元 $f_1, \ldots, f_n \in R$ によって
$U_i = D(f_i)$ として与えられる。性質 (1) により
$$
\mathcal{F}(U_i) =
\mathcal{F}(U) \otimes_R R_{f_i} =
\mathcal{F}(U)_{f_i}
$$
および
$$
\mathcal{F}(U_i \cap U_j) =
\mathcal{F}(U) \otimes_R R_{f_if_j} =
\mathcal{F}(U)_{f_if_j}
$$
を得る。したがって Algebra, Lemma \ref{algebra-lemma-cover-module}
から $\mathcal{F}$ は $(\textit{Aff}/S)_{Zar}$ 上の層である。
自由 $R$-加群による表示
$$
\bigoplus\nolimits_{k \in K} R
\longrightarrow
\bigoplus\nolimits_{l \in L} R
\longrightarrow
\mathcal{F}(U)
\longrightarrow 0
$$
を選ぶ。性質 (1) とテンソル積の右完全性により、
$(\textit{Aff}/S)_{Zar}$ の任意の射 $U' \to U$ に対して表示
$$
\bigoplus\nolimits_{k \in K} \mathcal{O}_{Aff}(U')
\longrightarrow
\bigoplus\nolimits_{l \in L} \mathcal{O}_{Aff}(U')
\longrightarrow
\mathcal{F}(U')
\longrightarrow 0
$$
を得る。言い換えれば、$\mathcal{F}$ の局所化された圏
$(\textit{Aff}/S)_{Zar}/U$ への制限は表示
$$
\bigoplus\nolimits_{k \in K} \mathcal{O}_{Aff}|_{(\textit{Aff}/S)_{Zar}/U}
\longrightarrow
\bigoplus\nolimits_{l \in L} \mathcal{O}_{Aff}|_{(\textit{Aff}/S)_{Zar}/U}
\longrightarrow
\mathcal{F}|_{(\textit{Aff}/S)_{Zar}/U}
\longrightarrow 0
$$
をもつ。この煩雑な記法をお詫びしつつ、これで証明を終える。
\end{proof}

\noindent
この節の導入で始めた議論を続ける。
$\tau \in \{Zariski, \etale\}$ とする。$S_{affine, \tau}$ は
$S_\tau$ の、対象をアフィン・スキームとする充満部分圏に、
被覆を標準 $\tau$-被覆と宣言してサイトの構造を入れたものであった。
Topologies, Definitions
\ref{topologies-definition-big-small-Zariski} および
\ref{topologies-definition-big-small-etale} を参照せよ。
Topologies, Lemmas \ref{topologies-lemma-alternative-zariski} および
\ref{topologies-lemma-alternative} により、それぞれトポスの同値
$g : \Sh(S_{affine, \tau}) \to \Sh(S_\tau)$ を得る。
その引き戻し関手は制限によって与えられる。
$\mathcal{O}$ が $S_\tau$ 上の構造層を表すことを思い出し、
一時的に厳密さを期して、$\mathcal{O}$ の $S_{affine, \tau}$ への制限を
$\mathcal{O}_{affine}$ と書く。このとき環付きトポスの同値
\begin{equation}
\label{descent-equation-alternative-small-ringed}
(\Sh(S_{affine, \tau}), \mathcal{O}_{affine})
\longrightarrow
(\Sh(S_\tau), \mathcal{O})
\end{equation}
を得る。これにより準連接加群も比較できる。

\begin{lemma}
\label{descent-lemma-quasi-coherent-alternative-small}
スキーム $S$ を取り、$\tau \in \{Zariski, \etale\}$ とする。
$\mathcal{F}$ を $S_{affine, \tau}$ 上の
$\mathcal{O}_{affine}$-加群の前層とする。次の条件は同値である。
\begin{enumerate}
\item $S_{affine, \tau}$ の任意の射 $U \to U'$ に対し、写像
$\mathcal{F}(U') \otimes_{\mathcal{O}(U')} \mathcal{O}(U) \to \mathcal{F}(U)$
は同型である。
\item $\mathcal{F}$ は $S_{affine, \tau}$ 上の層であり、かつ
Modules on Sites, Definition \ref{sites-modules-definition-site-local}
の意味で環付きサイト
$(S_{affine, \tau}, \mathcal{O}_{affine})$ 上の準連接加群である。
\item 同値 (\ref{descent-equation-alternative-small-ringed}) を介して、
$\mathcal{F}$ は $S_\tau$ 上の準連接加群に対応する。
\item $\mathcal{F}$ は、Section \ref{descent-section-quasi-coherent-sheaves}
で述べた手続きによって一意な準連接 $\mathcal{O}_S$-加群
$\mathcal{G}$ から得られる。
\end{enumerate}
\end{lemma}

\begin{proof}
\'etale 位相の場合に証明する。

\medskip\noindent
(1) が成り立つと仮定する。$\mathcal{F}$ が層であることを示すため、
$\mathcal{U} = \{U_i \to U\}_{i = 1, \ldots, n}$ を
$S_{affine, \etale}$ の被覆とする。$\mathcal{F}$ に関する仮定により、
$\mathcal{F}$ と $\mathcal{U}$ に対する層条件は
$$
0 \to \mathcal{F}(U) \to
\prod \mathcal{F}(U) \otimes_{\mathcal{O}(U)} \mathcal{O}(U_i) \to
\prod \mathcal{F}(U) \otimes_{\mathcal{O}(U)} \mathcal{O}(U_i \times_U U_j)
$$
が完全であることに帰着する。これは
$\mathcal{O}(U) \to \prod \mathcal{O}(U_i)$ が忠実平坦だから成り立つ
（Lemma \ref{descent-lemma-standard-covering-Cech} と、
$S_{affine, \etale}$ の被覆が標準 \'etale 被覆であることによる）。
実際、Lemma \ref{descent-lemma-ff-exact} を適用できる。
次に $\mathcal{F}$ が $S_{affine, \etale}$ 上準連接であることを示す。
$S_{affine, \etale}$ の対象 $U$ に対し $R = \mathcal{O}(U)$ と置き、
自由 $R$-加群による表示
$$
\bigoplus\nolimits_{k \in K} R
\longrightarrow
\bigoplus\nolimits_{l \in L} R
\longrightarrow
\mathcal{F}(U)
\longrightarrow 0
$$
を選ぶ。性質 (1) とテンソル積の右完全性により、
$S_{affine, \etale}$ の任意の射 $U' \to U$ に対して表示
$$
\bigoplus\nolimits_{k \in K} \mathcal{O}(U')
\longrightarrow
\bigoplus\nolimits_{l \in L} \mathcal{O}(U')
\longrightarrow
\mathcal{F}(U')
\longrightarrow 0
$$
を得る。言い換えれば、$\mathcal{F}$ の局所化された圏
$S_{affine, etale}/U$ への制限は表示
$$
\bigoplus\nolimits_{k \in K} \mathcal{O}_{affine}|_{S_{affine, \etale}/U}
\longrightarrow
\bigoplus\nolimits_{l \in L} \mathcal{O}_{affine}|_{S_{affine, \etale}/U}
\longrightarrow
\mathcal{F}|_{S_{affine, \etale}/U}
\longrightarrow 0
$$
をもつ。これは $\mathcal{F}$ が準連接であることを示すために必要な表示である。
この煩雑な記法をお詫びしつつ、これで (1) が (2) を導くことの証明を終える。

\medskip\noindent
準連接加群の概念は内在的なので
（Modules on Sites, Lemma \ref{sites-modules-lemma-special-locally-free}）、
同値 (\ref{descent-equation-alternative-small-ringed}) は準連接加群の圏の間の同値を誘導する。
したがって (2) と (3) は同値である。

\medskip\noindent
(3) と (4) の同値性は
Proposition \ref{descent-proposition-equivalence-quasi-coherent} から従う。

\medskip\noindent
(4) を仮定して (1) を示す。$\mathcal{G}$ を (4) のものとし、
$h : U \to U' \to S$ を $S_{affine, \etale}$ の射とする。
構造射を $f : U \to S$ および $f' : U' \to S$ と書けば、
$f = f' \circ h$ である。また
$\mathcal{F}(U') = \Gamma(U', (f')^*\mathcal{G})$ および
$\mathcal{F}(U) = \Gamma(U, f^*\mathcal{G}) = \Gamma(U, h^*(f')^*\mathcal{G})$.
である。したがって Schemes, Lemma
\ref{schemes-lemma-widetilde-pullback} により (1) が成り立つ。

\medskip\noindent
Zariski 位相の場合の証明は省略する。
\end{proof}





\section{寄生加群}
\label{descent-section-parasitic}

\noindent
寄生加群とは、基底スキーム上平坦なスキームへ制限すると零になる加群である。
以下が正式な定義である。

\begin{definition}
\label{descent-definition-parasitic}
スキーム $S$ と $\tau \in \{Zar, \etale,
smooth, syntomic, fppf\}$ を取る。$\mathcal{F}$ を
$(\Sch/S)_\tau$ 上の $\mathcal{O}$-加群の前層とする。
\begin{enumerate}
\item 任意の平坦射 $U \to S$ に対して $\mathcal{F}(U) = 0$ ならば、
$\mathcal{F}$ は {\it 寄生的}\footnote{これは非標準的な用語かもしれない。}
であるという。
\item 任意の $\tau$-被覆 $\{U_i \to S\}_{i \in I}$ に対して、
すべての $i$ について $\mathcal{F}(U_i) = 0$ ならば、
$\mathcal{F}$ は {\it $\tau$-位相に関して寄生的} であるという。
\end{enumerate}
\end{definition}

\noindent
$\tau = fppf$ の場合、これは $U \to S$ が平坦かつ局所有限表示ならば常に
$\mathcal{F}|_{U_{Zar}} = 0$ であることを意味する。他の場合も同様である。

\begin{lemma}
\label{descent-lemma-cohomology-parasitic}
スキーム $S$ と $\tau \in \{Zar, \etale, smooth,
syntomic, fppf\}$ を取る。$\mathcal{G}$ を
$(\Sch/S)_\tau$ 上の $\mathcal{O}$-加群の前層とする。
\begin{enumerate}
\item $\mathcal{G}$ が $\tau$-位相に関して寄生的ならば、
$S$ の開部分である、あるいはそれぞれ $S$ 上 \'etale、$S$ 上 smooth、
$S$ 上 syntomic、$S$ 上平坦かつ局所有限表示である任意の $U$ に対し
$H^p_\tau(U, \mathcal{G}) = 0$ である。
\item $\mathcal{G}$ が寄生的ならば、$S$ 上平坦な任意の $U$ に対し
$H^p_\tau(U, \mathcal{G}) = 0$ である。
\end{enumerate}
\end{lemma}

\begin{proof}
$\tau = fppf$ の場合を証明する。他の場合もまったく同様に証明される。
仮定は、平坦かつ局所有限表示な任意の $U \to S$ に対して
$\mathcal{G}(U) = 0$ であることを意味する。
Cohomology on Sites, Lemma
\ref{sites-cohomology-lemma-cech-vanish-collection} を、
平坦かつ局所有限表示な $U \to S$ からなる部分集合
$\mathcal{B} \subset \Ob((\Sch/S)_{fppf})$ と、
$\mathcal{B}$ の元のすべての fppf 被覆からなる族
$\text{Cov}$ に適用する。
\end{proof}

\begin{lemma}
\label{descent-lemma-direct-image-parasitic}
スキームの射 $f : T \to S$ を取る。$(\Sch/T)_\tau$ 上の任意の
寄生的 $\mathcal{O}$-加群に対し、その直像
$f_*\mathcal{F}$ および高次直像 $R^if_*\mathcal{F}$ は
$(\Sch/S)_\tau$ 上の寄生的 $\mathcal{O}$-加群である。
\end{lemma}

\begin{proof}
$R^if_*\mathcal{F}$ は前層
$$
U \mapsto H^i((\Sch/U \times_S T)_\tau, \mathcal{F})
$$
に付随する層であることを思い出そう。Cohomology on Sites, Lemma
\ref{sites-cohomology-lemma-higher-direct-images} を参照せよ。
$U \to S$ が平坦ならば、基底変換
$U \times_S T \to T$ も平坦である。したがって表示した群は
Lemma \ref{descent-lemma-cohomology-parasitic} により零である。
$\{U_i \to U\}$ が $\tau$-被覆ならば、
$U_i \times_S T \to T$ も平坦である。
したがって、表示した前層の層化は $S$ 上平坦なスキーム $U$ 上で
零であることが明らかである。
\end{proof}

\begin{lemma}
\label{descent-lemma-quasi-coherent-and-flat-base-change}
スキーム $S$ と $\tau \in \{Zar, \etale\}$ を取る。
$\mathcal{G}$ を $(\Sch/S)_{fppf}$ 上の $\mathcal{O}$-加群の層で、
次を満たすものとする。
\begin{enumerate}
\item $\mathcal{G}|_{S_\tau}$ は準連接である。
\item 任意の平坦かつ局所有限表示な射 $g : U \to S$ に対し、標準写像
$g_{\tau, small}^*(\mathcal{G}|_{S_\tau}) \to \mathcal{G}|_{U_\tau}$
は同型である。
\end{enumerate}
このとき $S$ 上平坦かつ局所有限表示な任意の $U$ に対し
$H^p(U, \mathcal{G}) = H^p(U, \mathcal{G}|_{U_\tau})$ である。
\end{lemma}

\begin{proof}
大 fppf サイト $(\Sch/S)_{fppf}$ への
$\mathcal{G}|_{S_\tau}$ の引き戻しを $\mathcal{F}$ とする。
$\mathcal{F}$ は準連接である。標準比較写像
$\varphi : \mathcal{F} \to \mathcal{G}$ が存在し、仮定 (1), (2) により
$\mathcal{F}|_{U_\tau} \to \mathcal{G}|_{U_\tau}$
は平坦かつ局所有限表示なすべての $g : U \to S$ に対して同型である。
したがって短完全列
$$
0 \to \Ker(\varphi) \to \mathcal{F} \to \Im(\varphi) \to 0
$$
および
$$
0 \to \Im(\varphi) \to \mathcal{G} \to \Coker(\varphi) \to 0
$$
において層 $\Ker(\varphi)$ と $\Coker(\varphi)$ は
fppf 位相に関して寄生的である。Lemma
\ref{descent-lemma-cohomology-parasitic} により、
$H^p(U, \mathcal{F}) \to H^p(U, \mathcal{G})$ は
平坦かつ局所有限表示な $g : U \to S$ に対して同型である。
Proposition \ref{descent-proposition-same-cohomology-quasi-coherent} により
結果は $\mathcal{F}$ について成り立つので、結論を得る。
\end{proof}













\section{Fpqc 被覆は普遍的有効全射である}
\label{descent-section-fpqc-universal-effective-epimorphisms}

\noindent
上の結果を応用して、興味深い結果 Lemma
\ref{descent-lemma-fpqc-universal-effective-epimorphisms} を証明する。
Sites, Section \ref{sites-section-representable-sheaves} により、
この補題から、$\tau \in \{Zariski, fppf, \etale, smooth, syntomic\}$
の各場合にサイト $(\Sch/S)_\tau$ 上の表現可能前層が層であることが従う。
まず補助補題を証明する。

\begin{lemma}
\label{descent-lemma-equiv-fibre-product}
スキーム $X$ に対し、その台集合を $|X|$ と書く。
スキームの射 $f : X \to S$ に対して、写像
$$
|X \times_S X| \to |X| \times_{|S|} |X|
$$
は全射である。
\end{lemma}

\begin{proof}
Schemes, Lemma \ref{schemes-lemma-points-fibre-product} における
ファイバー積の点の記述から直ちに従う。
\end{proof}


\begin{lemma}
\label{descent-lemma-universal-effective-epimorphism-affine}
アフィン・スキームの射の族
$\{f_i : X_i \to X\}_{i \in I}$ を取る。次の条件は同値である。
\begin{enumerate}
\item 任意の準連接 $\mathcal{O}_X$-加群 $\mathcal{F}$ に対し
$$
\Gamma(X, \mathcal{F}) =
\text{等化子}\left(
\xymatrix{
\prod\nolimits_{i \in I} \Gamma(X_i, f_i^*\mathcal{F})
\ar@<1ex>[r] \ar@<-1ex>[r] &
\prod\nolimits_{i, j \in I}
\Gamma(X_i \times_X X_j, (f_i \times f_j)^*\mathcal{F})
}
\right)
$$
が成り立つ。
\item $\{f_i : X_i \to X\}_{i \in I}$ はアフィン・スキームの圏における
普遍的有効全射である（Sites, Definition
\ref{sites-definition-universal-effective-epimorphisms}）。
\end{enumerate}
\end{lemma}

\begin{proof}
(2) が成り立つと仮定し、$\mathcal{F}$ を準連接
$\mathcal{O}_X$-加群とする。スキーム
（Constructions, Section \ref{constructions-section-spec}）
$$
X' = \underline{\Spec}_X(\mathcal{O}_X \oplus \mathcal{F})
$$
を考える。ここで $\mathcal{O}_X \oplus \mathcal{F}$ は、乗法
$(f, s)(f', s') = (ff', fs' + f's)$.
をもつ $\mathcal{O}_X$-代数である。
$s_i \in \Gamma(X_i, f_i^*\mathcal{F})$ が切断ならば、
$s_i$ は
$$
\Gamma(X' \times_X X_i, \mathcal{O}_{X' \times_X X_i}) =
\Gamma(X_i, \mathcal{O}_{X_i}) \oplus \Gamma(X_i, f_i^*\mathcal{F})
$$
の一意な元を定める。等式の証明は省略する。
$(s_i)_{i \in I}$ が (1) の等化子に属するとする。このとき、
任意のスキーム $T$ について成り立つ等式
$$
\Mor(T, \mathbf{A}^1_\mathbf{Z}) = \Gamma(T, \mathcal{O}_T)
$$
を用いると、これらの切断は射の族
$h_i : X' \times_X X_i \to \mathbf{A}^1_\mathbf{Z}$ を定め、
$h_i \circ \text{pr}_1 = h_j \circ \text{pr}_2$ が射
$(X' \times_X X_i) \times_{X'} (X' \times_X X_j) \to \mathbf{A}^1_\mathbf{Z}$.
として成り立つ。(2) を仮定したので、射 $h_i$ と両立する射
$h : X' \to \mathbf{A}^1_\mathbf{Z}$ を得る。これはさらに元
$s \in \Gamma(X, \mathcal{F})$ を定める。
$s$ が $\Gamma(X_i, f_i^*\mathcal{F})$ において $s_i$ へ写ることの確認は省略する。

\medskip\noindent
(1) を仮定する。$T$ をアフィン・スキームとし、
$h_i : X_i \to T$ を射の族で、すべての $i, j \in I$ に対し
$X_i \times_X X_j$ 上で
$h_i \circ \text{pr}_1 = h_j \circ \text{pr}_2$ を満たすものとする。このとき
$$
\prod h_i^\sharp :
\Gamma(T, \mathcal{O}_T)
\to
\prod \Gamma(X_i, \mathcal{O}_{X_i})
$$
は等化子へ写るので、補題の仮定を
$\mathcal{F} = \mathcal{O}_X$ に適用して環準同型
$\Gamma(T, \mathcal{O}_T) \to \Gamma(X, \mathcal{O}_X)$
を得る。この環準同型は $h_i = h \circ f_i$ を満たす射
$h : X \to T$ に対応する。したがってこの族は有効全射である。

\medskip\noindent
アフィン間の射 $p : Y \to X$ を取る。前段落の結果を適用して、
$f_i$ の基底変換 $g_i : Y_i \to Y$ が有効全射をなすことを示す。
実際、$\mathcal{G}$ が準連接 $\mathcal{O}_Y$-加群ならば
$$
\Gamma(Y, \mathcal{G}) = \Gamma(X, p_*\mathcal{G}),\quad
\Gamma(Y_i, g_i^*\mathcal{G}) = \Gamma(X, f_i^*p_*\mathcal{G}),
$$
および
$$
\Gamma(Y_i \times_Y Y_j, (g_i \times g_j)^*\mathcal{G}) =
\Gamma(X, (f_i \times f_j)^*p_*\mathcal{G})
$$
が自明な基底変換公式
（Cohomology of Schemes, Lemma \ref{coherent-lemma-affine-base-change}）
により成り立つ。したがって性質 (1) は族 $g_i$ について成り立つ。
\end{proof}

\begin{lemma}
\label{descent-lemma-universal-effective-epimorphism-surjective}
スキームの射の族 $\{f_i : X_i \to X\}_{i \in I}$ を取る。
\begin{enumerate}
\item この族がスキームの圏における普遍的有効全射ならば、
$\coprod f_i$ は全射である。
\item $X$ と $X_i$ がアフィンで、この族がアフィン・スキームの圏における
普遍的有効全射ならば、$\coprod f_i$ は全射である。
\end{enumerate}
\end{lemma}

\begin{proof}
省略する。ヒント：$\Spec(\kappa(x)) \to X$ により基底変換して、
任意の $x \in X$ が像に属さなければならないことを示せ。
\end{proof}

\begin{lemma}
\label{descent-lemma-check-universal-effective-epimorphism-affine}
スキームの射の族 $\{f_i : X_i \to X\}_{i \in I}$ を取る。
$Y$ がアフィンである任意の射 $Y \to X$ に対し、
基底変換の族 $g_i : Y_i \to Y$ が有効全射をなすならば、
$f_i$ の族はスキームの圏における普遍的有効全射をなす。
\end{lemma}

\begin{proof}
スキームの射 $Y \to X$ を取る。基底変換
$g_i : Y_i \to Y$ が有効全射をなすことを示せばよい。
そのため、スキーム $T$ と射 $h_i : Y_i \to T$ が与えられ、
$Y_i \times_Y Y_j$ 上で
$h_i \circ \text{pr}_1 = h_j \circ \text{pr}_2$ が成り立つと仮定する。
アフィン開被覆 $Y = \bigcup V_\alpha$ を選ぶ。
$Y_i$ における $V_\alpha$ の逆像を $V_{\alpha, i}$ と置く。
すると $V_{\alpha, i}  \to V_\alpha$ は
$V_\alpha \to X$ による $f_i$ の基底変換である。したがって仮定により、
制限の族 $h_i|_{V_{\alpha, i}}$ はスキームの射
$h_\alpha : V_\alpha \to T$ から生じる。
これらが交叉上で一致し、所望の射 $Y \to T$ を定めることの確認は読者に委ねる。
Schemes, Section \ref{schemes-section-glueing-schemes} の議論を参照せよ。
\end{proof}

\begin{lemma}
\label{descent-lemma-universal-effective-epimorphism}
アフィン・スキームの射の族
$\{f_i : X_i \to X\}_{i \in I}$ を取る。
Lemma \ref{descent-lemma-universal-effective-epimorphism-affine} の同値な条件が
成り立つと仮定する。さらに、アフィン間の任意の射 $Y \to X$ に対して写像
$$
\coprod X_i \times_X Y \longrightarrow Y
$$
が位相空間の商写像であると仮定する
（Topology, Definition \ref{topology-definition-submersive}）。
このとき、この射の族はスキームの圏における普遍的有効全射である。
\end{lemma}

\begin{proof}
Lemma \ref{descent-lemma-check-universal-effective-epimorphism-affine} により、
$Y$ がアフィンである射 $Y \to X$ でこの射の族を基底変換すれば十分である。
$Y_i = X_i \times_X Y$ と置く。$T$ をスキームとし、
$h_i : Y_i \to T$ を射の族で、
$Y_i \times_Y Y_j$ 上で
$h_i \circ \text{pr}_1 = h_j \circ \text{pr}_2$ を満たすものとする。
集合としての $Y$ は、$\coprod Y_i \times_Y Y_j$ から $\coprod Y_i$ への
二つの写像の余等化子である。実際、全射性は Lemma
\ref{descent-lemma-universal-effective-epimorphism-surjective} のアフィンの場合から、
単射性は Lemma \ref{descent-lemma-equiv-fibre-product} から従う。
したがって写像 $h_i$ と両立する台集合の写像 $h : Y \to T$ が存在する。
補題の第二の条件により $h$ は連続である。
ゆえに $y \in Y$ かつ $U \subset T$ が $h(y)$ のアフィン開近傍ならば、
$h(V) \subset U$ となるアフィン開部分 $V \subset Y$ を取れる。
$V_i = Y_i \times_Y V = X_i \times_X V$ と置けば、Lemma
\ref{descent-lemma-universal-effective-epimorphism-affine} で証明した結果により、
$h|_V : V \to U \subset T$ は、すべての $i$ についてスキームの射として
$h_i|_{V_i}$ と一致する一意なアフィン・スキームの射
$h_V : V \to U$ から生じる。これらの $h_V$ を貼り合わせると
（Schemes, Section \ref{schemes-section-glueing-schemes} を参照せよ）、
所望の射 $Y \to T$ が得られる。
\end{proof}

\begin{lemma}
\label{descent-lemma-open-fpqc-covering}
fpqc 被覆 $\{f_i : T_i \to T\}_{i \in I}$ を取る。
各 $i$ について開部分集合 $W_i \subset T_i$ があり、
すべての $i, j \in I$ に対して
$T_i \times_T T_j$ の開部分集合として
$\text{pr}_0^{-1}(W_i) = \text{pr}_1^{-1}(W_j)$ であると仮定する。
このとき、各 $i$ について $W_i = f_i^{-1}(W)$ となる一意な開部分集合
$W \subset T$ が存在する。
\end{lemma}

\begin{proof}
写像 $\coprod_{i \in I} T_i \to T$ に
Lemma \ref{descent-lemma-equiv-fibre-product} を適用する。
すると、各 $i$ について $W_i = f_i^{-1}(W)$ となる部分集合
$W \subset T$ が存在する。すなわち $W = \bigcup f_i(W_i)$ である。
$W$ が開であることを示すには、$T$ 上 Zariski 局所的に議論してよい。
したがって $T$ はアフィンと仮定してよい。Topologies, Definition
\ref{topologies-definition-fpqc-covering} により、
$\{T_i \to T\}_{i \in I}$ を細分する標準 fpqc 被覆
$\{g_j : V_j \to T\}_{j \in J}$ を選べる。$\alpha : J \to I$ と
$h_j : V_j \to T_{\alpha(j)}$ は Sites, Definition
\ref{sites-definition-morphism-coverings} のものとする。
すると $g_j^{-1}(W) = h_j^{-1}(W_{\alpha(j)})$ である。
したがって $\{f_i : T_i \to T\}$ が標準 fpqc 被覆である場合に帰着する。
この場合、射 $\coprod T_i \to T$ に Morphisms, Lemma
\ref{morphisms-lemma-fpqc-quotient-topology} を適用して、
$W$ が開であると結論できる。
\end{proof}

\begin{lemma}
\label{descent-lemma-fpqc-universal-effective-epimorphisms}
$\{T_i \to T\}$ を fpqc 被覆とする。Topologies, Definition
\ref{topologies-definition-fpqc-covering} を参照せよ。
このとき $\{T_i \to T\}$ はスキームの圏における普遍的有効全射である。
Sites, Definition \ref{sites-definition-universal-effective-epimorphisms}
を参照せよ。言い換えれば、スキームの圏上の任意の表現可能関手は
fpqc 位相に関する層条件を満たす。Topologies, Definition
\ref{topologies-definition-sheaf-property-fpqc} を参照せよ。
\end{lemma}

\begin{proof}
スキーム $S$ を取る。次を示さなければならない。
射 $\varphi_i : T_i \to S$ が与えられ、
$\varphi_i|_{T_i \times_T T_j} = \varphi_j|_{T_i \times_T T_j}$
を満たすならば、各 $T_i$ 上で $\varphi_i$ に制限される一意な射
$T \to S$ が存在する。言い換えれば、関手
$h_S = \Mor_{\Sch}( - , S)$ が fpqc 位相に関する層条件を満たすことを
示さなければならない。

\medskip\noindent
$\{T_i \to T\}$ が Zariski 被覆ならば、これは Schemes, Lemma
\ref{schemes-lemma-glue} から従う。したがって Topologies, Lemma
\ref{topologies-lemma-sheaf-property-fpqc} により、
アフィン間の単一の全射平坦射からなる被覆
$\{X \to Y\}$ の場合に帰着する。

\medskip\noindent
第一の証明。Lemma \ref{descent-lemma-sheaf-condition-holds} により、
$\{X \to Y\}$ に対して準連接加群の層条件が成り立つ。
Lemma \ref{descent-lemma-open-fpqc-covering} により、射 $X \to Y$ は
普遍的に商写像である。したがって Lemma
\ref{descent-lemma-universal-effective-epimorphism} を適用して、
$\{X \to Y\}$ が普遍的有効全射であることが分かる。

\medskip\noindent
第二の証明。$R \to A$ を、全射平坦射 $\pi : X \to Y$ に対応する
忠実平坦な環準同型とする。$f : X \to S$ を射で、
$X \times_Y X = \Spec(A \otimes_R A) \to S$ という射として
$f \circ \text{pr}_1 = f \circ \text{pr}_2$ を満たすものとする。
Lemma \ref{descent-lemma-equiv-fibre-product} により、台集合上の写像として
$f$ は、ある（集合論的）写像 $g : \Spec(R) \to S$ を用いて
$f = g \circ \pi$ と書ける。Morphisms, Lemma
\ref{morphisms-lemma-fpqc-quotient-topology} と $f$ の連続性により、
$g$ も連続である。

\medskip\noindent
$y \in Y = \Spec(R)$ を取る。$g(y)$ を含むアフィン開部分
$U \subset S$ を選び、$U = \Spec(B)$ とする。
上の議論により $y \in D(r) \subset g^{-1}(U)$ となる
$r \in R$ を選べる。$f$ の $\pi^{-1}(D(r))$ への制限を $U$ へ写す射は、
環準同型 $B \to A_r$ に対応する。$f$ に関する仮定により、
誘導される二つの環準同型
$B \to A_r \otimes_{R_r} A_r = (A \otimes_R A)_r$ は等しい。
$R_r \to A_r$ は忠実平坦である。Lemma \ref{descent-lemma-ff-exact} により、
二つの写像 $A_r \to A_r \otimes_{R_r} A_r$ の等化子は $R_r$ である。
したがって $B \to A_r$ は一意に写像 $B \to R_r$ を経由する。
この写像はさらにスキームの射 $D(r) \to U \to S$ を与える。
Schemes, Lemma \ref{schemes-lemma-morphism-into-affine} を参照せよ。

\medskip\noindent
ここまでに何を示したか。任意の素イデアル
$\mathfrak p \subset R$ に対し、標準アフィン開部分
$D(r) \subset \Spec(R)$ であって、射
$f|_{\pi^{-1}(D(r))} : \pi^{-1}(D(r)) \to S$ が、あるスキームの射
$D(r) \to S$ を一意に経由するものが存在することを示した。
これらの射が貼り合わさって所望の射 $\Spec(R) \to S$ を与えることの確認は省略する。
\end{proof}

\begin{lemma}
\label{descent-lemma-coequalizer-fpqc-local}
スキーム $X, Y, Z$、射 $a, b : X \to Y$、および
$c \circ a = c \circ b$ を満たす射 $c : Y \to Z$ を考える。
$d = c \circ a = c \circ b$ と置く。次を満たす fpqc 被覆
$\{Z_i \to Z\}$ が存在すると仮定する。
\begin{enumerate}
\item すべての $i$ について、射 $Y \times_{c, Z} Z_i \to Z_i$ は
$(a, 1) : X \times_{d, Z} Z_i \to Y \times_{c, Z} Z_i$ および
$(b, 1) : X \times_{d, Z} Z_i \to Y \times_{c, Z} Z_i$ の余等化子である。
\item すべての $i$ と $i'$ について、射
$Y \times_{c, Z} (Z_i \times_Z Z_{i'}) \to (Z_i \times_Z Z_{i'})$
は
$(a, 1) : X \times_{d, Z} (Z_i \times_Z Z_{i'}) \to
Y \times_{c, Z} (Z_i \times_Z Z_{i'})$ および
$(b, 1) : X \times_{d, Z} (Z_i \times_Z Z_{i'}) \to
Y \times_{c, Z} (Z_i \times_Z Z_{i'})$ の余等化子である。
\end{enumerate}
このとき $c$ は $a$ と $b$ の余等化子である。
\end{lemma}

\begin{proof}
実際、スキーム $T$ に対し、射 $Z \to T$ を与えることは、
交叉上で一致する射 $Z_i \to T$ の族を与えることと同値である。
これは Lemma \ref{descent-lemma-fpqc-universal-effective-epimorphisms} による。
\end{proof}















\section{射の有限性および平滑性に関する性質の降下}
\label{descent-section-descent-finiteness-morphisms}

\noindent
本節では、射のいくつかの性質（平滑であること、局所有限表示であること
など）が忠実平坦射に沿って降下することを示す。まず代数的な形から
始める。（Noetherian の場合だけに関心のある読者は、次の補題の代わりに
補題 \ref{descent-lemma-finite-type-local-source-fppf-algebra} を参照されたい。）

\begin{lemma}
\label{descent-lemma-flat-finitely-presented-permanence-algebra}
$R \to A \to B$ を環準同型とする。
$R \to B$ は有限表示であり、$A \to B$ は忠実平坦かつ
有限表示であると仮定する。このとき $R \to A$ は有限表示である。
\end{lemma}

\begin{proof}
代数 $C = B \otimes_A B$ と、$p(b) = b \otimes 1$ および
$q(b) = 1 \otimes b$ で定まる二つの準同型
$p, q : B \to C$ を考える。二つの合成 $A \to B \to C$ は
もちろん一致する。$p : B \to C$ は $A \to B$ の基底変換なので
平坦かつ有限表示であり、したがって環準同型 $R \to C$ は
$R \to B \to C$ の合成として有限表示であることに注意する。

\medskip\noindent
Algebra, Lemma \ref{algebra-lemma-characterize-finite-presentation}
の判定法を用いて $R \to A$ が有限表示であることを示す。
$S$ を任意の $R$-代数とし、
$S = \colim_{\lambda \in \Lambda} S_\lambda$ が $R$-代数の
有向余極限として書かれているとする。
$A \to S$ を $R$-代数準同型とする。$A \to S$ がいずれかの
$S_\lambda$ を経由することを示せばよい。
環 $B' = S \otimes_A B$ および
$C' = S \otimes_A C = B' \otimes_S B'$ を考える。
$B$ は $A$ 上忠実平坦かつ有限表示なので、$B'$ も
$S$ 上忠実平坦かつ有限表示である。
Algebra, Lemma \ref{algebra-lemma-flat-finite-presentation-limit-flat}
の (2) を組 $(S \to B', B')$ と系 $(S_\lambda)$ に適用すると、
ある $\lambda_0 \in \Lambda$ と、平坦かつ有限表示な
$S_{\lambda_0}$-代数 $B_{\lambda_0}$ が存在して
$B' = S \otimes_{S_{\lambda_0}} B_{\lambda_0}$ となる。
$\lambda \geq \lambda_0$ に対して
$B_\lambda = S_\lambda \otimes_{S_{\lambda_0}} B_{\lambda_0}$ および
$C_\lambda = B_\lambda \otimes_{S_\lambda} B_\lambda$ とおく。

\medskip\noindent
ここで議論をいったん中断し、十分大きな $\lambda$ に対して
$S_\lambda \to B_\lambda$ が忠実平坦であることを示す。
（これは本来、たとえば極限の章の別の補題にすべき主張である。）
$\Spec(B_{\lambda_0}) \to \Spec(S_{\lambda_0})$ は平坦かつ
有限表示なので開写像である（Morphisms,
Lemma \ref{morphisms-lemma-fppf-open} 参照）。
像の補集合が $V(I) \subset \Spec(S_{\lambda_0})$ となるような
イデアル $I \subset S_{\lambda_0}$ をとる。像の形成は基底変換と
可換であることに注意する。したがって
$\Spec(B') \to \Spec(S)$ が全射であり、かつ
$B' = B_{\lambda_0} \otimes_{S_{\lambda_0}} S$ であることから
$IS = S$ を得る。ゆえに、ある $\lambda \geq \lambda_0$ に対して
$IS_{\lambda} = S_\lambda$ である。この値およびそれ以上の
すべての $\lambda$ に対して、射
$\Spec(B_\lambda) \to \Spec(S_\lambda)$ は全射である。

\medskip\noindent
証明の第1段落の記法にならい、二つの標準準同型を
$p_\lambda, q_\lambda : B_\lambda \to C_\lambda$ と書く。
すると $B' = \colim_{\lambda \geq \lambda_0} B_\lambda$ および
$C' = \colim_{\lambda \geq \lambda_0} C_\lambda$ である。
$B$ と $C$ は $R$ 上有限表示なので、Algebra,
Lemma \ref{algebra-lemma-characterize-finite-presentation} を
数回適用すれば、ある $\lambda \geq \lambda_0$ と
$R$-代数準同型 $B \to B_\lambda$, $C \to C_\lambda$ が存在し、
次の図式は可換となる。
$$
\xymatrix{
C \ar[rr] & &
C_\lambda \\
B \ar[rr]
\ar@<1ex>[u]^-p
\ar@<-1ex>[u]_-q
& &
B_\lambda
\ar@<1ex>[u]^-{p_\lambda}
\ar@<-1ex>[u]_-{q_\lambda}
}
$$
したがって $A \to B \to B_\lambda$ の像は
$p_\lambda$ と $q_\lambda$ の等化子に含まれる。
補題 \ref{descent-lemma-ff-exact} により、$S_\lambda$ は
$p_\lambda$ と $q_\lambda$ の等化子である。
よって所望の環準同型 $A \to S_\lambda$ が得られ、証明が完了する。
\end{proof}

\noindent
有限型という性質だけを扱う、これより容易な形を次に述べる。

\begin{lemma}
\label{descent-lemma-finite-type-local-source-fppf-algebra}
$R \to A \to B$ を環準同型とする。
$R \to B$ は有限型であり、$A \to B$ は忠実平坦かつ
有限表示であると仮定する。このとき $R \to A$ は有限型である。
\end{lemma}

\begin{proof}
Algebra, Lemma \ref{algebra-lemma-descend-faithfully-flat-finite-presentation}
により、次の可換図式が存在する。
$$
\xymatrix{
R \ar[r] \ar@{=}[d] &
A_0 \ar[d] \ar[r] &
B_0 \ar[d] \\
R \ar[r] & A \ar[r] & B
}
$$
ここで $R \to A_0$ は有限表示、
$A_0 \to B_0$ は忠実平坦かつ有限表示であり、
$B = A \otimes_{A_0} B_0$ である。仮定により $R \to B$ は
有限型なので、$A_0$ にいくつかの元を加え、
準同型 $B_0 \to B$ が全射であると仮定してよい。
このとき $A_0 \to B_0$ は忠実平坦であり、
$$
(A_0 \to A) \otimes_{A_0} B_0 \cong (B_0 \to B)
$$
が全射なので、$A_0 \to A$ も全射である。したがって結論を得る。
\end{proof}

\begin{lemma}
\label{descent-lemma-flat-finitely-presented-permanence}
\begin{reference}
\cite[IV, 17.7.5 (i) and (ii)]{EGA}.
\end{reference}
次の可換図式を考える。
$$
\xymatrix{
X \ar[rr]_f \ar[rd]_p & &
Y \ar[dl]^q \\
& S
}
$$
ここで $f$ は全射、平坦、かつ局所有限表示であり、
$p$ は局所有限表示（resp.\ 局所有限型）であると仮定する。
このとき $q$ は局所有限表示（resp.\ 局所有限型）である。
\end{lemma}

\begin{proof}
問題は $S$ 上および $Y$ 上局所的なので、$S$ と $Y$ は
アフィンであると仮定してよい。$f$ は平坦かつ局所有限表示なので、
$f$ は開写像である
（Morphisms, Lemma \ref{morphisms-lemma-fppf-open}）。
したがって $Y$ は準コンパクトであるから、
$Y = \bigcup f(X_i)$ となる有限個のアフィン開集合
$X_i \subset X$ が存在する。明らかに $X$ を $\coprod X_i$ で
置き換えてよく、ゆえに $X$ もアフィンであると仮定できる。
この場合、本補題は上の
補題 \ref{descent-lemma-flat-finitely-presented-permanence-algebra}
（resp. 補題 \ref{descent-lemma-finite-type-local-source-fppf-algebra}）
と同値である。
\end{proof}

\noindent
これを用いて、先に得た射に関する結果のいくつかを強める。

\begin{lemma}
\label{descent-lemma-syntomic-smooth-etale-permanence}
次の可換図式を考える。
$$
\xymatrix{
X \ar[rr]_f \ar[rd]_p & &
Y \ar[dl]^q \\
& S
}
$$
次を仮定する。
\begin{enumerate}
\item $f$ は全射で、かつ syntomic（resp.\ 平滑、resp.\ \'etale）である。
\item $p$ は syntomic（resp.\ 平滑、resp.\ \'etale）である。
\end{enumerate}
このとき $q$ は syntomic（resp.\ 平滑、resp.\ \'etale）である。
\end{lemma}

\begin{proof}
Morphisms, Lemmas
\ref{morphisms-lemma-syntomic-permanence},
\ref{morphisms-lemma-smooth-permanence} および
\ref{morphisms-lemma-etale-permanence-two}
を上の補題 \ref{descent-lemma-flat-finitely-presented-permanence} と
組み合わせればよい。
\end{proof}

\noindent
実際、この結果は次のように強められる。

\begin{lemma}
\label{descent-lemma-smooth-permanence}
次の可換図式を考える。
$$
\xymatrix{
X \ar[rr]_f \ar[rd]_p & &
Y \ar[dl]^q \\
& S
}
$$
次を仮定する。
\begin{enumerate}
\item $f$ は全射、平坦、かつ局所有限表示である。
\item $p$ は平滑（resp.\ \'etale）である。
\end{enumerate}
このとき $q$ は平滑（resp.\ \'etale）である。
\end{lemma}

\begin{proof}
(1) を仮定し、$p$ は平滑であるとする。補題
\ref{descent-lemma-flat-finitely-presented-permanence} により、
$q$ は局所有限表示である。Morphisms,
Lemma \ref{morphisms-lemma-flat-permanence} により、$q$ は平坦である。
したがって、あとは $q$ のファイバーが平滑であることを示せば十分である。
Morphisms, Lemma \ref{morphisms-lemma-smooth-flat-smooth-fibres} を参照。
$s \in S$ に対する平坦全射 $X_s \to Y_s$ に
Varieties, Lemma \ref{varieties-lemma-flat-under-smooth} を
適用すれば結論を得る。\'etale の場合の証明は省略する。
\end{proof}

\begin{remark}
\label{descent-remark-smooth-permanence}
補題 \ref{descent-lemma-smooth-permanence} の仮定 (1) が成り立ち、
$p$ が平滑であっても、$X \to Y$ が自動的に平滑になるとは限らない。
反例は $S = \Spec(k)$、$X = \Spec(k[s])$、
$Y = \Spec(k[t])$ とし、$f$ を $t \mapsto s^2$ で定めれば得られる。
$f$ の構造に関する情報については
補題 \ref{descent-lemma-syntomic-permanence} も参照されたい。
\end{remark}

\begin{lemma}
\label{descent-lemma-syntomic-permanence}
次の可換図式を考える。
$$
\xymatrix{
X \ar[rr]_f \ar[rd]_p & &
Y \ar[dl]^q \\
& S
}
$$
次を仮定する。
\begin{enumerate}
\item $f$ は全射、平坦、かつ局所有限表示である。
\item $p$ は syntomic である。
\end{enumerate}
このとき $q$ と $f$ はともに syntomic である。
\end{lemma}

\begin{proof}
補題 \ref{descent-lemma-flat-finitely-presented-permanence} により
$q$ は有限表示である。Morphisms,
Lemma \ref{morphisms-lemma-flat-permanence} により $q$ は平坦である。
Morphisms, Lemma \ref{morphisms-lemma-syntomic-locally-standard-syntomic}
により、あとは $Y \to S$ のファイバーおよび $X \to Y$ の
ファイバーの局所環が局所完全交叉環であることを示せば十分である。
そのためには、点 $s \in S$ における $X \to Y \to S$ の
ファイバーをとってよい。すなわち $S$ は体のスペクトルであると
仮定してよい。$x \in X$ をとり、その像を $y \in Y$ として、
環準同型
$$
\mathcal{O}_{Y, y} \longrightarrow \mathcal{O}_{X, x}
$$
を考える。これは局所 Noether 環の平坦な局所準同型である。
局所環 $\mathcal{O}_{X, x}$ は完全交叉である。
そこで Avramov の結果
（Divided Power Algebra, Lemma \ref{dpa-lemma-avramov}）を用いると、
$\mathcal{O}_{Y, y}$ と
$\mathcal{O}_{X, x}/\mathfrak m_y\mathcal{O}_{X, x}$
はいずれも完全交叉環であると分かる。
\end{proof}

\noindent
次の型の補題は、ときに有用である。

\begin{lemma}
\label{descent-lemma-curiosity}
$X \to Y \to Z$ をスキームの射とする。$P$ をスキームの射の
次の性質のいずれかとする：平坦、局所有限型、局所有限表示。
$X \to Z$ が $P$ をもち、$\{X \to Y\}$ が $Y$ の fppf 被覆で
細分できると仮定する。このとき $Y \to Z$ は $P$ をもつ。
\end{lemma}

\begin{proof}
$\Spec(C) \subset Z$ をアフィン開集合とし、
$\Spec(B) \subset Y$ を $\Spec(C)$ に写るアフィン開集合とする。
$X \to Y$ に関する仮定から、標準アフィン fppf 被覆
$\{\Spec(B_j) \to \Spec(B)\}$ と持ち上げ
$x_j : \Spec(B_j) \to X$ が得られる。$\Spec(B_j)$ は
準コンパクトなので、$\Spec(B)$ 上にある有限個のアフィン開集合
$\Spec(A_i) \subset X$ を、各 $x_j$ の像が
$\bigcup \Spec(A_i)$ に含まれるようにとれる。
したがって各 $\Spec(B_j)$ をその標準アフィン Zariski 被覆で
置き換えれば、各 $\Spec(B_i) \to Y$ が $\Spec(B)$ 上にある
アフィン開集合 $\Spec(A_i) \subset X$ を経由するような
標準アフィン fppf 被覆 $\{\Spec(B_i) \to \Spec(B)\}$ が
あると仮定してよい。言い換えれば、各 $i$ に対して環準同型
$C \to B \to A_i \to B_i$ がある。さらに
$$
C \to B \to A = \prod A_i \to B' = \prod B_i
$$
も考えることができ、環準同型 $B \to \prod B_i$ は
忠実平坦かつ有限表示である。

\medskip\noindent
$P = flat$ の場合。この場合 $C \to A$ は平坦であり、
$C \to B$ が平坦であることを示す必要がある。
$N \to N' \to N''$ を $C$-加群の完全系列とする。
$N \otimes_C B \to N' \otimes_C B \to N'' \otimes_C B$
が完全であることを示したい。そのコホモロジーを $H$ とし、
$N \otimes_C B' \to N' \otimes_C B' \to N'' \otimes_C B'$
のコホモロジーを $H'$ とする。$B \to B'$ は平坦なので
$H' = H \otimes_B B'$ である。一方、
$N \otimes_C A \to N' \otimes_C A \to N'' \otimes_C A$
は完全であり、したがってそのコホモロジーは零である。
ゆえに準同型 $H \to H'$ は零加群を経由するので零であり、
$H' = 0$ となる。$B \to B'$ は忠実平坦だから、
所望のとおり $H = 0$ を得る。

\medskip\noindent
$P = locally\ finite\ type$ の場合。この場合 $C \to A$ は有限型であり、
$C \to B$ が有限型であることを示す必要がある。
$B \to B'$ は有限表示（したがって有限型）なので、
$A \to B'$ は有限型である。Algebra,
Lemma \ref{algebra-lemma-compose-finite-type} を参照。
したがって $C \to B'$ は有限型であり、
補題 \ref{descent-lemma-finite-type-local-source-fppf-algebra} から結論を得る。

\medskip\noindent
$P = locally\ finite\ presentation$ の場合。この場合 $C \to A$ は
有限表示であり、$C \to B$ が有限表示であることを示す必要がある。
$B \to B'$ は有限表示であり、$B \to A$ は有限型なので、
$A \to B'$ は有限表示である。Algebra,
Lemma \ref{algebra-lemma-compose-finite-type} を参照。
したがって $C \to B'$ は有限表示であり、
補題 \ref{descent-lemma-flat-finitely-presented-permanence-algebra} から
結論を得る。
\end{proof}










\section{スキームの局所的性質}
\label{descent-section-descending-properties}

\noindent
スキームの被覆の各要素がある性質をもつことを証明できる場合が
しばしばある。多くの場合、これはスキーム自身もその性質をもつことを
意味する。たとえば $S$ をスキームとし、$f : S' \to S$ を
全射平坦な射で $S'$ が被約スキームであるものとすれば、$S$ も
被約である。これは局所環を調べ、Algebra,
Lemma \ref{algebra-lemma-descent-reduced} を用いれば証明できる。
このとき、被約であるという性質は
{\it 全射平坦射を通じて降下する}という。
この型の結果の一部は Algebra,
Section \ref{algebra-section-descending-properties} に、スキームの場合は
Section \ref{descent-section-variants} にまとめられている。
射の性質の降下に関する類似の結果は
Section \ref{descent-section-descent-finiteness-morphisms} にある。

\medskip\noindent
他方、$S$ は被約だが $S'$ は被約でない全射平坦射
$f : S' \to S$ の例がある。たとえば射
$\Spec(k[x]/(x^2)) \to \Spec(k)$ がそうである。したがって、
被約であるという性質は{\it 平坦射に沿って上昇しない}。
剰余体が無限であるという性質は平坦射に沿って上昇する
（ただし、もちろん全射平坦射に沿って降下しない）。
この型の結果の一部は Algebra,
Section \ref{algebra-section-ascending-properties} にまとめられている。

\medskip\noindent
最後に、ある性質が平坦射に沿って上昇し、全射平坦射に沿って
降下するとき、その性質は{\it 平坦位相について局所的}であるという。
やや自明な例は、剰余体が所定の標数をもつという性質である。
より正確に述べ、これをスキーム上のさまざまな位相と結び付けるため、
次の形式的な定義を置く。

\begin{definition}
\label{descent-definition-property-local}
$\mathcal{P}$ をスキームの性質とする。また
$\tau \in \{fpqc, \linebreak[0] fppf, \linebreak[0] syntomic, \linebreak[0]
smooth, \linebreak[0] \etale, \linebreak[0] Zariski\}$.
任意の $\tau$-被覆 $\{S_i \to S\}_{i \in I}$
（Topologies, Section \ref{topologies-section-procedure} 参照）に対して
次が成り立つとき、$\mathcal{P}$ は
{\it $\tau$-位相について局所的}であるという。
$$
S \text{ は }\mathcal{P}\text{ をもつ}
\Leftrightarrow
\text{各 }S_i \text{ は }\mathcal{P}\text{ をもつ}.
$$
\end{definition}

\noindent
同型射は常に被覆なので、性質 $\mathcal{P}$ が $S$ について
成り立つことと、$S$ と同型な任意のスキーム $S'$ について
成り立つこととは同値であると分かる（あるいは、そう要請する）。
実際、$\tau = fpqc, \linebreak[0] fppf, \linebreak[0] syntomic,
\linebreak[0] smooth, \linebreak[0] \etale$ または $Zariski$
のいずれかであるとする。$S$ が $\mathcal{P}$ をもち、
$S' \to S$ がそれぞれ平坦、平坦かつ局所有限表示、syntomic、
平滑、\'etale、または開埋め込みであれば、$S'$ も
$\mathcal{P}$ をもつ。これは $\{S' \to S\}$ を常に
$\tau$-被覆へ拡張できるからである。

\medskip\noindent
次の含意が成り立つ：
$\mathcal{P}$ は fpqc 位相について局所的
$\Rightarrow$
$\mathcal{P}$ は fppf 位相について局所的
$\Rightarrow$
$\mathcal{P}$ は syntomic 位相について局所的
$\Rightarrow$
$\mathcal{P}$ は smooth 位相について局所的
$\Rightarrow$
$\mathcal{P}$ は \'etale 位相について局所的
$\Rightarrow$
$\mathcal{P}$ は Zariski 位相について局所的。
これは Topologies, Lemmas
\ref{topologies-lemma-zariski-etale},
\ref{topologies-lemma-zariski-etale-smooth},
\ref{topologies-lemma-zariski-etale-smooth-syntomic},
\ref{topologies-lemma-zariski-etale-smooth-syntomic-fppf} および
\ref{topologies-lemma-zariski-etale-smooth-syntomic-fppf-fpqc}
から従う。

\begin{lemma}
\label{descent-lemma-descending-properties}
$\mathcal{P}$ をスキームの性質とする。
$\tau \in \{fpqc, \linebreak[0] fppf, \linebreak[0]
\etale, \linebreak[0] smooth, \linebreak[0] syntomic\}$.
次を仮定する。
\begin{enumerate}
\item この性質は Zariski 位相について局所的である。
\item アフィン・スキームの任意の射 $S' \to S$ が、
$\tau$ が fpqc、fppf、\'etale、smooth、syntomic のいずれであるかに
応じて、それぞれ平坦、有限表示の平坦射、\'etale、平滑、syntomic
であるとする。$\mathcal{P}$ が $S$ について成り立つならば、
$\mathcal{P}$ は $S'$ について成り立つ。
\item アフィン・スキームの任意の全射 $S' \to S$ が、
$\tau$ が fpqc、fppf、\'etale、smooth、syntomic のいずれであるかに
応じて、それぞれ平坦、有限表示の平坦射、\'etale、平滑、syntomic
であるとする。$\mathcal{P}$ が $S'$ について成り立つならば、
$\mathcal{P}$ は $S$ について成り立つ。
\end{enumerate}
このとき $\mathcal{P}$ は基底上 $\tau$-局所的である。
\end{lemma}

\begin{proof}
これは $\tau$-被覆の定義からほとんど直ちに従う。
Topologies, Definition
\ref{topologies-definition-fpqc-covering}
\ref{topologies-definition-fppf-covering}
\ref{topologies-definition-etale-covering}
\ref{topologies-definition-smooth-covering} または
\ref{topologies-definition-syntomic-covering}
と、Topologies, Lemma
\ref{topologies-lemma-fpqc-affine},
\ref{topologies-lemma-fppf-affine},
\ref{topologies-lemma-etale-affine},
\ref{topologies-lemma-smooth-affine} または
\ref{topologies-lemma-syntomic-affine} を参照。
詳細は省略する。
\end{proof}

\begin{remark}
\label{descent-remark-descending-properties-standard}
上の補題 \ref{descent-lemma-descending-properties} で
$\tau = smooth$ ならば、条件 (3) の射は（全射な）標準平滑射であると
仮定してよい。$\tau = syntomic$ または $\tau = \etale$ の場合も
同様である。
\end{remark}





\section{fppf 位相について局所的なスキームの性質}
\label{descent-section-descending-properties-fppf}

\noindent
本節では、fppf 位相において基底上局所的なスキームの性質をいくつか求める。

\begin{lemma}
\label{descent-lemma-Noetherian-local-fppf}
性質 $\mathcal{P}(S) =$「$S$ は局所 Noether である」は
fppf 位相について局所的である。
\end{lemma}

\begin{proof}
補題 \ref{descent-lemma-descending-properties} を用いる。
まず、「局所 Noether であること」は Zariski 位相について局所的である。
これは定義から明らかである。Properties,
Definition \ref{properties-definition-noetherian} を参照。
次に、$S' \to S$ がアフィン・スキームの平坦かつ有限表示な射で、
$S$ が局所 Noether ならば、$S'$ も局所 Noether であることを示す。
これは Morphisms, Lemma
\ref{morphisms-lemma-finite-type-noetherian} である。
最後に、$S' \to S$ がアフィン・スキームの全射、平坦、かつ
有限表示な射で、$S'$ が局所 Noether ならば、$S$ も局所 Noether
であることを示す必要がある。これは Algebra,
Lemma \ref{algebra-lemma-descent-Noetherian} から従う。
したがって補題 \ref{descent-lemma-descending-properties} の (1)、(2)、(3) が
成り立ち、結論を得る。
\end{proof}

\begin{lemma}
\label{descent-lemma-Jacobson-local-fppf}
性質 $\mathcal{P}(S) =$「$S$ は Jacobson である」は
fppf 位相について局所的である。
\end{lemma}

\begin{proof}
補題 \ref{descent-lemma-descending-properties} を用いる。
まず、「Jacobson であること」は Zariski 位相について局所的である。
これは Properties, Lemma
\ref{properties-lemma-locally-jacobson} である。
次に、$S' \to S$ がアフィン・スキームの平坦かつ有限表示な射で、
$S$ が Jacobson ならば $S'$ も Jacobson であることを示す。
これは Morphisms, Lemma
\ref{morphisms-lemma-Jacobson-universally-Jacobson} である。
最後に、$f : S' \to S$ がアフィン・スキームの全射、平坦、かつ
有限表示な射で、$S'$ が Jacobson ならば $S$ も Jacobson
であることを示す必要がある。$S = \Spec(A)$、
$S' = \Spec(B)$ と書き、$S' \to S$ は $A \to B$ により
与えられるとする。このとき $A \to B$ は有限表示かつ忠実平坦である。
さらに環 $B$ は Jacobson である。Properties,
Lemma \ref{properties-lemma-locally-jacobson} を参照。

\medskip\noindent
Algebra, Lemma \ref{algebra-lemma-fppf-fpqf} により、次の図式が存在する。
$$
\xymatrix{
B \ar[rr] & & B' \\
& A \ar[ru] \ar[lu] &
}
$$
ここで $A \to B'$ は有限表示、忠実平坦、かつ準有限である。
特に $B \to B'$ は有限型であり、Algebra,
Proposition \ref{algebra-proposition-Jacobson-permanence} から
$B'$ は Jacobson であると分かる。したがって $A \to B$ は
忠実平坦かつ有限表示であることに加えて準有限であると仮定してよい。

\medskip\noindent
矛盾を導くため、$A$ は Jacobson でないと仮定する。
Algebra, Lemma \ref{algebra-lemma-characterize-jacobson} により、
極大でない素イデアル $\mathfrak p \subset A$ と元
$f \in A$, $f \not \in \mathfrak p$ で
$V(\mathfrak p) \cap D(f) = \{\mathfrak p\}$.
を満たすものが存在する。

\medskip\noindent
次のように矛盾が生じる。まず $\mathfrak p \subset \mathfrak m$ となる
$A$ の極大イデアルをとる。
$\mathfrak m$ の上にある素イデアル $\mathfrak m' \subset B$ をとる
（$A \to B$ は忠実平坦なので存在する。Algebra,
Lemma \ref{algebra-lemma-ff-rings} 参照）。
$A \to B$ は平坦なので、下降定理
（Algebra, Lemma \ref{algebra-lemma-flat-going-down}）により、
$\mathfrak p$ の上にある素イデアル
$\mathfrak q \subset \mathfrak m'$ をとれる。
特に $\mathfrak q$ は極大でない。したがって再び Algebra,
Lemma \ref{algebra-lemma-characterize-jacobson} により、集合
$V(\mathfrak q) \cap D(f)$ は無限である
（ここで初めて $B$ が Jacobson であることを用いる）。
$V(\mathfrak q) \cap D(f)$ のすべての点は
$V(\mathfrak p) \cap D(f) = \{\mathfrak p\}$ に写る。
したがって $\mathfrak p$ 上のファイバーは無限である。
これは $A \to B$ が準有限であることに反する（
Algebra, Lemma \ref{algebra-lemma-quasi-finite}
または、より明示的には
Morphisms, Lemma \ref{morphisms-lemma-quasi-finite} 参照）。
以上で補題が証明された。
\end{proof}

\begin{lemma}
\label{descent-lemma-locally-finite-nr-irred-local-fppf}
性質 $\mathcal{P}(S) =$「$S$ の各準コンパクト開集合は有限個の
既約成分をもつ」は fppf 位相について局所的である。
\end{lemma}

\begin{proof}
補題 \ref{descent-lemma-descending-properties} を用いる。まず
$\mathcal{P}$ は Zariski 位相について局所的である。
次に、$T \to S$ がアフィン・スキームの平坦かつ有限表示な射で、
$S$ が有限個の既約成分をもつならば、$T$ もそうであることを示す。
実際、$T \to S$ は平坦なので、$T$ の生成点は $S$ の生成点に写る。
Morphisms, Lemma
\ref{morphisms-lemma-generalizations-lift-flat} を参照。
したがって、$s \in S$ に対してファイバー $T_s$ が有限個の
生成点をもつことを示せば十分である。$T_s$ は $\kappa(s)$ 上
有限型なアフィン・スキームであることに注意する。Morphisms,
Lemma \ref{morphisms-lemma-base-change-finite-type} を参照。
ゆえに $T_s$ は Noether で、有限個の既約成分をもつ
（Morphisms, Lemma \ref{morphisms-lemma-finite-type-noetherian} および
Properties, Lemma
\ref{properties-lemma-Noetherian-irreducible-components}）。
最後に、$T \to S$ がアフィン・スキームの全射、平坦、かつ有限表示な
射で、$T$ が有限個の既約成分をもつならば、$S$ もそうであることを
示す必要がある。これは Topology, Lemma
\ref{topology-lemma-surjective-continuous-irreducible-components} から従う。
したがって補題 \ref{descent-lemma-descending-properties} の (1)、(2)、(3) が
成り立ち、結論を得る。
\end{proof}




\section{syntomic 位相について局所的なスキームの性質}
\label{descent-section-descending-properties-syntomic}

\noindent
本節では、syntomic 位相において基底上局所的なスキームの性質を
いくつか求める。

\begin{lemma}
\label{descent-lemma-Sk-local-syntomic}
性質 $\mathcal{P}(S) =$「$S$ は局所 Noether であり $(S_k)$ を満たす」は
syntomic 位相について局所的である。
\end{lemma}

\begin{proof}
補題 \ref{descent-lemma-descending-properties} の (1)、(2)、(3) を確認する。
syntomic 射は有限表示の平坦射なので
（Morphisms, Lemmas \ref{morphisms-lemma-syntomic-flat} および
\ref{morphisms-lemma-syntomic-locally-finite-presentation}）、
「局所 Noether であること」については補題
\ref{descent-lemma-Noetherian-local-fppf} の証明ですでに確認済みである。
以下ではこれを断りなく用いる。
まず $\mathcal{P}$ は Zariski 位相について局所的である。
これは定義から明らかである。Cohomology of Schemes,
Definition \ref{coherent-definition-depth} を参照。
次に、$S' \to S$ がアフィン・スキームの syntomic 射で、
$S$ が $\mathcal{P}$ をもつならば $S'$ も $\mathcal{P}$ を
もつことを示す。これは Algebra, Lemma
\ref{algebra-lemma-Sk-goes-up} である
（Morphisms, Lemma \ref{morphisms-lemma-syntomic-characterize}、
Algebra, Definition \ref{algebra-definition-lci} および
Lemma \ref{algebra-lemma-lci-CM} を用いる）。
最後に、$S' \to S$ がアフィン・スキームの全射 syntomic 射で、
$S'$ が $\mathcal{P}$ をもつならば $S$ も $\mathcal{P}$ を
もつことを示す。これは Algebra, Lemma
\ref{algebra-lemma-descent-Sk} である。
したがって補題 \ref{descent-lemma-descending-properties} の (1)、(2)、(3) が
成り立ち、結論を得る。
\end{proof}

\begin{lemma}
\label{descent-lemma-CM-local-syntomic}
性質 $\mathcal{P}(S) =$「$S$ は Cohen--Macaulay である」は
syntomic 位相について局所的である。
\end{lemma}

\begin{proof}
これは上の補題 \ref{descent-lemma-Sk-local-syntomic} から明らかである。
実際、スキームが Cohen--Macaulay であることと、局所 Noether であり、
すべての $k \geq 0$ に対して $(S_k)$ を満たすこととは同値である。
Properties, Lemma \ref{properties-lemma-scheme-CM-iff-all-Sk} を参照。
\end{proof}







\section{smooth 位相について局所的なスキームの性質}
\label{descent-section-descending-properties-smooth}

\noindent
本節では、smooth 位相において基底上局所的なスキームの性質を
いくつか求める。

\begin{lemma}
\label{descent-lemma-reduced-local-smooth}
性質 $\mathcal{P}(S) =$「$S$ は被約である」は
smooth 位相について局所的である。
\end{lemma}

\begin{proof}
補題 \ref{descent-lemma-descending-properties} を用いる。
まず、「被約であること」は Zariski 位相について局所的である。
これは定義から明らかである。Schemes,
Definition \ref{schemes-definition-reduced} を参照。
次に、$S' \to S$ がアフィン・スキームの平滑射で、
$S$ が被約ならば $S'$ も被約であることを示す。
これは Algebra, Lemma \ref{algebra-lemma-reduced-goes-up} である。
最後に、$S' \to S$ がアフィン・スキームの全射平滑射で、
$S'$ が被約ならば $S$ も被約であることを示す。
これは Algebra, Lemma \ref{algebra-lemma-descent-reduced} である。
したがって補題 \ref{descent-lemma-descending-properties} の (1)、(2)、(3) が
成り立ち、結論を得る。
\end{proof}

\begin{lemma}
\label{descent-lemma-normal-local-smooth}
\begin{slogan}
正規性は smooth 位相について局所的である。
\end{slogan}
性質 $\mathcal{P}(S) =$「$S$ は正規である」は
smooth 位相について局所的である。
\end{lemma}

\begin{proof}
補題 \ref{descent-lemma-descending-properties} を用いる。
まず、「正規であること」は Zariski 位相について局所的である。
これは定義から明らかである。Properties,
Definition \ref{properties-definition-normal} を参照。
次に、$S' \to S$ がアフィン・スキームの平滑射で、
$S$ が正規ならば $S'$ も正規であることを示す。
これは Algebra, Lemma \ref{algebra-lemma-normal-goes-up} である。
最後に、$S' \to S$ がアフィン・スキームの全射平滑射で、
$S'$ が正規ならば $S$ も正規であることを示す。
これは Algebra, Lemma \ref{algebra-lemma-descent-normal} である。
したがって補題 \ref{descent-lemma-descending-properties} の (1)、(2)、(3) が
成り立ち、結論を得る。
\end{proof}

\begin{lemma}
\label{descent-lemma-Rk-local-smooth}
性質 $\mathcal{P}(S) =$「$S$ は局所 Noether であり $(R_k)$ を満たす」は
smooth 位相について局所的である。
\end{lemma}

\begin{proof}
補題 \ref{descent-lemma-descending-properties} の (1)、(2)、(3) を確認する。
平滑射は有限表示の平坦射なので
（Morphisms, Lemmas \ref{morphisms-lemma-smooth-flat} および
\ref{morphisms-lemma-smooth-locally-finite-presentation}）、
「局所 Noether であること」については補題
\ref{descent-lemma-Noetherian-local-fppf} の証明ですでに確認済みである。
以下ではこれを断りなく用いる。
まず $\mathcal{P}$ は Zariski 位相について局所的である。
これは定義から明らかである。Properties,
Definition \ref{properties-definition-Rk} を参照。
次に、$S' \to S$ がアフィン・スキームの平滑射で、
$S$ が $\mathcal{P}$ をもつならば $S'$ も $\mathcal{P}$ を
もつことを示す。これは Algebra, Lemmas
\ref{algebra-lemma-Rk-goes-up} である
（Morphisms, Lemma \ref{morphisms-lemma-smooth-characterize}、
Algebra, Lemmas \ref{algebra-lemma-base-change-smooth} および
\ref{algebra-lemma-characterize-smooth-over-field} を用いる）。
最後に、$S' \to S$ がアフィン・スキームの全射平滑射で、
$S'$ が $\mathcal{P}$ をもつならば $S$ も $\mathcal{P}$ を
もつことを示す。これは Algebra, Lemma
\ref{algebra-lemma-descent-Rk} である。
したがって補題 \ref{descent-lemma-descending-properties} の (1)、(2)、(3) が
成り立ち、結論を得る。
\end{proof}

\begin{lemma}
\label{descent-lemma-regular-local-smooth}
性質 $\mathcal{P}(S) =$「$S$ は正則である」は
smooth 位相について局所的である。
\end{lemma}

\begin{proof}
これは上の補題 \ref{descent-lemma-Rk-local-smooth} から明らかである。
実際、局所 Noether スキームが正則であることと、局所 Noether であり、
すべての $k \geq 0$ に対して $(R_k)$ を満たすこととは同値である。
\end{proof}

\begin{lemma}
\label{descent-lemma-Nagata-local-smooth}
性質 $\mathcal{P}(S) =$「$S$ は Nagata である」は
smooth 位相について局所的である。
\end{lemma}

\begin{proof}
補題 \ref{descent-lemma-descending-properties} の (1)、(2)、(3) を確認する。
まず Nagata であることは Zariski 位相について局所的である。
これは Properties, Lemma \ref{properties-lemma-locally-nagata} である。
次に、$S' \to S$ がアフィン・スキームの平滑射で、
$S$ が Nagata ならば $S'$ も Nagata であることを示す。
これは Morphisms, Lemma
\ref{morphisms-lemma-finite-type-nagata} である。
最後に、$S' \to S$ がアフィン・スキームの全射平滑射で、
$S'$ が Nagata ならば $S$ も Nagata であることを示す。
これは Algebra, Lemma \ref{algebra-lemma-descent-nagata} である。
したがって補題 \ref{descent-lemma-descending-properties} の (1)、(2)、(3) が
成り立ち、結論を得る。
\end{proof}





\section{性質の降下に関する変種}
\label{descent-section-variants}

\noindent
局所的ではない性質でも降下できることがある。この型の結果を
本節にまとめる。平滑性など、射の性質の降下については
Section \ref{descent-section-descent-finiteness-morphisms} も参照されたい。

\begin{lemma}
\label{descent-lemma-descend-reduced}
$f : X \to Y$ をスキームの平坦全射とする。
$X$ が被約ならば $Y$ も被約である。
\end{lemma}

\begin{proof}
局所環を調べ（Schemes, Definition
\ref{schemes-definition-reduced}）、Algebra,
Lemma \ref{algebra-lemma-descent-reduced} を用いれば従う。
\end{proof}

\begin{lemma}
\label{descent-lemma-descend-regular}
$f : X \to Y$ を代数空間の射とする。$f$ が局所有限表示、平坦、
かつ全射であり、$X$ が正則ならば $Y$ も正則である。
\end{lemma}

\begin{proof}
本補題は次の代数的主張に帰着する：$A \to B$ が忠実平坦かつ
有限表示な環準同型で、$B$ が Noether かつ正則ならば、
$A$ も Noether かつ正則である。$A$ が Noether であることは
Algebra, Lemma \ref{algebra-lemma-descent-Noetherian} から、
正則であることは Algebra,
Lemma \ref{algebra-lemma-flat-under-regular} から分かる。
\end{proof}









\section{スキームの芽}
\label{descent-section-germs}

\begin{definition}
\label{descent-definition-germs}
スキームの芽を次のように定義する。
\begin{enumerate}
\item スキーム $X$ と点 $x \in X$ からなる組 $(X, x)$ を
{\it $x$ における $X$ の芽}という。
\item {\it 芽の射} $f : (X, x) \to (S, s)$ とは、$U \subset X$ が
$x$ の開近傍であるときの、$f(x) = s$ を満たすスキームの射
$f : U \to S$ の同値類である。このような二つの射 $f$, $f'$ が
同値であるとは、$x$ のある開近傍上で $f$ と $f'$ が一致することをいう。
\item {\it 芽の射の合成}は代表元を合成することにより定義する
（これは良定義である）。
\end{enumerate}
\end{definition}

\noindent
先へ進む前に、もう一つ定義が必要である。

\begin{definition}
\label{descent-definition-etale-morphism-germs}
$f : (X, x) \to (S, s)$ を芽の射とする。$f$ の代表元
$f : U \to S$ で、スキームの \'etale 射
（resp.\ 平滑射）であるものが存在するとき、$f$ は
{\it \'etale}（resp.\ {\it 平滑}）であるという。
\end{definition}








\section{芽の局所的性質}
\label{descent-section-properties-germs-local}

\begin{definition}
\label{descent-definition-local-at-point}
$\mathcal{P}$ をスキームの芽の性質とする。
任意の \'etale（resp.\ 平滑）な芽の射 $(U', u') \to (U, u)$ に対して
$\mathcal{P}(U, u) \Leftrightarrow \mathcal{P}(U', u')$ が成り立つとき、
$\mathcal{P}$ は {\it \'etale 局所的}
（resp.\ {\it smooth 局所的}）であるという。
\end{definition}

\noindent
$(X, x)$ をスキームの芽とする。$x$ における $X$ の次元とは、
$X$ における $x$ の開近傍の次元の最小値であり、十分小さい任意の
開近傍はこの次元をもつ。したがって、これは芽の同型類の不変量である。
これを単に $\dim_x(X)$ と書く。次の補題は、主張
$\dim_x(X) = d$ が芽の \'etale 局所的性質であることを示す。

\begin{lemma}
\label{descent-lemma-dimension-at-point-local}
$f : U \to V$ をスキームの \'etale 射とする。
$u \in U$ とし、$v = f(u)$ とおく。このとき
$\dim_u(U) = \dim_v(V)$ である。
\end{lemma}

\begin{proof}
主張に現れる $\dim_u(U)$ は Topology, Definition
\ref{topology-definition-Krull} で定義された $u$ における $U$ の
次元、すなわち $U$ における $u$ の開近傍の Krull 次元の最小値である。
$\dim_v(V)$ についても同様である。

\medskip\noindent
まず $\dim_v(V) \geq \dim_u(U)$ を示す。
$V'$ を $V$ における $v$ の開近傍とする。
$f^{-1}(V')$ に含まれる $U$ における $u$ の開近傍 $U'$ で
$\dim_u(U) = \dim(U')$ を満たすものが存在する。
$Z_0 \subset Z_1 \subset \ldots \subset Z_n$ を $U'$ の
既約閉部分スキームの鎖とする。$\xi_i \in Z_i$ を生成点とすれば、
特殊化の列
$\xi_n \leadsto \xi_{n - 1} \leadsto \ldots \leadsto \xi_0$.
がある。これにより $V'$ における特殊化の列
$f(\xi_n) \leadsto f(\xi_{n - 1}) \leadsto \ldots \leadsto f(\xi_0)$
を得る。$f$ のファイバーは離散的なので、$i \not = j$ ならば
$f(\xi_j) \not = f(\xi_i)$ であることに注意する
（Morphisms, Lemma \ref{morphisms-lemma-etale-over-field} 参照）。
したがって $\dim(V') \geq n$ である。これから形式的に
$\dim_v(V) \geq \dim_u(U)$ が従う。

\medskip\noindent
次に $\dim_u(U) \geq \dim_v(V)$ を示す。
$U'$ を $U$ における $u$ の開近傍とする。Morphisms,
Lemma \ref{morphisms-lemma-fppf-open} により
$V' = f(U')$ は $v$ の開近傍であることに注意する。
したがって $\dim(V') \geq \dim_v(V)$ である。
$Z_0 \subset Z_1 \subset \ldots \subset Z_n$ を $V'$ の
既約閉部分スキームの鎖とする。$\xi_i \in Z_i$ を生成点とすれば、
特殊化の列
$\xi_n \leadsto \xi_{n - 1} \leadsto \ldots \leadsto \xi_0$.
がある。$\xi_0 \in f(U')$ なので、$f(\eta_0) = \xi_0$ を満たす点
$\eta_0 \in U'$ をとれる。局所環の準同型
$$
\mathcal{O}_{V', \xi_0} \longrightarrow \mathcal{O}_{U', \eta_0}
$$
を考える。これは Morphisms, Lemma
\ref{morphisms-lemma-etale-flat} により平坦な局所環準同型である。
Schemes, Lemma \ref{schemes-lemma-specialize-points} により、
点 $\xi_i$ は左辺の環の素イデアルに対応する。
したがって、表示した環準同型に下降定理
（Algebra, Section \ref{algebra-section-going-up} 参照）を適用すると、
$U'$ における特殊化の列
$\eta_n \leadsto \eta_{n - 1} \leadsto \ldots \leadsto \eta_0$
で、$f$ により列
$\xi_n \leadsto \xi_{n - 1} \leadsto \ldots \leadsto \xi_0$
へ写るものが得られる。これより
$\dim_u(U) \geq \dim_v(V)$ が従う。
\end{proof}

\noindent
$(X, x)$ をスキームの芽とする。局所環 $\mathcal{O}_{X, x}$ の
同型類は芽の不変量である。次の補題は、性質
$\dim(\mathcal{O}_{X, x}) = d$ が芽の \'etale 局所的性質であることを
述べる。

\begin{lemma}
\label{descent-lemma-dimension-local-ring-local}
$f : U \to V$ をスキームの \'etale 射とする。
$u \in U$ とし、$v = f(u)$ とおく。このとき
$\dim(\mathcal{O}_{U, u}) = \dim(\mathcal{O}_{V, v})$.
\end{lemma}

\begin{proof}
証明すべき代数的主張は次のものである：
$A \to B$ が \'etale 環準同型で、$\mathfrak q$ が
$\mathfrak p \subset A$ の上にある $B$ の素イデアルならば
$\dim(A_{\mathfrak p}) = \dim(B_{\mathfrak q})$.
これは More on Algebra,
Lemma \ref{more-algebra-lemma-dimension-etale-extension} である。
\end{proof}

\noindent
$(X, x)$ をスキームの芽とする。局所環 $\mathcal{O}_{X, x}$ の
同型類は芽の不変量である。次の補題は、「$\mathcal{O}_{X, x}$ は
正則である」という性質が芽の \'etale 局所的性質であることを述べる。

\begin{lemma}
\label{descent-lemma-regular-local-ring-local}
$f : U \to V$ をスキームの \'etale 射とする。
$u \in U$ とし、$v = f(u)$ とおく。このとき
$\mathcal{O}_{U, u}$ が正則局所環であることと、
$\mathcal{O}_{V, v}$ が正則局所環であることとは同値である。
\end{lemma}

\begin{proof}
証明すべき代数的主張は次のものである：
$A \to B$ が \'etale 環準同型で、$\mathfrak q$ が
$\mathfrak p \subset A$ の上にある $B$ の素イデアルならば、
$A_{\mathfrak p}$ が正則であることと $B_{\mathfrak q}$ が
正則であることとは同値である。これは More on Algebra,
Lemma \ref{more-algebra-lemma-regular-etale-extension} である。
\end{proof}







\section{標的上局所的な射の性質}
\label{descent-section-descending-properties-morphisms}

\noindent
$f : X \to Y$ をスキームの射とする。
$g : Y' \to Y$ をスキームの射とし、
$f' : X' \to Y'$ を $g$ による $f$ の基底変換とする：
$$
\xymatrix{
X' \ar[d]_{f'} \ar[r]_{g'} & X \ar[d]^f \\
Y' \ar[r]^g & Y
}
$$
$\mathcal{P}$ をスキームの射の性質とする。このとき
(a) $\mathcal{P}(f) \Rightarrow \mathcal{P}(f')$ が成り立つか、
また逆向きの (b) $\mathcal{P}(f') \Rightarrow \mathcal{P}(f)$ が
成り立つかを問うことができる。$g$ が平坦なとき常に (a) が
成り立つならば、$\mathcal{P}$ は平坦基底変換で保存されるという。
$g$ が全射平坦なとき常に (b) が成り立つならば、
$\mathcal{P}$ は全射平坦基底変換を通じて降下するという。
$\mathcal{P}$ が平坦基底変換で保存され、かつ全射平坦基底変換を
通じて降下するとき、$\mathcal{P}$ は標的上平坦局所的であるという。
Section \ref{descent-section-descending-properties} の議論と比較されたい。
これは非常に重要な概念であることが分かるので、次の定義で形式化する。

\begin{definition}
\label{descent-definition-property-morphisms-local}
$\mathcal{P}$ を基底上のスキームの射の性質とする。\newline
$\tau \in \{fpqc, fppf, syntomic, smooth, \etale, Zariski\}$ とする。
任意の $\tau$-被覆 $\{Y_i \to Y\}_{i \in I}$
（Topologies, Section \ref{topologies-section-procedure} 参照）と、
$S$ 上の任意のスキームの射 $f : X \to Y$ に対して次が成り立つとき、
$\mathcal{P}$ は{\it 基底上 $\tau$-局所的}、または
{\it 標的上 $\tau$-局所的}、または
{\it $\tau$-位相について基底上局所的}であるという。
$$
f \text{ は }\mathcal{P}\text{ をもつ}
\Leftrightarrow
\text{各 }Y_i \times_Y X \to Y_i\text{ は }\mathcal{P}\text{ をもつ}.
$$
\end{definition}

\noindent
同型射は常に被覆なので、性質 $\mathcal{P}$ が $X \to Y$ について
成り立つことと、$X \to Y$ と同型な任意の矢 $X' \to Y'$ について
成り立つこととは同値であると分かる（あるいは、そう要請する）。
ある性質が標的上 $\tau$-局所的ならば、$\tau$-被覆に現れる射による
基底変換で保存される。これを形式的に述べる。

\begin{lemma}
\label{descent-lemma-pullback-property-local-target}
$\tau \in \{fpqc, fppf, syntomic, smooth, \etale, Zariski\}$ とする。
$\mathcal{P}$ を標的上 $\tau$-局所的な射の性質とし、
$f : X \to Y$ が性質 $\mathcal{P}$ をもつとする。
$Y' \to Y$ が平坦、resp.\ 平坦かつ局所有限表示、
resp.\ syntomic、resp.\ \'etale、resp.\ 開埋め込みである任意の射ならば、
$f$ の基底変換 $f' : Y' \times_Y X \to Y'$ は性質
$\mathcal{P}$ をもつ。
\end{lemma}

\begin{proof}
$Y' \to Y$ を $\tau$-被覆をなす射の族に組み込めるからである。
\end{proof}

\noindent
上の結果から得られる、しばしば用いられる簡単な帰結は次である。
$f : X \to Y$ が標的上 $\tau$-局所的な性質 $\mathcal{P}$ をもち、
ある開部分スキーム $V \subset Y$ に対して $f(X) \subset V$ ならば、
誘導される射 $X \to V$ も $\mathcal{P}$ をもつ。
証明：$V \to Y$ による $f$ の基底変換は $X \to V$ である。

\begin{lemma}
\label{descent-lemma-largest-open-of-the-base}
$\tau \in \{fppf, syntomic, smooth, \etale\}$ とする。
$\mathcal{P}$ を標的上 $\tau$-局所的な射の性質とする。
スキームの任意の射 $f : X \to Y$ に対し、制限
$X_{W(f)} \to W(f)$ が $\mathcal{P}$ をもつような最大の開集合
$W(f) \subset Y$ が存在する。さらに、
\begin{enumerate}
\item $g : Y' \to Y$ が平坦かつ局所有限表示、syntomic、平滑、
または \'etale であり、基底変換 $f' : X_{Y'} \to Y'$ が
$\mathcal{P}$ をもつならば、$g(Y') \subset W(f)$ である。
\item $g : Y' \to Y$ が平坦かつ局所有限表示、syntomic、平滑、
または \'etale ならば、$W(f') = g^{-1}(W(f))$ である。
\item $\{g_i : Y_i \to Y\}$ が $\tau$-被覆ならば
$g_i^{-1}(W(f)) = W(f_i)$ である。ここで $f_i$ は
$Y_i \to Y$ による $f$ の基底変換である。
\end{enumerate}
\end{lemma}

\begin{proof}
次の性質をもつ射 $g : Y' \to Y$ の像 $g(Y') \subset Y$ の
合併を $W$ とする：
\begin{enumerate}
\item $g$ は平坦かつ局所有限表示、syntomic、平滑、または \'etale
である。
\item 基底変換 $Y' \times_{g, Y} X \to Y'$ は性質 $\mathcal{P}$ をもつ。
\end{enumerate}
このような射 $g$ は開写像なので
（Morphisms, Lemma \ref{morphisms-lemma-fppf-open} 参照）、
$W \subset Y$ は $Y$ の開部分集合である。$\mathcal{P}$ は
$\tau$ 位相について局所的であり、各引き戻しが $\mathcal{P}$ をもつ
$W$ の被覆 $\{Y' \to W\}$ が与えられているので、制限
$X_W \to W$ は性質 $\mathcal{P}$ をもつ。これで存在が示され、
$W(f)$ が性質 (1) をもつことも示された。(2) を見るため、
$\mathcal{P}$ は平坦かつ局所有限表示、syntomic、平滑、または
\'etale な射による基底変換で保存されるので
（補題 \ref{descent-lemma-pullback-property-local-target} 参照）、
$W(f') \supset g^{-1}(W(f))$ であることに注意する。
他方、$Y'' \subset Y'$ が開集合で $X_{Y''} \to Y''$ が性質
$\mathcal{P}$ をもつならば、構成により $Y'' \to Y$ は $W$ を経由する。
すなわち $Y'' \subset g^{-1}(W(f))$ である。これで (2) が示された。
主張 (3) は (2) から従う。実際、$\tau$-被覆の定義により各射
$Y_i \to Y$ は平坦かつ局所有限表示、syntomic、平滑、または
\'etale である。
\end{proof}

\begin{lemma}
\label{descent-lemma-descending-properties-morphisms}
$\mathcal{P}$ を基底上のスキームの射の性質とする。
$\tau \in \{fpqc, fppf, \etale, smooth, syntomic\}$ とし、次を仮定する。
\begin{enumerate}
\item この性質は、$\tau$ が fpqc、fppf、\'etale、smooth、syntomic
のいずれであるかに応じて、それぞれ平坦、平坦かつ局所有限表示、
\'etale、平滑、syntomic な基底変換で保存される
（Schemes, Definition
\ref{schemes-definition-preserved-by-base-change} と比較せよ）。
\item この性質は基底上 Zariski 局所的である。
\item アフィン・スキームの任意の全射 $S' \to S$ が、
$\tau$ が fpqc、fppf、\'etale、smooth、syntomic のいずれであるかに
応じて、それぞれ平坦、有限表示の平坦射、\'etale、平滑、syntomic
であるとする。スキームの任意の射 $f : X \to S$ に対し、
基底変換 $f' : X' = S' \times_S X \to S'$ が性質 $\mathcal{P}$ を
もつならば、$f$ も性質 $\mathcal{P}$ をもつ。
\end{enumerate}
このとき $\mathcal{P}$ は基底上 $\tau$-局所的である。
\end{lemma}

\begin{proof}
これは $\tau$-被覆の定義からほとんど直ちに従う。
Topologies, Definition
\ref{topologies-definition-fpqc-covering}
\ref{topologies-definition-fppf-covering}
\ref{topologies-definition-etale-covering}
\ref{topologies-definition-smooth-covering} または
\ref{topologies-definition-syntomic-covering}
と、Topologies, Lemma
\ref{topologies-lemma-fpqc-affine},
\ref{topologies-lemma-fppf-affine},
\ref{topologies-lemma-etale-affine},
\ref{topologies-lemma-smooth-affine} または
\ref{topologies-lemma-syntomic-affine} を参照。
詳細は省略する。
\end{proof}

\begin{remark}
\label{descent-remark-descending-properties-morphisms-standard}
（これは上の Remark \ref{descent-remark-descending-properties-standard} の
繰り返しである。）上の補題
\ref{descent-lemma-descending-properties-morphisms} で
$\tau = smooth$ ならば、条件 (3) の射は（全射な）標準平滑射であると
仮定してよい。$\tau = syntomic$ または $\tau = \etale$ の場合も
同様である。
\end{remark}




\section{標的上の fpqc 位相について局所的な射の性質}
\label{descent-section-descending-properties-morphisms-fpqc}

\noindent
本節では、fpqc 位相について基底上局所的なスキームの射の性質を
多数求める。これとは対照的に、
Examples, Section \ref{examples-section-non-descending-property-projective}
では「射影的」および「準射影的」という性質が Zariski 位相においてさえ
基底上局所的でないことを示す。

\begin{lemma}
\label{descent-lemma-descending-property-quasi-compact}
性質 $\mathcal{P}(f) =$「$f$ は準コンパクトである」は
基底上 fpqc 局所的である。
\end{lemma}

\begin{proof}
準コンパクト射の基底変換は準コンパクトである。Schemes,
Lemma \ref{schemes-lemma-quasi-compact-preserved-base-change} を参照。
準コンパクトであることは基底上 Zariski 局所的である。Schemes,
Lemma \ref{schemes-lemma-quasi-compact-affine} を参照。
最後に、$S' \to S$ をアフィン・スキームの平坦全射とし、
$f : X \to S$ を射とする。基底変換 $f' : X' \to S'$ が
準コンパクトであると仮定する。このとき $X'$ は準コンパクトであり、
$X' \to X$ は全射である。したがって $X$ は準コンパクトであり、
$f$ も準コンパクトである。ゆえに補題
\ref{descent-lemma-descending-properties-morphisms} を適用でき、結論を得る。
\end{proof}

\begin{lemma}
\label{descent-lemma-descending-property-quasi-separated}
性質 $\mathcal{P}(f) =$「$f$ は準分離である」は
基底上 fpqc 局所的である。
\end{lemma}

\begin{proof}
準分離射の任意の基底変換は準分離である。Schemes,
Lemma \ref{schemes-lemma-separated-permanence} を参照。
準分離であることは基底上 Zariski 局所的である
（定義、または Schemes,
Lemma \ref{schemes-lemma-characterize-quasi-separated} による）。
最後に、$S' \to S$ をアフィン・スキームの平坦全射とし、
$f : X \to S$ を射とする。基底変換 $f' : X' \to S'$ が
準分離であると仮定する。これは
$\Delta' : X' \to X'\times_{S'} X'$ が準コンパクトであることを意味する。
$\Delta'$ は $S' \to S$ による
$\Delta : X \to X \times_S X$ の基底変換であることに注意する。
補題 \ref{descent-lemma-descending-property-quasi-compact} により
$\Delta$ は準コンパクトであり、したがって $f$ は準分離である。
ゆえに補題 \ref{descent-lemma-descending-properties-morphisms} を適用でき、
結論を得る。
\end{proof}

\begin{lemma}
\label{descent-lemma-descending-property-universally-closed}
性質 $\mathcal{P}(f) =$「$f$ は普遍閉である」は
基底上 fpqc 局所的である。
\end{lemma}

\begin{proof}
定義により、普遍閉射の基底変換は普遍閉である。
普遍閉であることは基底上 Zariski 局所的である
（定義、または Morphisms, Lemma
\ref{morphisms-lemma-universally-closed-local-on-the-base} による）。
最後に、$S' \to S$ をアフィン・スキームの平坦全射とし、
$f : X \to S$ を射とする。基底変換 $f' : X' \to S'$ が
普遍閉であると仮定する。$T \to S$ を任意の射とし、図式
$$
\xymatrix{
X' \ar[d] &
S' \times_S T \times_S X \ar[d] \ar[r] \ar[l] &
T \times_S X \ar[d] \\
S' &
S' \times_S T \ar[r] \ar[l] &
T
}
$$
を考える。二つの正方形はいずれも Cartesian である。
したがって仮定により中央の垂直射は閉写像である。
右側の水平射は $S' \to S$ の基底変換なので、平坦、準コンパクト、
かつ全射である。ゆえに $T$ の部分集合が閉であることと、
$S' \times_S T$ における逆像が閉であることとは同値である。
Morphisms, Lemma \ref{morphisms-lemma-fpqc-quotient-topology} を参照。
容易な図式追跡により右側の垂直射も閉写像であることが分かり、
$X \to S$ は普遍閉であると結論する。
したがって補題 \ref{descent-lemma-descending-properties-morphisms} を
適用でき、結論を得る。
\end{proof}

\begin{lemma}
\label{descent-lemma-descending-property-universally-open}
性質 $\mathcal{P}(f) =$「$f$ は普遍開である」は
基底上 fpqc 局所的である。
\end{lemma}

\begin{proof}
証明は補題 \ref{descent-lemma-descending-property-universally-closed} の
証明と同じである。
\end{proof}

\begin{lemma}
\label{descent-lemma-descending-property-universally-submersive}
性質 $\mathcal{P}(f) =$「$f$ は普遍的に商写像である」は
基底上 fpqc 局所的である。
\end{lemma}

\begin{proof}
証明は補題 \ref{descent-lemma-descending-property-universally-closed} の
証明と同じである。ただし、Morphisms,
Lemma \ref{morphisms-lemma-fpqc-quotient-topology} により
準コンパクト平坦全射が普遍的に商写像であることを用いる。
\end{proof}

\begin{lemma}
\label{descent-lemma-descending-property-separated}
性質 $\mathcal{P}(f) =$「$f$ は分離的である」は
基底上 fpqc 局所的である。
\end{lemma}

\begin{proof}
分離射の基底変換は分離的である。Schemes,
Lemma \ref{schemes-lemma-separated-permanence} を参照。
分離的であることは基底上 Zariski 局所的である
（定義、または Schemes,
Lemma \ref{schemes-lemma-characterize-separated} による）。
最後に、$S' \to S$ をアフィン・スキームの平坦全射とし、
$f : X \to S$ を射とする。基底変換 $f' : X' \to S'$ が
分離的であると仮定する。これは
$\Delta' : X' \to X'\times_{S'} X'$ が閉埋め込みであり、
したがって普遍閉であることを意味する。
$\Delta'$ は $S' \to S$ による
$\Delta : X \to X \times_S X$ の基底変換であることに注意する。
補題 \ref{descent-lemma-descending-property-universally-closed} により
$\Delta$ は普遍閉である。また、これは埋め込みなので
（Schemes, Lemma \ref{schemes-lemma-diagonal-immersion}）、
$\Delta$ は閉埋め込みである。したがって $f$ は分離的である。
ゆえに補題 \ref{descent-lemma-descending-properties-morphisms} を適用でき、
結論を得る。
\end{proof}

\begin{lemma}
\label{descent-lemma-descending-property-surjective}
性質 $\mathcal{P}(f) =$「$f$ は全射である」は
基底上 fpqc 局所的である。
\end{lemma}

\begin{proof}
明らかである。
\end{proof}

\begin{lemma}
\label{descent-lemma-descending-property-dominant}
性質 $\mathcal{P}(f) =$「$f$ は準コンパクトかつ支配的である」は
基底上 fpqc 局所的である。
\end{lemma}

\begin{proof}
Morphisms, Lemma \ref{morphisms-lemma-base-change-dominant} により、
準コンパクトな支配的射は平坦引き戻しで保存される。
逆向きはより容易である。実際、補題
\ref{descent-lemma-descending-property-quasi-compact} により、
準コンパクト性は基底上 fpqc 局所的である。
支配的であることは明らかに基底上 Zariski 局所的である。
最後に、$S' \to S$ をスキームの全射（平坦とは限らない）とし、
$f : X \to S$ を射とする。基底変換 $f' : X' \to S'$ が
支配的であると仮定する。このとき $X' \to S'\to S$ は二つの
支配的射の合成なので支配的である。これは合成
$X' \to X \to S$ でもあるから、$X\to S$ は支配的である。
\end{proof}

\noindent
支配的射は一般には平坦引き戻しで保存されないが、次の特別な場合は
述べておく価値がある。補題
\ref{descent-lemma-descending-fppf-property-dominant} も参照されたい。

\begin{lemma}
\label{descent-lemma-descending-property-dominant-over-field}
$E/k$ を体拡大とする。このとき $k$ 上の射 $X \to Y$ が
支配的であることと、引き戻し $X_E \to Y_E$ が支配的であることとは
同値である。
\end{lemma}

\begin{proof}
Morphisms,
Lemma \ref{morphisms-lemma-scheme-over-field-universally-open} により、
射 $\Spec(E) \to \Spec(k)$ は普遍開である。したがって
$Y_E \to Y$ は開写像である。ゆえに $X \to Y$ が支配的ならば、
Morphisms, Lemma \ref{morphisms-lemma-open-base-change-dominant} により
$X_E \to Y_E$ は支配的である。
逆に $X_E \to Y_E$ が支配的であると仮定する。
Morphisms, Lemma \ref{morphisms-lemma-base-change-surjective} により
射 $Y_E \to Y$ は全射なので、合成 $X_E \to Y_E \to Y$ は
支配的である。これは合成 $X_E \to X \to Y$ でもあるから、
$X \to Y$ は支配的である。
\end{proof}

\begin{lemma}
\label{descent-lemma-descending-property-universally-injective}
性質 $\mathcal{P}(f) =$「$f$ は普遍単射である」は
基底上 fpqc 局所的である。
\end{lemma}

\begin{proof}
普遍単射の基底変換は普遍単射である（これは形式的である）。
普遍単射であることは基底上 Zariski 局所的であり、定義から明らかである。
最後に、$S' \to S$ をアフィン・スキームの平坦全射とし、
$f : X \to S$ を射とする。基底変換 $f' : X' \to S'$ が
普遍単射であると仮定する。$K$ を体とし、
$a, b : \Spec(K) \to X$ を $f \circ a = f \circ b$ を満たす
二つの射とする。$S' \to S$ は全射なので、Schemes,
Section \ref{schemes-section-points} の議論により、体拡大 $K'/K$ と
射 $\Spec(K') \to S'$ で、次の実線の図式を可換にするものが存在する。
$$
\xymatrix{
\Spec(K') \ar[rrd] \ar@{-->}[rd]_{a', b'} \ar[dd] \\
 &
X' \ar[r] \ar[d] &
S' \ar[d] \\
\Spec(K) \ar[r]^{a, b} &
X \ar[r] &
S
}
$$
正方形は Cartesian なので、図式を可換にする二つの点線の矢
$a'$, $b'$ が得られる。$X' \to S'$ は普遍単射なので、Morphisms,
Lemma \ref{morphisms-lemma-universally-injective} により $a' = b'$ である。
これは明らかに $a = b$ を強制する
（Schemes, Section \ref{schemes-section-points} の議論による）。
したがって補題 \ref{descent-lemma-descending-properties-morphisms} を
適用でき、結論を得る。

\medskip\noindent
別の証明として、普遍単射とは対角射が全射である射であるという特徴付け
（Morphisms, Lemma \ref{morphisms-lemma-universally-injective}）を
用いることもできる。このとき本補題は、全射であるという性質が基底上
fpqc 局所的であることから従う。補題
\ref{descent-lemma-descending-property-surjective} を参照。
（ヒント：対角射の基底変換は基底変換の対角射であることを用いる。）
\end{proof}

\begin{lemma}
\label{descent-lemma-descending-property-universal-homeomorphism}
性質 $\mathcal{P}(f) =$「$f$ は普遍同相である」は
基底上 fpqc 局所的である。
\end{lemma}

\begin{proof}
補題 \ref{descent-lemma-descending-property-universally-closed} とまったく
同じ方法で証明できる。別法として、位相空間の写像が同相写像であることと、
単射、全射、かつ開写像であることとの同値性を用いてもよい。
したがって普遍同相であることは、全射、普遍単射、かつ普遍開であることと
同値である。ゆえに本補題は
Lemmas \ref{descent-lemma-descending-property-surjective},
\ref{descent-lemma-descending-property-universally-injective} および
\ref{descent-lemma-descending-property-universally-open} から従う。
\end{proof}

\begin{lemma}
\label{descent-lemma-descending-property-locally-finite-type}
性質 $\mathcal{P}(f) =$「$f$ は局所有限型である」は
基底上 fpqc 局所的である。
\end{lemma}

\begin{proof}
局所有限型であることは基底変換で保存される。Morphisms,
Lemma \ref{morphisms-lemma-base-change-finite-type} を参照。
局所有限型であることは基底上 Zariski 局所的である。Morphisms,
Lemma \ref{morphisms-lemma-locally-finite-type-characterize} を参照。
最後に、$S' \to S$ をアフィン・スキームの平坦全射とし、
$f : X \to S$ を射とする。基底変換 $f' : X' \to S'$ が
局所有限型であると仮定する。$U \subset X$ をアフィン開集合とする。
すると $U' = S' \times_S U$ はアフィンで、$S'$ 上有限型である。
$S = \Spec(R)$,
$S' = \Spec(R')$,
$U = \Spec(A)$ および
$U' = \Spec(A')$ と書く。
$R \to R'$ は忠実平坦であり、$A' = R' \otimes_R A$、
$R' \to A'$ は有限型である。$R \to A$ が有限型であることを
示す必要がある。これは Algebra,
Lemma \ref{algebra-lemma-finite-type-descends} の結果である。
したがって $f$ は局所有限型である。ゆえに補題
\ref{descent-lemma-descending-properties-morphisms} を適用でき、結論を得る。
\end{proof}

\begin{lemma}
\label{descent-lemma-descending-property-locally-finite-presentation}
性質 $\mathcal{P}(f) =$「$f$ は局所有限表示である」は
基底上 fpqc 局所的である。
\end{lemma}

\begin{proof}
局所有限表示であることは基底変換で保存される。Morphisms,
Lemma \ref{morphisms-lemma-base-change-finite-presentation} を参照。
局所有限型であることは基底上 Zariski 局所的である。Morphisms,
Lemma \ref{morphisms-lemma-locally-finite-presentation-characterize} を参照。
最後に、$S' \to S$ をアフィン・スキームの平坦全射とし、
$f : X \to S$ を射とする。基底変換 $f' : X' \to S'$ が
局所有限表示であると仮定する。$U \subset X$ をアフィン開集合とする。
すると $U' = S' \times_S U$ はアフィンで、$S'$ 上有限型である。
$S = \Spec(R)$,
$S' = \Spec(R')$,
$U = \Spec(A)$ および
$U' = \Spec(A')$ と書く。
$R \to R'$ は忠実平坦であり、$A' = R' \otimes_R A$、
$R' \to A'$ は有限表示である。$R \to A$ が有限表示であることを
示す必要がある。これは Algebra,
Lemma \ref{algebra-lemma-finite-presentation-descends} の結果である。
したがって $f$ は局所有限表示である。ゆえに補題
\ref{descent-lemma-descending-properties-morphisms} を適用でき、結論を得る。
\end{proof}

\begin{lemma}
\label{descent-lemma-descending-property-finite-type}
性質 $\mathcal{P}(f) =$「$f$ は有限型である」は
基底上 fpqc 局所的である。
\end{lemma}

\begin{proof}
Lemmas \ref{descent-lemma-descending-property-quasi-compact} と
\ref{descent-lemma-descending-property-locally-finite-type} を組み合わせる。
\end{proof}

\begin{lemma}
\label{descent-lemma-descending-property-finite-presentation}
性質 $\mathcal{P}(f) =$「$f$ は有限表示である」は
基底上 fpqc 局所的である。
\end{lemma}

\begin{proof}
Lemmas \ref{descent-lemma-descending-property-quasi-compact}、
\ref{descent-lemma-descending-property-quasi-separated} および
\ref{descent-lemma-descending-property-locally-finite-presentation} を
組み合わせる。
\end{proof}

\begin{lemma}
\label{descent-lemma-descending-property-proper}
性質 $\mathcal{P}(f) =$「$f$ は固有である」は
基底上 fpqc 局所的である。
\end{lemma}

\begin{proof}
Lemmas \ref{descent-lemma-descending-property-universally-closed}、
\ref{descent-lemma-descending-property-separated} および
\ref{descent-lemma-descending-property-finite-type} を組み合わせれば従う。
\end{proof}

\begin{lemma}
\label{descent-lemma-descending-property-flat}
性質 $\mathcal{P}(f) =$「$f$ は平坦である」は
基底上 fpqc 局所的である。
\end{lemma}

\begin{proof}
平坦性は任意の基底変換で保存される。Morphisms,
Lemma \ref{morphisms-lemma-base-change-flat} を参照。
定義により、平坦性は基底上 Zariski 局所的である。
最後に、$S' \to S$ をアフィン・スキームの平坦全射とし、
$f : X \to S$ を射とする。基底変換 $f' : X' \to S'$ が
平坦であると仮定する。$U \subset X$ をアフィン開集合とする。
すると $U' = S' \times_S U$ はアフィンである。
$S = \Spec(R)$,
$S' = \Spec(R')$,
$U = \Spec(A)$ および
$U' = \Spec(A')$ と書く。
$R \to R'$ は忠実平坦であり、$A' = R' \otimes_R A$、
$R' \to A'$ は平坦である。目標は $R \to A$ が平坦であることを
示すことである。これは Algebra,
Lemma \ref{algebra-lemma-flatness-descends} から直ちに従う。
ゆえに $f$ は平坦である。したがって補題
\ref{descent-lemma-descending-properties-morphisms} を適用でき、結論を得る。
\end{proof}

\begin{lemma}
\label{descent-lemma-descending-property-open-immersion}
性質 $\mathcal{P}(f) =$「$f$ は開埋め込みである」は
基底上 fpqc 局所的である。
\end{lemma}

\begin{proof}
開埋め込みであるという性質は基底変換で保存される。Schemes,
Lemma \ref{schemes-lemma-base-change-immersion} を参照。
開埋め込みであるという性質は基底上 Zariski 局所的である
（これは明らかである）。

\medskip\noindent
$S' \to S$ をアフィン・スキームの平坦全射とし、
$f : X \to S$ を射とする。基底変換 $f' : X' \to S'$ が
開埋め込みであると仮定する。$f$ が開埋め込みであると主張する。
このとき $f'$ は普遍開かつ普遍単射である。したがって補題
\ref{descent-lemma-descending-property-universally-open} により $f$ は普遍開、
補題 \ref{descent-lemma-descending-property-universally-injective} により
普遍単射である。特に $f(X) \subset S$ は開である。
任意のアフィン開集合 $U \subset f(X)$ に対して
$f^{-1}(U) \to U$ が同型であることを示せれば、$f$ は開埋め込みとなり
証明が完了する。
$U' \subset S'$ を $U$ の逆像とする。このとき $U' \to U$ は
アフィン・スキームの忠実平坦射であり、
$(f')^{-1}(U') \to U'$ は同型である
（$U$ の選び方により $f'(X')$ は $U'$ を含む）。
したがって次段落で論じる場合に帰着する。

\medskip\noindent
$S' \to S$ をアフィン・スキームの平坦全射、
$f : X \to S$ を射とし、基底変換 $f' : X' \to S'$ が
同型であると仮定する。$f$ も同型であることを示す必要がある。
$f$ が全射、普遍単射、かつ普遍開であることは明らかである
（後二者については上の議論を参照）。したがって $f$ は全単射、
すなわち $f$ は同相写像である。ゆえに Morphisms,
Lemma \ref{morphisms-lemma-homeomorphism-affine} により $f$ は
アフィンである。また
$$
\mathcal{O}(S') \to
\mathcal{O}(X') =
\mathcal{O}(S') \otimes_{\mathcal{O}(S)} \mathcal{O}(X)
$$
は同型であり、$\mathcal{O}(S) \to \mathcal{O}(S')$ は忠実平坦なので、
$\mathcal{O}(S) \to \mathcal{O}(X)$ は同型であることが従う。
したがって $f$ は同型であり、上の主張の証明が完了する。
ゆえに補題 \ref{descent-lemma-descending-properties-morphisms} を適用でき、
結論を得る。
\end{proof}

\begin{lemma}
\label{descent-lemma-descending-property-isomorphism}
性質 $\mathcal{P}(f) =$「$f$ は同型である」は
基底上 fpqc 局所的である。
\end{lemma}

\begin{proof}
Lemmas \ref{descent-lemma-descending-property-surjective} と
\ref{descent-lemma-descending-property-open-immersion} を組み合わせる。
\end{proof}

\begin{lemma}
\label{descent-lemma-descending-property-affine}
性質 $\mathcal{P}(f) =$「$f$ はアフィンである」は
基底上 fpqc 局所的である。
\end{lemma}

\begin{proof}
アフィン射の基底変換はアフィンである。Morphisms,
Lemma \ref{morphisms-lemma-base-change-affine} を参照。
アフィンであることは基底上 Zariski 局所的である。Morphisms,
Lemma \ref{morphisms-lemma-characterize-affine} を参照。
最後に、$g : S' \to S$ をアフィン・スキームの平坦全射とし、
$f : X \to S$ を射とする。基底変換 $f' : X' \to S'$ が
アフィンであると仮定する。言い換えれば $X'$ はアフィンであり、
$X' = \Spec(A')$ と書ける。また $S = \Spec(R)$ および
$S' = \Spec(R')$ と書く。$X$ がアフィンであることを示す必要がある。

\medskip\noindent
Lemmas \ref{descent-lemma-descending-property-quasi-compact} と
\ref{descent-lemma-descending-property-separated} により、
$X \to S$ は分離的かつ準コンパクトである。したがって
$f_*\mathcal{O}_X$ は $\mathcal{O}_S$-代数の準連接層である。
Schemes, Lemma \ref{schemes-lemma-push-forward-quasi-coherent} を参照。
ゆえに、ある $R$-代数 $A$ に対して
$f_*\mathcal{O}_X = \widetilde{A}$ である。もちろん実際には
$A = \Gamma(X, \mathcal{O}_X)$ である。
また、平坦基底変換により
（たとえば Cohomology of Schemes,
Lemma \ref{coherent-lemma-flat-base-change-cohomology} 参照）、
$g^*f_*\mathcal{O}_X = f'_*\mathcal{O}_{X'}$ である。
言い換えれば $A' = R' \otimes_R A$ である。標準射
$$
X \longrightarrow \Spec(A)
$$
を考える。これは Schemes, Lemma
\ref{schemes-lemma-morphism-into-affine} の $S$ 上の射である。
上の議論により、この射の $S'$ への基底変換は同型である。
したがって補題 \ref{descent-lemma-descending-property-isomorphism} により
この射自身も同型である。ゆえに補題
\ref{descent-lemma-descending-properties-morphisms} を適用でき、結論を得る。
\end{proof}

\begin{lemma}
\label{descent-lemma-descending-property-closed-immersion}
性質 $\mathcal{P}(f) =$「$f$ は閉埋め込みである」は
基底上 fpqc 局所的である。
\end{lemma}

\begin{proof}
$f : X \to Y$ をスキームの射とし、
$\{Y_i \to Y\}$ を fpqc 被覆とする。各
$f_i : Y_i \times_Y X \to Y_i$ が閉埋め込みであると仮定する。
このとき各 $f_i$ はアフィンである。Morphisms,
Lemma \ref{morphisms-lemma-closed-immersion-affine} を参照。
補題 \ref{descent-lemma-descending-property-affine} により $f$ は
アフィンである。あとは $\mathcal{O}_Y \to f_*\mathcal{O}_X$ が
全射であることを示せばよい。任意の $y \in Y$ に対し、ある $i$ と
$y$ に写る点 $y_i \in Y_i$ が存在する。
Cohomology of Schemes, Lemma
\ref{coherent-lemma-flat-base-change-cohomology} により、層
$f_{i, *}(\mathcal{O}_{Y_i \times_Y X})$ は
$f_*\mathcal{O}_X$ の引き戻しである。仮定により、これは
$\mathcal{O}_{Y_i}$ の商である。したがって
$$
\Big(
\mathcal{O}_{Y, y} \longrightarrow (f_*\mathcal{O}_X)_y
\Big)
\otimes_{\mathcal{O}_{Y, y}} \mathcal{O}_{Y_i, y_i}
$$
は全射である。$\mathcal{O}_{Y_i, y_i}$ は
$\mathcal{O}_{Y, y}$ 上忠実平坦なので、所望のとおり
$\mathcal{O}_{Y, y} \longrightarrow (f_*\mathcal{O}_X)_y$ が
全射であることが従う。
\end{proof}

\begin{lemma}
\label{descent-lemma-descending-property-quasi-affine}
性質 $\mathcal{P}(f) =$「$f$ は準アフィンである」は
基底上 fpqc 局所的である。
\end{lemma}

\begin{proof}
$f : X \to Y$ をスキームの射とし、
$\{g_i : Y_i \to Y\}$ を fpqc 被覆とする。各
$f_i : Y_i \times_Y X \to Y_i$ が準アフィンであると仮定する。
このとき各 $f_i$ は準コンパクトかつ分離的である。
Lemmas \ref{descent-lemma-descending-property-quasi-compact} と
\ref{descent-lemma-descending-property-separated} により、
$f$ は準コンパクトかつ分離的である。
$\mathcal{O}_Y$-代数の層 $\mathcal{A} = f_*\mathcal{O}_X$ を考える。
Schemes, Lemma \ref{schemes-lemma-push-forward-quasi-coherent} により、
これは準連接 $\mathcal{O}_Y$-代数である。標準射
$$
j : X \longrightarrow \underline{\Spec}_Y(\mathcal{A})
$$
を考える。Constructions, Lemma
\ref{constructions-lemma-canonical-morphism} を参照。
平坦基底変換により
（たとえば Cohomology of Schemes,
Lemma \ref{coherent-lemma-flat-base-change-cohomology} 参照）、
$g_i^*f_*\mathcal{O}_X = f_{i, *}\mathcal{O}_{X'}$ である。
ここで $g_i : Y_i \to Y$ は与えられた平坦射である。
したがって $g_i$ による $j$ の基底変換 $j_i$ は、射 $f_i$ に対する
Constructions, Lemma \ref{constructions-lemma-canonical-morphism} の
標準射である。仮定と Morphisms,
Lemma \ref{morphisms-lemma-characterize-quasi-affine} により、
これらの射 $j_i$ はすべて準コンパクトな開埋め込みである。
ゆえに Lemmas \ref{descent-lemma-descending-property-quasi-compact} と
\ref{descent-lemma-descending-property-open-immersion} により、
$j$ は準コンパクトな開埋め込みである。したがって再び Morphisms,
Lemma \ref{morphisms-lemma-characterize-quasi-affine} により
$f$ は準アフィンである。
\end{proof}

\begin{lemma}
\label{descent-lemma-descending-property-quasi-compact-immersion}
性質 $\mathcal{P}(f) =$「$f$ は準コンパクトな埋め込みである」は
基底上 fpqc 局所的である。
\end{lemma}

\begin{proof}
$f : X \to Y$ をスキームの射とし、
$\{Y_i \to Y\}$ を fpqc 被覆とする。
$X_i = Y_i \times_Y X$ と書き、$f_i : X_i \to Y_i$ を
$f$ の基底変換とする。また、与えられた平坦射を
$q_i : Y_i \to Y$ と書く。各 $f_i$ が準コンパクトな埋め込みであると
仮定する。Schemes, Lemma
\ref{schemes-lemma-immersions-monomorphisms} により各 $f_i$ は
分離的である。Lemmas \ref{descent-lemma-descending-property-quasi-compact} と
\ref{descent-lemma-descending-property-separated} により、
$f$ は準コンパクトかつ分離的である。
$X \to Z \to Y$ を $f$ のスキーム論的像を通る分解とする。
Morphisms, Lemma
\ref{morphisms-lemma-quasi-compact-scheme-theoretic-image} により、
閉部分スキーム $Z \subset Y$ は、$f$ が準コンパクトであることから
準連接イデアル層
$\mathcal{I} = \Ker(\mathcal{O}_Y \to f_*\mathcal{O}_X)$
によって切り出される。平坦基底変換により
（たとえば Cohomology of Schemes,
Lemma \ref{coherent-lemma-flat-base-change-cohomology} 参照。
ここで $f$ が分離的であることを用いる）、
$f_{i, *}\mathcal{O}_{X_i}$ は $q_i^*f_*\mathcal{O}_X$ という
引き戻しである。したがって $Y_i \times_Y Z$ は準連接イデアル層
$q_i^*\mathcal{I} =
\Ker(\mathcal{O}_{Y_i} \to f_{i, *}\mathcal{O}_{X_i})$.
によって切り出される。Morphisms,
Lemma \ref{morphisms-lemma-quasi-compact-immersion} により、
射 $X_i \to Y_i \times_Y Z$ は開埋め込みである。
したがって補題 \ref{descent-lemma-descending-property-open-immersion} により
$X \to Z$ は開埋め込みであり、ゆえに所望のとおり $f$ は
埋め込みである（準コンパクトであることはすでに示した）。
\end{proof}

\begin{lemma}
\label{descent-lemma-descending-property-integral}
性質 $\mathcal{P}(f) =$「$f$ は整である」は
基底上 fpqc 局所的である。
\end{lemma}

\begin{proof}
整射であることは、アフィンかつ普遍閉であることと同値である。
Morphisms, Lemma
\ref{morphisms-lemma-integral-universally-closed} を参照。
したがって Lemmas
\ref{descent-lemma-descending-property-universally-closed} と
\ref{descent-lemma-descending-property-affine} を組み合わせれば本補題を得る。
\end{proof}

\begin{lemma}
\label{descent-lemma-descending-property-finite}
性質 $\mathcal{P}(f) =$「$f$ は有限である」は
基底上 fpqc 局所的である。
\end{lemma}

\begin{proof}
有限射であることは、整かつ局所有限型であることと同値である。
Morphisms, Lemma \ref{morphisms-lemma-finite-integral} を参照。
したがって Lemmas
\ref{descent-lemma-descending-property-locally-finite-type} と
\ref{descent-lemma-descending-property-integral} を組み合わせれば従う。
\end{proof}

\begin{lemma}
\label{descent-lemma-descending-property-quasi-finite}
性質 $\mathcal{P}(f) =$「$f$ は局所準有限である」と
$\mathcal{P}(f) =$「$f$ は準有限である」は、
いずれも基底上 fpqc 局所的である。
\end{lemma}

\begin{proof}
$f : X \to S$ をスキームの射とし、$\{S_i \to S\}$ を、
各基底変換 $f_i : X_i \to S_i$ が局所準有限となる fpqc 被覆とする。
「局所有限型であること」は基底上 fpqc 局所的であることをすでに見た
（補題 \ref{descent-lemma-descending-property-locally-finite-type}）。
したがって $f$ は局所有限型である。このとき Morphisms,
Lemma \ref{morphisms-lemma-base-change-quasi-finite} により
$f$ は局所準有限である。準有限の場合は、「準コンパクトであること」が
基底上 fpqc 局所的であることをすでに見たことから従う
（補題 \ref{descent-lemma-descending-property-quasi-compact}）。
\end{proof}

\begin{lemma}
\label{descent-lemma-descending-property-relative-dimension-d}
性質 $\mathcal{P}(f) =$「$f$ は相対次元 $d$ の局所有限型である」は
基底上 fpqc 局所的である。
\end{lemma}

\begin{proof}
局所有限型であることが基底上 fpqc 局所的であるという事実と
Morphisms, Lemma
\ref{morphisms-lemma-dimension-fibre-after-base-change} から直ちに従う。
\end{proof}

\begin{lemma}
\label{descent-lemma-descending-property-syntomic}
性質 $\mathcal{P}(f) =$「$f$ は syntomic である」は
基底上 fpqc 局所的である。
\end{lemma}

\begin{proof}
射が syntomic であることと、局所有限表示かつ平坦であり、
ファイバーが局所完全交叉であることとは同値である。
平坦であることと局所有限表示であることが基底上 fpqc 局所的で
あることはすでに見た（Lemmas
\ref{descent-lemma-descending-property-flat} および
\ref{descent-lemma-descending-property-locally-finite-presentation}）。
したがって syntomic の場合の結果は Morphisms,
Lemma \ref{morphisms-lemma-set-points-where-fibres-lci} から従う。
\end{proof}

\begin{lemma}
\label{descent-lemma-descending-property-smooth}
性質 $\mathcal{P}(f) =$「$f$ は平滑である」は
基底上 fpqc 局所的である。
\end{lemma}

\begin{proof}
射が平滑であることと、局所有限表示かつ平坦であり、
ファイバーが平滑であることとは同値である。
平坦であることと局所有限表示であることが基底上 fpqc 局所的で
あることはすでに見た（Lemmas
\ref{descent-lemma-descending-property-flat} および
\ref{descent-lemma-descending-property-locally-finite-presentation}）。
したがって平滑の場合の結果は Morphisms,
Lemma \ref{morphisms-lemma-set-points-where-fibres-smooth} から従う。
\end{proof}

\begin{lemma}
\label{descent-lemma-descending-property-unramified}
性質 $\mathcal{P}(f) =$「$f$ は不分岐である」は
基底上 fpqc 局所的である。
性質 $\mathcal{P}(f) =$「$f$ は G-不分岐である」も
基底上 fpqc 局所的である。
\end{lemma}

\begin{proof}
射が不分岐（resp.\ G-不分岐）であることと、局所有限型
（resp.\ 有限表示）であり、かつその対角射が開埋め込みであることとは
同値である（Morphisms,
Lemma \ref{morphisms-lemma-diagonal-unramified-morphism} 参照）。
局所有限型（resp.\ 局所有限表示）であることと開埋め込みであることが
基底上 fpqc 局所的であることはすでに見た（Lemmas
\ref{descent-lemma-descending-property-locally-finite-presentation},
\ref{descent-lemma-descending-property-locally-finite-type} および
\ref{descent-lemma-descending-property-open-immersion}）。
したがって結果は形式的に従う。
\end{proof}

\begin{lemma}
\label{descent-lemma-descending-property-etale}
性質 $\mathcal{P}(f) =$「$f$ は \'etale である」は
基底上 fpqc 局所的である。
\end{lemma}

\begin{proof}
射が \'etale であることと、平坦かつ G-不分岐であることとは同値である。
Morphisms, Lemma \ref{morphisms-lemma-flat-unramified-etale} を参照。
平坦であることと G-不分岐であることが基底上 fpqc 局所的であることは
すでに見た（Lemmas \ref{descent-lemma-descending-property-flat} および
\ref{descent-lemma-descending-property-unramified}）。したがって結果を得る。
\end{proof}

\begin{lemma}
\label{descent-lemma-descending-property-finite-locally-free}
性質 $\mathcal{P}(f) =$「$f$ は有限局所自由である」は
基底上 fpqc 局所的である。$d \geq 0$ とする。
性質 $\mathcal{P}(f) =$「$f$ は次数 $d$ の有限局所自由である」も
基底上 fpqc 局所的である。
\end{lemma}

\begin{proof}
有限局所自由であることは、有限、平坦、かつ局所有限表示であることと
同値である（Morphisms, Lemma
\ref{morphisms-lemma-finite-flat}）。したがって Lemmas
\ref{descent-lemma-descending-property-finite},
\ref{descent-lemma-descending-property-flat} および
\ref{descent-lemma-descending-property-locally-finite-presentation} から従う。
$f : Z \to U$ が有限局所自由で、$\{U_i \to U\}$ が射の全射族であり、
各引き戻し $Z \times_U U_i \to U_i$ が次数 $d$ をもつならば、
$Z \to U$ も次数 $d$ をもつ。たとえば、点 $u \in U$ における次数は
ファイバー
$(f_*\mathcal{O}_Z)_u \otimes_{\mathcal{O}_{U, u}} \kappa(u)$.
から読み取れる。
\end{proof}

\begin{lemma}
\label{descent-lemma-descending-property-monomorphism}
性質 $\mathcal{P}(f) =$「$f$ はモノ射である」は
基底上 fpqc 局所的である。
\end{lemma}

\begin{proof}
$f : X \to S$ をスキームの射とする。$\{S_i \to S\}$ を fpqc 被覆とし、
$f$ の各基底変換 $f_i : X_i \to S_i$ がモノ射であると仮定する。
$a, b : T \to X$ を $f \circ a = f \circ b$ を満たす二つの射とする。
$a = b$ を示す必要がある。$f_i$ はモノ射なので $a_i = b_i$ である。
ここで $a_i, b_i : S_i \times_S T \to X_i$ は基底変換である。
特に合成 $S_i \times_S T \to T \to X$ は等しい。
$\coprod S_i \times_S T \to T$ はエピ射なので
（たとえば補題 \ref{descent-lemma-fpqc-universal-effective-epimorphisms} 参照）、
$a = b$ を得る。
\end{proof}

\begin{lemma}
\label{descent-lemma-descending-property-regular-immersion}
次の性質
\begin{enumerate}
\item[] $\mathcal{P}(f) =$「$f$ は Koszul-正則埋め込みである」、
\item[] $\mathcal{P}(f) =$「$f$ は $H_1$-正則埋め込みである」、
\item[] $\mathcal{P}(f) =$「$f$ は準正則埋め込みである」
\end{enumerate}
はいずれも基底上 fpqc 局所的である。
\end{lemma}

\begin{proof}
補題 \ref{descent-lemma-descending-properties-morphisms} の判定法を用いる。
Divisors, Definition \ref{divisors-definition-regular-immersion} により、
Koszul-正則（resp.\ $H_1$-正則、準正則）埋め込みであることは基底上
Zariski 局所的である。Divisors, Lemma
\ref{divisors-lemma-flat-base-change-regular-immersion} により、
Koszul-正則（resp.\ $H_1$-正則、準正則）埋め込みであることは
平坦基底変換で保存される。補題
\ref{descent-lemma-descending-properties-morphisms} の最後の仮定 (3) は、
次の代数的主張に翻訳される：$A \to B$ を忠実平坦な環準同型とし、
$I \subset A$ をイデアルとする。$IB$ が $\Spec(B)$ 上局所的に
$B$ の Koszul-正則（resp.\ $H_1$-正則、準正則）列で生成されるならば、
$I \subset A$ は $\Spec(A)$ 上局所的に $A$ の Koszul-正則
（resp.\ $H_1$-正則、準正則）列で生成される。これは More on Algebra,
Lemma \ref{more-algebra-lemma-flat-descent-regular-ideal} である。
\end{proof}





\section{標的上の fppf 位相について局所的な射の性質}
\label{descent-section-descending-properties-morphisms-fppf}

\noindent
本節では、fpqc 位相について基底上局所的であることは（まだ）示せないが、
fppf 位相については基底上局所的であるスキームの射の性質をいくつか求める。

\begin{lemma}
\label{descent-lemma-descending-fppf-property-immersion}
性質 $\mathcal{P}(f) =$「$f$ は埋め込みである」は
基底上 fppf 局所的である。
\end{lemma}

\begin{proof}
埋め込みであるという性質は基底変換で保存される。Schemes,
Lemma \ref{schemes-lemma-base-change-immersion} を参照。
埋め込みであるという性質は基底上 Zariski 局所的である。
最後に、$\pi : S' \to S$ をアフィン・スキームの全射で、
平坦かつ局所有限表示なものとする。Morphisms,
Lemma \ref{morphisms-lemma-fppf-open} により
$\pi : S' \to S$ は開写像であることに注意する。
$f : X \to S$ を射とし、基底変換 $f' : X' \to S'$ が
埋め込みであると仮定する。特に
$f'(X') = \pi^{-1}(f(X))$ は局所閉である。
したがって Topology, Lemma
\ref{topology-lemma-open-morphism-quotient-topology} により
$f(X) \subset S$ は局所閉である。$Z \subset S$ を閉部分集合
$Z = \overline{f(X)} \setminus f(X)$ とする。再び Topology,
Lemma \ref{topology-lemma-open-morphism-quotient-topology} により、
$f'(X')$ は $S' \setminus Z'$ において閉である。
ゆえに fpqc 被覆 $\{S' \setminus Z' \to S \setminus Z\}$ に
補題 \ref{descent-lemma-descending-property-closed-immersion} を適用でき、
$f : X \to S \setminus Z$ は閉埋め込みであると結論する。
言い換えれば $f$ は埋め込みである。したがって補題
\ref{descent-lemma-descending-properties-morphisms} を適用でき、結論を得る。
\end{proof}

\begin{lemma}
\label{descent-lemma-descending-fppf-property-dominant}
性質 $\mathcal{P}(f) =$「$f$ は支配的である」は
基底上 fppf 局所的である。
\end{lemma}

\begin{proof}
Morphisms, Lemma \ref{morphisms-lemma-open-base-change-dominant} により、
支配的射は開写像による引き戻しで保存され、したがって局所有限表示な
平坦射による引き戻しで保存される
（Morphisms, Lemma \ref{morphisms-lemma-fppf-open}）。
逆向きはより容易である。実際、支配的であることは明らかに基底上
Zariski 局所的である。次に $S' \to S$ をスキームの全射
（平坦とも局所有限表示とも限らない）とし、$f : X \to S$ を射とする。
基底変換 $f' : X' \to S'$ が支配的であると仮定する。このとき
$X' \to S'\to S$ は二つの支配的射の合成なので支配的である。
これは合成 $X' \to X \to S$ でもあるから、$X\to S$ は支配的である。
\end{proof}






\section{射の性質の fpqc 降下の応用}
\label{descent-section-application-descending-properties-morphisms}

\noindent
次の補題はやや軽い主張に見えるかもしれないが、\'etale 射と不分岐射を
研究するうえで有用な道具となる。

\begin{lemma}
\label{descent-lemma-flat-surjective-quasi-compact-monomorphism-isomorphism}
$f : X \to Y$ を平坦、準コンパクト、かつ全射なモノ射とする。
このとき f は同型である。
\end{lemma}

\begin{proof}
$f$ は平坦、準コンパクト、かつ全射なので、
$\{X \to Y\}$ は $Y$ の fpqc 被覆である。
対角射 $\Delta : X \to X \times_Y X$ は同型である
（Schemes, Lemma \ref{schemes-lemma-monomorphism}）。
これは $f$ による $f$ の基底変換が同型であることを意味する。
したがって補題 \ref{descent-lemma-descending-property-isomorphism} により
$f$ は同型である。
\end{proof}

\noindent
この補題を用いて次の重要な結果を示せる。fpqc 降下を用いない証明も
与える。この結果と関連する結果は \'Etale Morphisms,
Section \ref{etale-section-topological-etale} でさらに詳しく論じる。

\begin{lemma}
\label{descent-lemma-universally-injective-etale-open-immersion}
普遍単射な \'etale 射は開埋め込みである。
\end{lemma}

\begin{proof}[第1の証明]
$f : X \to Y$ を普遍単射な \'etale 射とする。このとき $f$ は開写像である
（Morphisms, Lemma \ref{morphisms-lemma-etale-open}）。
したがって $Y$ を $f(X)$ で置き換え、$f$ は全射であると仮定してよい。
すると $f$ は全単射かつ開写像なので同相写像であり、したがって
$f$ は準コンパクトである。ゆえに補題
\ref{descent-lemma-flat-surjective-quasi-compact-monomorphism-isomorphism} により、
$f$ がモノ射であることを示せば十分である。
$X \to Y$ は \'etale なので、Morphisms,
Lemma \ref{morphisms-lemma-diagonal-unramified-morphism}
（および Morphisms,
Lemma \ref{morphisms-lemma-flat-unramified-etale}）により、射
$\Delta_{X/Y} : X \to X \times_Y X$ は開埋め込みである。
$f$ は普遍単射なので、Morphisms,
Lemma \ref{morphisms-lemma-universally-injective} により
$\Delta_{X/Y}$ は全射でもある。したがって $\Delta_{X/Y}$ は同型、
すなわち $X \to Y$ はモノ射である。
\end{proof}

\begin{proof}[第2の証明]
$f : X \to Y$ を普遍単射な \'etale 射とする。このとき $f$ は開写像である
（Morphisms, Lemma \ref{morphisms-lemma-etale-open}）。
したがって $Y$ を $f(X)$ で置き換え、$f$ は全射であると仮定してよい。
仮定は任意の基底変換後も成り立つので、$f$ は普遍同相である。
したがって $f$ は整である。Morphisms,
Lemma \ref{morphisms-lemma-universal-homeomorphism} を参照。
さらに Morphisms, Lemma \ref{morphisms-lemma-finite-integral} により
$f$ は有限であり、Morphisms, Lemma
\ref{morphisms-lemma-finite-flat} により $f$ は有限局所自由である。
証明を終えるには、$f$ が次数 $1$ の有限局所自由であることを
示せば十分である（次数 $1$ の有限局所自由射は同型である）。
$Y$ は開かつ閉な部分スキーム $V_d$ に分解され、その上で
$f^{-1}(V_d) \to V_d$ は次数 $d$ の有限局所自由となる。
Morphisms, Lemma \ref{morphisms-lemma-finite-locally-free} を参照。
$V_d$ が空でなければ、$k$ を代数閉体として射
$\Spec(k) \to V_d \subset Y$ をとれる
（$V_d$ のある点の剰余体の代数閉包をとればよい）。
このとき Morphisms, Lemma \ref{morphisms-lemma-etale-over-field} と
$k$ が代数閉であることにより、
$\Spec(k) \times_Y X \to \Spec(k)$ は $\Spec(k)$ のコピーの
非交和である。しかし $f$ は普遍単射なのでコピーは一つしかなく、
所望のとおり $d = 1$ である。
\end{proof}

\noindent
平坦な普遍単射の次の特徴付けを用いると、上の補題の仮定を少し
言い換えることができる。

\begin{lemma}
\label{descent-lemma-flat-universally-injective}
$f : X \to Y$ をスキームの射とする。$X^0$ を $X$ の既約成分の
生成点全体とする。次を仮定する。
\begin{enumerate}
\item $f$ は平坦かつ分離的である。
\item $\xi \in X^0$ に対して $\kappa(f(\xi)) = \kappa(\xi)$ である。
\item $\xi, \xi' \in X^0$, $\xi \not = \xi'$ ならば
$f(\xi) \not = f(\xi')$ である。
\end{enumerate}
このとき $f$ は普遍単射である。
\end{lemma}

\begin{proof}
$\Delta : X \to X \times_Y X$ が全射であることを示す必要がある。
Morphisms, Lemma \ref{morphisms-lemma-universally-injective} を参照。
$X \to Y$ は分離的なので、$\Delta$ の像は閉である。
したがって $\Delta$ が全射でなければ、$\Delta$ の像に含まれない
$X \times_S X$ の既約成分の生成点
$\eta \in X \times_S X$ をとれる。射影
$\text{pr}_1 : X \times_Y X \to X$ は平坦射 $X \to Y$ の基底変換なので
平坦である。Morphisms, Lemma
\ref{morphisms-lemma-base-change-flat} を参照。したがって
$\text{pr}_1$ に沿って一般化を持ち上げられる。Morphisms,
Lemma \ref{morphisms-lemma-generalizations-lift-flat} を参照。
ゆえに $\xi = \text{pr}_1(\eta) \in X^0$ である。
しかし仮定 (2) と (3) により、任意の $\xi \in X^0$ に対して
スキーム $(X \times_Y X)_{f(\xi)}$ は高々一つの点しかもたない。
言い換えれば $\Delta(\xi) = \eta$ であり、矛盾である。
\end{proof}

\noindent
したがって補題
\ref{descent-lemma-universally-injective-etale-open-immersion} を次のように
言い換えられる。

\begin{lemma}
\label{descent-lemma-characterize-open-immersion}
$f : X \to Y$ をスキームの射とする。$X^0$ を $X$ の既約成分の
生成点全体とする。次を仮定する。
\begin{enumerate}
\item $f$ は \'etale かつ分離的である。
\item $\xi \in X^0$ に対して $\kappa(f(\xi)) = \kappa(\xi)$ である。
\item $\xi, \xi' \in X^0$, $\xi \not = \xi'$ ならば
$f(\xi) \not = f(\xi')$ である。
\end{enumerate}
このとき $f$ は開埋め込みである。
\end{lemma}

\begin{proof}
Lemmas \ref{descent-lemma-flat-universally-injective} と
\ref{descent-lemma-universally-injective-etale-open-immersion} から直ちに従う。
\end{proof}

\begin{lemma}
\label{descent-lemma-descending-property-proper-over-base}
$f : X \to Y$ を局所有限型なスキームの射とし、
$Z$ を $X$ の閉部分集合とする。逆像 $Z_i \subset Y_i \times_Y X$ が
$Y_i$ 上固有となる fpqc 被覆 $\{Y_i \to Y\}$ が存在するならば
（Cohomology of Schemes, Definition
\ref{coherent-definition-proper-over-base}）、$Z$ は $Y$ 上固有である。
\end{lemma}

\begin{proof}
$Z$ に誘導された被約閉部分スキーム構造を入れる。Schemes,
Definition \ref{schemes-definition-reduced-induced-scheme} を参照。
各 $i$ に対し、基底変換 $Y_i \times_Y Z$ は
$Y_i \times_Y X$ の閉部分スキームで、その台となる閉部分集合は
$Z_i$ である。定義により（Cohomology of Schemes,
Lemma \ref{coherent-lemma-closed-proper-over-base} を介して）、
射影 $Y_i \times_Y Z \to Y_i$ は固有射である。
したがって補題 \ref{descent-lemma-descending-property-proper} により
$Z \to Y$ は固有射である。ゆえに定義により $Z$ は $Y$ 上固有である。
\end{proof}

\begin{lemma}
\label{descent-lemma-descending-property-ample}
$f : X \to S$ をスキームの射とし、
$\mathcal{L}$ を可逆 $\mathcal{O}_X$-加群とする。
$\{g_i : S_i \to S\}_{i \in I}$ を fpqc 被覆とする。
$f_i : X_i \to S_i$ を $f$ の基底変換とし、$\mathcal{L}_i$ を
$X_i$ への $\mathcal{L}$ の引き戻しとする。次は同値である。
\begin{enumerate}
\item $\mathcal{L}$ は $X/S$ 上豊富である。
\item すべての $i \in I$ に対して $\mathcal{L}_i$ は
$X_i/S_i$ 上豊富である。
\end{enumerate}
\end{lemma}

\begin{proof}
(1) $\Rightarrow$ (2) は Morphisms,
Lemma \ref{morphisms-lemma-ample-base-change} から従う。
すべての $i \in I$ に対して $\mathcal{L}_i$ が $X_i/S_i$ 上
豊富であると仮定する。Morphisms, Definition
\ref{morphisms-definition-relatively-ample} により
$X_i \to S_i$ は準コンパクトであり、Morphisms,
Lemma \ref{morphisms-lemma-relatively-ample-separated} により
$X_i \to S$ は分離的である。したがって Lemmas
\ref{descent-lemma-descending-property-quasi-compact} と
\ref{descent-lemma-descending-property-separated} により、$f$ は
準コンパクトかつ分離的である。

\medskip\noindent
これは
$\mathcal{A} = \bigoplus_{d \geq 0} f_*\mathcal{L}^{\otimes d}$
が準連接次数付き $\mathcal{O}_S$-代数であることを意味する
（Schemes, Lemma \ref{schemes-lemma-push-forward-quasi-coherent}）。
さらに Cohomology of Schemes, Lemma
\ref{coherent-lemma-flat-base-change-cohomology} により、
$\mathcal{A}$ の形成は平坦基底変換と可換である。特に
$\mathcal{A}_i = \bigoplus_{d \geq 0} f_{i, *}\mathcal{L}_i^{\otimes d}$
とおけば $\mathcal{A}_i = g_i^*\mathcal{A}$ である。したがって
$\mathcal{O}_X$ の標準準同型
$\psi_d : f^*\mathcal{A}_d \to \mathcal{L}^{\otimes d}$
を引き戻すと、$\mathcal{O}_{X_i}$-加群の標準準同型
$\psi_{i, d} : f_i^*(\mathcal{A}_i)_d \to \mathcal{L}_i^{\otimes d}$
が得られる。$\mathcal{L}_i$ は $X_i/S_i$ 上豊富なので、任意の点
$x_i \in X_i$ に対して、準同型
$f_i^*(\mathcal{A}_i)_d \to \mathcal{L}_i^{\otimes d}$ が
$x_i$ における茎上全射となる $d \geq 1$ が存在する。
これは相対的に豊富な加群の定義から直接、または Morphisms,
Lemma \ref{morphisms-lemma-characterize-relatively-ample} から従う。
$x \in X$ ならば、$x$ に写る $i$ と $x_i \in X_i$ をとれる。
$\mathcal{O}_{X, x} \to \mathcal{O}_{X_i, x_i}$ は平坦、したがって
忠実平坦なので、任意の $x \in X$ に対して、準同型
$f^*\mathcal{A}_d \to \mathcal{L}^{\otimes d}$
が $x$ における茎上全射となる $d \geq 1$ が存在することが分かる。
これは Constructions, Lemma
\ref{constructions-lemma-invertible-map-into-relative-proj} の開部分集合
$U(\psi) \subset X$、すなわち次数付き $\mathcal{O}_X$-代数の準同型
$\psi : f^*\mathcal{A} \to \bigoplus_{d \geq 0} \mathcal{L}^{\otimes d}$
に対応する開部分集合が $X$ に等しいことを意味する。
対応する射
$$
r_{\mathcal{L}, \psi} : X \longrightarrow \underline{\text{Proj}}_S(\mathcal{A})
$$
を考える。上の議論から、$r_{\mathcal{L}, \psi}$ の $S_i$ への
基底変換は射 $r_{\mathcal{L}_i, \psi_i}$ であり、これは Morphisms,
Lemma \ref{morphisms-lemma-characterize-relatively-ample} により
開埋め込みである。したがって補題
\ref{descent-lemma-descending-property-open-immersion} により
$r_{\mathcal{L}, \psi}$ は開埋め込みであり、再び Morphisms,
Lemma \ref{morphisms-lemma-characterize-relatively-ample} により
$\mathcal{L}$ は $X/S$ 上豊富であると結論する。
\end{proof}















\section{始域上局所的な射の性質}
\label{descent-section-properties-morphisms-local-source}

\noindent
別の射を前合成した後で、ある射が所定の性質をもつことを証明できる場合が
しばしばある。多くの場合、これは元の射もその性質をもつことを含意する。
以下の定義でこれを定式化する。

\begin{definition}
\label{descent-definition-property-morphisms-local-source}
スキームの射の性質 $\mathcal{P}$ をとる。
$\tau \in \{Zariski, \linebreak[0] fpqc, \linebreak[0] fppf, \linebreak[0]
\etale, \linebreak[0] smooth, \linebreak[0] syntomic\}$.
$S$ 上の任意のスキームの射 $f : X \to Y$ と任意の
$\tau$-被覆 $\{X_i \to X\}_{i \in I}$ に対して
$$
f \text{ は }\mathcal{P}\text{ をもつ}
\Leftrightarrow
\text{各 }X_i \to Y\text{ は }\mathcal{P}\text{ をもつ}.
$$
が成り立つとき、$\mathcal{P}$ は
{\it 始域上 $\tau$-局所的}、または
{\it $\tau$-位相について始域上局所的}であるという。
\end{definition}

\noindent
同型射は常に被覆なので、性質 $\mathcal{P}$ が $X \to Y$ について
成り立つことと、$X \to Y$ に同型な任意の矢 $X' \to Y'$ について
成り立つこととは同値であることが分かる（あるいは、そう要請する）。
ある性質が始域上 $\tau$-局所的ならば、$\tau$-被覆に現れる射を
前合成しても保たれる。これを形式的に述べると次のようになる。

\begin{lemma}
\label{descent-lemma-precompose-property-local-source}
$\tau \in \{fpqc, fppf, syntomic, smooth, \etale, Zariski\}$ とする。
$\mathcal{P}$ を始域上 $\tau$ 局所的な射の性質とし、
$f : X \to Y$ が性質 $\mathcal{P}$ をもつとする。
平坦、resp.\ 平坦かつ局所有限表示、resp.\ syntomic、
resp.\ \'etale、resp.\ 開埋め込みである任意の射
$a : X' \to X$ に対し、合成
$f \circ a : X' \to Y$ は性質 $\mathcal{P}$ をもつ。
\end{lemma}

\begin{proof}
$X' \to X$ を $\tau$-被覆をなす射の族に組み込めるからである。
\end{proof}

\begin{lemma}
\label{descent-lemma-largest-open-of-the-source}
$\tau \in \{fppf, syntomic, smooth, \etale\}$ とする。
$\mathcal{P}$ を始域上 $\tau$ 局所的な射の性質とする。
任意のスキームの射 $f : X \to Y$ に対して、制限
$f|_{W(f)} : W(f) \to Y$ が $\mathcal{P}$ をもつような最大の開集合
$W(f) \subset X$ が存在する。さらに、$g : X' \to X$ が
平坦かつ局所有限表示、syntomic、smooth、または \'etale であり、
$f' = f \circ g : X' \to Y$ ならば
$g^{-1}(W(f)) = W(f')$.
\end{lemma}

\begin{proof}
次の性質をもつ射 $g : X' \to X$ の像 $g(X') \subset X$ の和集合を
$W$ とする。
\begin{enumerate}
\item $g$ は平坦かつ局所有限表示、syntomic、smooth、または \'etale であり、
\item 合成 $X' \to X \to Y$ は性質 $\mathcal{P}$ をもつ。
\end{enumerate}
このような射 $g$ は開射なので（Morphisms, Lemma
\ref{morphisms-lemma-fppf-open}）、$W \subset X$ は $X$ の開部分集合である。
$\mathcal{P}$ は $\tau$ 位相について局所的であり、引き戻しが
$\mathcal{P}$ をもつような $W$ の $\tau$ 被覆 $\{X' \to W\}$ が
与えられているから、制限 $f|_W : W \to Y$ は性質 $\mathcal{P}$ をもつ。
これで $W(f)$ の存在が示された。最後の文で述べた整合性は
$W(f)$ の構成から直ちに従う。
\end{proof}

\begin{lemma}
\label{descent-lemma-properties-morphisms-local-source}
スキームの射の性質 $\mathcal{P}$ をとる。
$\tau \in \{fpqc, \linebreak[0] fppf, \linebreak[0]
\etale, \linebreak[0] smooth, \linebreak[0] syntomic\}$.
次を仮定する。
\begin{enumerate}
\item $\tau$ が fpqc、fppf、\'etale、smooth、または syntomic のいずれで
あるかに応じて、平坦、平坦かつ局所有限表示、\'etale、smooth、または
syntomic な射を前合成しても、この性質は保たれる。
\item この性質は始域上 Zariski 局所的である。
\item この性質は標的上 Zariski 局所的である。
\item アフィン・スキームの任意の射 $f : X \to Y$ と、$\tau$ が
fpqc、fppf、\'etale、smooth、または syntomic のいずれであるかに応じて
平坦、平坦かつ有限表示、\'etale、smooth、または syntomic である
アフィン・スキームの任意の全射 $X' \to X$ に対し、合成
$f' : X' \to Y$ について性質 $\mathcal{P}$ が成り立つならば、
$f$ についても性質 $\mathcal{P}$ が成り立つ。
\end{enumerate}
このとき $\mathcal{P}$ は始域上 $\tau$ 局所的である。
\end{lemma}

\begin{proof}
これは $\tau$-被覆の定義からほとんど直ちに従う。Topologies, Definition
\ref{topologies-definition-fpqc-covering}
\ref{topologies-definition-fppf-covering}
\ref{topologies-definition-etale-covering}
\ref{topologies-definition-smooth-covering}、または
\ref{topologies-definition-syntomic-covering}
および Topologies, Lemma
\ref{topologies-lemma-fpqc-affine},
\ref{topologies-lemma-fppf-affine},
\ref{topologies-lemma-etale-affine},
\ref{topologies-lemma-smooth-affine}、または
\ref{topologies-lemma-syntomic-affine} を参照。
詳細は省略する。（ヒント：始域と標的上の局所性を用いて、性質
$\mathcal{P}$ の検証をアフィン間の射の場合に帰着し、その後 (1) と (4) を適用する。）
\end{proof}

\begin{remark}
\label{descent-remark-properties-morphisms-local-source-standard}
（これは上の Remarks \ref{descent-remark-descending-properties-standard} と
\ref{descent-remark-descending-properties-morphisms-standard} の繰り返しである。）
上の Lemma \ref{descent-lemma-properties-morphisms-local-source} で
$\tau = smooth$ ならば、条件 (4) の射は（全射な）標準 smooth 射であると
仮定してよい。$\tau = syntomic$ または $\tau = \etale$ の場合も同様である。
\end{remark}



\section{始域上の fpqc 位相について局所的な射の性質}
\label{descent-section-fpqc-local-source}

\noindent
始域上 fpqc 局所的な射の性質をいくつか挙げる。

\begin{lemma}
\label{descent-lemma-flat-fpqc-local-source}
性質 $\mathcal{P}(f)=$「$f$ は平坦である」は始域上 fpqc 局所的である。
\end{lemma}

\begin{proof}
平坦性は局所環の写像によって定義されるので（Morphisms, Definition
\ref{morphisms-definition-flat}）、示すべきことは次の代数的事実である。
$A \to B \to C$ を局所環の局所準同型とし、$B \to C$ が平坦であると
仮定する。このとき $A \to B$ が平坦であることと
$A \to C$ が平坦であることは同値である。
$A \to B$ が平坦ならば、Algebra, Lemma
\ref{algebra-lemma-composition-flat} により $A \to C$ は平坦である。
逆に $A \to C$ が平坦であると仮定する。Algebra, Lemma
\ref{algebra-lemma-local-flat-ff} により $B \to C$ は忠実平坦である。
したがって Algebra, Lemma \ref{algebra-lemma-flat-permanence} により
$A \to B$ は平坦である。
（直接の証明については Morphisms, Lemma
\ref{morphisms-lemma-flat-permanence} も参照。）
\end{proof}

\begin{lemma}
\label{descent-lemma-injective-local-rings-fpqc-local-source}
このとき、性質
$\mathcal{P}(f : X \to Y)=$「すべての $x \in X$ に対して局所環の写像
$\mathcal{O}_{Y, f(x)} \to \mathcal{O}_{X, x}$ は単射である」は
始域上 fpqc 局所的である。
\end{lemma}

\begin{proof}
省略する。これは単に、（おそらく見当違いな）遊び心を出そうとしただけである。
\end{proof}





\section{始域上の fppf 位相について局所的な射の性質}
\label{descent-section-fppf-local-source}

\noindent
始域上 fppf 局所的な射の性質をいくつか挙げる。

\begin{lemma}
\label{descent-lemma-locally-finite-presentation-fppf-local-source}
性質 $\mathcal{P}(f)=$「$f$ は局所有限表示である」は
始域上 fppf 局所的である。
\end{lemma}

\begin{proof}
局所有限表示であることは始域と標的上 Zariski 局所的である。
Morphisms, Lemma
\ref{morphisms-lemma-locally-finite-presentation-characterize} を参照。
また、合成で保たれる性質である。Morphisms, Lemma
\ref{morphisms-lemma-composition-finite-presentation} を参照。
これで Lemma \ref{descent-lemma-properties-morphisms-local-source} の
(1)、(2)、(3) が示された。最後の条件 (4) は
Lemma \ref{descent-lemma-flat-finitely-presented-permanence-algebra} である。
これで証明が終わる。
\end{proof}

\begin{lemma}
\label{descent-lemma-locally-finite-type-fppf-local-source}
性質 $\mathcal{P}(f)=$「$f$ は局所有限型である」は
始域上 fppf 局所的である。
\end{lemma}

\begin{proof}
局所有限型であることは始域と標的上 Zariski 局所的である。
Morphisms, Lemma
\ref{morphisms-lemma-locally-finite-type-characterize} を参照。
また、合成で保たれる性質であり（Morphisms, Lemma
\ref{morphisms-lemma-composition-finite-type}）、局所有限表示な平坦射は
局所有限型である（Morphisms, Lemma
\ref{morphisms-lemma-finite-presentation-finite-type}）。
これで Lemma \ref{descent-lemma-properties-morphisms-local-source} の
(1)、(2)、(3) が示された。最後の条件 (4) は
Lemma \ref{descent-lemma-finite-type-local-source-fppf-algebra} である。
これで証明が終わる。
\end{proof}

\begin{lemma}
\label{descent-lemma-open-fppf-local-source}
性質 $\mathcal{P}(f)=$「$f$ は開射である」は始域上 fppf 局所的である。
\end{lemma}

\begin{proof}
開射であることは明らかに始域と標的上 Zariski 局所的である。
また、合成で保たれる性質であり（Morphisms, Lemma
\ref{morphisms-lemma-composition-open}）、有限表示な平坦射は開射である。
Morphisms, Lemma \ref{morphisms-lemma-fppf-open} を参照。
これで Lemma \ref{descent-lemma-properties-morphisms-local-source} の
(1)、(2)、(3) が示された。最後の条件 (4) は Morphisms, Lemma
\ref{morphisms-lemma-fpqc-quotient-topology} から従う。
これで証明が終わる。
\end{proof}

\begin{lemma}
\label{descent-lemma-universally-open-fppf-local-source}
性質 $\mathcal{P}(f)=$「$f$ は普遍開である」は始域上 fppf 局所的である。
\end{lemma}

\begin{proof}
$f : X \to Y$ をスキームの射とし、
$\{X_i \to X\}_{i \in I}$ を fppf 被覆とする。
合成を $f_i : X_i \to X$ と記す。示すべきことは、
$f$ が普遍開であることと各 $f_i$ が普遍開であることが同値なことである。
$f$ が普遍開ならば、写像 $X_i \to X$ は普遍開であり、普遍開射の合成も
普遍開なので（Morphisms, Lemmas \ref{morphisms-lemma-fppf-open} および
\ref{morphisms-lemma-composition-open}）、各 $f_i$ も普遍開である。
逆に、各 $f_i$ が普遍開であると仮定する。
$Y' \to Y$ をスキームの射とする。
$X' = Y' \times_Y X$ および $X'_i = Y' \times_Y X_i$ と記す。
$\{X_i' \to X'\}_{i \in I}$ も fppf 被覆である。
仮定により射 $f'_i : X_i' \to Y'$ は開射である。
したがって上の Lemma \ref{descent-lemma-open-fppf-local-source} により、
望むとおり $f' : X' \to Y'$ は開射である。
\end{proof}



\section{始域上の syntomic 位相について局所的な射の性質}
\label{descent-section-syntomic-local-source}

\noindent
始域上 syntomic 局所的な射の性質をいくつか挙げる。

\begin{lemma}
\label{descent-lemma-syntomic-syntomic-local-source}
性質 $\mathcal{P}(f)=$「$f$ は syntomic である」は
始域上 syntomic 局所的である。
\end{lemma}

\begin{proof}
Lemma \ref{descent-lemma-properties-morphisms-local-source} と
Morphisms, Lemma \ref{morphisms-lemma-syntomic-characterize}
（始域と標的上 Zariski 局所的）、
Morphisms, Lemma \ref{morphisms-lemma-composition-syntomic}（前合成）、
および Lemma \ref{descent-lemma-syntomic-smooth-etale-permanence}（(4) の部分）
を組み合わせる。
\end{proof}




\section{始域上の smooth 位相について局所的な射の性質}
\label{descent-section-smooth-local-source}

\noindent
始域上 smooth 局所的な射の性質をいくつか挙げる。
平坦射による smooth 性の降下については、いくつかの点でより強い結果である
Lemma \ref{descent-lemma-smooth-permanence} にも注意せよ。

\begin{lemma}
\label{descent-lemma-smooth-smooth-local-source}
性質 $\mathcal{P}(f)=$「$f$ は smooth である」は
始域上 smooth 局所的である。
\end{lemma}

\begin{proof}
Lemma \ref{descent-lemma-properties-morphisms-local-source} と
Morphisms, Lemma \ref{morphisms-lemma-smooth-characterize}
（始域と標的上 Zariski 局所的）、
Morphisms, Lemma \ref{morphisms-lemma-composition-smooth}（前合成）、
および Lemma \ref{descent-lemma-syntomic-smooth-etale-permanence}（(4) の部分）
を組み合わせる。
\end{proof}



\section{始域上の \'etale 位相について局所的な射の性質}
\label{descent-section-etale-local-source}

\noindent
始域上 \'etale 局所的な射の性質をいくつか挙げる。

\begin{lemma}
\label{descent-lemma-etale-etale-local-source}
性質 $\mathcal{P}(f)=$「$f$ は \'etale である」は
始域上 \'etale 局所的である。
\end{lemma}

\begin{proof}
Lemma \ref{descent-lemma-properties-morphisms-local-source} と
Morphisms, Lemma \ref{morphisms-lemma-etale-characterize}
（始域と標的上 Zariski 局所的）、
Morphisms, Lemma \ref{morphisms-lemma-composition-etale}（前合成）、
および Lemma \ref{descent-lemma-syntomic-smooth-etale-permanence}（(4) の部分）
を組み合わせる。
\end{proof}

\begin{lemma}
\label{descent-lemma-locally-quasi-finite-etale-local-source}
性質 $\mathcal{P}(f)=$「$f$ は局所準有限である」は
始域上 \'etale 局所的である。
\end{lemma}

\begin{proof}
Lemma \ref{descent-lemma-properties-morphisms-local-source} を用いる。
Morphisms, Lemma
\ref{morphisms-lemma-locally-quasi-finite-characterize} により、
局所準有限であることは始域と標的上 Zariski 局所的である。
Morphisms, Lemmas
\ref{morphisms-lemma-composition-quasi-finite} および
\ref{morphisms-lemma-etale-locally-quasi-finite}
により、局所準有限射に \'etale 射を前合成したものは局所準有限である。
最後に、$X \to Y$ をアフィン・スキームの射とし、
$X' \to X$ をアフィン・スキームの全射 \'etale 射であって
$X' \to Y$ が局所準有限となるものとする。このとき $X' \to Y$ は
有限型であり、Lemma \ref{descent-lemma-finite-type-local-source-fppf-algebra}
により $X \to Y$ も有限型である。さらに仮定により
$X' \to Y$ のファイバーは有限であり、したがって $X \to Y$ の
ファイバーも有限である。Morphisms, Lemma
\ref{morphisms-lemma-quasi-finite} により $X \to Y$ は準有限である。
これで Lemma \ref{descent-lemma-properties-morphisms-local-source} の最後の仮定が
示され、証明が終わる。
\end{proof}

\begin{lemma}
\label{descent-lemma-unramified-etale-local-source}
性質 $\mathcal{P}(f)=$「$f$ は不分岐である」は
始域上 \'etale 局所的である。
性質 $\mathcal{P}(f)=$「$f$ は G-不分岐である」も
始域上 \'etale 局所的である。
\end{lemma}

\begin{proof}
Lemma \ref{descent-lemma-properties-morphisms-local-source} を用いる。
Morphisms, Lemma \ref{morphisms-lemma-unramified-characterize} により、
不分岐（resp.\ G-不分岐）であることは始域と標的上 Zariski 局所的である。
Morphisms, Lemmas \ref{morphisms-lemma-composition-unramified} および
\ref{morphisms-lemma-etale-smooth-unramified}
により、不分岐（resp.\ G-不分岐）射に \'etale 射を前合成したものは
不分岐（resp.\ G-不分岐）である。
最後に、$X \to Y$ をアフィン・スキームの射とし、
$f : X' \to X$ をアフィン・スキームの全射 \'etale 射であって
$X' \to Y$ が不分岐（resp.\ G-不分岐）となるものとする。
このとき $X' \to Y$ は有限型（resp.\ 有限表示）であり、
Lemma \ref{descent-lemma-finite-type-local-source-fppf-algebra}
（resp.\ Lemma \ref{descent-lemma-flat-finitely-presented-permanence-algebra}）により、
$X \to Y$ も有限型（resp.\ 有限表示）である。Morphisms, Lemma
\ref{morphisms-lemma-triangle-differentials-smooth}
により短完全列
$$
0 \to f^*\Omega_{X/Y} \to \Omega_{X'/Y} \to \Omega_{X'/X} \to 0.
$$
を得る。$X' \to Y$ は不分岐なので中項は零である。
さらに $f$ は忠実平坦だから $\Omega_{X/Y} = 0$ である。
したがって Morphisms, Lemma
\ref{morphisms-lemma-unramified-omega-zero} により
$X \to Y$ は不分岐（resp.\ G-不分岐）である。
これで Lemma \ref{descent-lemma-properties-morphisms-local-source} の最後の仮定が
示され、証明が終わる。
\end{proof}




\section{始域・標的上 \'etale 局所的な射の性質}
\label{descent-section-properties-etale-local-source-target}

\noindent
スキームの射の性質 $\mathcal{P}$ をとる。「$\mathcal{P}$ は始域と
標的上 \'etale 局所的である」という語句には直観的な意味がある。
しかし、この概念は $\mathcal{P}$ が始域上 \'etale 局所的かつ
標的上 \'etale 局所的であることを要求するのとは同じでない。
これをさらに論じる前に、二つのばかげた例を挙げる。

\begin{example}
\label{descent-example-silly-one}
規則 $\mathcal{P}(X \to Y) = $「$Y$ は局所 Noether である」によって
定義されるスキームの射の性質 $\mathcal{P}$ を考える。
これが始域上 \'etale 局所的かつ標的上 \'etale 局所的であることは
読者が確かめられる（省略する。Lemma \ref{descent-lemma-Noetherian-local-fppf}
を参照）。しかし、$f : X \to Y$ が $\mathcal{P}$ をもち、
$g : Y \to Z$ が \'etale ならば $g \circ f$ が $\mathcal{P}$ をもつ、
ということは {\bf 成り立たない}。実際、$f$ は $Y$ 上の恒等射であり、
$g$ は局所 Noether スキーム $Y$ から局所 Noether でないスキーム
$Z$ への開埋め込みであり得る。
\end{example}

\noindent
次の例は、ある意味でさらに悪い。

\begin{example}
\label{descent-example-silly-two}
規則 $\mathcal{P}(f : X \to Y) = $「ある $f(x)$（$x \in X$）の
特殊化であるすべての $y \in Y$ に対して、局所環
$\mathcal{O}_{Y, y}$ は Noether である」によって定義される
スキームの射の性質 $\mathcal{P}$ を考える。これが始域上
\'etale 局所的かつ標的上 \'etale 局所的であることを確かめよう。
Schemes, Lemma \ref{schemes-lemma-specialize-points} を自由に用いる。

\medskip\noindent
標的上局所的であること：
$\{g_i : Y_i \to Y\}$ を \'etale 被覆とする。
$f_i : X_i \to Y_i$ を $f$ の基底変換とし、射影を
$h_i : X_i \to X$ と記す。$\mathcal{P}(f)$ を仮定する。
$f(x_i) \leadsto y_i$ を特殊化とする。このとき
$f(h_i(x_i)) \leadsto g_i(y_i)$ なので、$\mathcal{P}(f)$ により
$\mathcal{O}_{Y, g_i(y_i)}$ は Noether である。また
$\mathcal{O}_{Y, g_i(y_i)} \to \mathcal{O}_{Y_i, y_i}$ は
\'etale 環写像の局所化である。したがって Algebra, Lemma
\ref{algebra-lemma-Noetherian-permanence} により
$\mathcal{O}_{Y_i, y_i}$ は Noether である。
逆に、すべての $i$ に対して $\mathcal{P}(f_i)$ を仮定する。
$f(x) \leadsto y$ を特殊化とし、$y$ に写る $i$ と
$y_i \in Y_i$ を選ぶ。$x$ は
$\Spec(\mathcal{O}_{Y, y}) \times_Y X$ の点とみなせ、また
$\mathcal{O}_{Y, y} \to \mathcal{O}_{Y_i, y_i}$ は忠実平坦なので、$x$ に写る点
$x_i \in \Spec(\mathcal{O}_{Y_i, y_i}) \times_Y X$
が存在する。このとき $x_i \in X_i$ であり、$f_i(x_i)$ は $y_i$ に
特殊化する。したがって $\mathcal{P}(f_i)$ により
$\mathcal{O}_{Y_i, y_i}$ は Noether であり、Algebra, Lemma
\ref{algebra-lemma-descent-Noetherian} により
$\mathcal{O}_{Y, y}$ も Noether である。

\medskip\noindent
始域上局所的であること：
$\{h_i : X_i \to X\}$ を \'etale 被覆とし、$f_i : X_i \to Y$ を
合成 $f \circ h_i$ とする。$\mathcal{P}(f)$ を仮定する。
$f(x_i) \leadsto y$ を特殊化とする。このとき
$f(h_i(x_i)) \leadsto y$ なので、$\mathcal{P}(f)$ により
$\mathcal{O}_{Y, y}$ は Noether である。したがって
$\mathcal{P}(f_i)$ が成り立つ。
逆に、すべての $i$ に対して $\mathcal{P}(f_i)$ を仮定する。
$f(x) \leadsto y$ を特殊化とし、$x$ に写る $i$ と
$x_i \in X_i$ を選ぶ。このとき $y$ は
$f_i(x_i) = f(x)$ の特殊化である。ゆえに $\mathcal{P}(f_i)$ により、
望むとおり $\mathcal{O}_{Y, y}$ は Noether である。

\medskip\noindent
次の可換図式
$$
\xymatrix{
U \ar[d]_a \ar[r]_h & V \ar[d]^b \\
X \ar[r]^f & Y
}
$$
で、縦の矢が全射 \'etale であり、$h$ は $\mathcal{P}$ をもつが
$f$ は $\mathcal{P}$ をもたないものが存在すると主張する。実際、
$$
Y =
\Spec\Big(
\mathbf{C}[x_n; n \in \mathbf{Z}]/(x_n x_m; n \not = m)
\Big)
$$
とし、$X \subset Y$ を、すべての座標が $x_n = 0$ である点の補集合
として得られる開部分スキームとする。$U = X$、$V = X \amalg Y$ とし、
$a, b$ を明らかな写像、$h : U \to V$ を $U = X$ から $V$ の
第一成分への包含とする。$U$ は局所 Noether だが $Y$ はそうでないので、
上の主張が成り立つ。
\end{example}

\noindent
始域・標的上 \'etale 局所的な性質の正しい概念は何であるべきだろうか。
Morphisms, Definition \ref{morphisms-definition-property-local} との
類似から、次のものだと考える。

\begin{definition}
\label{descent-definition-local-source-target}
スキームの射の性質 $\mathcal{P}$ をとる。次を満たすとき
$\mathcal{P}$ は {\it 始域・標的上 \'etale 局所的}であるという。
\begin{enumerate}
\item （\'etale 写像の前合成で安定）
$f : X \to Y$ が \'etale で $g : Y \to Z$ が $\mathcal{P}$ をもつならば、
$g \circ f$ は $\mathcal{P}$ をもつ。
\item （\'etale 基底変換で安定）
$f : X \to Y$ が $\mathcal{P}$ をもち $Y' \to Y$ が \'etale ならば、
基底変換 $f' : Y' \times_Y X \to Y'$ は $\mathcal{P}$ をもつ。
\item （局所性）射 $f : X \to Y$ に対して次は同値である。
\begin{enumerate}
\item $f$ は $\mathcal{P}$ をもつ。
\item すべての $x \in X$ に対して可換図式
$$
\xymatrix{
U \ar[d]_a \ar[r]_h & V \ar[d]^b \\
X \ar[r]^f & Y
}
$$
で、縦の矢が \'etale であり、$a(u) = x$ を満たす $u \in U$ が存在し、
$h$ が $\mathcal{P}$ をもつものが存在する。
\end{enumerate}
\end{enumerate}
\end{definition}

\noindent
この定義は Examples \ref{descent-example-silly-one} と
\ref{descent-example-silly-two} で見た挙動を排除することが分かる。
Remark \ref{descent-remark-compare-definitions} で、これを Deligne--Mumford の論文
\cite{DM} における定義と比較する。さらに、始域・標的上 \'etale 局所的な
性質は、始域上 \'etale 局所的かつ標的上 \'etale 局所的である。
最後に、Lemma \ref{descent-lemma-etale-local-source-target} で見るように、
その逆もほとんど成り立つ。

\begin{lemma}
\label{descent-lemma-local-source-target-implies}
$\mathcal{P}$ を始域・標的上 \'etale 局所的なスキームの射の性質とする。
このとき
\begin{enumerate}
\item $\mathcal{P}$ は始域上 \'etale 局所的である。
\item $\mathcal{P}$ は標的上 \'etale 局所的である。
\item $\mathcal{P}$ は \'etale 射の後合成で安定である。すなわち、
$f : X \to Y$ が $\mathcal{P}$ をもち $g : Y \to Z$ が \'etale ならば、
$g \circ f$ は $\mathcal{P}$ をもつ。
\item $\mathcal{P}$ は次の保持性をもつ。$f : X \to Y$ と
\'etale 射 $g : Y \to Z$ が与えられ、$g \circ f$ が
$\mathcal{P}$ をもつならば、$f$ は $\mathcal{P}$ をもつ。
\end{enumerate}
\end{lemma}

\begin{proof}
すべてを詳しく書き下す。

\medskip\noindent
(1) の証明。$f : X \to Y$ をスキームの射とし、
$\{X_i \to X\}_{i \in I}$ を $X$ の \'etale 被覆とする。
各合成 $h_i : X_i \to Y$ が $\mathcal{P}$ をもつならば、各
$x \in X$ に対して、$x$ に写る $i \in I$ と点 $x_i \in X_i$ をとれる。
このとき $(X_i, x_i) \to (X, x)$ は芽の \'etale 射であり、
$\text{id}_Y : Y \to Y$ は \'etale 射であり、$h_i$ は Definition
\ref{descent-definition-local-source-target} の (3) にある射である。
したがって $f$ は $\mathcal{P}$ をもつ。
逆に $f$ が $\mathcal{P}$ をもつならば、Definition
\ref{descent-definition-local-source-target} の (1) により
各 $X_i \to Y$ は $\mathcal{P}$ をもつ。

\medskip\noindent
(2) の証明。$f : X \to Y$ をスキームの射とし、
$\{Y_i \to Y\}_{i \in I}$ を $Y$ の \'etale 被覆とする。
$X_i = Y_i \times_Y X$ とおき、$h_i : X_i \to Y_i$ を $f$ の
基底変換とする。各 $h_i : X_i \to Y_i$ が $\mathcal{P}$ をもつならば、
各 $x \in X$ に対して、$x$ に写る $i \in I$ と点 $x_i \in X_i$ をとる。
このとき $(X_i, x_i) \to (X, x)$ は芽の \'etale 射であり、
$Y_i \to Y$ は \'etale であり、$h_i$ は Definition
\ref{descent-definition-local-source-target} の (3) にある射である。
したがって $f$ は $\mathcal{P}$ をもつ。
逆に $f$ が $\mathcal{P}$ をもつならば、Definition
\ref{descent-definition-local-source-target} の (2) により
各 $X_i \to Y_i$ は $\mathcal{P}$ をもつ。

\medskip\noindent
(3) の証明。$f : X \to Y$ が $\mathcal{P}$ をもち、
$g : Y \to Z$ が \'etale であると仮定する。すべての $x \in X$ に対して、
$(X, x) \to (X, x)$ を芽の \'etale 射とみなせ、$Y \to Z$ は
\'etale 射であり、$h = f$ は Definition
\ref{descent-definition-local-source-target} の (3) にある射である。
したがって $g \circ f$ は $\mathcal{P}$ をもつ。

\medskip\noindent
(4) の証明。$f : X \to Y$ を射、$g : Y \to Z$ を \'etale 射とし、
$g \circ f$ が $\mathcal{P}$ をもつとする。Definition
\ref{descent-definition-local-source-target} の (2) により
$\text{pr}_Y : Y \times_Z X \to Y$ は $\mathcal{P}$ をもつ。
一方、射 $(f, 1) : X \to Y \times_Z X$ は \'etale 射影
$\text{pr}_X : Y \times_Z X \to X$ の切断なので \'etale である。
Morphisms, Lemma \ref{morphisms-lemma-etale-permanence} を参照。
したがって Definition \ref{descent-definition-local-source-target} の (1) により
$f = \text{pr}_Y \circ (f, 1)$ は $\mathcal{P}$ をもつ。
\end{proof}

\noindent
次の補題は Morphisms, Lemma
\ref{morphisms-lemma-locally-P-characterize} の類似である。

\begin{lemma}
\label{descent-lemma-local-source-target-characterize}
$\mathcal{P}$ を始域・標的上 \'etale 局所的なスキームの射の性質とし、
$f : X \to Y$ をスキームの射とする。次は同値である。
\begin{enumerate}
\item[(a)] $f$ は性質 $\mathcal{P}$ をもつ。
\item[(b)] すべての $x \in X$ に対して、芽の \'etale 射
$a : (U, u) \to (X, x)$、\'etale 射 $b : V \to Y$、および射
$h : U \to V$ で、$f \circ a = b \circ h$ を満たし
$h$ が $\mathcal{P}$ をもつものが存在する。
\item[(c)]
任意の可換図式
$$
\xymatrix{
U \ar[d]_a \ar[r]_h & V \ar[d]^b \\
X \ar[r]^f & Y
}
$$
で $a$, $b$ が \'etale ならば、射 $h$ は $\mathcal{P}$ をもつ。
\item[(d)] (c) にある形のある図式で、$a : U \to X$ が全射であり、
$h$ が $\mathcal{P}$ をもつものが存在する。
\item[(e)] \'etale 被覆 $\{Y_i \to Y\}_{i \in I}$ で、各基底変換
$Y_i \times_Y X \to Y_i$ が $\mathcal{P}$ をもつものが存在する。
\item[(f)] \'etale 被覆 $\{X_i \to X\}_{i \in I}$ で、各合成
$X_i \to Y$ が $\mathcal{P}$ をもつものが存在する。
\item[(g)] \'etale 被覆 $\{Y_i \to Y\}_{i \in I}$ と、各
$i \in I$ に対する \'etale 被覆
$\{X_{ij} \to Y_i \times_Y X\}_{j \in J_i}$ で、各射
$X_{ij} \to Y_i$ が $\mathcal{P}$ をもつものが存在する。
\end{enumerate}
\end{lemma}

\begin{proof}
(a) と (b) の同値性は Definition
\ref{descent-definition-local-source-target} の一部である。
(a) と (e) の同値性は Lemma
\ref{descent-lemma-local-source-target-implies} の (2) であり、
(a) と (f) の同値性は Lemma
\ref{descent-lemma-local-source-target-implies} の (1) である。
(a) は (e) と (f) のそれぞれと同値だから、(a) は (g) と同値である。

\medskip\noindent
(c) が (a) を含意することは明らかである。(a) が成り立つならば、
(c) にある任意の図式に対して、Definition
\ref{descent-definition-local-source-target} の (1) により射
$f \circ a$ は $\mathcal{P}$ をもつ。すると Lemma
\ref{descent-lemma-local-source-target-implies} の (4) により
$h$ は $\mathcal{P}$ をもつ。したがって (a) と (c) は同値である。
(c) が (d) を含意することは明らかである。(d) が (a) を含意することを
見るため、(c) にある形の図式で、$a : U \to X$ が全射であり
$h$ が $\mathcal{P}$ をもつものを仮定する。Lemma
\ref{descent-lemma-local-source-target-implies} の (3) により
$b \circ h$ は $\mathcal{P}$ をもつ。$\{a : U \to X\}$ は
\'etale 被覆なので、Lemma
\ref{descent-lemma-local-source-target-implies} の (1) により
$f$ は $\mathcal{P}$ をもつ。
\end{proof}

\noindent
次の補題の結果は形式的なものではなく、実際に \'etale 射の幾何に
関する何らかの事実を用いるように思われる。

\begin{lemma}
\label{descent-lemma-etale-local-source-target}
$\mathcal{P}$ をスキームの射の性質とする。次を仮定する。
\begin{enumerate}
\item $\mathcal{P}$ は始域上 \'etale 局所的である。
\item $\mathcal{P}$ は標的上 \'etale 局所的である。
\item $\mathcal{P}$ は開埋め込みの後合成で安定である。すなわち、
$f : X \to Y$ が $\mathcal{P}$ をもち、$Y \subset Z$ が
開部分スキームならば、$X \to Z$ は $\mathcal{P}$ をもつ。
\end{enumerate}
このとき $\mathcal{P}$ は始域・標的上 \'etale 局所的である。
\end{lemma}

\begin{proof}
$\mathcal{P}$ を補題の条件 (1)、(2)、(3) を満たすスキームの射の
性質とする。Lemma \ref{descent-lemma-precompose-property-local-source} により、
$\mathcal{P}$ は \'etale 射の前合成で安定である。Lemma
\ref{descent-lemma-pullback-property-local-target} により、$\mathcal{P}$ は
\'etale 基底変換で安定である。したがって Definition
\ref{descent-definition-local-source-target} の (3) が成り立つことを
示せば十分である。

\medskip\noindent
より正確には、$f : X \to Y$ を Definition
\ref{descent-definition-local-source-target} の (3)(b) を満たすスキームの射とする。
言い換えると、すべての $x \in X$ に対して、\'etale 射
$a_x : U_x \to X$、$x$ に写る点 $u_x \in U_x$、\'etale 射
$b_x : V_x \to Y$、および射 $h_x : U_x \to V_x$ で、
$f \circ a_x = b_x \circ h_x$ を満たし $h_x$ が $\mathcal{P}$ を
もつものが存在する。$f$ が $\mathcal{P}$ をもつことを示せば、
補題の証明は完了する。
$U = \coprod U_x$、$a = \coprod a_x$、$V = \coprod V_x$、
$b = \coprod b_x$、$h = \coprod h_x$ とおく。可換図式
$$
\xymatrix{
U \ar[d]_a \ar[r]_h & V \ar[d]^b \\
X \ar[r]^f & Y
}
$$
を得る。ここで $a$, $b$ は \'etale であり、$a$ は全射である。
$h_x$ はすべて同様であり、$\mathcal{P}$ は標的上
\'etale 局所的なので、$h$ は $\mathcal{P}$ をもつ。
$a$ は全射で $\mathcal{P}$ は始域上 \'etale 局所的だから、
$b \circ h$ が $\mathcal{P}$ をもつことを示せば十分である。
したがって補題は、$\mathcal{P}$ が \'etale 射の後合成で安定であることの
証明に帰着される。

\medskip\noindent
証明の残りでは、$f : X \to Y$ を性質 $\mathcal{P}$ をもつ射とし、
$g : Y \to Z$ を \'etale 射とする。次の命題を考える。
\begin{enumerate}
\item[(-)] 追加の仮定なしに $g \circ f$ は性質 $\mathcal{P}$ をもつ。
\item[(A)] $Z$ がアフィンならば $g \circ f$ は性質 $\mathcal{P}$ をもつ。
\item[(AA)] $X$ と $Z$ がアフィンならば
$g \circ f$ は性質 $\mathcal{P}$ をもつ。
\item[(AAA)] $X$, $Y$, $Z$ がアフィンならば
$g \circ f$ は性質 $\mathcal{P}$ をもつ。
\end{enumerate}
(-) を示せば補題の証明は完了する。

\medskip\noindent
主張 1：(AAA) $\Rightarrow$ (AA)。
実際、$f : X \to Y$、$g : Y \to Z$ を上のとおりとし、
$X$, $Z$ がアフィンであるとする。$X$ はアフィン、したがって
準コンパクトなので、アフィン開集合 $Y_i \subset Y$、
$i = 1, \ldots, n$ を有限個とり、
$X = \bigcup_{i = 1, \ldots, n} f^{-1}(Y_i)$ とできる。
$X_i = f^{-1}(Y_i)$ とおく。Lemma
\ref{descent-lemma-pullback-property-local-target} により、各射
$X_i \to Y_i$ は $\mathcal{P}$ をもつ。$\mathcal{P}$ は標的上
\'etale 局所的だから、
$\coprod_{i = 1, \ldots, n} X_i \to \coprod_{i = 1, \ldots, n} Y_i$
は $\mathcal{P}$ をもつ。(AAA) を
$\coprod_{i = 1, \ldots, n} X_i \to \coprod_{i = 1, \ldots, n} Y_i$
と \'etale 射 $\coprod_{i = 1, \ldots, n} Y_i \to Z$ に適用すると、
$\coprod_{i = 1, \ldots, n} X_i \to Z$ は $\mathcal{P}$ をもつ。
ここで $\{\coprod_{i = 1, \ldots, n} X_i \to X\}$ は \'etale 被覆であり、
$\mathcal{P}$ は始域上 \'etale 局所的なので、望むとおり
$X \to Z$ は $\mathcal{P}$ をもつ。

\medskip\noindent
主張 2：(AAA) $\Rightarrow$ (A)。
実際、$f : X \to Y$、$g : Y \to Z$ を上のとおりとし、
$Z$ がアフィンであるとする。アフィン開被覆 $X = \bigcup X_i$ を選ぶ。
$\mathcal{P}$ は始域上 \'etale 局所的なので、各
$f|_{X_i} : X_i \to Y$ は $\mathcal{P}$ をもつ。
主張 1 により (AAA) から従う (AA) を用いると、各 $i$ に対して
$X_i \to Z$ は $\mathcal{P}$ をもつ。
$\{X_i \to X\}$ は \'etale 被覆であり、$\mathcal{P}$ は始域上
\'etale 局所的なので、$X \to Z$ は $\mathcal{P}$ をもつ。

\medskip\noindent
主張 3：(AAA) $\Rightarrow$ (-)。
実際、$f : X \to Y$、$g : Y \to Z$ を上のとおりとする。
アフィン開被覆 $Z = \bigcup Z_i$ を選び、
$Y_i = g^{-1}(Z_i)$、$X_i = f^{-1}(Y_i)$ とおく。Lemma
\ref{descent-lemma-pullback-property-local-target} により、各射
$X_i \to Y_i$ は $\mathcal{P}$ をもつ。主張 2 により (AAA) から従う
(A) を用いると、各 $i$ に対して $X_i \to Z_i$ は $\mathcal{P}$ をもつ。
$\mathcal{P}$ は標的上局所的であり
$X_i = (g \circ f)^{-1}(Z_i)$ だから、
$X \to Z$ は $\mathcal{P}$ をもつ。

\medskip\noindent
したがって補題を示すには (AAA) を示せば十分である。
$f : X \to Y$ と $g : Y \to Z$ を上のとおりとし、
$X, Y, Z$ がアフィンであるとする。アフィン間の \'etale 射の
ファイバーは普遍的に有界であることに注意する。Morphisms,
Lemma \ref{morphisms-lemma-etale-locally-quasi-finite} および
Lemma \ref{morphisms-lemma-locally-quasi-finite-qc-source-universally-bounded}
を参照。したがって $Y \to Z$ のファイバーの次数を上から押さえる整数
$n$ について帰納法を行える。\'etale 射の場合のこの整数の記述については
Morphisms, Lemma \ref{morphisms-lemma-etale-universally-bounded} を参照。
$n = 1$ ならば、Lemma
\ref{descent-lemma-universally-injective-etale-open-immersion} により
$Y \to Z$ は開埋め込みであり、結果は補題の仮定 (3) から従う。
$n > 1$ と仮定する。

\medskip\noindent
次の可換図式を考える。
$$
\xymatrix{
X \times_Z Y \ar[d] \ar[r]_{f_Y} &
Y \times_Z Y \ar[d] \ar[r]_-{\text{pr}} &
Y \ar[d] \\
X \ar[r]^f &
Y \ar[r]^g &
Z
}
$$
開閉部分スキームへの分解
$Y \times_Z Y = \Delta_{Y/Z}(Y) \amalg Y'$ があることに注意する。
Morphisms, Lemma
\ref{morphisms-lemma-diagonal-unramified-morphism} を参照。
基底変換なので、第二射影
$\text{pr} : Y \times_Z Y \to Y$ のファイバーの次数は $n$ 以下である。
Morphisms, Lemma
\ref{morphisms-lemma-base-change-universally-bounded} を参照。
一方、$\text{pr}|_{\Delta(Y)} : \Delta(Y) \to Y$ は同型射であり、
各ファイバーはちょうど一点をもつ。したがって Morphisms, Lemma
\ref{morphisms-lemma-etale-universally-bounded} を適用すると、制限
$\text{pr}|_{Y'} : Y' \to Y$ のファイバーの次数は $n - 1$ 以下である。
$X' = f_Y^{-1}(Y')$ とおく。図式は
$$
\xymatrix{
X \amalg X' \ar@{=}[d] \ar[r]_-{f \amalg f'} &
\Delta(Y) \amalg Y' \ar@{=}[d] \ar[r] &
Y \ar@{=}[d] \\
X \times_Z Y \ar[r]^{f_Y} &
Y \times_Z Y \ar[r]^-{\text{pr}} &
Y
}
$$
となる。$\mathcal{P}$ は標的上 \'etale 局所的であり、したがって
\'etale 基底変換で安定なので（Lemma
\ref{descent-lemma-pullback-property-local-target}）、$f_Y$ は
$\mathcal{P}$ をもつ。さらに $\mathcal{P}$ は始域上 \'etale 局所的だから、
$f' = f_Y|_{X'}$ は $\mathcal{P}$ をもつ。帰納法の仮定により
$X' \to Y$ は $\mathcal{P}$ をもつ。
$\mathcal{P}$ は始域上局所的であり、
$\{X \to X \times_Z Y, X' \to X \times_Y Z\}$ は \'etale 被覆なので、
$\text{pr} \circ f_Y$ は $\mathcal{P}$ をもつ。
$g \circ f$ は射 $g \circ f : X \to g(Y)$ とみなせることに注意する。
$\text{pr} \circ f_Y$ は \'etale 被覆 $\{Y \to g(Y)\}$ による
$g \circ f : X \to g(Y)$ の引き戻しであり、$\mathcal{P}$ は標的上
\'etale 局所的だから、$g \circ f : X \to g(Y)$ は性質
$\mathcal{P}$ をもつ。最後に、補題の仮定 (3) をもう一度適用すると、
$g \circ f : X \to Z$ は性質 $\mathcal{P}$ をもつ。
\end{proof}

\begin{remark}
\label{descent-remark-list-local-source-target}
Lemma \ref{descent-lemma-etale-local-source-target} と本章の先行する節で行った
作業を用いると、始域・標的上 \'etale 局所的な射の型の一覧を
容易に作れる。各場合について、その性質が始域上 \'etale 局所的であることを
含意する補題と、標的上 \'etale 局所的であることを含意する補題を挙げる。
どの場合も Lemma \ref{descent-lemma-etale-local-source-target} の第三の仮定は
自明に確かめられるので省略する。一覧は次のとおりである。
\begin{enumerate}
\item 平坦：Lemmas \ref{descent-lemma-flat-fpqc-local-source} および
\ref{descent-lemma-descending-property-flat} を参照。
\item 局所有限表示：Lemmas
\ref{descent-lemma-locally-finite-presentation-fppf-local-source} および
\ref{descent-lemma-descending-property-locally-finite-presentation} を参照。
\item 局所有限型：Lemmas
\ref{descent-lemma-locally-finite-type-fppf-local-source} および
\ref{descent-lemma-descending-property-locally-finite-type} を参照。
\item 普遍開：Lemmas \ref{descent-lemma-universally-open-fppf-local-source} および
\ref{descent-lemma-descending-property-universally-open} を参照。
\item syntomic：Lemmas \ref{descent-lemma-syntomic-syntomic-local-source} および
\ref{descent-lemma-descending-property-syntomic} を参照。
\item smooth：Lemmas \ref{descent-lemma-smooth-smooth-local-source} および
\ref{descent-lemma-descending-property-smooth} を参照。
\item \'etale：Lemmas \ref{descent-lemma-etale-etale-local-source} および
\ref{descent-lemma-descending-property-etale} を参照。
\item 局所準有限：Lemmas
\ref{descent-lemma-locally-quasi-finite-etale-local-source} および
\ref{descent-lemma-descending-property-quasi-finite} を参照。
\item 不分岐：Lemmas \ref{descent-lemma-unramified-etale-local-source} および
\ref{descent-lemma-descending-property-unramified} を参照。
\item G-不分岐：Lemmas \ref{descent-lemma-unramified-etale-local-source} および
\ref{descent-lemma-descending-property-unramified} を参照。
\item 必要に応じてここにさらに追加する。
\end{enumerate}
\end{remark}

\begin{remark}
\label{descent-remark-compare-definitions}
この時点で、射の性質 $\mathcal{P}$ が「始域と標的上 \'etale 局所的」
であることの意味について、三つの定義が考えられる。
\begin{enumerate}
\item[(ST)] $\mathcal{P}$ は始域上 \'etale 局所的であり、かつ
$\mathcal{P}$ は標的上 \'etale 局所的である。
\item[(DM)] （Deligne--Mumford の論文 \cite[Page 100]{DM} における定義）
任意の図式
$$
\xymatrix{
U \ar[d]_a \ar[r]_h & V \ar[d]^b \\
X \ar[r]^f & Y
}
$$
で縦の矢が全射 \'etale ならば
$\mathcal{P}(h) \Leftrightarrow \mathcal{P}(f)$ が成り立つ。
\item[(SP)] $\mathcal{P}$ は始域・標的上 \'etale 局所的である。
\end{enumerate}
本節で (SP) $\Rightarrow$ (DM) $\Rightarrow$ (ST) を見た。
Examples \ref{descent-example-silly-one} と \ref{descent-example-silly-two} は、
どちらの含意も逆向きにはできないことを示す。最後に Lemma
\ref{descent-lemma-etale-local-source-target} は、開埋め込みの後合成で安定な
射の性質を考えると差が消えることを示す。実際上は常にこの場合になる。
\end{remark}

\begin{lemma}
\label{descent-lemma-etale-etale-local-source-target}
$\mathcal{P}$ を始域・標的上 \'etale 局所的なスキームの射の性質とする。
スキームの可換図式
$$
\vcenter{
\xymatrix{
X' \ar[d]_{g'} \ar[r]_{f'} & Y' \ar[d]^g \\
X \ar[r]^f & Y
}
}
\quad\text{点を伴い}\quad
\vcenter{
\xymatrix{
x' \ar[d] \ar[r] & y' \ar[d] \\
x \ar[r] & y
}
}
$$
が与えられ、$g'$ が $x'$ で \'etale、$g$ が $y'$ で \'etale ならば、
$x \in W(f) \Leftrightarrow x' \in W(f')$
である。ここで $W(-)$ は Lemma
\ref{descent-lemma-largest-open-of-the-source} におけるものとする。
\end{lemma}

\begin{proof}
Lemma \ref{descent-lemma-local-source-target-implies} により $\mathcal{P}$ は
始域上 \'etale 局所的だから、Lemma
\ref{descent-lemma-largest-open-of-the-source} を適用できる。

\medskip\noindent
$x \in W(f)$ と仮定する。$U' \subset X'$ と $V' \subset Y'$ を
$x'$ と $y'$ の開近傍で、$f'(U') \subset V'$、
$g'(U') \subset W(f)$ を満たし、$g'|_{U'}$ と $g|_{V'}$ が
\'etale となるものとする。Definition
\ref{descent-definition-local-source-target} の性質 (1) により
$f \circ g'|_{U'} = g \circ f'|_{U'}$ は $\mathcal{P}$ をもつ。
Lemma \ref{descent-lemma-local-source-target-implies} の (4) により
$f'|_{U'} : U' \to V'$ は性質 $\mathcal{P}$ をもつ。
さらに Lemma \ref{descent-lemma-local-source-target-implies} の (3) により
$f'_{U'} : U' \to Y'$ は $\mathcal{P}$ をもつ。
ゆえに定義により $U' \subset W(f')$ であり、$x' \in W(f')$ である。

\medskip\noindent
$x' \in W(f')$ と仮定する。$U' \subset X'$ と $V' \subset Y'$ を
$x'$ と $y'$ の開近傍で、$f'(U') \subset V'$、
$U' \subset W(f')$ を満たし、$g'|_{U'}$ と $g|_{V'}$ が
\'etale となるものとする。$W(f')$ の定義により
$U' \to Y'$ は $\mathcal{P}$ をもつ。Lemma
\ref{descent-lemma-local-source-target-implies} の (4) により
$U' \to V'$ は $\mathcal{P}$ をもち、Lemma
\ref{descent-lemma-local-source-target-implies} の (3) により
$U' \to Y$ は $\mathcal{P}$ をもつ。
$U \subset X$ を \'etale（したがって開）射
$g'|_U' : U' \to X$ の像とする。このとき $\{U' \to U\}$ は
\'etale 被覆なので、Lemma \ref{descent-lemma-local-source-target-implies} の
(1) により $U \to Y$ は $\mathcal{P}$ をもつ。
したがって定義により $U \subset W(f)$ であり、$x \in W(f)$ である。
\end{proof}

\begin{lemma}
\label{descent-lemma-orbits}
$k$ を体とし、$n \geq 2$ とする。$1 \leq i, j \leq n$、
$i \not = j$、$d \geq 0$ に対して、座標で
$$
(x_1, \ldots, x_n) \longmapsto
(x_1, \ldots, x_{i - 1}, x_i + x_j^d, x_{i + 1}, \ldots, x_n)
$$
で与えられる $\mathbf{A}^n_k$ の自己同型を $T_{i, j, d}$ と記す。
$W \subset \mathbf{A}^n_k$ を空でない開部分スキームで、上のすべての
$i, j, d$ に対して $T_{i, j, d}(W) = W$ を満たすものとする。
このとき、$W = \mathbf{A}^n_k$ であるか、または $k$ の標数が
$p > 0$ であり $\mathbf{A}^n_k \setminus W$ は座標が
$\mathbf{F}_p$ 上代数的な閉点の有限集合である。
\end{lemma}

\begin{proof}
これを示す際、$k$ を任意の拡大体で置き換えてよい。
$Z$ を $\mathbf{A}^n_k \setminus W$ の既約成分とする。
矛盾を導くため $\dim(Z) \geq 1$ と仮定する。
このとき拡大体 $k'/k$ と、
$Z_{k'} \subset \mathbf{A}^n_{k'}$
の $k'$-値点 $\xi = (\xi_1, \ldots, \xi_n) \in (k')^n$ で、
$x_1, \ldots, x_n$ の少なくとも一つが素体上超越的なものが存在する。
主張：変換 $T_{i, j, d}$ が生成する群による $\xi$ の軌道は
$\mathbf{A}^n_{k'}$ で Zariski 稠密である。この主張から望む矛盾が従う。

\medskip\noindent
$k'$ の標数が零ならば、作用素 $T_{i, j, 0}$ だけで十分である。
実際、これらは $\xi$ を点
$$
(\xi_1 + a_1, \ldots, \xi_n + a_n)
$$
へ移す。ここで $(a_1, \ldots, a_n) \in \mathbf{Z}_{\geq 0}^n$ は任意である。
標数が $p > 0$ ならば、番号を付け替えることにより
$\xi_n$ が $\mathbf{F}_p$ 上超越的であると仮定してよい。
$i < n$ に対する作用素 $T_{i, n, d}$ を順に適用すると、
$\xi$ の軌道は元
$$
(\xi_1 + P_1(\xi_n), \ldots, \xi_{n - 1} + P_{n - 1}(\xi_n), \xi_n)
$$
を含む。ここで $(P_1, \ldots, P_{n - 1}) \in \mathbf{F}_p[t]$ は任意である。
したがって軌道の Zariski 閉包は座標超平面 $x_n = \xi_n$ を含む。
別の座標について議論を繰り返すと、Zariski 閉包は
$x_i = \xi_i + P(\xi_n)$
を含む。ここで $P \in \mathbf{F}_p[t]$ は
$\xi_i + P(\xi_n)$ が $\mathbf{F}_p$ 上超越的となる任意の多項式である。
このような $P$ は無限個あるから主張が従う。

\medskip\noindent
もちろん、$Z = \{z\}$ が次元 $0$ で、$\kappa(z)$ における $z$ の座標が
$\mathbf{F}_p$ 上代数的でない場合にも、前段落の議論は適用できる。
これで補題が従う。
\end{proof}

\begin{lemma}
\label{descent-lemma-etale-tau-local-source-target}
$\mathcal{P}$ をスキームの射の性質とする。次を仮定する。
\begin{enumerate}
\item $\mathcal{P}$ は始域上 \'etale 局所的である。
\item $\mathcal{P}$ は標的上 smooth 局所的である。
\item $\mathcal{P}$ は開埋め込みの後合成で安定である。すなわち、
$f : X \to Y$ が $\mathcal{P}$ をもち、$Y \subset Z$ が
開部分スキームならば、$X \to Z$ は $\mathcal{P}$ をもつ。
\end{enumerate}
スキームの可換図式
$$
\vcenter{
\xymatrix{
X' \ar[d]_{g'} \ar[r]_{f'} & Y' \ar[d]^g \\
X \ar[r]^f & Y
}
}
\quad\text{点を伴い}\quad
\vcenter{
\xymatrix{
x' \ar[d] \ar[r] & y' \ar[d] \\
x \ar[r] & y
}
}
$$
が与えられ、$g$ が smooth $y'$ であり
$X' \to X \times_Y Y'$ が $x'$ で \'etale ならば、
$x \in W(f) \Leftrightarrow x' \in W(f')$ である。
ここで $W(-)$ は Lemma \ref{descent-lemma-largest-open-of-the-source}
におけるものとする。
\end{lemma}

\begin{proof}
$\mathcal{P}$ は始域上 \'etale 局所的なので、$x \in W(f)$ であることと、
$X \times_Y Y'$ における $x$ の像が
$W(X \times_Y Y' \to Y')$ に属することは同値である。
したがって補題の図式は Cartesian であると仮定してよい。

\medskip\noindent
$x \in W(f)$ と仮定する。$\mathcal{P}$ は標的上 smooth 局所的なので、
$(g')^{-1}W(f) = W(f) \times_Y Y' \to Y'$ は $\mathcal{P}$ をもつ。
したがって $(g')^{-1}W(f) \subset W(f')$ であり、
$x' \in W(f')$ と結論する。

\medskip\noindent
$x' \in W(f')$ と仮定する。$y'$ の任意の開近傍
$V' \subset Y'$ に対して、$Y'$ を $V'$ で、$X'$ を
$U' = (f')^{-1}V'$ で置き換えてよい。実際、$V' \to Y'$ は smooth であり、
したがって $W(f') \to Y'$ の基底変換
$W(f') \cap U' \to V'$ は性質 $\mathcal{P}$ をもつ。
そこで $Y$ 上の \'etale 射 $Y' \to \mathbf{A}^n_Y$ が存在すると
仮定してよい。Morphisms, Lemma
\ref{morphisms-lemma-smooth-etale-over-affine-space} を参照。図式は
$$
\xymatrix{
X' \ar[r] \ar[d] & Y' \ar[d] \\
\mathbf{A}^n_X \ar[r]_{f_n} \ar[d] & \mathbf{A}^n_Y \ar[d] \\
X \ar[r]^f & Y
}
$$
となる。Lemma \ref{descent-lemma-etale-local-source-target}
（および \'etale 被覆は smooth 被覆でもあること）により、
$\mathcal{P}$ は始域・標的上 \'etale 局所的である。Lemma
\ref{descent-lemma-etale-etale-local-source-target} により、$W(f')$ は
開集合 $W(f_n) \subset \mathbf{A}^n_X$ の逆像である。
特に $W(f_n)$ は $x$ 上の点を含む。$X$ を $W(f_n)$ の像
（これは開である）で置き換えることにより、
$W(f_n) \to X$ は全射であると仮定してよい。
主張：$W(f_n) = \mathbf{A}^n_X$。
$\mathcal{P}$ は smooth 位相について局所的であり、
$\{\mathbf{A}^n_Y \to Y\}$ は smooth 被覆なので、この主張から
$f$ が $\mathcal{P}$ をもつことが従う。

\medskip\noindent
本質的には、$W(f_n) \subset \mathbf{A}^n_X$ が
$\mathbf{A}^n_X \to X$ の各ファイバーと交わる「平行移動不変」な
開集合であることから主張は従う。しかし、この方針で議論するには、
十分な平行移動を作るため $Y$ 上で \'etale 局所化を行わねばならず、
少し面倒になる。代わりに Lemma \ref{descent-lemma-orbits} の自己同型と
アフィン空間の \'etale 射を用いる。$n \geq 2$ と仮定してよい。
実際、$n = 0$ ならば証明は終わっている。$n = 1$ ならば図式
$$
\xymatrix{
\mathbf{A}^2_X \ar[r]_{f_2} \ar[d]_p & \mathbf{A}^2_Y \ar[d] \\
\mathbf{A}^1_X \ar[r]^{f_1} & \mathbf{A}^1_Y
}
$$
を考える。$p^{-1}(W(f_1)) \subset W(f_2)$ である
（証明の第一段落を参照）。したがって $W(f_2) \to X$ は依然として
全射であり、$f_2$ を用いてよい。$n \geq 2$ と仮定する。

\medskip\noindent
$1 \leq i, j \leq n$、$i \not = j$、$d \geq 0$ に対して、
Lemma \ref{descent-lemma-orbits} で定義した $\mathbf{A}^n$ の自己同型を
$T_{i, j, d}$ と記す。このとき可換図式
$$
\xymatrix{
\mathbf{A}^n_X \ar[r]_{f_n} \ar[d]_{T_{i, j, d}} &
\mathbf{A}^n_Y \ar[d]^{T_{i, j, d}} \\
\mathbf{A}^n_X \ar[r]^{f_n} & \mathbf{A}^n_Y
}
$$
を得る。縦の矢は同型射である。したがって
$T_{i, j, d}(W(f_n)) = W(f_n)$ である。Lemma \ref{descent-lemma-orbits} を
適用すると、任意の $x \in X$ に対してファイバー
$W(f_n)_x \subset \mathbf{A}^n_x$ は $\mathbf{A}^n_x$
（これが望む場合）に等しいか、または $\kappa(x)$ の標数が
$p > 0$ で、$W(f_n)_x$ は閉点の有限集合
$Z_x \subset \mathbf{A}^n_x$ の補集合である。
第二の場合は起こり得ない。実際、次で与えられる射
$T_p : \mathbf{A}^n \to \mathbf{A}^n$ を考える。
$$
(x_1, \ldots, x_n) \mapsto (x_1 - x_1^p, \ldots, x_n - x_n^p)
$$
上と同様に可換図式
$$
\xymatrix{
\mathbf{A}^n_X \ar[r]_{f_n} \ar[d]_{T_p} &
\mathbf{A}^n_Y \ar[d]^{T_p} \\
\mathbf{A}^n_X \ar[r]^{f_n} & \mathbf{A}^n_Y
}
$$
を得る。射 $T_p : \mathbf{A}^n_X \to \mathbf{A}^n_X$ は
$x$ 上のすべての点で \'etale であり、射
$T_p : \mathbf{A}^n_Y \to \mathbf{A}^n_Y$ は
$Y$ における $x$ の像の上のすべての点で \'etale である。
（詳細は省略する。ヒント：導関数を計算せよ。）
したがって
$$
T_p^{-1}(W) \cap \mathbf{A}^n_x = W \cap \mathbf{A}^n_x
$$
である。これは Lemma \ref{descent-lemma-etale-etale-local-source-target} による
（$\mathcal{P}$ が始域・標的上 \'etale 局所的であることは既に見た）。
$T_p : \mathbf{A}^n_x \to \mathbf{A}^n_x$ は次数
$p^n > 1$ の有限 \'etale 射なので、$Z_x$ が空でなければ、
それは自身より大きい $T_p^{-1}(Z_x)$ を含む。この矛盾で証明が終わる。
\end{proof}





\section{始域・標的上局所的な芽の射の性質}
\label{descent-section-local-source-target-at-point}

\noindent
本節では、スキームの芽の射について、Section
\ref{descent-section-properties-etale-local-source-target} の内容の類似を論じる。

\begin{definition}
\label{descent-definition-local-source-target-at-point}
$\mathcal{Q}$ をスキームの芽の射の性質とする。芽の任意の可換図式
$$
\xymatrix{
(U', u') \ar[d]_a \ar[r]_{h'} & (V', v') \ar[d]^b \\
(U, u) \ar[r]^h & (V, v)
}
$$
で、縦の矢が \'etale ならば
$\mathcal{Q}(h) \Leftrightarrow \mathcal{Q}(h')$ が成り立つとき、
$\mathcal{Q}$ は {\it 始域・標的上 \'etale 局所的}であるという。
\end{definition}

\begin{lemma}
\label{descent-lemma-local-source-target-global-implies-local}
$\mathcal{P}$ を始域・標的上 \'etale 局所的なスキームの射の性質とする。
次の規則で定義される芽の射の性質 $\mathcal{Q}$ を考える。
$$
\mathcal{Q}((X, x) \to (S, s))
\Leftrightarrow
\text{代表射 }U \to S\text{ で }\mathcal{P}\text{ をもつものが存在する}
$$
このとき $\mathcal{Q}$ は Definition
\ref{descent-definition-local-source-target-at-point} の意味で
始域・標的上 \'etale 局所的である。
\end{lemma}

\begin{proof}
芽の射 $(X, x) \to (S, s)$ が $\mathcal{Q}$ をもつならば、
$x$ の任意に小さい近傍 $U \subset X$ と $s$ の任意に小さい近傍
$V \subset S$ で、$(X, x) \to (S, s)$ の代表射
$U \to V$ が $\mathcal{P}$ をもつものが存在する。
これは Lemma \ref{descent-lemma-local-source-target-implies} から従う。
$$
\xymatrix{
(U', u') \ar[r]_{h'} \ar[d]_a & (V', v') \ar[d]^b \\
(U, u) \ar[r]^h & (V, v)
}
$$
を Definition \ref{descent-definition-local-source-target-at-point} にある図式とする。
$U_1 \subset U$ と $h$ の代表射 $h_1 : U_1 \to V$ を選ぶ。
$V'_1 \subset V'$ と $b$ の \'etale 代表射 $b_1 : V'_1 \to V$ を選ぶ
（Definition \ref{descent-definition-etale-morphism-germs}）。
$U'_1 \subset U'$ と、$a$, $h'$ の代表射
$a_1 : U'_1 \to U_1$、$h'_1 : U'_1 \to V'_1$ を、$a_1$ が
\'etale となるように選ぶ。$U'_1$ を縮小して
$h_1 \circ a_1 = b_1 \circ h'_1$ と仮定してよい。
証明冒頭の注意により、示すべきことは
$u' \in W(h'_1) \Leftrightarrow u \in W(h_1)$ である。
ここで $W(-)$ は Lemma \ref{descent-lemma-largest-open-of-the-source} に
おけるものとする。したがって補題は Lemma
\ref{descent-lemma-etale-etale-local-source-target} から従う。
\end{proof}

\begin{lemma}
\label{descent-lemma-local-source-target-local-implies-global}
$\mathcal{P}$ を始域・標的上 \'etale 局所的なスキームの射の性質とし、
$Q$ を対応する芽の射の性質とする。Lemma
\ref{descent-lemma-local-source-target-global-implies-local} を参照。
$f : X \to Y$ をスキームの射とする。次は同値である。
\begin{enumerate}
\item $f$ は性質 $\mathcal{P}$ をもつ。
\item すべての $x \in X$ に対して芽の射
$(X, x) \to (Y, f(x))$ は性質 $\mathcal{Q}$ をもつ。
\end{enumerate}
\end{lemma}

\begin{proof}
含意 (1) $\Rightarrow$ (2) は定義から直接従う。
含意 (2) $\Rightarrow$ (1) も Definition
\ref{descent-definition-local-source-target} の (3) から従う。
\end{proof}

\noindent
芽の射 $(X, x) \to (S, s)$ は局所環の矛盾なく定まる写像を決定する。
したがって次の補題には意味がある。

\begin{lemma}
\label{descent-lemma-flat-at-point}
芽の射の性質
$$
\mathcal{P}((X, x) \to (S, s)) =
\mathcal{O}_{S, s} \to \mathcal{O}_{X, x}\text{ は平坦である}
$$
は始域・標的上 \'etale 局所的である。
\end{lemma}

\begin{proof}
Definition \ref{descent-definition-local-source-target-at-point} にある図式から、
局所環の局所準同型の図式
$$
\xymatrix{
\mathcal{O}_{U', u'} & \mathcal{O}_{V', v'} \ar[l] \\
\mathcal{O}_{U, u} \ar[u] & \mathcal{O}_{V, v} \ar[l] \ar[u]
}
$$
を得る。縦の矢は \'etale 環写像の局所化であり、特に本質的有限表示、
平坦、かつ不分岐であることに注意する（Algebra, Section
\ref{algebra-section-etale}）。特に縦の写像は忠実平坦である。
Algebra, Lemma \ref{algebra-lemma-local-flat-ff} を参照。
上の水平矢が平坦ならば、$R = \mathcal{O}_{V, v}$、
$S = \mathcal{O}_{U, u}$、$M = \mathcal{O}_{U', u'}$ として
Algebra, Lemma \ref{algebra-lemma-flat-permanence} を適用することにより、
下の水平矢は平坦である。下の水平矢が平坦ならば、環写像
$$
\mathcal{O}_{V', v'} \otimes_{\mathcal{O}_{V, v}} \mathcal{O}_{U, u}
\longleftarrow
\mathcal{O}_{V', v'}
$$
は Algebra, Lemma \ref{algebra-lemma-flat-base-change} により平坦である。
また環写像
$$
\mathcal{O}_{U', u'}
\longleftarrow
\mathcal{O}_{V', v'} \otimes_{\mathcal{O}_{V, v}} \mathcal{O}_{U, u}
$$
は $\mathcal{O}_{U, u}$ の \'etale 環拡大の間の写像の局所化なので、
Algebra, Lemma \ref{algebra-lemma-map-between-etale} により平坦である。
\end{proof}

\begin{lemma}
\label{descent-lemma-etale-on-fiber}
スキームの射の可換図式
$$
\xymatrix{
U' \ar[r] \ar[d] & V' \ar[d] \\
U \ar[r] & V
}
$$
で、縦の矢が \'etale であり、点 $v' \in V'$ が $v \in V$ に写るものを
考える。このときファイバーの射 $U'_{v'} \to U_v$ は \'etale である。
\end{lemma}

\begin{proof}
$U'_v \to U_v$ は \'etale 射 $U' \to U$ の基底変換なので
\'etale であることに注意する。スキーム $U'_v$ は $V'_v$ 上の
スキームである。Morphisms, Lemma
\ref{morphisms-lemma-etale-over-field} により、スキーム $V'_v$ は
$\kappa(v)$ の有限分離拡大のスペクトルの非交和である。
その一つが $v' = \Spec(\kappa(v'))$ である。したがって
$U'_{v'}$ は $U'_v$ の開閉部分スキームであり、
$U'_{v'} \to U'_v \to U_v$ は \'etale である
（開埋め込みと \'etale 射の合成である。Morphisms, Section
\ref{morphisms-section-etale} を参照）。
\end{proof}

\noindent
スキームの芽の射 $(X, x) \to (S, s)$ が与えられたとき、
{\it ファイバー}を芽 $(U_s, x)$ の同型類として定義できる。
ここで $U \to S$ は任意の代表射である。しばしば記号を濫用して
単に $(X_s, x)$ と書く。

\begin{lemma}
\label{descent-lemma-dimension-local-ring-fibre}
$d \in \{0, 1, 2, \ldots, \infty\}$ とする。芽の射の性質
$$
\mathcal{P}_d((X, x) \to (S, s)) =
\text{ファイバーの局所環 }
\mathcal{O}_{X_s, x}
\text{ の次元は }d\text{ である}
$$
は始域・標的上 \'etale 局所的である。
\end{lemma}

\begin{proof}
Definition \ref{descent-definition-local-source-target-at-point} にある図式から、
$u'$ を $u$ に写すファイバーの \'etale 射
$U'_{v'} \to U_v$ を得る。Lemma \ref{descent-lemma-etale-on-fiber} を参照。
したがって結果は Lemma \ref{descent-lemma-dimension-local-ring-local} から従う。
\end{proof}

\begin{lemma}
\label{descent-lemma-transcendence-degree-at-point}
$r \in \{0, 1, 2, \ldots, \infty\}$ とする。芽の射の性質
$$
\mathcal{P}_r((X, x) \to (S, s))
\Leftrightarrow
\text{trdeg}_{\kappa(s)} \kappa(x) = r
$$
は始域・標的上 \'etale 局所的である。
\end{lemma}

\begin{proof}
Definition \ref{descent-definition-local-source-target-at-point} にある図式から、
局所環の局所準同型の図式
$$
\xymatrix{
\mathcal{O}_{U', u'} & \mathcal{O}_{V', v'} \ar[l] \\
\mathcal{O}_{U, u} \ar[u] & \mathcal{O}_{V, v} \ar[l] \ar[u]
}
$$
を得る。縦の矢は \'etale 環写像の局所化であり、特に不分岐である
（Algebra, Section \ref{algebra-section-etale}）。
したがって $\kappa(u')/\kappa(u)$ と $\kappa(v')/\kappa(v)$ は
有限分離体拡大である。ゆえに
$\text{trdeg}_{\kappa(v)} \kappa(u) = \text{trdeg}_{\kappa(v')} \kappa(u)$
であり、補題が示された。
\end{proof}

\noindent
$(X, x)$ をスキームの芽とする。$x$ における $X$ の次元とは、
$X$ における $x$ の開近傍の次元の最小値であり、十分小さい任意の
開近傍はこの次元をもつ。したがってこれは芽の同型類の不変量である。
これを単に $\dim_x(X)$ と記す。

\begin{lemma}
\label{descent-lemma-dimension-at-point}
$d \in \{0, 1, 2, \ldots, \infty\}$ とする。芽の射の性質
$$
\mathcal{P}_d((X, x) \to (S, s))
\Leftrightarrow
\dim_x (X_s) = d
$$
は始域・標的上 \'etale 局所的である。
\end{lemma}

\begin{proof}
Definition \ref{descent-definition-local-source-target-at-point} にある図式から、
$u'$ を $u$ に写すファイバーの \'etale 射
$U'_{v'} \to U_v$ を得る。Lemma \ref{descent-lemma-etale-on-fiber} を参照。
このとき等式 $\dim_u(U_v) = \dim_{u'}(U'_{v'})$ は
Lemma \ref{descent-lemma-dimension-at-point-local} から従う。
\end{proof}







\section{スキーム上のスキームの降下データ}
\label{descent-section-descent-datum}

\noindent
本節の議論の大部分は形式的であり、降下データの定義のみに依存する。
Simplicial Spaces, Section \ref{spaces-simplicial-section-simplicial-descent}
では単体スキームとの関係を検討し、これによって状況はいくらか明確になる。

\begin{definition}
\label{descent-definition-descent-datum}
スキームの射 $f : X \to S$ をとる。
\begin{enumerate}
\item $V \to X$ を $X$ 上のスキームとする。
{\it $V/X/S$ の降下データ}とは、$X \times_S X$ 上のスキームの同型
$\varphi : V \times_S X \to X \times_S V$ であって、
図式
$$
\xymatrix{
V \times_S X \times_S X \ar[rd]^{\varphi_{01}} \ar[rr]_{\varphi_{02}} &
&
X \times_S X \times_S V\\
&
X \times_S V \times_S X \ar[ru]^{\varphi_{12}}
}
$$
が可換となるという{\it コサイクル条件}を満たすものをいう（記法は明らかである）。
\item 組 $(V/X, \varphi)$ を
{\it $X \to S$ に関する降下データ}ともいう。
\item {\it $X \to S$ に関する降下データの射
$g : (V/X, \varphi) \to (V'/X, \varphi')$}とは、
図式
$$
\xymatrix{
V \times_S X \ar[r]_{\varphi} \ar[d]_{g \times \text{id}_X} &
X \times_S V \ar[d]^{\text{id}_X \times g} \\
V' \times_S X \ar[r]^{\varphi'} & X \times_S V'
}
$$
が可換となるような、$X$ 上のスキームの射 $g : V \to V'$ をいう。
\end{enumerate}
\end{definition}

\noindent
上の定義からは、さまざまな「奇跡的な」恒等式が生じる。
たとえば、対角射 $\Delta : X \to X \times_S X$ による
$\varphi$ の引き戻しは、射 $\Delta^*\varphi : V \to V$ とみなせる。
これは $X \times_{\Delta, X \times_S X} (V \times_S X) = V$ であり、
かつ $X \times_{\Delta, X \times_S X} (X \times_S V) = V$
だからである。実際、$\Delta^*\varphi$ は恒等射に等しい。
この内容に不慣れなら、これはよい演習となる。

\begin{remark}
\label{descent-remark-easier}
スキームの射 $X \to S$ をとり、$(V/X, \varphi)$ を
$X \to S$ に関する降下データとする。同型 $\varphi$ は、
$$
(X \times_S X) \times_{\text{pr}_0, X} V
\longrightarrow
(X \times_S X) \times_{\text{pr}_1, X} V
$$
$X \times_S X$ 上のスキームの同型とみなせる。したがって大まかには、
$\varphi$ を写像
$\varphi : \text{pr}_0^*V \to \text{pr}_1^*V$\footnote{残念ながら、
ここでは矢印の「誤った」向きを選んでいる。「関数」と「空間」は双対である
という一般原理により、Definitions \ref{descent-definition-descent-datum} と
\ref{descent-definition-descent-datum-for-family-of-morphisms} では、
Definition \ref{descent-definition-descent-datum-quasi-coherent}
で採用した向きとは反対にすべきであった。}
と考えられる。このときコサイクル条件は
$\text{pr}_{02}^*\varphi =
\text{pr}_{12}^*\varphi \circ \text{pr}_{01}^*\varphi$.
と述べている。この形では、準連接層上の降下データの場合と非常によく似ている。
\end{remark}

\noindent
次に、標的を固定した射の族に対する定義を与える。

\begin{definition}
\label{descent-definition-descent-datum-for-family-of-morphisms}
スキーム $S$ をとる。
$\{X_i \to S\}_{i \in I}$ を、標的 $S$ をもつ射の族とする。
\begin{enumerate}
\item {\it 族 $\{X_i \to S\}$ に関する降下データ
$(V_i, \varphi_{ij})$}とは、各 $i \in I$ に対する
$X_i$ 上のスキーム $V_i$ と、各組 $(i, j) \in I^2$ に対する
$X_i \times_S X_j$ 上のスキームの同型
$\varphi_{ij} : V_i \times_S X_j \to X_i \times_S V_j$
であって、任意の添字の三つ組 $(i, j, k) \in I^3$ に対して図式
$$
\xymatrix{
V_i \times_S X_j \times_S X_k
\ar[rd]^{\text{pr}_{01}^*\varphi_{ij}}
\ar[rr]_{\text{pr}_{02}^*\varphi_{ik}} &
&
X_i \times_S X_j \times_S V_k\\
&
X_i \times_S V_j \times_S X_k
\ar[ru]^{\text{pr}_{12}^*\varphi_{jk}}
}
$$
が $X_i \times_S X_j \times_S X_k$ 上のスキームの図式として
可換となるものをいう（記法は明らかである）。
\item {\it 降下データの射
$\psi : (V_i, \varphi_{ij}) \to (V'_i, \varphi'_{ij})$}とは、
すべての図式
$$
\xymatrix{
V_i \times_S X_j \ar[r]_{\varphi_{ij}} \ar[d]_{\psi_i \times \text{id}} &
X_i \times_S V_j \ar[d]^{\text{id} \times \psi_j} \\
V'_i \times_S X_j \ar[r]^{\varphi'_{ij}} & X_i \times_S V'_j
}
$$
が可換となるような、$X_i$-スキームの射
$\psi_i : V_i \to V'_i$ の族
$\psi = (\psi_i)_{i \in I}$ をいう。
\end{enumerate}
\end{definition}

\noindent
この概念は、たとえば相対曲線のファイバー圏が fpqc 位相に関して
スタックであるかを問う際に自然に現れる
（少なくともスキームに限定するかぎり、スタックではない）。

\begin{remark}
\label{descent-remark-easier-family}
スキーム $S$ をとる。
$\{X_i \to S\}_{i \in I}$ を、標的 $S$ をもつ射の族とする。
$(V_i, \varphi_{ij})$ を $\{X_i \to S\}$ に関する降下データとする。
同型 $\varphi_{ij}$ は、
$$
(X_i \times_S X_j) \times_{\text{pr}_0, X_i} V_i
\longrightarrow
(X_i \times_S X_j) \times_{\text{pr}_1, X_j} V_j
$$
$X_i \times_S X_j$ 上のスキームの同型とみなせる。
したがって大まかには、$\varphi_{ij}$ を $X_i \times_S X_j$ 上の同型
$\text{pr}_0^*V_i \to \text{pr}_1^*V_j$ と考えられる。
このときコサイクル条件は
$\text{pr}_{02}^*\varphi_{ik} =
\text{pr}_{12}^*\varphi_{jk} \circ \text{pr}_{01}^*\varphi_{ij}$.
と述べている。この形では、準連接層上の降下データの場合と非常によく似ている。
\end{remark}

\noindent
通常、射を一つだけ含む族の形で扱う理由は、次の補題にある。

\begin{lemma}
\label{descent-lemma-family-is-one}
スキーム $S$ をとる。
$\{X_i \to S\}_{i \in I}$ を、標的 $S$ をもつ射の族とする。
$X = \coprod_{i \in I} X_i$ とおき、これを $S$-スキームとみなす。
次の標準的な圏同値が存在する。
$$
\begin{matrix}
\text{降下データの圏} \\
\text{族 } \{X_i \to S\}_{i \in I}\text{ に関する}
\end{matrix}
\longrightarrow
\begin{matrix}
\text{降下データの圏} \\
X/S \text{ に関する}
\end{matrix}
$$
これは $(V_i, \varphi_{ij})$ を、$V = \coprod_{i\in I} V_i$ および
$\varphi = \coprod \varphi_{ij}$ により定まる $(V, \varphi)$ に写す。
\end{lemma}

\begin{proof}
$X \times_S X = \coprod_{ij} X_i \times_S X_j$ であり、
高次のファイバー積についても同様であることに注意する。
射 $V \to X$ を与えることは、族 $V_i \to X_i$ を与えることと
まったく同じである。また、降下データ $\varphi$ を与えることは、
族 $\varphi_{ij}$ を与えることとまったく同じである。
\end{proof}

\begin{lemma}
\label{descent-lemma-pullback}
スキーム上のスキームの降下データの引き戻し。
\begin{enumerate}
\item
$$
\xymatrix{
X' \ar[r]_f \ar[d]_{a'} & X \ar[d]^a \\
S' \ar[r]^h & S
}
$$
をスキームの射の可換図式とする。構成
$$
(V \to X, \varphi) \longmapsto f^*(V \to X, \varphi) = (V' \to X', \varphi')
$$
において $V' = X' \times_X V$ とし、$\varphi'$ を合成
$$
\xymatrix{
V' \times_{S'} X' \ar@{=}[r] &
(X' \times_X V) \times_{S'} X' \ar@{=}[r] &
(X' \times_{S'} X') \times_{X \times_S X} (V \times_S X)
\ar[d]^{\text{id} \times \varphi} \\
X' \times_{S'} V' \ar@{=}[r] &
X' \times_{S'} (X' \times_X V) &
(X' \times_{S'} X') \times_{X \times_S X} (X \times_S V) \ar@{=}[l]
}
$$
として定めると、これは $X \to S$ に関する降下データの圏から
$X' \to S'$ に関する降下データの圏への関手を定める。
\item 図式を可換にする二つの射 $f_i : X' \to X$, $i = 0, 1$
が与えられると、関手 $f_0^*$ と $f_1^*$ は標準的に同型である。
\end{enumerate}
\end{lemma}

\begin{proof}
(1) の証明は省略するが、射 $\varphi'$ は
Remark \ref{descent-remark-easier} で導入した記法では
射 $(f \times f)^*\varphi$ であることに注意する。
(2) について、関手的同型を与える射
$f_0^*V \to f_1^*V$ を明示する。実際、$f_0$ と $f_1$ はともに
可換図式に適合するので、$f_i = \text{pr}_i \circ r$ を満たす
一意な射 $r : X' \to X \times_S X$ が存在する。そこで
\begin{eqnarray*}
f_0^*V & = &
X' \times_{f_0, X} V \\
& = &
X' \times_{\text{pr}_0 \circ r, X} V \\
& = &
X' \times_{r, X \times_S X} (X \times_S X) \times_{\text{pr}_0, X} V \\
& \xrightarrow{\varphi} &
X' \times_{r, X \times_S X} (X \times_S X) \times_{\text{pr}_1, X} V \\
& = &
X' \times_{\text{pr}_1 \circ r, X} V \\
& = &
X' \times_{f_1, X} V \\
& = & f_1^*V
\end{eqnarray*}
これが機能することの確認は省略する。
\end{proof}

\begin{definition}
\label{descent-definition-pullback-functor}
Lemma \ref{descent-lemma-pullback} の $S, S', X, X', f, a, a', h$ に対し、
同補題で構成した降下データ上の関手
$$
(V, \varphi) \longmapsto f^*(V, \varphi)
$$
を{\it 引き戻し関手}という。
\end{definition}

\begin{lemma}[族上のスキームの降下データの引き戻し]
\label{descent-lemma-pullback-family}
標的を固定した射の族
$\mathcal{U} = \{U_i \to S'\}_{i \in I}$ および
$\mathcal{V} = \{V_j \to S\}_{j \in J}$ をとる。
$\alpha : I \to J$, $h : S' \to S$ および
$g_i : U_i \to V_{\alpha(i)}$ を、標的を固定した写像の族の射とする。
Sites, Definition \ref{sites-definition-morphism-coverings} を参照せよ。
\begin{enumerate}
\item $(Y_j, \varphi_{jj'})$ を族 $\{V_j \to S'\}$ に関する
降下データとする。系
$$
\left(
g_i^*Y_{\alpha(i)},
(g_i \times g_{i'})^*\varphi_{\alpha(i)\alpha(i')}
\right)
$$
は（Remark \ref{descent-remark-easier-family} の記法により）
$\mathcal{V}$ に関する降下データである。
\item この構成は、$\mathcal{U}$ に関する降下データと
$\mathcal{V}$ に関する降下データの間の関手を定める。
\item 標的を固定した写像の族の第二の射
$\alpha' : I \to J$, $h' : S' \to S$ および
$g'_i : U_i \to V_{\alpha'(i)}$ が与えられたとき、$h = h'$ ならば、
得られる二つの降下データ間の関手は標準的に同型である。
\item これらの関手は Lemma \ref{descent-lemma-family-is-one} を介して、
Lemma \ref{descent-lemma-pullback} で構成した引き戻し関手と一致する。
\end{enumerate}
\end{lemma}

\begin{proof}
これは Lemma \ref{descent-lemma-family-is-one} の対応を介して
Lemma \ref{descent-lemma-pullback} から従う。
\end{proof}

\begin{definition}
\label{descent-definition-pullback-functor-family}
Lemma \ref{descent-lemma-pullback-family} の
$\mathcal{U} = \{U_i \to S'\}_{i \in I}$,
$\mathcal{V} = \{V_j \to S\}_{j \in J}$, $\alpha : I \to J$, $h : S' \to S$,
および $g_i : U_i \to V_{\alpha(i)}$ に対し、同補題で構成した
降下データ上の関手
$$
(Y_j, \varphi_{jj'}) \longmapsto
(g_i^*Y_{\alpha(i)}, (g_i \times g_{i'})^*\varphi_{\alpha(i)\alpha(i')})
$$
を{\it 引き戻し関手}という。
\end{definition}

\noindent
もし $\mathcal{U}$ と $\mathcal{V}$ が同じ標的 $S$ をもち、
$\mathcal{U}$ が $\mathcal{V}$ の細分であるならば
（Sites, Definition \ref{sites-definition-morphism-coverings}
を参照）、組 $(\alpha, g_i)$ が明示的に与えられていなくても、
引き戻し関手について論じられる。実際、Lemma
\ref{descent-lemma-pullback-family} により、この組の選択は
（標準的同型を除いて）影響しない。


\begin{definition}
\label{descent-definition-effective}
スキーム $S$ とスキームの射 $f : X \to S$ をとる。
\begin{enumerate}
\item $S$ 上のスキーム $U$ が与えられると、
$\text{id} : S \to S$ に関する $U$ の{\it 自明な降下データ}、
すなわち $U$ 上の恒等射が得られる。
\item Lemma \ref{descent-lemma-pullback} により、この自明な降下データを
$f$ で引き戻すことで、$X \to S$ に関する $X \times_S U$ 上の
{\it 標準降下データ}が得られる。この降下データをしばしば
$(X \times_S U, can)$ と表す。
\item $X/S$ に関する降下データ $(V, \varphi)$ が{\it 有効}であるとは、
ある $S$ 上のスキーム $U$ に対し、$(V, \varphi)$ が標準降下データ
$(X \times_S U, can)$ と同型であることをいう。
\end{enumerate}
\end{definition}

\noindent
したがって、有効であるとは、$S$ 上のスキーム $U$ と
$X$-スキームの同型 $\psi : V \to X \times_S U$ が存在し、
$\varphi$ が合成
$$
V \times_S X \xrightarrow{\psi \times \text{id}_X}
X \times_S U \times_S X =
X \times_S X \times_S U
\xrightarrow{\text{id}_X \times \psi^{-1}}
X \times_S V
$$

\begin{definition}
\label{descent-definition-effective-family}
スキーム $S$ をとる。
$\{X_i \to S\}$ を、標的 $S$ をもつ射の族とする。
\begin{enumerate}
\item $S$ 上のスキーム $U$ が与えられると、
$\{\text{id} : S \to S\}$ に関する $U$ の自明な降下データを
引き戻すことにより、スキームの族 $X_i \times_S U$ 上の
{\it 標準降下データ}が得られる。この降下データを
$(X_i \times_S U, can)$ と表す。
\item $\{X_i \to S\}$ に関する降下データ
$(V_i, \varphi_{ij})$ が{\it 有効}であるとは、
ある $S$ 上のスキーム $U$ が存在して
$(V_i, \varphi_{ij})$ が $(X_i \times_S U, can)$ と同型であることをいう。
\end{enumerate}
\end{definition}












\section{引き戻し関手の充満忠実性}
\label{descent-section-fully-faithful}

\noindent
fpqc 被覆に対する降下データ間の引き戻し関手は充満忠実である。
言い換えれば、スキームの射は fpqc 降下を満たす。
本節の目的はこれを証明することであるが、読者自身で証明することを勧める。
鍵となるのは Lemma \ref{descent-lemma-fpqc-universal-effective-epimorphisms}
を用いることである。

\begin{lemma}
\label{descent-lemma-surjective-flat-epi}
全射かつ平坦な射は、スキームの圏におけるエピ射である。
\end{lemma}

\begin{proof}
$h : X' \to X$ を全射かつ平坦とし、
$a \circ h = b \circ h$ を満たす射 $a, b : X \to Y$ をとる。
$h$ は全射なので、$a$ と $b$ は台位相空間上で一致する。
$x' \in X'$ をとり、$x = h(x')$ および $y = a(x) = b(x)$ とおく。
局所環の写像
$$
a^\sharp_x, b^\sharp_x : \mathcal{O}_{Y, y} \to \mathcal{O}_{X, x}
$$
を考える。これらは、平坦な局所準同型
$h^\sharp_{x'} : \mathcal{O}_{X, x} \to \mathcal{O}_{X', x'}$
と合成すると等しくなる。平坦な局所準同型は忠実平坦なので
(Algebra, Lemma \ref{algebra-lemma-local-flat-ff})
$h^\sharp_{x'}$ は単射である。したがって
$a^\sharp_x = b^\sharp_x$ であり、所望の $a = b$ が従う。
\end{proof}

\begin{lemma}
\label{descent-lemma-ff-base-change-faithful}
$h : S' \to S$ をスキームの全射かつ平坦な射とする。
基底変換関手
$$
\Sch/S \longrightarrow \Sch/S', \quad
X \longmapsto S' \times_S X
$$
は忠実である。
\end{lemma}

\begin{proof}
$X_1$, $X_2$ を $S$ 上のスキームとし、
$\alpha, \beta : X_2 \to X_1$ を $S$ 上の射とする。
$\alpha$, $\beta$ の基底変換が同じ射になるならば、次の可換図式を得る。
$$
\xymatrix{
X_2 \ar[d]^\alpha &
S' \times_S X_2 \ar[l] \ar[d] \ar[r] &
X_2 \ar[d]^\beta \\
X_1 &
S' \times_S X_1 \ar[l] \ar[r] &
X_1
}
$$
したがって、$S' \times_S X_2 \to X_2$ がエピ射であることを
示せば十分である。これは全射かつ平坦な射の基底変換なので、
全射かつ平坦である（
Morphisms, Lemmas \ref{morphisms-lemma-base-change-surjective}
および \ref{morphisms-lemma-base-change-flat} を参照）。
ゆえに Lemma \ref{descent-lemma-surjective-flat-epi} から従う。
\end{proof}

\begin{lemma}
\label{descent-lemma-faithful}
Lemma \ref{descent-lemma-pullback} の状況で、
$f : X' \to X$ が全射かつ平坦であると仮定する。
このとき引き戻し関手は忠実である。
\end{lemma}

\begin{proof}
$(V_i, \varphi_i)$, $i = 1, 2$ を $X \to S$ に対する降下データとし、
$\alpha, \beta : V_1 \to V_2$ を降下データの射とする。
$f^*\alpha = f^*\beta$ と仮定し、$\alpha = \beta$ を示す。
$\alpha$, $\beta$ は $X$ 上のスキームの射であり、
$f^*\alpha$, $f^*\beta$ は単に $\alpha$, $\beta$ を
$X'$ 上の射へ基底変換したものだから、
Lemma \ref{descent-lemma-ff-base-change-faithful} から従う。
\end{proof}

\noindent
次が本節の鍵となる補題である。

\begin{lemma}
\label{descent-lemma-fully-faithful}
Lemma \ref{descent-lemma-pullback} の状況で、次を仮定する。
\begin{enumerate}
\item $\{f : X' \to X\}$ は fpqc 被覆である
（たとえば $f$ が全射、平坦かつ準コンパクトである場合）。
\item $S = S'$.
\end{enumerate}
このとき引き戻し関手は充満忠実である。
\end{lemma}

\begin{proof}
仮定 (1) より $f$ は全射かつ平坦である。したがって
Lemma \ref{descent-lemma-faithful} により引き戻し関手は忠実である。
$(V, \varphi)$ と $(W, \psi)$ を $X \to S$ に関する二つの
降下データとする。$(V', \varphi') = f^*(V, \varphi)$ および
$(W', \psi') = f^*(W, \psi)$ とおく。
$\alpha' : V' \to W'$ を $S$ 上の $X'$ に対する降下データの射とする。
その引き戻しが $\alpha'$ となる、$S$ 上の $X$ に対する降下データの射
$\alpha : V \to W$ が存在することを示さなければならない。

\medskip\noindent
$V'$ は $f$ による $V$ の基底変換であり、$\varphi'$ は
$f \times f$ による $\varphi$ の基底変換であることを思い出そう
（Remark \ref{descent-remark-easier} を参照）。仮定により図式
$$
\xymatrix{
V' \times_S X' \ar[r]_{\varphi'} \ar[d]_{\alpha' \times \text{id}} &
X' \times_S V' \ar[d]^{\text{id} \times \alpha'} \\
W' \times_S X' \ar[r]^{\psi'} &
X' \times_S W'
}
$$
は可換である。二つの合成
$$
\xymatrix{
V' \times_V V' \ar[r]^-{\text{pr}_i} &
V' \ar[r]^{\alpha'} &
W' \ar[r] &
W
}
, \quad i = 0, 1
$$
が等しいと主張する。以下を読む前に、読者自身でこの主張を証明することを
勧める（簡潔でよい議論を見つけたらメールで知らせてほしい）。
$v_0, v_1$ を同じ点 $v \in V$ に写る $V'$ の点とする。
$x_i \in X'$ を $v_i$ の像とし、$x$ を $X$ の点、
すなわち $v$ の $X$ における像とする。
言い換えれば、$V' = X' \times_X V$ において
$v_i = (x_i, v)$ である。ある $V$ の点 $v'$ に対し
$\varphi(v, x) = (x, v')$ と書く。$\varphi$ は
$X \times_S X$ 上の射なので、これは可能である。
$V'$ の点 $v_i' = (x_i, v')$ とおく。
すると、$\varphi'$ の定義を用いた計算により
$\varphi'(v_i, x_j) = (x_i, v'_j)$ が分かる。
$w_i = \alpha'(v_i)$ および $w'_i = \alpha'(v_i')$ とおく。
このとき、ある $W$ の点 $u_i$ と、ある $W$ の点 $u_i'$ に対して
$w_i = (x_i, u_i)$ および $w_i' = (x_i, u'_i)$ と書ける。
主張は $u_0 = u_1$ と同値である。
$\psi'$ の定義を用いた形式的計算
（Lemma \ref{descent-lemma-pullback} を参照）によれば、
上の図式の可換性は、すべての $i, j \in \{0, 1\}$ に対して
$$
((x_i, x_j), \psi(u_i, x)) = ((x_i, x_j), (x, u'_j))
$$
が
$(X' \times_S X') \times_{X \times_S X} (X \times_S W)$
の点として成り立つことをいう。これより
$\psi(u_0, x) = \psi(u_1, x)$、したがって $\psi^{-1}$ をとって
$u_0 = u_1$ である。上の議論は形式的であり、スキーム点をとれる
（言い換えれば $(v_0, v_1) = \text{id}_{V' \times_V V'}$
ととれる）ので、主張が証明された。

\medskip\noindent
ここで Lemma \ref{descent-lemma-fpqc-universal-effective-epimorphisms}
を用いることができる。実際、$\{V' \to V\}$ は
射 $f : X' \to X$ の基底変換なので fpqc 被覆である。
したがって Lemma \ref{descent-lemma-fpqc-universal-effective-epimorphisms}
により、射 $\alpha' : V' \to W' \to W$ は、
基底変換が必然的に $\alpha'$ となる一意な射
$\alpha : V \to W$ を経由する。最後に、図式
$$
\xymatrix{
V \times_S X \ar[r]_{\varphi} \ar[d]_{\alpha \times \text{id}} &
X \times_S V \ar[d]^{\text{id} \times \alpha} \\
W \times_S X \ar[r]^{\psi} & X \times_S W
}
$$
は可換である。その $X' \times_S X'$ への基底変換が可換であり、
射 $X' \times_S X' \to X \times_S X$ が全射かつ平坦だからである
（Lemma \ref{descent-lemma-ff-base-change-faithful} を用いる）。
したがって $\alpha$ は、所望の降下データの射
$(V, \varphi) \to (W, \psi)$ である。
\end{proof}

\noindent
以下の二つの補題は、先行する内容の記述が改善されたため不要になった。
しかし、なお正しい。

\begin{lemma}
\label{descent-lemma-pullback-selfmap}
$X \to S$ をスキームの射とし、$f : X \to X$ を
$S$ 上の $X$ の自己写像とする。このとき、$f$ による引き戻しは、
$X \to S$ に関する降下データの圏上の恒等関手と同型である。
\end{lemma}

\begin{proof}
Lemma \ref{descent-lemma-pullback} が
$f^* \cong \text{id}^*$ を与えるので明らかである。
\end{proof}

\begin{lemma}
\label{descent-lemma-morphism-with-section-equivalence}
$f : X' \to X$ を基底スキーム $S$ 上のスキームの射とする。
たとえば $f$ が切断をもつ場合のように、$S$ 上の射
$g : X \to X'$ が存在すると仮定する。このとき
Lemma \ref{descent-lemma-pullback} の引き戻し関手は、
$X/S$ に関する降下データの圏と $X'/S$ に関する降下データの圏との
圏同値を定める。
\end{lemma}

\begin{proof}
射 $g : X \to X'$ を $S$ 上でとる。
上の Lemma \ref{descent-lemma-pullback-selfmap} により、関手
$f^* \circ g^* = (g \circ f)^*$ および
$g^* \circ f^* = (f \circ g)^*$
はそれぞれ対応する恒等関手と同型であり、主張が従う。
\end{proof}

\begin{lemma}
\label{descent-lemma-morphism-source-faithfully-flat}
$f : X \to X'$ を基底スキーム $S$ 上のスキームの射とする。
$X \to S$ が全射かつ平坦であると仮定する。このとき
Lemma \ref{descent-lemma-pullback} の引き戻し関手は、
$X'/S$ に関する降下データの圏から $X/S$ に関する
降下データの圏への忠実関手である。
\end{lemma}

\begin{proof}
$X \to X'$ を $X \to X \times_S X' \to X'$ と分解できる。
最初の射は切断をもつので、
Lemma \ref{descent-lemma-morphism-with-section-equivalence} により
降下データの圏の圏同値を誘導する。第二の射は全射かつ平坦なので、
Lemma \ref{descent-lemma-faithful} により忠実関手を誘導する。
\end{proof}

\begin{lemma}
\label{descent-lemma-morphism-source-fpqc-covering}
$f : X \to X'$ を基底スキーム $S$ 上のスキームの射とする。
$\{X \to S\}$ が fpqc 被覆であると仮定する
（たとえば $f$ が全射、平坦かつ準コンパクトである場合）。
このとき Lemma \ref{descent-lemma-pullback} の引き戻し関手は、
$X'/S$ に関する降下データの圏から $X/S$ に関する降下データの圏への
充満忠実関手である。
\end{lemma}

\begin{proof}
$X \to X'$ を $X \to X \times_S X' \to X'$ と分解できる。
最初の射は切断をもつので、
Lemma \ref{descent-lemma-morphism-with-section-equivalence} により
降下データの圏の圏同値を誘導する。第二の射は fpqc 被覆なので、
Lemma \ref{descent-lemma-fully-faithful} により充満忠実関手を誘導する。
\end{proof}

\begin{lemma}
\label{descent-lemma-fpqc-refinement-coverings-fully-faithful}
スキーム $S$ をとる。
$\mathcal{U} = \{U_i \to S\}_{i \in I}$ および
$\mathcal{V} = \{V_j \to S\}_{j \in J}$,
を標的 $S$ をもつ射の族とする。
$\alpha : I \to J$, $\text{id} : S \to S$ および
$g_i : U_i \to V_{\alpha(i)}$ を、標的を固定した写像の族の射とする。
Sites, Definition \ref{sites-definition-morphism-coverings} を参照せよ。
各 $j \in J$ に対し、族
$\{g_i : U_i \to V_j\}_{\alpha(i) = j}$ が $V_j$ の fpqc 被覆であると
仮定する。このとき Lemma \ref{descent-lemma-pullback-family} の引き戻し関手
$$
\text{降下データの基準族：}
\mathcal{V}
\longrightarrow
\text{降下データの基準族：}
\mathcal{U}
$$
は充満忠実である。
\end{lemma}

\begin{proof}
スキームの射
$$
g :
X = \coprod\nolimits_{i \in I} U_i
\longrightarrow
Y = \coprod\nolimits_{j \in J} V_j
$$
を $S$ 上で考える。これは第 $i$ 成分を射 $g_{\alpha(i)}$ により
第 $\alpha(i)$ 成分へ写す。$\{g : X \to Y\}$ はスキームの
fpqc 被覆であると主張する。実際、
Topologies, Lemma \ref{topologies-lemma-disjoint-union-is-fpqc-covering}
により、各 $j$ について射
$\{\coprod_{\alpha(i) = j} U_i \to V_j\}$ は fpqc 被覆である。
したがって、任意のアフィン開集合 $V \subset V_j$
（これを $Y$ のアフィン開集合とみなせる）に対し、有限個のアフィン開集合
$W_1, \ldots, W_n \subset \coprod_{\alpha(i) = j} U_i$
（これらを $X$ のアフィン開集合とみなせる）で
$V = \bigcup_{i = 1, \ldots, n} g(W_i)$ を満たすものが存在する。
これにより、$X$ の有限個のアフィン開集合で被覆できる $Y$ の
アフィン開集合が十分に得られるので、
Topologies, Lemma \ref{topologies-lemma-recognize-fpqc-covering} (3)
を適用でき、主張が従う。$X/S$ に関する降下データの圏を
$DD(X/S)$、$\mathcal{U}$ に関するものを $DD(\mathcal{U})$ と書き、
$Y/S$ と $\mathcal{V}$ についても同様に書く。図式
$$
\xymatrix{
DD(Y/S) \ar[r] & DD(X/S) \\
DD(\mathcal{V}) \ar[u]^{\text{補題 }\ref{descent-lemma-family-is-one}} \ar[r] &
DD(\mathcal{U}) \ar[u]_{\text{補題 }\ref{descent-lemma-family-is-one}}
}
$$
を考える。この図式は可換である
（Lemma \ref{descent-lemma-pullback-family} の証明を参照）。
縦の矢印は圏同値である。したがって、図式の上側の横矢印が
充満忠実であることを示す Lemma \ref{descent-lemma-fully-faithful} から従う。
\end{proof}

\noindent
次の補題は、有効性を確認する際、標的 $S$ をもつ所与の射の族を
常に Zariski 細分してよいことを示す。

\begin{lemma}
\label{descent-lemma-Zariski-refinement-coverings-equivalence}
スキーム $S$ をとる。
$\mathcal{U} = \{U_i \to S\}_{i \in I}$ および
$\mathcal{V} = \{V_j \to S\}_{j \in J}$,
を標的 $S$ をもつ射の族とする。
$\alpha : I \to J$, $\text{id} : S \to S$ および
$g_i : U_i \to V_{\alpha(i)}$ を、標的を固定した写像の族の射とする。
Sites, Definition \ref{sites-definition-morphism-coverings} を参照せよ。
各 $j \in J$ に対し、族
$\{g_i : U_i \to V_j\}_{\alpha(i) = j}$ が $V_j$ の Zariski 被覆であると
仮定する（
Topologies, Definition \ref{topologies-definition-zariski-covering}
を参照）。このとき Lemma \ref{descent-lemma-pullback-family} の引き戻し関手
$$
\text{降下データの基準族：}
\mathcal{V}
\longrightarrow
\text{降下データの基準族：}
\mathcal{U}
$$
は圏同値である。特に、$S$ 上のスキームの圏は、
$S$ の任意の Zariski 被覆に関する降下データの圏と同値である。
\end{lemma}

\begin{proof}
Lemma \ref{descent-lemma-fpqc-refinement-coverings-fully-faithful} により、
この関手は忠実かつ充満忠実である。本質的全射性の証明を説明する。
$(X_i, \varphi_{ii'})$ を $\mathcal{U}$ に関する降下データとする。
$j \in J$ を固定し、$I_j = \{i \in I \mid \alpha(i) = j\}$ とおく。
$i, i' \in I_j$ に対し、標準射
$$
c_{ii'} : U_i \times_{g_i, V_j, g_{i'}} U_{i'} \to U_i \times_S U_{i'}.
$$
が存在することに注意する。したがって、この射で
$\varphi_{ii'}$ を引き戻し、$i, i' \in I_j$ に対して
$\psi_{ii'} = c_{ii'}^*\varphi_{ii'}$ とおける。
これにより、Zariski 被覆
$\{g_i : U_i \to V_j\}_{i \in I_j}$ に関する降下データ
$(X_i, \psi_{ii'})$ を得る。$\psi_{ii'}$ は、$X_i$ の開集合
$X_{i, U_i \times_{V_j} U_{i'}}$ から $X_{i'}$ の対応する
開集合への同型であることに注意する。
Schemes, Section \ref{schemes-section-glueing-schemes}
により、$(X_i, \psi_{ii'})$ を貼り合わせて $V_j$ 上の
スキーム $Y_j$ を得られる。さらに、$i \in I_j$ および
$i' \in I_{j'}$ に対する射 $\varphi_{ii'}$ は貼り合わされて、
コサイクル条件を満たす射
$\varphi_{jj'} : Y_j \times_S V_{j'} \to V_j \times_S Y_{j'}$
となる（詳細は省略する）。したがって、$\mathcal{V}$ に関する
所望の降下データ $(Y_j, \varphi_{jj'})$ を得る。
\end{proof}

\begin{lemma}
\label{descent-lemma-refine-coverings-fully-faithful}
スキーム $S$ をとる。
$\mathcal{U} = \{U_i \to S\}_{i \in I}$ および
$\mathcal{V} = \{V_j \to S\}_{j \in J}$,
を $S$ の fpqc 被覆とする。$\mathcal{U}$ が
$\mathcal{V}$ の細分であるならば、引き戻し関手
$$
\text{降下データの基準族：}
\mathcal{V}
\longrightarrow
\text{降下データの基準族：}
\mathcal{U}
$$
は充満忠実である。特に、$S$ 上のスキームの圏は、
$S$ の任意の fpqc 被覆に関する降下データの圏の充満部分圏と同一視される。
\end{lemma}

\begin{proof}
$S$ の fpqc 被覆
$\mathcal{W} = \{U_i \times_S V_j \to S\}_{(i, j) \in I \times J}$
を考える。これは $\mathcal{U}$ と $\mathcal{V}$ の両方の細分である。
したがって、関手と圏の $2$-可換図式
$$
\xymatrix{
DD(\mathcal{V}) \ar[rd] \ar[rr] & & DD(\mathcal{U}) \ar[ld] \\
& DD(\mathcal{W}) &
}
$$
を得る。記法は
Lemma \ref{descent-lemma-fpqc-refinement-coverings-fully-faithful} の証明と同じであり、
可換性は Lemma \ref{descent-lemma-pullback-family} (3) による。
したがって、関手
$DD(\mathcal{V}) \to DD(\mathcal{W})$ および
$DD(\mathcal{U}) \to DD(\mathcal{W})$ が充満忠実であることを
示せば十分なのは明らかである。これは
Lemma \ref{descent-lemma-fpqc-refinement-coverings-fully-faithful}
から従い、証明が完了する。
\end{proof}

\begin{remark}
\label{descent-remark-morphisms-of-schemes-satisfy-fpqc-descent}
Lemma \ref{descent-lemma-refine-coverings-fully-faithful} は、
スキームの射が fpqc 降下を満たすと述べている。
言い換えれば、スキーム $S$ と $S$ 上のスキーム $X$, $Y$ が
与えられたとき、関手
$$
(\Sch/S)^{opp} \longrightarrow \textit{Sets},
\quad
T \longmapsto \Mor_T(X_T, Y_T)
$$
は fpqc 位相に対する層条件を満たす。
最も簡単な場合は次のとおりである。$T \to S$ をアフィン間の
全射平坦な射とし、$\psi_0 : X_T \to Y_T$ を標準降下データと
両立する $T$ 上のスキームの射とする。このとき、$T$ への基底変換が
$\psi_0$ となる一意な射 $\psi : X \to Y$ が存在する。
実際、この特殊な場合は Lemma \ref{descent-lemma-fully-faithful} から
直接従う。そして同補題は、次の二つの事実の形式的帰結である。
(a) 忠実平坦な射による基底変換関手は忠実である
（Lemma \ref{descent-lemma-ff-base-change-faithful} を参照）。
(b) スキームは fpqc 位相に対する層条件を満たす
（Lemma \ref{descent-lemma-fpqc-universal-effective-epimorphisms} を参照）。
\end{remark}

\begin{lemma}
\label{descent-lemma-effective-for-fpqc-is-local-upstairs}
$X \to S$ をスキームの全射、準コンパクトかつ平坦な射とし、
$(V, \varphi)$ を $X/S$ に関する降下データとする。
任意の $v \in V$ に対し、
$\varphi(W \times_S X) \subset X \times_S W$ を満たし、かつ
降下データ $(W, \varphi|_{W \times_S X})$ が有効となる開部分スキーム
$v \in W \subset V$ が存在すると仮定する。このとき
$(V, \varphi)$ は有効である。
\end{lemma}

\begin{proof}
開被覆 $V = \bigcup W_i$ が
$\varphi(W_i \times_S X) \subset X \times_S W_i$ を満たし、
降下データ $(W_i, \varphi|_{W_i \times_S X})$ が有効であるとする。
$U_i \to S$ をスキームとし、
$\alpha_i : (X \times_S U_i, can) \to (W_i, \varphi|_{W_i \times_S X})$
を降下データの同型とする。各添字の組 $(i, j)$ に対し、開集合
$\alpha_i^{-1}(W_i \cap W_j) \subset X \times_S U_i$ を考える。
すべてが降下データと両立し、$\{X \to S\}$ は fpqc 被覆なので、
Lemma \ref{descent-lemma-open-fpqc-covering} を適用して、
$\alpha_i^{-1}(W_i \cap W_j) = X \times_S U_{ij}$ を満たす
開集合 $U_{ij} \subset U_i$ を得る。いま $W_i \cap W_j$ 上の恒等射は
降下データと両立するので、$S$ 上の一意な射
$\varphi_{ij} : U_{ij} \to U_{ji}$ から生じる
（Remark \ref{descent-remark-morphisms-of-schemes-satisfy-fpqc-descent} を参照）。
このとき $(U_i, U_{ij}, \varphi_{ij})$ は
Schemes, Section \ref{schemes-section-glueing-schemes} の貼り合わせデータ
である（証明は省略する）。したがって、$S$ 上のスキーム $U$ が存在し、
$U_i \subset U$ が開、$U_{ij} = U_i \cap U_j$、かつ
$\varphi_{ij} = \text{id}_{U_i \cap U_j}$ であると仮定してよい。
以下を参照せよ。
Schemes, Lemma \ref{schemes-lemma-glue}。
$X$ へ引き戻すと、$\alpha_i$ を用いて所望の同型
$\alpha : X \times_S U \to V$ を得る。
\end{proof}







\section{射の型の降下}
\label{descent-section-descending-types-morphisms}

\noindent
以下では、fpqc 被覆に関するスキームの降下データが有効であるかを調べる。
まず注意すべきことは、常に有効とは限らないことである。これは
Algebraic Spaces, Example
\ref{spaces-example-non-representable-descent} で見る。
Examples, Lemma
\ref{examples-lemma-non-effective-descent-projective} によれば、
射影射でさえ fpqc 被覆に対する降下を常に満たすわけではない。

\medskip\noindent
一方、降下させようとするスキームが特に単純ならば、スキームのあるクラス
全体について降下データが有効となる場合がある。ここでは、この現象を
抽象的に記述する用語を導入する。ただし、後で正しく用いなければ
混乱を招きうる。

\begin{definition}
\label{descent-definition-descending-types-morphisms}
$\mathcal{P}$ を基底上のスキームの射の性質とする。\newline
$\tau \in \{Zariski, fpqc, fppf, \etale, smooth, syntomic\}$ とする。
任意の $\tau$-被覆
$\mathcal{U} : \{U_i \to S\}_{i \in I}$
（Topologies, Section \ref{topologies-section-procedure} を参照）と、
各射 $X_i \to U_i$ が性質 $\mathcal{P}$ をもつ
$\mathcal{U}$ に関する任意の降下データ $(X_i, \varphi_{ij})$ が
有効であるとき、{\it 型 $\mathcal{P}$ の射は
$\tau$-被覆に対する降下を満たす}という。
\end{definition}

\noindent
いずれの場合にも、$S$ 上のスキームから $\mathcal{U}$ 上の
降下データへの関手が充満忠実であることはすでに見た
（Lemma \ref{descent-lemma-refine-coverings-fully-faithful} と、
任意の $\tau$-被覆が fpqc 被覆でもあるという Topologies の結果を組み合わせる）。
また、Zariski 被覆に関して降下データは常に有効であることも見た
（Lemma \ref{descent-lemma-Zariski-refinement-coverings-equivalence}）。
誤った議論、すなわち降下の途中で $\mathcal{P}$ に依存する議論を
避けるには、$\mathcal{P}$ が任意の基底変換で安定であるか、
少なくとも $\tau$-位相で基底上局所的である場合に限って
上の概念を検討するのが賢明であろう
（Definition \ref{descent-definition-property-morphisms-local} を参照）。

\medskip\noindent
次は、この問題をアフィン間の一つの射で与えられる被覆の場合に
帰着するための標準的な補題である。

\begin{lemma}
\label{descent-lemma-descending-types-morphisms}
$\mathcal{P}$ を基底上のスキームの射の性質とし、
$\tau \in \{fpqc, fppf, \etale, smooth, syntomic\}$ とする。
次を仮定する。
\begin{enumerate}
\item $\mathcal{P}$ は任意の基底変換で安定である
（Schemes, Definition \ref{schemes-definition-preserved-by-base-change} を参照）。
\item $Y_j \to V_j$, $j = 1, \ldots, m$ が $\mathcal{P}$ をもつならば、
$\coprod Y_j \to \coprod V_j$ も同じ性質をもつ。
\item アフィン間の任意の全射 $X \to S$ であって、
$\tau$ が fpqc, fppf, \'etale, smooth, syntomic のいずれであるかに応じて、
平坦、有限表示平坦、\'etale、smooth、syntomic であるものに対し、
$V \to X$ が $\mathcal{P}$ をもつような $S$ 上の $X$ に関する
任意の降下データ $(V, \varphi)$ は有効である。
\end{enumerate}
このとき、型 $\mathcal{P}$ の射は $\tau$-被覆に対する降下を満たす。
\end{lemma}

\begin{proof}
スキーム $S$ をとる。
$\mathcal{U} = \{\varphi_i : U_i \to S\}_{i \in I}$ を
$S$ の $\tau$-被覆とする。$(X_i, \varphi_{ii'})$ を
$\mathcal{U}$ に関する降下データとし、各射 $X_i \to U_i$ が
性質 $\mathcal{P}$ をもつと仮定する。
$(X_i, \varphi_{ii'}) \cong (U_i \times_S X, can)$ を満たす
スキーム $X \to S$ が存在することを示さなければならない。

\medskip\noindent
本来の証明を始める前に次を注意しておく。
任意の射の族 $\mathcal{V} : \{V_j \to S\}$ と任意の族の射
$\mathcal{V} \to \mathcal{U}$ に対し、降下データ
$(X_i, \varphi_{ii'})$ を $\mathcal{V}$ 上の降下データ
$(Y_j, \varphi_{jj'})$ へ引き戻すと、各射 $Y_j \to V_j$ も
性質 $\mathcal{P}$ をもつ。これは $\mathcal{P}$ が任意の基底変換で
安定であるという仮定 (1) と引き戻しの定義による
（Definition \ref{descent-definition-pullback-functor-family} を参照）。
以下では断りなくこれを用いる。

\medskip\noindent
まず $S$ がアフィンの場合に補題を証明する。
Topologies, Lemma
\ref{topologies-lemma-fpqc-affine},
\ref{topologies-lemma-fppf-affine},
\ref{topologies-lemma-etale-affine},
\ref{topologies-lemma-smooth-affine}, または
\ref{topologies-lemma-syntomic-affine}
のいずれかにより、$\mathcal{U}$ を細分する標準 $\tau$-被覆
$\mathcal{V} : \{V_j \to S\}_{j = 1, \ldots, m}$ が存在する。
降下データの圏の間の引き戻し関手
$DD(\mathcal{U}) \to DD(\mathcal{V})$
は Lemma \ref{descent-lemma-refine-coverings-fully-faithful} により充満忠実である。
したがって、標準 $\tau$-被覆 $\mathcal{V}$ 上の降下データが
有効であることを示せば十分である。仮定 (2) により
$\coprod Y_j \to \coprod V_j$ は性質 $\mathcal{P}$ をもつ。
Lemma \ref{descent-lemma-family-is-one} により、問題は被覆
$\{\coprod_{j = 1, \ldots, m} V_j \to S\}$ に帰着されるが、
これは補題の仮定 (3) で結果を仮定した場合である。
したがって $S$ がアフィンならば補題は成り立つ。

\medskip\noindent
一般の $S$ を考える。$V \subset S$ をアフィン開集合とする。
サイトの性質により、族
$\mathcal{U}_V = \{V \times_S U_i \to V\}_{i \in I}$ は
$V$ の $\tau$-被覆である。所与の降下データを
$\mathcal{U}_V$ へ制限（または引き戻し）したものを
$(X_i, \varphi_{ii'})_V$ と表す。直前の結果により、
$\mathcal{U}_V$ に関する標準降下データが
$(X_i, \varphi_{ii'})_V$ と同型となる $V$ 上のスキーム $X_V$ を得る。
$V' \subset V$ を $V$ のアフィン開集合とする。このとき
$X_{V'}$ と $V' \times_V X_V$ の標準降下データはともに
$(X_i, \varphi_{ii'})_{V'}$ と同型である。したがって再び
Lemma \ref{descent-lemma-refine-coverings-fully-faithful} により、
$X_{V'}$ を $X_V$ における $V'$ の逆像と同一視する
$S$ 上の標準射 $\rho^V_{V'} : X_{V'} \to X_V$ を得る。
$S$ のアフィン開集合 $V'' \subset V' \subset V$ に対して
$\rho^V_{V''} = \rho^V_{V'} \circ \rho^{V'}_{V''}$ となることの確認は省略する。

\medskip\noindent
Constructions, Lemma \ref{constructions-lemma-relative-glueing} により、
データ $(X_V, \rho^V_{V'})$ はスキーム $X \to S$ に貼り合わされる。
さらに、射 $\rho^V_{V'}$ を復元する同型 $V \times_S X \to X_V$
が与えられる。スキーム $X_V$ の構成を展開すると、同型
$$
V \times_S U_i \times_S X
\longrightarrow
V \times_S X_i
$$
を得る。これは射 $\varphi_{ii'}$ と両立し、$X$ のより小さい
アフィン開集合への制限とも両立する。したがって
$U_i \times_S X$ 上の標準降下データは所与の降下データと同型であり、
証明が完了する。
\end{proof}










\section{アフィン射の降下}
\label{descent-section-affine}

\noindent
本節では「アフィン射は fpqc 被覆に対する降下を満たす」ことを示す。
厳密な主張は次のとおりである。

\begin{lemma}
\label{descent-lemma-affine}
スキーム $S$ をとる。
$\{X_i \to S\}_{i\in I}$ を fpqc 被覆とする
（Topologies, Definition \ref{topologies-definition-fpqc-covering} を参照せよ）。
$(V_i/X_i, \varphi_{ij})$ を $\{X_i \to S\}$ に関する
降下データとする。各射 $V_i \to X_i$ がアフィンならば、
この降下データは有効である。
\end{lemma}

\begin{proof}
アフィンであることは、基底上局所的で任意の基底変換により保たれる
スキームの射の性質である。Morphisms, Lemmas
\ref{morphisms-lemma-characterize-affine} および
\ref{morphisms-lemma-base-change-affine} を参照せよ。
したがって Lemma \ref{descent-lemma-descending-types-morphisms} を適用でき、
fpqc 被覆がアフィン間の一つの平坦全射
$\{X \to S\}$ で与えられる場合に補題の主張を証明すれば十分である。
$X = \Spec(A)$ および $S = \Spec(R)$ と書くと、
$R \to A$ は忠実平坦な環準同型である。
$(V, \varphi)$ を $S$ 上の $X$ に関する降下データとし、
$V \to X$ がアフィンであると仮定する。
このとき、$V \to X$ がアフィンであることから、
ある $A$-代数 $B$ に対して $V = \Spec(B)$ である
（Morphisms, Definition \ref{morphisms-definition-affine} を参照）。
同型 $\varphi$ は、環の同型
$$
\varphi^\sharp :
B \otimes_R A \longleftarrow A \otimes_R B
$$
に対応し、これは $A \otimes_R A$-代数の同型である。
$\varphi$ のコサイクル条件は、図式
$$
\xymatrix{
B \otimes_R A \otimes_R A & &
A \otimes_R A \otimes_R B \ar[ll] \ar[ld]\\
& A \otimes_R B \otimes_R A \ar[lu] &
}
$$
が可換であると述べている。これらの矢印を逆にすると、
Definition \ref{descent-definition-descent-datum-modules} の意味で
$R \to A$ に関する加群の降下データが得られる。
したがって Proposition \ref{descent-proposition-descent-module} を適用して、
$R$-加群
$C = \Ker(B \to A \otimes_R B)$
と、降下データを尊重する同型 $A \otimes_R C \cong B$ を得る。
任意の $c, c' \in C$ に対し、$\varphi$ は代数準同型なので
$B$ における積 $cc'$ は $C$ に属する。したがって $C$ は
$R$-代数であり、その $A$ への基底変換は降下データと両立して
$B$ と同型である。$\Spec$ を適用すると、
$(V, \varphi) \cong (X \times_S U, can)$ を満たす
$S$ 上のスキーム $U$ を得て、主張が従う。
\end{proof}

\begin{lemma}
\label{descent-lemma-closed-immersion}
スキーム $S$ をとる。
$\{X_i \to S\}_{i\in I}$ を fpqc 被覆とする
（Topologies, Definition \ref{topologies-definition-fpqc-covering} を参照せよ）。
$(V_i/X_i, \varphi_{ij})$ を $\{X_i \to S\}$ に関する
降下データとする。各射 $V_i \to X_i$ が閉埋め込みならば、
この降下データは有効である。
\end{lemma}

\begin{proof}
閉埋め込みはアフィン射なので
（Morphisms, Lemma \ref{morphisms-lemma-closed-immersion-affine}）、
Lemma \ref{descent-lemma-affine} を適用できる。
\end{proof}


\section{準アフィン射の降下}
\label{descent-section-quasi-affine}

\noindent
本節では「準アフィン射は fpqc 被覆に対する降下を満たす」ことを示す。
厳密な主張は次のとおりである。

\begin{lemma}
\label{descent-lemma-quasi-affine}
スキーム $S$ をとる。
$\{X_i \to S\}_{i\in I}$ を fpqc 被覆とする
（Topologies, Definition \ref{topologies-definition-fpqc-covering} を参照せよ）。
$(V_i/X_i, \varphi_{ij})$ を $\{X_i \to S\}$ に関する
降下データとする。各射 $V_i \to X_i$ が準アフィンならば、
この降下データは有効である。
\end{lemma}

\begin{proof}
準アフィンであることは、任意の基底変換により保たれる
スキームの射の性質である。Morphisms, Lemmas
\ref{morphisms-lemma-characterize-quasi-affine} および
\ref{morphisms-lemma-base-change-quasi-affine} を参照せよ。
したがって Lemma \ref{descent-lemma-descending-types-morphisms} を適用でき、
fpqc 被覆がアフィン間の一つの平坦全射
$\{X \to S\}$ で与えられる場合に補題の主張を証明すれば十分である。
$X = \Spec(A)$ および $S = \Spec(R)$ と書くと、
$R \to A$ は忠実平坦な環準同型である。
$(V, \varphi)$ を $S$ 上の $X$ に関する降下データとし、
$\pi : V \to X$ が準アフィンであると仮定する。

\medskip\noindent
Morphisms, Lemma \ref{morphisms-lemma-characterize-quasi-affine} によれば、
これは
$$
V \longrightarrow \underline{\Spec}_X(\pi_*\mathcal{O}_V) = W
$$
が $X$ 上のスキームの準コンパクト開埋め込みであることを意味する。
射影 $\text{pr}_i : X \times_S X \to X$ は平坦なので、
$$
\text{pr}_0^*\pi_*\mathcal{O}_V =
(\pi \times \text{id}_X)_*\mathcal{O}_{V \times_S X}, \quad
\text{pr}_1^*\pi_*\mathcal{O}_V =
(\text{id}_X \times \pi)_*\mathcal{O}_{X \times_S V}
$$
を得る。これは平坦基底変換による
（Cohomology of Schemes, Lemma
\ref{coherent-lemma-flat-base-change-cohomology}）。
したがって、$X \times_S X$ 上の同型である
$\varphi : V \times_S X \to X \times_S V$ は、
$X \times_S X$ 上の準連接代数層の同型
$$
\varphi^\sharp :
\text{pr}_0^*\pi_*\mathcal{O}_V
\longrightarrow
\text{pr}_1^*\pi_*\mathcal{O}_V
$$
を誘導する。$\varphi$ のコサイクル条件は
$\varphi^\sharp$ のコサイクル条件を導く。言い換えれば、
これは $S$ 上の $X$ に関するアフィン・スキーム $W$ 上の
降下データ $\varphi'$ を生じ、さらに射 $V \to W$ は
降下データの射となる。したがって Lemma \ref{descent-lemma-affine}
（または準連接代数の降下の有効性）により、
降下データの同型 $(W, \varphi') \cong (X \times_S U', can)$ をもつ
スキーム $U' \to S$ を得る。なお、
Lemma \ref{descent-lemma-descending-property-affine} により $U'$ はアフィンである。

\medskip\noindent
いま $V$ は、降下データ
$$
can : X \times_S U' \times_S X \to X \times_S X \times_S U',
(x_0, u', x_1) \mapsto (x_0, x_1, u').
$$
のもとで安定な（準コンパクト）開集合
$V \subset X \times_S U'$ とみなせる。言い換えれば、
$S$ の同じ点に写る任意の $x_0, x_1, u'$ に対して
$(x_0, u') \in V \Rightarrow (x_1, u') \in V$ である。
$X \to S$ は全射なので、$V$ は射 $X \times_S U' \to U'$ による
ある部分集合 $U \subset U'$ の逆像であることが直ちに分かる。
$X \to S$ は準コンパクト、平坦かつ全射なので、
$X \times_S U' \to U'$ も準コンパクト、平坦かつ全射である。
したがって Morphisms, Lemma
\ref{morphisms-lemma-fpqc-quotient-topology} により、
部分集合 $U \subset U'$ は開であり、証明が完了する。
\end{proof}












\section{層による降下データの記述}
\label{descent-section-descent-data-sheaves}


\noindent
サイト上の被覆に対する降下データを考える別の方法を示す。

\begin{lemma}
\label{descent-lemma-descent-data-sheaves}
$\tau \in \{Zariski, fppf, \etale, smooth, syntomic\}$ とする\footnote{
fpqc が含まれていないのは誤植ではない。Topologies, Section
\ref{topologies-section-fpqc} の議論を参照せよ。}。
$\Sch_\tau$ を大 $\tau$-サイトとし、$S \in \Ob(\Sch_\tau)$ とする。
$\{S_i \to S\}_{i \in I}$ をサイト $(\Sch/S)_\tau$ における被覆とする。
次の圏同値が存在する。
$$
\left\{
\begin{matrix}
\text{降下データ }(X_i, \varphi_{ii'})\text{ であって}\\
\text{各 }X_i \in \Ob((\Sch/S)_\tau)\text{ となるもの}
\end{matrix}
\right\}
\leftrightarrow
\left\{
\begin{matrix}
(\Sch/S)_\tau\text{ 上の層 }F\text{ であって}\\
\text{各 }h_{S_i} \times F\text{ が表現可能となるもの}
\end{matrix}
\right\}.
$$
さらに、
\begin{enumerate}
\item 右辺で $h_{S_i} \times F$ を表現する対象は、
左辺のスキーム $X_i$ に対応する。
\item 層 $F$ が表現可能であるための必要十分条件は、
対応する降下データ $(X_i, \varphi_{ii'})$ が有効であることである。
\end{enumerate}
\end{lemma}

\begin{proof}
Section \ref{descent-section-fpqc-universal-effective-epimorphisms} で、
表現可能な前層はサイト $(\Sch/S)_\tau$ 上の層であることを見た。
さらに Yoneda の補題
（Categories, Lemma \ref{categories-lemma-yoneda}）により、
表現可能層の間の写像は表現対象の間の写像と一対一に対応する。
証明では以下、これらを断りなく用いる。

\medskip\noindent
右から左への関手を構成する。$F$ を $(\Sch/S)_\tau$ 上の層で、
各 $h_{S_i} \times F$ が表現可能なものとする。
この場合、$X_i$ を $(\Sch/S)_\tau$ における表現対象とする。
これには射 $X_i \to S_i$ が備わる。このとき
$X_i \times_S S_{i'}$ と $S_i \times_S X_{i'}$ はともに層
$h_{S_i} \times F \times h_{S_{i'}}$ を表現するので、同型
$$
\varphi_{ii'} : X_i \times_S S_{i'} \to S_i \times_S X_{i'}
$$
を得る。写像 $\varphi_{ii'}$ が $S_i \times_S S_{i'}$ 上の射であり、
コサイクル条件を満たすことは直ちに分かる。右から左への関手は、
構成 $F \mapsto (X_i, \varphi_{ii'})$ によって与えられる。

\medskip\noindent
左から右への関手を構成する。各 $i$ について、層 $h_{X_i}$ を
$F_i$ と表す。同型 $\varphi_{ii'}$ は同型
$$
\varphi_{ii'} :
F_i \times h_{S_{i'}}
\longrightarrow
h_{S_i} \times F_{i'}
$$
を $h_{S_i} \times h_{S_{i'}}$ 上で与える。$F$ を次の図式の余等化子とする。
$$
\xymatrix{
\coprod_{i, i'} F_i \times h_{S_{i'}}
\ar@<1ex>[rr]^-{\text{pr}_0}
\ar@<-1ex>[rr]_-{\text{pr}_1 \circ \varphi_{ii'}}
& &
\coprod_i F_i \ar[r]
&
F
}
$$
コサイクル条件により $h_{S_i} \times F$ は $F_i$ と同型であり、
したがって表現可能である。左から右への関手は、構成
$(X_i, \varphi_{ii'}) \mapsto F$ によって与えられる。

\medskip\noindent
これらの構成が互いに擬逆関手であることの確認は省略する。
最後の主張 (1), (2) は構成から従う。
\end{proof}

\begin{remark}
\label{descent-remark-what-product-means}
Lemma \ref{descent-lemma-descent-data-sheaves} の主張において、
$h_{S_i} \times F$ が表現可能であるという条件は、
$F$ の $(\Sch/S_i)_\tau$ への制限が表現可能であるという条件と同値である。
\end{remark}
